\documentclass[10pt,a4paper]{article}
\usepackage[margin=17mm,headheight=14pt]{geometry}
\usepackage{fontspec}\setmainfont{Latin Modern Roman}
\setsansfont{DejaVu Sans}[Scale=0.91]\setmonofont{DejaVu Sans Mono}[Scale=0.82]
\usepackage{amsmath,amssymb,mathtools,xcolor,tikz,pdflscape,fvextra}
\usetikzlibrary{arrows.meta,shapes.geometric}
\usepackage[most]{tcolorbox}\usepackage{longtable,booktabs,array,tabularx,fancyhdr,hyperref}
\definecolor{navy}{HTML}{203650}\definecolor{ink}{HTML}{243449}
\definecolor{muted}{HTML}{617389}\definecolor{panel}{HTML}{F3F6FA}
\definecolor{core}{HTML}{244766}\definecolor{api}{HTML}{58687C}
\definecolor{assembly}{HTML}{7562A3}\definecolor{compressed}{HTML}{76509E}
\hypersetup{colorlinks=true,linkcolor=navy,urlcolor=navy,
 pdftitle={Multi-scale convex Wolff: Theorem 5.2 Dependency Atlas},bookmarksopen=true}
\pagestyle{fancy}\fancyhf{}
\lhead{\sffamily\small Multi-scale convex Wolff | Theorem 5.2}
\rhead{\sffamily\small Dependency atlas}\cfoot{\sffamily\small\thepage}
\renewcommand{\headrulewidth}{0.25pt}
\setlength{\parindent}{0pt}\setlength{\parskip}{4pt}\setlength{\emergencystretch}{3em}
\newcommand{\BB}{\mathbb}\newcommand{\RR}{\mathbb R}\newcommand{\TT}{\mathcal T}\newcommand{\CC}{\mathcal C}
\newcommand{\AD}{\operatorname{AD}}\newcommand{\vol}{\operatorname{vol}}\newcommand{\norm}[1]{\left\|#1\right\|}
\newcommand{\nodecard}[1]{\hyperlink{node-#1}{\texttt{\detokenize{#1}}}}
\newcommand{\sectitle}[1]{\clearpage{\sffamily\Large\bfseries\color{navy}#1}\par\medskip}
\definecolor{setupbg}{HTML}{eef3f8}
\definecolor{setupfg}{HTML}{60748c}
\definecolor{planebg}{HTML}{e4f2ee}
\definecolor{planefg}{HTML}{257966}
\definecolor{grainsbg}{HTML}{e6f4ee}
\definecolor{grainsfg}{HTML}{23845b}
\definecolor{regularitybg}{HTML}{fff0db}
\definecolor{regularityfg}{HTML}{b66a0e}
\definecolor{slopebg}{HTML}{e7effc}
\definecolor{slopefg}{HTML}{356db0}
\definecolor{cinematicbg}{HTML}{eee8ff}
\definecolor{cinematicfg}{HTML}{7345b3}
\definecolor{radialbg}{HTML}{eaf4ff}
\definecolor{radialfg}{HTML}{2e779d}
\definecolor{externalbg}{HTML}{fff0f2}
\definecolor{externalfg}{HTML}{b24a68}
\definecolor{apibg}{HTML}{eef3f8}
\definecolor{apifg}{HTML}{60748c}

\begin{document}\thispagestyle{empty}
{\sffamily\small\color{muted}WANG--ZAHL / WZ1 + WZ2 / OCTOBER 11, 2026}\par\vspace{5mm}
{\sffamily\Huge\bfseries\color{navy}Multi-scale convex Wolff}\par
{\sffamily\Large\color{navy}Theorem 5.2 dependency atlas and Lean statement blueprint}\par\vspace{5mm}
\begin{tcolorbox}[colback=panel,colframe=navy,boxrule=0.4pt,arc=2mm]
\textbf{Sole root.} The discrete volume estimate in WZ2 Theorem 5.2, not the dimension conclusions of either paper.
\textbf{WZ1:} Wang--Zahl, \emph{Sticky Kakeya sets and the sticky Kakeya conjecture}, supplied 71-page JAMS version, DOI \href{https://doi.org/10.1090/jams/1067}{10.1090/jams/1067}.
\textbf{WZ2:} Wang--Zahl, \emph{The Assouad dimension of Kakeya sets in $\mathbb R^3$}, supplied 54-page Inventiones version, DOI \href{https://doi.org/10.1007/s00222-025-01336-x}{10.1007/s00222-025-01336-x}.
WZ1 page references use its printed 1--71 pagination. WZ2 references use the one-based PDF page, because the supplied version does not display a numerical running page.
\end{tcolorbox}
\textbf{Coverage.} 102 nodes, 237 recorded edges, one core overview and nine detailed views. Every node is an ancestor of the requested root. The HTML adds the complete unified graph, zoom, pan, search and immediate-neighbour tracing. Each card supplies its mathematical statement, Lean statement sketch, support areas and source/transfer note.
\textbf{What survives from the reference atlas.} The WZ1 multiscale, planiness, graininess, slope regularity and internal Section 8 projection branches. What changes is the extremal class, normalized tube counts, nonflatness endpoint and final twisted projection argument.

\textbf{Reading order.} Begin at Theorem 5.2 and read downwards through the overview. Follow the section views into the WZ1 proof interfaces and the new WZ2 constructions. Every diagram label is clickable and opens its node card.

\textbf{Arrow semantics.} Arrows point from prerequisites to consequences, generally upwards. Solid blue-grey edges are proof dependencies; dashed grey edges expose definitions or standard operations; dash-dot purple edges assemble an output into a later hypothesis. Long-dashed purple edges occur only in the overview and compress a verified path. Crossings are not vertices. Definitions/API use rectangles, results use ellipses, and external sources use dashed rounded boxes. Colours identify proof branches, not formalization status.

\textbf{Source status.} WZ1 configuration cards marked ``/ CW'' are transferred statements with WZ2 extremality and loss conventions. They do not assert the original direction-separated statements verbatim. Unnumbered construction cards are proof operations in the cited argument, not new external axioms. The mathematical proof is described by the papers; the Lean declarations are uncompiled statement sketches.

\textbf{Excluded.} WZ1 QStick, Lemma 2.2, Proposition 2.1 and the Hausdorff-dimension roots; WZ2 Sections 1, 3, 4, the deductions surrounding Theorem 5.2 in Section 5, and Section 8. WZ2 self-similar convex Wolff axioms (Definition 2.15) are also unnecessary here. Historical provenance leaves are removed when they supply no independent input.
\sectitle{Conventions and porting contracts}
In the structural branch write $\sigma=\sigma_{\rm CW}>0$ and $C=(\mathcal T,Y)_\delta$. This deficit belongs to the WZ2 admissible class; it is not identified with WZ1's $\sigma_3$. The contradiction needs only a small projected image of order $\delta^{\sigma-O(\varepsilon)}$.

\textbf{Cardinality normalization.} Retain $\#\mathcal T$ in every multiplicity calculation. A dense shading has mass $\gtrsim\delta^\eta\delta^2\#\mathcal T$, whereas an extremal union has size at most $\delta^{\sigma-\eta}$. Hence typical multiplicity has scale $\delta^{2-\sigma}\#\mathcal T$ with reserved power losses. No upper bound $\#\mathcal T\lesssim\delta^{-2}$ is imported. In the transferred WZ1 calculations the matching cardinality factors cancel.

\textbf{Refinement normalization.} Shading refinement controls total incidence mass. Obtaining a uniformly dense family or changing tubes can alter normalized convex counts. The partitioning cover and each child fibre are rebalanced with explicit loss witnesses; neither point selection nor arbitrary tube deletion is treated as automatic multi-scale inheritance.

\textbf{Three entropy conditions.} ADsets have an upper bound at every coarser mesh and pass to subsets. Frostman sets are normalized by the total mesh entropy. Katz--Tao sets have an absolute $(r/\delta)^s$ bound. WZ1 hypergraph uniform density and WZ2 occupied-cell uniformity are separate notions.

\textbf{Section 7 scales.} Use $q$ for a local parameter cube, $h=\delta/q$ for its rescaled mesh, and $r=q$ for the coarse target neighbourhood. Fixed-$c$ slicing means a $\delta$-slab followed by freezing with $O(\delta)$ error. The common $ct$ term is removed by an area-preserving target shear.

\textbf{External stopping rule.} Open every invoked WZ1 proof to the reference atlas's granularity. Stop at the other papers' cited results: endpoint multilinear Kakeya, Lipschitz extension, Orponen--Shmerkin--Wang radial bootstrap, Pramanik--Yang--Zahl cinematic maximal, F\"assler--Orponen content extraction, Orponen--Shmerkin uniform refinement, and Shmerkin--Wang scale decomposition/merging.
\medskip\textbf{WZ1 to WZ2 crosswalk.}\par
\begin{longtable}{@{}p{.23\linewidth}p{.25\linewidth}p{.47\linewidth}@{}}\toprule WZ1 interface & WZ2 interface & Required change\\\midrule\endhead
Admissibility / Def. 3.1 & Sec. 6.2 / Def. 6.1 & New deficit; normalized shading density; no delta\textasciicircum{}(-2) upper cardinality.\\
Def. 3.2 & Def. 2.8 and Sec. 6.1 & John normalization plus compatible shaded re-cubulation.\\
Prop. 3.2 & Prop. 6.2 / Lem. 2.13 & Multi-scale CW coarse/fine witnesses and normalized multiplicities.\\
Prop. 4.1 & Prop. 6.3 & Expanded WZ1 planiness and graininess proof.\\
Prop. 5.1 & Prop. 6.4 & Expanded WZ1 C2 construction, including its Section 8 branch.\\
Lem. 6.3 & Lem. 6.6 & Steps 1-3 retained; new convex slab count in Step 4.\\
Prop. 6.1 & Prop. 6.5 (1 double-prime) & Comparable slope slab and anisotropic normalization retained.\\
Cor. 6.5 & Lem. 6.7 & Small projected area is the needed output; the old large-norm branch is removed.\\
Prop. 7.1 & Prop. 6.8 via Prop. 7.1 & New Frostman extraction, Katz-Tao replication and scale iteration.\\
Thm. 7.2 / Lem. 7.3 & Thm. 7.13 / Lem. 7.12 proof & Cinematic theorem and model geometry retained.\\
\bottomrule\end{longtable}
\sectitle{Source discrepancies and loss bookkeeping}
These notes affect an interface actually used by Theorem 5.2. They distinguish the printed source from the weaker or loss-budget form carried by this atlas; they are not a claim to have formalized or independently reconstructed every omitted proof.
\begin{longtable}{@{}p{.29\linewidth}p{.66\linewidth}@{}}\toprule Source location & Roadmap treatment\\\midrule\endhead
WZ2 (6.11), PDF p. 43 & The printed exponent is sigma+epsilon. The direct AD/Fubini consequence is sigma-epsilon with losses; this weaker conclusion proves the contradiction. Card W2L6\_7 states and derives it.\\
WZ2 Def. 7.2, PDF p. 44 & The printed ratio (delta/r)\textasciicircum{}s is reversed. Use (r/delta)\textasciicircum{}s, as required by Lemma 7.11 and the cinematic absolute ball condition.\\
WZ2 Lem. 7.11 / Lem. 7.12 & Katz-Tao extraction is only nearly saturated. The replication card retains its density deficit kappa instead of silently using the saturated cardinality hypothesis.\\
WZ2 Lem. 7.12 scale bookkeeping & q is the parameter cube size, h=delta/q its rescaled mesh, and r=q the coarse target scale. Reserve a separate cinematic exponent so a delta-loss becomes an h-loss.\\
WZ2 Prop. 7.1 iteration, PDF p. 51 & Use a finite bound depending on epsilon with the required logarithmic factor, and telescope to the terminal r\_N. The literal N <= epsilon\textasciicircum{}(-2) and last rho\_2 are unnecessary.\\
\bottomrule\end{longtable}
\textbf{Why the weaker small-image estimate is enough.} Choose the lower projection exponent $\sigma/4$. Make the loss in Proposition 6.5 smaller than both the lower theorem's admissibility threshold and $\sigma/4$. Then $|X_f(C)|\lesssim\delta^{3\sigma/4}$ while $|X_f(C)|\ge\delta^{\sigma/4}$, impossible at an arbitrarily small scale.

\textbf{A fixed iteration bound.} For $r_{j+1}>r_j^{1-\theta}$, the exponent of $\delta$ at the $j$th scale is at most $(1-\theta)^j$. A fixed $N\ge\log\theta/\log(1-\theta)$ reaches $r_N\ge\delta^\theta$. Reserve the losses for this finite $N(\theta)$ before choosing the initial scale.

\textbf{Formalization boundary.} The ported WZ1 statements and the WZ2 sketches still require complete implementations of cover rebuilding, affine re-cubulation and parameter losses. The witness cards expose these obligations. No node is certified as already formalized.
\sectitle{Lean interface conventions and support}
All Lean blocks are statement sketches. Use \texttt{import Mathlib} for an initial prototype. The suggested modules identify support areas, not a tested minimal import list. Local \texttt{Hyp} and data-packet names stand for the full mathematical hypotheses on their own cards.

Use \texttt{Config 3 delta} for a finite shaded configuration, \texttt{vol S} only after proving finite volume, and \texttt{coverNumR} for a real coercion of a chosen finite covering number. The affine line chart, covering number, grids, convex-Wolff hierarchy, hypergraph contracts and $C^2$ jet metric need project definitions. A declaration marked \texttt{axiom} is an external mathematical interface, not a proposed additional axiom for the internal proof.
{\footnotesize\textbf{G: Euclidean geometry.} \texttt{Mathlib.Analysis.InnerProductSpace.PiL2}\par EuclideanSpace, inner, norm, Metric.ball; the affine-line chart is project-defined.}\par
{\footnotesize\textbf{M: Measure and integration.} \texttt{Mathlib.MeasureTheory.Measure.Lebesgue.Basic}\par volume, measurable sets, product measures, integrals; real volumes need finiteness.}\par
{\footnotesize\textbf{F: Finite combinatorics.} \texttt{Mathlib.Data.Finset.Basic}\par Finset.card, filter, product, sum; finite grids and hypergraphs need project definitions.}\par
{\footnotesize\textbf{Q: Powers and suprema.} \texttt{Mathlib.Analysis.SpecialFunctions.Pow.Real}\par Real.rpow, Real.log, sSup; positive-mesh and boundedness certificates.}\par
{\footnotesize\textbf{L: Lipschitz maps.} \texttt{Mathlib.Topology.MetricSpace.Lipschitz}\par LipschitzWith, LipschitzOnWith and scalar extension; constants use NNReal.}\par
{\footnotesize\textbf{C: Calculus and patches.} \texttt{Mathlib.Analysis.Calculus.ContDiff.Basic}\par ContDiff, ContDiffOn, deriv, HasDerivAt; also Deriv.Basic and BumpFunction.Basic.}\par
{\footnotesize\textbf{P: Lp estimates.} \texttt{Mathlib.MeasureTheory.Function.LpSeminorm.Defs}\par eLpNorm, MemLp and Holder; exponent 3/2 is a real/ENNReal value.}\par
{\footnotesize\textbf{A: Affine transport.} \texttt{Mathlib.Analysis.Normed.Module.Basic}\par ContinuousLinearEquiv, affine maps and determinants; transport and re-cubulation need project proofs.}\par
\sectitle{Node index}
\begin{longtable}{@{}p{.18\linewidth}p{.30\linewidth}p{.47\linewidth}@{}}\toprule ID / origin & Paper label & Interface\\\midrule\endhead
\nodecard{Igeom} / WZ2 & Geometry API & Lines, tubes, cubes, shadings\\
\nodecard{Ientropy} / API & Entropy API & Finite covering numbers\\
\nodecard{Imass} / WZ2 & Refinement API & Mass and constant multiplicity\\
\nodecard{Icover} / WZ2 & Cover API & Balanced coarse/fine configurations\\
\nodecard{Iaffine} / API & Affine transport API & Jacobian, direction and grid transport\\
\nodecard{Icalc} / API & Calculus API & C2 bounds and compactly supported patches\\
\nodecard{Iparams} / API & Quantifier API & Small-parameter and finite-iteration contracts\\
\nodecard{Iadgap} / API & AD exponent gap & Why the expansion alternative is impossible\\
\nodecard{L3_3} / WZ1 & WZ1 Lemma 3.3 / CW & Local rescaling and multiplicity control\\
\nodecard{L3_4} / WZ1 & WZ1 Lemma 3.4 / CW & Robust transversality\\
\nodecard{L3_5} / WZ1 & WZ1 Lemma 3.5 / CW & Mild anisotropic rescaling preserves extremality\\
\nodecard{Dplane} / WZ2 & WZ2 Section 6.1: plane map & Plane map\\
\nodecard{DAD} / WZ2 & WZ2 Section 6.1: ADset & Upper AD covering condition\\
\nodecard{T4_2} / WZ1 & WZ1 Theorem 4.2 & Endpoint multilinear Kakeya input\\
\nodecard{L4_3} / WZ1 & WZ1 Lemma 4.3 / CW & Weak planiness\\
\nodecard{L4_4} / WZ1 & WZ1 Lemma 4.4 / CW & Coarsen weak planiness to a plane map\\
\nodecard{L4_5} / WZ1 & WZ1 Lemma 4.5 / CW & Small oscillation on a fixed cube scale\\
\nodecard{C4_6} / WZ1 & WZ1 Corollary 4.6 / CW & Small oscillation for nearby points\\
\nodecard{L4_7} / WZ1 & WZ1 Lemma 4.7 / CW & Iterate to a Lipschitz plane map\\
\nodecard{L4_8} / WZ1 & WZ1 Lemma 4.8 / CW & Existence of Lipschitz plane data\\
\nodecard{L4_9} / WZ1 & WZ1 Lemma 4.9 / CW & Grain count in tau-cubes\\
\nodecard{L4_10} / WZ1 & WZ1 Lemma 4.10 / CW & Grain count in thin prisms\\
\nodecard{L4_11} / WZ1 & WZ1 Lemma 4.11 / CW & AD spacing at one grain scale\\
\nodecard{L4_12} / WZ1 & WZ1 Lemma 4.12 / CW & AD spacing at every grain scale\\
\nodecard{T5_2} / WZ1 & WZ1 Theorem 5.2 & Dot-product expansion or orthogonal lines\\
\nodecard{L5_3} / WZ1 & WZ1 Lemma 5.3 & Local linearity from the two grain structures\\
\nodecard{L5_4} / WZ1 & WZ1 Lemma 5.4 / CW & Piecewise linear approximation at one scale\\
\nodecard{L5_5} / WZ1 & WZ1 Lemma 5.5 / CW & Compatible approximations at many scales\\
\nodecard{Dtrap} / WZ1 & WZ1 Definition 5.1 & Vertical trapezoids and parents\\
\nodecard{C5_6} / WZ1 & WZ1 Corollary 5.6 / CW & A nested trapezoid tree for the slope graph\\
\nodecard{L5_8} / WZ1 & WZ1 Lemma 5.8 & Smooth extension of an affine segment\\
\nodecard{P5_7} / WZ1 & WZ1 Proposition 5.7 & A trapezoid tree produces a C2 approximation\\
\nodecard{L6_2} / WZ1 & WZ1 Lemma 6.2 & A thin projection forces near orthogonality\\
\nodecard{L6_4} / WZ1 & WZ1 Lemma 6.4 / CW & A thick slab with comparable nonzero derivative\\
\nodecard{Dtwist} / WZ1 & WZ1 Def. 6.1 / WZ2 Sec. 6.5 & Twisted projection\\
\nodecard{Dcinematic} / WZ1 & WZ1 Definition 7.1 & Cinematic families of C2 functions\\
\nodecard{T7_2} / WZ1/WZ2 & WZ1 Thm. 7.2 = WZ2 Thm. 7.13 & Cinematic maximal estimate\\
\nodecard{L7_3} / WZ1 & WZ1 Lemma 7.3 & The twisted-line family is cinematic\\
\nodecard{DFrost} / WZ1 & WZ1 Definition 8.1 & Normalized discretized Frostman sets\\
\nodecard{Dhyp} / WZ1 & WZ1 Definitions 8.2-8.4 & Hypergraph fibres and uniform density\\
\nodecard{L8_1} / WZ1 & WZ1 Lemma 8.1 & Hypergraph refinement\\
\nodecard{T8_2} / WZ1 & WZ1 Theorem 8.2 & A discretized Kaufman estimate with dense fibres\\
\nodecard{Dthin} / WZ1 & WZ1 Definition 8.5 & Thin tubes for probability measures\\
\nodecard{DthinD} / WZ1 & WZ1 Definition 8.5' & Discretized thin tubes\\
\nodecard{L8_3} / WZ1 & WZ1 Lemma 8.3 & Pass from measure thin tubes to grid thin tubes\\
\nodecard{Dsep} / WZ1 & WZ1 Definition 8.6 & Standard separation\\
\nodecard{L8_4} / WZ1 & WZ1 Lemma 8.4 & Thin tubes force a large dot-product image\\
\nodecard{P8_6} / WZ1 & WZ1 Proposition 8.6 & Radial bootstrap for measures\\
\nodecard{P8_6p} / WZ1 & WZ1 Proposition 8.6' & Radial bootstrap for grid sets\\
\nodecard{C8_7} / WZ1 & WZ1 Corollary 8.7 & Iterate the bootstrap up to exponent 1\\
\nodecard{L8_8} / WZ1 & WZ1 Lemma 8.8 & Obtain a fixed 1/4-thin-tube seed\\
\nodecard{P8_5} / WZ1 & WZ1 Proposition 8.5 & Strip nonconcentration forces expansion\\
\nodecard{L8_10} / WZ1 & WZ1 Lemma 8.10 & Two-ends localization to a strip\\
\nodecard{L8_11} / WZ1 & WZ1 Lemma 8.11 & Restore Frostman after anisotropic rescaling\\
\nodecard{L8_12} / WZ1 & WZ1 Lemma 8.12 & Frostman coarsening after pigeonholing\\
\nodecard{L8_13} / WZ1 & WZ1 Lemma 8.13 & Comparable orthogonal strips or expansion\\
\nodecard{P8_9} / WZ1 & WZ1 Proposition 8.9 & Separated, uniform version of the dichotomy\\
\nodecard{T5_2p} / WZ1 & WZ1 Theorem 5.2' & Frostman form, without separation or uniformity\\
\nodecard{XBCT} / External & External [1] / [8] & Endpoint multilinear Kakeya\\
\nodecard{XExt} / External & External [29] / Mathlib & Scalar Lipschitz extension\\
\nodecard{XOSW} / External & External [35, Lemma 2.9] & Radial thin-tube bootstrap\\
\nodecard{XPYZ} / External & External: Pramanik-Yang-Zahl & Cinematic L3/2 maximal theorem\\
\nodecard{W2CWA} / WZ2 & WZ2 Definition 2.5 & Normalized convex-containment count\\
\nodecard{W2Unit} / WZ2 & WZ2 Definition 2.8 / Sec. 6.1 & Geometric and shaded unit rescaling\\
\nodecard{W2MW} / WZ2 & WZ2 Definition 2.12 & Convex Wolff axioms at every scale\\
\nodecard{W2L2_13} / WZ2 & WZ2 Lemma 2.13 & Coarse and fine multi-scale inheritance\\
\nodecard{Iwolffbalance} / WZ2 & WZ2 proof operation & Restriction and rebuilding of cover witnesses\\
\nodecard{W2Sigma} / WZ2 & WZ2 Section 6.2 & The counterexample deficit sigma\_CW\\
\nodecard{Ienvelope} / WZ2 & WZ2 Section 6.2: consequence & Uniform lower envelope for admissible unions\\
\nodecard{W2D6_1} / WZ2 & WZ2 Definition 6.1 & Extremal configurations under convex Wolff\\
\nodecard{W2Exists} / WZ2 & WZ2 Section 6.2 / WZ1 Lem. 3.1 & Extremizers exist at arbitrarily small scales\\
\nodecard{W2P6_2} / WZ2 & WZ2 Proposition 6.2 & Balanced extremal coarse and fine families\\
\nodecard{W2P6_3} / WZ2 & WZ2 Proposition 6.3 & Global and local grains with Lipschitz data\\
\nodecard{W2P6_4} / WZ2 & WZ2 Proposition 6.4 & Global grains with a C2 slope\\
\nodecard{W1SlabPrep} / WZ1 & WZ1 Lemma 6.3: Steps 1-3 & Popular grains in a thin horizontal slab\\
\nodecard{W2SlabCount} / WZ2 & WZ2 Lemma 6.6: new Step 4 & Convex Wolff contradicts a popular thin slab\\
\nodecard{W2L6_6} / WZ2 & WZ2 Lemma 6.6 & The slope cannot be flat on a popular slab\\
\nodecard{W2P6_5} / WZ2 & WZ2 Proposition 6.5, part 1 double-prime & Nonsingular C2 global grains\\
\nodecard{W2L6_7} / WZ2 & WZ2 Lemma 6.7: small-image interface & Extremality produces a small twisted image\\
\nodecard{W2TW} / WZ2 & WZ2 Section 7: Tube Wolff & The weaker radius-squared tube count\\
\nodecard{W2DFrost} / WZ2 & WZ2 Definition 7.3 & Normalized Frostman entropy\\
\nodecard{W2DKT} / WZ2 & WZ2 Definition 7.2 & Absolute Katz-Tao entropy bound\\
\nodecard{W2DUnif} / WZ2 & WZ2 Definition 7.8 & Uniform occupied-cell counts\\
\nodecard{W2L7_6} / WZ2 & WZ2 Lemma 7.6 & Frostman sets contain large Katz-Tao subsets\\
\nodecard{XFO} / External & External: Fässler-Orponen Lem. 3.13 & Content-to-Katz-Tao extraction\\
\nodecard{W2L7_9} / WZ2 & WZ2 Lemma 7.9 & Sparse-scale uniform refinement\\
\nodecard{XOSuniform} / External & External: Orponen-Shmerkin Lem. 2.15 & Uniform refinement black box\\
\nodecard{XSWdecomp} / External & External: Shmerkin-Wang Lem. 5.21 & Decompose the scale-entropy profile\\
\nodecard{XSWmerge} / External & External: Shmerkin-Wang Lem. 5.20 & Merge adjacent entropy intervals\\
\nodecard{W2L7_10} / WZ2 & WZ2 Lemma 7.10 & Zoom to Frostman fibres of a uniform set\\
\nodecard{W2L7_11} / WZ2 & WZ2 Lemma 7.11 & Replicate a local Katz-Tao cluster\\
\nodecard{W2ImageUniform} / WZ2 & WZ2 Lemma 7.12: preparation & Uniformize the projected incidence set\\
\nodecard{W2Slice} / WZ2 & WZ2 Lemma 7.12: parameter slicing & A rich fixed-c slice of R4 parameters\\
\nodecard{W2JetCount} / WZ2 & WZ2 Lemma 7.12: cinematic assembly & Convert Katz-Tao triples to separated curves\\
\nodecard{W2AreaCompare} / WZ2 & WZ2 Lemma 7.12: Holder and translations & From cinematic norm to a local area ratio\\
\nodecard{W2L7_12} / WZ2 & WZ2 Lemma 7.12 & One-step fine-to-coarse projected-area comparison\\
\nodecard{IcoarseTW} / WZ2 & WZ2 proof operation & Uniform Tube-Wolff coarsening\\
\nodecard{W2Iterate} / WZ2 & WZ2 Proposition 7.1: iteration & Telescope the ratios up to a coarse scale\\
\nodecard{W2P7_1} / WZ2 & WZ2 Proposition 7.1 & Twisted-image lower bound under Tube Wolff\\
\nodecard{W2P6_8} / WZ2 & WZ2 Proposition 6.8 & Twisted-image lower bound under convex Wolff\\
\nodecard{W2Zero} / WZ2 & WZ2 Section 6.5: contradiction & The multi-scale convex Wolff deficit is zero\\
\nodecard{W2T5_2} / WZ2 & WZ2 Theorem 5.2 & Discrete union estimate for multi-scale convex Wolff tubes\\
\bottomrule\end{longtable}
\sectitle{Core proof DAG}
\hypertarget{view-overview}{}
{\sffamily\small\color{muted}Theorem 5.2 alone: conclusions above, supporting geometry and analytic interfaces below.}\par
\begin{center}\resizebox{493.00bp}{!}{\begin{tikzpicture}[x=1bp,y=1bp, >={Stealth[length=4bp,width=3.2bp]}]
\path[use as bounding box] (0,0) rectangle (757.46,971.00);
\fill[panel,rounded corners=5bp] (0,0) rectangle (757.46,971.00);
\draw[line width=3.5bp,panel] (349.23,135.00) .. controls (349.23,150.41) and (327.99,153.09) .. (327.99,168.50);
\draw[->,line width=1.0bp,api,dash pattern=on 5bp off 4bp] (349.23,135.00) .. controls (349.23,150.41) and (327.99,153.09) .. (327.99,168.50);
\draw[line width=3.5bp,panel] (390.89,135.00) .. controls (390.89,143.10) and (47.50,145.80) .. (47.50,145.80) .. controls (47.50,174.68) and (47.50,213.18) .. (47.50,242.06) .. controls (47.50,242.06) and (449.28,244.76) .. (449.28,252.86);
\draw[->,line width=1.0bp,core] (390.89,135.00) .. controls (390.89,143.10) and (47.50,145.80) .. (47.50,145.80) .. controls (47.50,174.68) and (47.50,213.18) .. (47.50,242.06) .. controls (47.50,242.06) and (449.28,244.76) .. (449.28,252.86);
\draw[line width=3.5bp,panel] (432.15,211.50) .. controls (432.15,230.53) and (311.67,233.84) .. (311.67,252.86);
\draw[->,line width=1.0bp,core] (432.15,211.50) .. controls (432.15,230.53) and (311.67,233.84) .. (311.67,252.86);
\draw[line width=3.5bp,panel] (328.40,135.00) .. controls (328.40,143.10) and (44.00,145.80) .. (44.00,145.80) .. controls (44.00,177.74) and (44.00,220.32) .. (44.00,252.26) .. controls (44.00,252.26) and (203.58,254.96) .. (203.58,263.06);
\draw[->,line width=1.0bp,core] (328.40,135.00) .. controls (328.40,143.10) and (44.00,145.80) .. (44.00,145.80) .. controls (44.00,177.74) and (44.00,220.32) .. (44.00,252.26) .. controls (44.00,252.26) and (203.58,254.96) .. (203.58,263.06);
\draw[line width=3.5bp,panel] (327.99,211.50) .. controls (327.99,232.89) and (151.50,236.61) .. (151.50,258.00);
\draw[->,line width=1.0bp,core] (327.99,211.50) .. controls (327.99,232.89) and (151.50,236.61) .. (151.50,258.00);
\draw[line width=3.5bp,panel] (151.50,305.00) .. controls (151.50,328.72) and (126.40,332.84) .. (126.40,356.56);
\draw[->,line width=1.0bp,core] (151.50,305.00) .. controls (151.50,328.72) and (126.40,332.84) .. (126.40,356.56);
\draw[line width=3.5bp,panel] (178.48,398.50) .. controls (178.48,406.60) and (40.50,409.30) .. (40.50,409.30) .. controls (40.50,442.44) and (40.50,486.62) .. (40.50,519.76) .. controls (40.50,519.76) and (126.40,522.46) .. (126.40,530.56);
\draw[->,line width=1.0bp,core] (178.48,398.50) .. controls (178.48,406.60) and (40.50,409.30) .. (40.50,409.30) .. controls (40.50,442.44) and (40.50,486.62) .. (40.50,519.76) .. controls (40.50,519.76) and (126.40,522.46) .. (126.40,530.56);
\draw[line width=3.5bp,panel] (279.48,727.00) .. controls (279.48,741.75) and (227.40,744.31) .. (227.40,759.06);
\draw[->,line width=1.0bp,core] (279.48,727.00) .. controls (279.48,741.75) and (227.40,744.31) .. (227.40,759.06);
\draw[line width=3.5bp,panel] (609.46,316.00) .. controls (609.46,324.10) and (713.46,326.80) .. (713.46,326.80) .. controls (713.46,383.17) and (713.46,458.33) .. (713.46,514.70) .. controls (713.46,514.70) and (380.48,517.40) .. (380.48,525.50);
\draw[->,line width=1.0bp,core] (609.46,316.00) .. controls (609.46,324.10) and (713.46,326.80) .. (713.46,326.80) .. controls (713.46,383.17) and (713.46,458.33) .. (713.46,514.70) .. controls (713.46,514.70) and (380.48,517.40) .. (380.48,525.50);
\draw[line width=3.5bp,panel] (279.48,801.00) .. controls (279.48,815.75) and (328.40,818.31) .. (328.40,833.06);
\draw[->,line width=1.0bp,core] (279.48,801.00) .. controls (279.48,815.75) and (328.40,818.31) .. (328.40,833.06);
\draw[line width=3.5bp,panel] (411.73,135.00) .. controls (411.73,143.10) and (716.96,145.80) .. (716.96,145.80) .. controls (716.96,347.22) and (716.96,615.78) .. (716.96,817.20) .. controls (716.96,817.20) and (380.48,819.90) .. (380.48,828.00);
\draw[->,line width=1.0bp,api,dash pattern=on 5bp off 4bp] (411.73,135.00) .. controls (411.73,143.10) and (716.96,145.80) .. (716.96,145.80) .. controls (716.96,347.22) and (716.96,615.78) .. (716.96,817.20) .. controls (716.96,817.20) and (380.48,819.90) .. (380.48,828.00);
\draw[line width=3.5bp,panel] (126.40,567.44) .. controls (126.40,585.18) and (279.48,588.26) .. (279.48,606.00);
\draw[->,line width=1.0bp,assembly,dash pattern=on 7bp off 3bp on 1bp off 3bp] (126.40,567.44) .. controls (126.40,585.18) and (279.48,588.26) .. (279.48,606.00);
\draw[line width=3.5bp,panel] (328.40,65.00) .. controls (328.40,73.10) and (37.00,75.80) .. (37.00,75.80) .. controls (37.00,99.77) and (37.00,131.73) .. (37.00,155.70) .. controls (37.00,155.70) and (178.07,158.40) .. (178.07,166.50);
\draw[->,line width=1.0bp,core] (328.40,65.00) .. controls (328.40,73.10) and (37.00,75.80) .. (37.00,75.80) .. controls (37.00,99.77) and (37.00,131.73) .. (37.00,155.70) .. controls (37.00,155.70) and (178.07,158.40) .. (178.07,166.50);
\draw[line width=3.5bp,panel] (432.56,65.00) .. controls (432.56,77.42) and (380.48,79.58) .. (380.48,92.00);
\draw[->,line width=1.0bp,core] (432.56,65.00) .. controls (432.56,77.42) and (380.48,79.58) .. (380.48,92.00);
\draw[line width=3.5bp,panel] (397.84,65.00) .. controls (397.84,73.10) and (33.50,75.80) .. (33.50,75.80) .. controls (33.50,100.37) and (33.50,133.13) .. (33.50,157.70) .. controls (33.50,157.70) and (380.07,160.40) .. (380.07,168.50);
\draw[->,line width=1.0bp,core] (397.84,65.00) .. controls (397.84,73.10) and (33.50,75.80) .. (33.50,75.80) .. controls (33.50,100.37) and (33.50,133.13) .. (33.50,157.70) .. controls (33.50,157.70) and (380.07,160.40) .. (380.07,168.50);
\draw[line width=3.5bp,panel] (370.06,135.00) .. controls (370.06,150.41) and (432.15,153.09) .. (432.15,168.50);
\draw[->,line width=1.0bp,core] (370.06,135.00) .. controls (370.06,150.41) and (432.15,153.09) .. (432.15,168.50);
\draw[line width=3.5bp,panel] (178.07,213.50) .. controls (178.07,236.30) and (99.42,240.26) .. (99.42,263.06);
\draw[->,line width=1.0bp,core] (178.07,213.50) .. controls (178.07,236.30) and (99.42,240.26) .. (99.42,263.06);
\draw[line width=3.5bp,panel] (230.56,567.44) .. controls (230.56,575.54) and (30.00,578.24) .. (30.00,578.24) .. controls (30.00,607.05) and (30.00,645.45) .. (30.00,674.26) .. controls (30.00,674.26) and (227.40,676.96) .. (227.40,685.06);
\draw[->,line width=1.0bp,assembly,dash pattern=on 7bp off 3bp on 1bp off 3bp] (230.56,567.44) .. controls (230.56,575.54) and (30.00,578.24) .. (30.00,578.24) .. controls (30.00,607.05) and (30.00,645.45) .. (30.00,674.26) .. controls (30.00,674.26) and (227.40,676.96) .. (227.40,685.06);
\draw[line width=3.5bp,panel] (363.12,65.00) .. controls (363.12,73.10) and (26.50,75.80) .. (26.50,75.80) .. controls (26.50,277.54) and (26.50,546.52) .. (26.50,748.26) .. controls (26.50,748.26) and (331.56,750.96) .. (331.56,759.06);
\draw[->,line width=1.0bp,api,dash pattern=on 5bp off 4bp] (363.12,65.00) .. controls (363.12,73.10) and (26.50,75.80) .. (26.50,75.80) .. controls (26.50,277.54) and (26.50,546.52) .. (26.50,748.26) .. controls (26.50,748.26) and (331.56,750.96) .. (331.56,759.06);
\draw[line width=3.5bp,panel] (530.40,403.00) .. controls (530.40,411.10) and (720.46,413.80) .. (720.46,413.80) .. controls (720.46,468.27) and (720.46,540.91) .. (720.46,595.38) .. controls (720.46,595.38) and (491.89,598.08) .. (491.89,606.18);
\draw[->,line width=1.0bp,api,dash pattern=on 5bp off 4bp] (530.40,403.00) .. controls (530.40,411.10) and (720.46,413.80) .. (720.46,413.80) .. controls (720.46,468.27) and (720.46,540.91) .. (720.46,595.38) .. controls (720.46,595.38) and (491.89,598.08) .. (491.89,606.18);
\draw[line width=3.5bp,panel] (481.48,727.00) .. controls (481.48,741.75) and (429.40,744.31) .. (429.40,759.06);
\draw[->,line width=1.0bp,core] (481.48,727.00) .. controls (481.48,741.75) and (429.40,744.31) .. (429.40,759.06);
\draw[line width=3.5bp,panel] (634.56,403.00) .. controls (634.56,411.10) and (723.96,413.80) .. (723.96,413.80) .. controls (723.96,514.14) and (723.96,647.92) .. (723.96,748.26) .. controls (723.96,748.26) and (533.56,750.96) .. (533.56,759.06);
\draw[->,line width=1.0bp,core] (634.56,403.00) .. controls (634.56,411.10) and (723.96,413.80) .. (723.96,413.80) .. controls (723.96,514.14) and (723.96,647.92) .. (723.96,748.26) .. controls (723.96,748.26) and (533.56,750.96) .. (533.56,759.06);
\draw[line width=3.5bp,panel] (481.48,801.00) .. controls (481.48,815.75) and (432.56,818.31) .. (432.56,833.06);
\draw[->,line width=1.0bp,core] (481.48,801.00) .. controls (481.48,815.75) and (432.56,818.31) .. (432.56,833.06);
\draw[line width=3.5bp,panel] (380.48,875.00) .. controls (380.48,889.75) and (328.40,892.31) .. (328.40,907.06);
\draw[->,line width=1.0bp,core] (380.48,875.00) .. controls (380.48,889.75) and (328.40,892.31) .. (328.40,907.06);
\draw[line width=3.5bp,panel] (432.56,135.00) .. controls (432.56,143.10) and (727.46,145.80) .. (727.46,145.80) .. controls (727.46,370.94) and (727.46,671.12) .. (727.46,896.26) .. controls (727.46,896.26) and (432.56,898.96) .. (432.56,907.06);
\draw[->,line width=1.0bp,core] (432.56,135.00) .. controls (432.56,143.10) and (727.46,145.80) .. (727.46,145.80) .. controls (727.46,370.94) and (727.46,671.12) .. (727.46,896.26) .. controls (727.46,896.26) and (432.56,898.96) .. (432.56,907.06);
\draw[line width=3.5bp,panel] (380.48,405.00) .. controls (380.48,413.10) and (23.00,415.80) .. (23.00,415.80) .. controls (23.00,446.99) and (23.00,488.57) .. (23.00,519.76) .. controls (23.00,519.76) and (230.56,522.46) .. (230.56,530.56);
\draw[->,line width=1.0bp,compressed,dash pattern=on 10bp off 4bp] (380.48,405.00) .. controls (380.48,413.10) and (23.00,415.80) .. (23.00,415.80) .. controls (23.00,446.99) and (23.00,488.57) .. (23.00,519.76) .. controls (23.00,519.76) and (230.56,522.46) .. (230.56,530.56);
\draw[line width=3.5bp,panel] (582.48,218.00) .. controls (582.48,226.10) and (730.96,228.80) .. (730.96,228.80) .. controls (730.96,260.42) and (730.96,302.58) .. (730.96,334.20) .. controls (730.96,334.20) and (380.48,336.90) .. (380.48,345.00);
\draw[->,line width=1.0bp,compressed,dash pattern=on 10bp off 4bp] (582.48,218.00) .. controls (582.48,226.10) and (730.96,228.80) .. (730.96,228.80) .. controls (730.96,260.42) and (730.96,302.58) .. (730.96,334.20) .. controls (730.96,334.20) and (380.48,336.90) .. (380.48,345.00);
\draw[line width=3.5bp,panel] (279.48,653.00) .. controls (279.48,667.75) and (331.56,670.31) .. (331.56,685.06);
\draw[->,line width=1.0bp,compressed,dash pattern=on 10bp off 4bp] (279.48,653.00) .. controls (279.48,667.75) and (331.56,670.31) .. (331.56,685.06);
\draw[line width=3.5bp,panel] (178.48,492.00) .. controls (178.48,500.10) and (19.50,502.80) .. (19.50,502.80) .. controls (19.50,532.04) and (19.50,571.02) .. (19.50,600.26) .. controls (19.50,600.26) and (429.40,602.96) .. (429.40,611.06);
\draw[->,line width=1.0bp,compressed,dash pattern=on 10bp off 4bp] (178.48,492.00) .. controls (178.48,500.10) and (19.50,502.80) .. (19.50,502.80) .. controls (19.50,532.04) and (19.50,571.02) .. (19.50,600.26) .. controls (19.50,600.26) and (429.40,602.96) .. (429.40,611.06);
\draw[line width=3.5bp,panel] (380.48,492.00) .. controls (380.48,500.10) and (734.46,502.80) .. (734.46,502.80) .. controls (734.46,531.03) and (734.46,568.66) .. (734.46,596.89) .. controls (734.46,596.89) and (450.23,599.59) .. (450.23,607.69);
\draw[->,line width=1.0bp,compressed,dash pattern=on 10bp off 4bp] (380.48,492.00) .. controls (380.48,500.10) and (734.46,502.80) .. (734.46,502.80) .. controls (734.46,531.03) and (734.46,568.66) .. (734.46,596.89) .. controls (734.46,596.89) and (450.23,599.59) .. (450.23,607.69);
\draw[line width=3.5bp,panel] (582.48,492.00) .. controls (582.48,500.10) and (737.96,502.80) .. (737.96,502.80) .. controls (737.96,531.03) and (737.96,568.66) .. (737.96,596.89) .. controls (737.96,596.89) and (512.73,599.59) .. (512.73,607.69);
\draw[->,line width=1.0bp,compressed,dash pattern=on 10bp off 4bp] (582.48,492.00) .. controls (582.48,500.10) and (737.96,502.80) .. (737.96,502.80) .. controls (737.96,531.03) and (737.96,568.66) .. (737.96,596.89) .. controls (737.96,596.89) and (512.73,599.59) .. (512.73,607.69);
\draw[line width=3.5bp,panel] (582.48,579.00) .. controls (582.48,593.75) and (533.56,596.31) .. (533.56,611.06);
\draw[->,line width=1.0bp,compressed,dash pattern=on 10bp off 4bp] (582.48,579.00) .. controls (582.48,593.75) and (533.56,596.31) .. (533.56,611.06);
\draw[line width=3.5bp,panel] (380.48,572.50) .. controls (380.48,587.99) and (471.06,590.69) .. (471.06,606.18);
\draw[->,line width=1.0bp,compressed,dash pattern=on 10bp off 4bp] (380.48,572.50) .. controls (380.48,587.99) and (471.06,590.69) .. (471.06,606.18);
\draw[line width=3.5bp,panel] (481.48,653.00) .. controls (481.48,665.42) and (481.48,667.58) .. (481.48,680.00);
\draw[->,line width=1.0bp,compressed,dash pattern=on 10bp off 4bp] (481.48,653.00) .. controls (481.48,665.42) and (481.48,667.58) .. (481.48,680.00);
\draw[line width=3.5bp,panel] (380.48,318.00) .. controls (380.48,335.74) and (230.56,338.82) .. (230.56,356.56);
\draw[->,line width=1.0bp,compressed,dash pattern=on 10bp off 4bp] (380.48,318.00) .. controls (380.48,335.74) and (230.56,338.82) .. (230.56,356.56);
\filldraw[fill=setupbg,draw=setupfg,line width=1.0bp] (380.48,925.50) ellipse (84.00bp and 23.50bp);
\node[inner sep=0pt,font=\sffamily\fontsize{10.7bp}{12bp}\selectfont\bfseries,text=ink] at (380.48,932.00) {\hyperlink{node-W2T5_2}{WZ2 Theorem 5.2}};
\node[inner sep=0pt,font=\sffamily\fontsize{10.7bp}{12bp}\selectfont,text=ink] at (380.48,919.00) {\hyperlink{node-W2T5_2}{Multi-scale CW union}};
\filldraw[fill=setupbg,draw=setupfg,line width=1.0bp] (380.48,851.50) ellipse (84.00bp and 23.50bp);
\node[inner sep=0pt,font=\sffamily\fontsize{10.7bp}{12bp}\selectfont\bfseries,text=ink] at (380.48,858.00) {\hyperlink{node-W2Zero}{WZ2 Contradiction}};
\node[inner sep=0pt,font=\sffamily\fontsize{10.7bp}{12bp}\selectfont,text=ink] at (380.48,845.00) {\hyperlink{node-W2Zero}{sigma\_CW = 0}};
\filldraw[fill=slopebg,draw=slopefg,line width=1.0bp] (279.48,777.50) ellipse (84.00bp and 23.50bp);
\node[inner sep=0pt,font=\sffamily\fontsize{10.7bp}{12bp}\selectfont\bfseries,text=ink] at (279.48,784.00) {\hyperlink{node-W2L6_7}{WZ2 Lemma 6.7}};
\node[inner sep=0pt,font=\sffamily\fontsize{10.7bp}{12bp}\selectfont,text=ink] at (279.48,771.00) {\hyperlink{node-W2L6_7}{Small twisted image}};
\filldraw[fill=cinematicbg,draw=cinematicfg,line width=1.0bp] (481.48,777.50) ellipse (84.00bp and 23.50bp);
\node[inner sep=0pt,font=\sffamily\fontsize{10.7bp}{12bp}\selectfont\bfseries,text=ink] at (481.48,784.00) {\hyperlink{node-W2P6_8}{WZ2 Proposition 6.8}};
\node[inner sep=0pt,font=\sffamily\fontsize{10.7bp}{12bp}\selectfont,text=ink] at (481.48,771.00) {\hyperlink{node-W2P6_8}{Convex CW projection}};
\filldraw[fill=slopebg,draw=slopefg,line width=1.0bp] (279.48,703.50) ellipse (84.00bp and 23.50bp);
\node[inner sep=0pt,font=\sffamily\fontsize{10.7bp}{12bp}\selectfont\bfseries,text=ink] at (279.48,710.00) {\hyperlink{node-W2P6_5}{WZ2 Prop. 6.5 (1'')}};
\node[inner sep=0pt,font=\sffamily\fontsize{10.7bp}{12bp}\selectfont,text=ink] at (279.48,697.00) {\hyperlink{node-W2P6_5}{Nonsingular C2 slope}};
\filldraw[fill=cinematicbg,draw=cinematicfg,line width=1.0bp] (481.48,703.50) ellipse (84.00bp and 23.50bp);
\node[inner sep=0pt,font=\sffamily\fontsize{10.7bp}{12bp}\selectfont\bfseries,text=ink] at (481.48,710.00) {\hyperlink{node-W2P7_1}{WZ2 Proposition 7.1}};
\node[inner sep=0pt,font=\sffamily\fontsize{10.7bp}{12bp}\selectfont,text=ink] at (481.48,697.00) {\hyperlink{node-W2P7_1}{Tube-Wolff projection}};
\filldraw[fill=regularitybg,draw=regularityfg,line width=1.0bp] (178.48,549.00) ellipse (84.00bp and 23.50bp);
\node[inner sep=0pt,font=\sffamily\fontsize{10.7bp}{12bp}\selectfont\bfseries,text=ink] at (178.48,555.50) {\hyperlink{node-W2P6_4}{WZ2 Proposition 6.4}};
\node[inner sep=0pt,font=\sffamily\fontsize{10.7bp}{12bp}\selectfont,text=ink] at (178.48,542.50) {\hyperlink{node-W2P6_4}{C2 grains}};
\filldraw[fill=grainsbg,draw=grainsfg,line width=1.0bp] (178.48,375.00) ellipse (84.00bp and 23.50bp);
\node[inner sep=0pt,font=\sffamily\fontsize{10.7bp}{12bp}\selectfont\bfseries,text=ink] at (178.48,381.50) {\hyperlink{node-W2P6_3}{WZ2 Proposition 6.3}};
\node[inner sep=0pt,font=\sffamily\fontsize{10.7bp}{12bp}\selectfont,text=ink] at (178.48,368.50) {\hyperlink{node-W2P6_3}{Lipschitz grains}};
\filldraw[fill=setupbg,draw=setupfg,line width=1.0bp] (151.50,281.50) ellipse (84.00bp and 23.50bp);
\node[inner sep=0pt,font=\sffamily\fontsize{10.7bp}{12bp}\selectfont\bfseries,text=ink] at (151.50,288.00) {\hyperlink{node-W2P6_2}{WZ2 Proposition 6.2}};
\node[inner sep=0pt,font=\sffamily\fontsize{10.7bp}{12bp}\selectfont,text=ink] at (151.50,275.00) {\hyperlink{node-W2P6_2}{Coarse/fine extremality}};
\filldraw[fill=setupbg,draw=setupfg,line width=1.0bp] (380.48,281.50) ellipse (110.98bp and 36.50bp);
\node[inner sep=0pt,font=\sffamily\fontsize{10.7bp}{12bp}\selectfont\bfseries,text=ink] at (380.48,301.00) {\hyperlink{node-W2Exists}{WZ2 Sec. 6.2 / WZ1 Lem.}};
\node[inner sep=0pt,font=\sffamily\fontsize{10.7bp}{12bp}\selectfont,text=ink] at (380.48,288.00) {\hyperlink{node-W2Exists}{3.1}};
\node[inner sep=0pt,font=\sffamily\fontsize{10.7bp}{12bp}\selectfont,text=ink] at (380.48,275.00) {\hyperlink{node-W2Exists}{Extremizers exist at}};
\node[inner sep=0pt,font=\sffamily\fontsize{10.7bp}{12bp}\selectfont,text=ink] at (380.48,262.00) {\hyperlink{node-W2Exists}{arbitrarily small scales}};
\filldraw[fill=setupbg,draw=setupfg,line width=1.0bp] (178.07,190.00) ellipse (84.00bp and 23.50bp);
\node[inner sep=0pt,font=\sffamily\fontsize{10.7bp}{12bp}\selectfont\bfseries,text=ink] at (178.07,196.50) {\hyperlink{node-W2L2_13}{WZ2 Lemma 2.13}};
\node[inner sep=0pt,font=\sffamily\fontsize{10.7bp}{12bp}\selectfont,text=ink] at (178.07,183.50) {\hyperlink{node-W2L2_13}{Multi-scale inheritance}};
\filldraw[fill=setupbg,draw=setupfg,line width=1.0bp,rounded corners=5bp] (296.48,22.00) rectangle (464.48,65.00);
\node[inner sep=0pt,font=\sffamily\fontsize{10.7bp}{12bp}\selectfont\bfseries,text=ink] at (380.48,50.00) {\hyperlink{node-W2MW}{WZ2 Definition 2.12}};
\node[inner sep=0pt,font=\sffamily\fontsize{10.7bp}{12bp}\selectfont,text=ink] at (380.48,37.00) {\hyperlink{node-W2MW}{Multi-scale convex CW}};
\filldraw[fill=setupbg,draw=setupfg,line width=1.0bp,rounded corners=5bp] (296.07,168.50) rectangle (464.07,211.50);
\node[inner sep=0pt,font=\sffamily\fontsize{10.7bp}{12bp}\selectfont\bfseries,text=ink] at (380.07,196.50) {\hyperlink{node-W2D6_1}{WZ2 Definition 6.1}};
\node[inner sep=0pt,font=\sffamily\fontsize{10.7bp}{12bp}\selectfont,text=ink] at (380.07,183.50) {\hyperlink{node-W2D6_1}{CW-extremal pair}};
\filldraw[fill=setupbg,draw=setupfg,line width=1.0bp,rounded corners=5bp] (296.48,92.00) rectangle (464.48,135.00);
\node[inner sep=0pt,font=\sffamily\fontsize{10.7bp}{12bp}\selectfont\bfseries,text=ink] at (380.48,120.00) {\hyperlink{node-W2Sigma}{WZ2 Section 6.2}};
\node[inner sep=0pt,font=\sffamily\fontsize{10.7bp}{12bp}\selectfont,text=ink] at (380.48,107.00) {\hyperlink{node-W2Sigma}{Counterexample deficit}};
\filldraw[fill=slopebg,draw=slopefg,line width=1.0bp] (279.48,629.50) ellipse (84.00bp and 23.50bp);
\node[inner sep=0pt,font=\sffamily\fontsize{10.7bp}{12bp}\selectfont\bfseries,text=ink] at (279.48,636.00) {\hyperlink{node-W2L6_6}{WZ2 Lemma 6.6}};
\node[inner sep=0pt,font=\sffamily\fontsize{10.7bp}{12bp}\selectfont,text=ink] at (279.48,623.00) {\hyperlink{node-W2L6_6}{Popular slab is nonflat}};
\filldraw[fill=regularitybg,draw=regularityfg,line width=1.0bp] (380.48,375.00) ellipse (84.00bp and 30.00bp);
\node[inner sep=0pt,font=\sffamily\fontsize{10.7bp}{12bp}\selectfont\bfseries,text=ink] at (380.48,388.00) {\hyperlink{node-T5_2}{WZ1 Thm. 5.2}};
\node[inner sep=0pt,font=\sffamily\fontsize{10.7bp}{12bp}\selectfont,text=ink] at (380.48,375.00) {\hyperlink{node-T5_2}{Dot-product expansion or}};
\node[inner sep=0pt,font=\sffamily\fontsize{10.7bp}{12bp}\selectfont,text=ink] at (380.48,362.00) {\hyperlink{node-T5_2}{orthogonal lines}};
\filldraw[fill=cinematicbg,draw=cinematicfg,line width=1.0bp] (481.48,629.50) ellipse (84.00bp and 23.50bp);
\node[inner sep=0pt,font=\sffamily\fontsize{10.7bp}{12bp}\selectfont\bfseries,text=ink] at (481.48,636.00) {\hyperlink{node-W2L7_12}{WZ2 Lemma 7.12}};
\node[inner sep=0pt,font=\sffamily\fontsize{10.7bp}{12bp}\selectfont,text=ink] at (481.48,623.00) {\hyperlink{node-W2L7_12}{One-step area ratio}};
\filldraw[fill=cinematicbg,draw=cinematicfg,line width=1.0bp] (178.48,462.00) ellipse (84.00bp and 30.00bp);
\node[inner sep=0pt,font=\sffamily\fontsize{10.7bp}{12bp}\selectfont\bfseries,text=ink] at (178.48,475.00) {\hyperlink{node-W2L7_10}{WZ2 Lemma 7.10}};
\node[inner sep=0pt,font=\sffamily\fontsize{10.7bp}{12bp}\selectfont,text=ink] at (178.48,462.00) {\hyperlink{node-W2L7_10}{Zoom to Frostman fibres}};
\node[inner sep=0pt,font=\sffamily\fontsize{10.7bp}{12bp}\selectfont,text=ink] at (178.48,449.00) {\hyperlink{node-W2L7_10}{of a uniform set}};
\filldraw[fill=cinematicbg,draw=cinematicfg,line width=1.0bp] (380.48,462.00) ellipse (84.00bp and 30.00bp);
\node[inner sep=0pt,font=\sffamily\fontsize{10.7bp}{12bp}\selectfont\bfseries,text=ink] at (380.48,475.00) {\hyperlink{node-W2L7_6}{WZ2 Lemma 7.6}};
\node[inner sep=0pt,font=\sffamily\fontsize{10.7bp}{12bp}\selectfont,text=ink] at (380.48,462.00) {\hyperlink{node-W2L7_6}{Frostman sets contain}};
\node[inner sep=0pt,font=\sffamily\fontsize{10.7bp}{12bp}\selectfont,text=ink] at (380.48,449.00) {\hyperlink{node-W2L7_6}{large Katz-Tao subsets}};
\filldraw[fill=cinematicbg,draw=cinematicfg,line width=1.0bp] (582.48,462.00) ellipse (84.00bp and 30.00bp);
\node[inner sep=0pt,font=\sffamily\fontsize{10.7bp}{12bp}\selectfont\bfseries,text=ink] at (582.48,475.00) {\hyperlink{node-W2L7_11}{WZ2 Lemma 7.11}};
\node[inner sep=0pt,font=\sffamily\fontsize{10.7bp}{12bp}\selectfont,text=ink] at (582.48,462.00) {\hyperlink{node-W2L7_11}{Replicate a local}};
\node[inner sep=0pt,font=\sffamily\fontsize{10.7bp}{12bp}\selectfont,text=ink] at (582.48,449.00) {\hyperlink{node-W2L7_11}{Katz-Tao cluster}};
\filldraw[fill=cinematicbg,draw=cinematicfg,line width=1.0bp] (380.48,549.00) ellipse (84.00bp and 23.50bp);
\node[inner sep=0pt,font=\sffamily\fontsize{10.7bp}{12bp}\selectfont\bfseries,text=ink] at (380.48,555.50) {\hyperlink{node-T7_2}{WZ1 7.2 = WZ2 7.13}};
\node[inner sep=0pt,font=\sffamily\fontsize{10.7bp}{12bp}\selectfont,text=ink] at (380.48,542.50) {\hyperlink{node-T7_2}{Cinematic maximal}};
\filldraw[fill=cinematicbg,draw=cinematicfg,line width=1.0bp] (582.48,549.00) ellipse (84.00bp and 30.00bp);
\node[inner sep=0pt,font=\sffamily\fontsize{10.7bp}{12bp}\selectfont\bfseries,text=ink] at (582.48,562.00) {\hyperlink{node-L7_3}{WZ1 Lemma 7.3}};
\node[inner sep=0pt,font=\sffamily\fontsize{10.7bp}{12bp}\selectfont,text=ink] at (582.48,549.00) {\hyperlink{node-L7_3}{The twisted-line family}};
\node[inner sep=0pt,font=\sffamily\fontsize{10.7bp}{12bp}\selectfont,text=ink] at (582.48,536.00) {\hyperlink{node-L7_3}{is cinematic}};
\filldraw[fill=externalbg,draw=externalfg,line width=1.0bp,rounded corners=5bp,dashed] (525.46,247.00) rectangle (693.46,316.00);
\node[inner sep=0pt,font=\sffamily\fontsize{10.7bp}{12bp}\selectfont\bfseries,text=ink] at (609.46,301.00) {\hyperlink{node-XPYZ}{External:}};
\node[inner sep=0pt,font=\sffamily\fontsize{10.7bp}{12bp}\selectfont,text=ink] at (609.46,288.00) {\hyperlink{node-XPYZ}{Pramanik-Yang-Zahl}};
\node[inner sep=0pt,font=\sffamily\fontsize{10.7bp}{12bp}\selectfont,text=ink] at (609.46,275.00) {\hyperlink{node-XPYZ}{Cinematic L3/2 maximal}};
\node[inner sep=0pt,font=\sffamily\fontsize{10.7bp}{12bp}\selectfont,text=ink] at (609.46,262.00) {\hyperlink{node-XPYZ}{theorem}};
\filldraw[fill=externalbg,draw=externalfg,line width=1.0bp,rounded corners=5bp,dashed] (498.07,162.00) rectangle (666.89,218.00);
\node[inner sep=0pt,font=\sffamily\fontsize{10.7bp}{12bp}\selectfont\bfseries,text=ink] at (582.48,203.00) {\hyperlink{node-XOSW}{External [35, Lemma 2.9]}};
\node[inner sep=0pt,font=\sffamily\fontsize{10.7bp}{12bp}\selectfont,text=ink] at (582.48,190.00) {\hyperlink{node-XOSW}{Radial thin-tube}};
\node[inner sep=0pt,font=\sffamily\fontsize{10.7bp}{12bp}\selectfont,text=ink] at (582.48,177.00) {\hyperlink{node-XOSW}{bootstrap}};
\filldraw[fill=cinematicbg,draw=cinematicfg,line width=1.0bp,rounded corners=5bp] (498.48,347.00) rectangle (666.48,403.00);
\node[inner sep=0pt,font=\sffamily\fontsize{10.7bp}{12bp}\selectfont\bfseries,text=ink] at (582.48,388.00) {\hyperlink{node-W2TW}{WZ2 Sec. 7: Tube Wolff}};
\node[inner sep=0pt,font=\sffamily\fontsize{10.7bp}{12bp}\selectfont,text=ink] at (582.48,375.00) {\hyperlink{node-W2TW}{The weaker radius-squared}};
\node[inner sep=0pt,font=\sffamily\fontsize{10.7bp}{12bp}\selectfont,text=ink] at (582.48,362.00) {\hyperlink{node-W2TW}{tube count}};
\end{tikzpicture}}\end{center}
{\footnotesize Solid: proof input. Dashed: API. Dash-dot: assembly. Long-dashed: overview path. Read from the top down; arrows point to consumers. Full prerequisites appear on the cards and in the ledger.}\par
\sectitle{Wolff infrastructure and self-similarity}
\hypertarget{view-scales}{}
{\sffamily\small\color{muted}WZ2 definitions replace sticky admissibility; multiplicities keep the total cardinality visible.}\par
\begin{center}\resizebox{493.00bp}{!}{\begin{tikzpicture}[x=1bp,y=1bp, >={Stealth[length=4bp,width=3.2bp]}]
\path[use as bounding box] (0,0) rectangle (792.33,754.00);
\fill[panel,rounded corners=5bp] (0,0) rectangle (792.33,754.00);
\draw[line width=3.5bp,panel] (368.42,78.00) .. controls (368.42,90.42) and (298.67,92.58) .. (298.67,105.00);
\draw[->,line width=1.0bp,api,dash pattern=on 5bp off 4bp] (368.42,78.00) .. controls (368.42,90.42) and (298.67,92.58) .. (298.67,105.00);
\draw[line width=3.5bp,panel] (389.25,78.00) .. controls (389.25,86.10) and (61.50,88.80) .. (61.50,88.80) .. controls (61.50,161.22) and (61.50,257.78) .. (61.50,330.20) .. controls (61.50,330.20) and (382.31,332.90) .. (382.31,341.00);
\draw[->,line width=1.0bp,api,dash pattern=on 5bp off 4bp] (389.25,78.00) .. controls (389.25,86.10) and (61.50,88.80) .. (61.50,88.80) .. controls (61.50,161.22) and (61.50,257.78) .. (61.50,330.20) .. controls (61.50,330.20) and (382.31,332.90) .. (382.31,341.00);
\draw[line width=3.5bp,panel] (272.63,161.00) .. controls (272.63,169.10) and (58.00,171.80) .. (58.00,171.80) .. controls (58.00,219.32) and (58.00,282.68) .. (58.00,330.20) .. controls (58.00,330.20) and (347.59,332.90) .. (347.59,341.00);
\draw[->,line width=1.0bp,api,dash pattern=on 5bp off 4bp] (272.63,161.00) .. controls (272.63,169.10) and (58.00,171.80) .. (58.00,171.80) .. controls (58.00,219.32) and (58.00,282.68) .. (58.00,330.20) .. controls (58.00,330.20) and (347.59,332.90) .. (347.59,341.00);
\draw[line width=3.5bp,panel] (347.59,384.00) .. controls (347.59,392.10) and (54.50,394.80) .. (54.50,394.80) .. controls (54.50,425.07) and (54.50,465.43) .. (54.50,495.70) .. controls (54.50,495.70) and (140.40,498.40) .. (140.40,506.50);
\draw[->,line width=1.0bp,api,dash pattern=on 5bp off 4bp] (347.59,384.00) .. controls (347.59,392.10) and (54.50,394.80) .. (54.50,394.80) .. controls (54.50,425.07) and (54.50,465.43) .. (54.50,495.70) .. controls (54.50,495.70) and (140.40,498.40) .. (140.40,506.50);
\draw[line width=3.5bp,panel] (448.59,161.00) .. controls (448.59,173.42) and (350.75,175.58) .. (350.75,188.00);
\draw[->,line width=1.0bp,api,dash pattern=on 5bp off 4bp] (448.59,161.00) .. controls (448.59,173.42) and (350.75,175.58) .. (350.75,188.00);
\draw[line width=3.5bp,panel] (246.59,161.00) .. controls (246.59,173.42) and (246.59,175.58) .. (246.59,188.00);
\draw[->,line width=1.0bp,api,dash pattern=on 5bp off 4bp] (246.59,161.00) .. controls (246.59,173.42) and (246.59,175.58) .. (246.59,188.00);
\draw[line width=3.5bp,panel] (389.25,384.00) .. controls (389.25,392.10) and (51.00,394.80) .. (51.00,394.80) .. controls (51.00,473.18) and (51.00,577.68) .. (51.00,656.06) .. controls (51.00,656.06) and (261.28,658.76) .. (261.28,666.86);
\draw[->,line width=1.0bp,core] (389.25,384.00) .. controls (389.25,392.10) and (51.00,394.80) .. (51.00,394.80) .. controls (51.00,473.18) and (51.00,577.68) .. (51.00,656.06) .. controls (51.00,656.06) and (261.28,658.76) .. (261.28,666.86);
\draw[line width=3.5bp,panel] (140.40,549.50) .. controls (140.40,557.60) and (47.50,560.30) .. (47.50,560.30) .. controls (47.50,589.03) and (47.50,627.33) .. (47.50,656.06) .. controls (47.50,656.06) and (123.67,658.76) .. (123.67,666.86);
\draw[->,line width=1.0bp,core] (140.40,549.50) .. controls (140.40,557.60) and (47.50,560.30) .. (47.50,560.30) .. controls (47.50,589.03) and (47.50,627.33) .. (47.50,656.06) .. controls (47.50,656.06) and (123.67,658.76) .. (123.67,666.86);
\draw[line width=3.5bp,panel] (430.91,384.00) .. controls (430.91,392.10) and (737.83,394.80) .. (737.83,394.80) .. controls (737.83,422.71) and (737.83,459.93) .. (737.83,487.85) .. controls (737.83,487.85) and (418.10,490.55) .. (418.10,498.65);
\draw[->,line width=1.0bp,core] (430.91,384.00) .. controls (430.91,392.10) and (737.83,394.80) .. (737.83,394.80) .. controls (737.83,422.71) and (737.83,459.93) .. (737.83,487.85) .. controls (737.83,487.85) and (418.10,490.55) .. (418.10,498.65);
\draw[line width=3.5bp,panel] (246.59,244.00) .. controls (246.59,252.10) and (44.00,254.80) .. (44.00,254.80) .. controls (44.00,324.71) and (44.00,417.93) .. (44.00,487.85) .. controls (44.00,487.85) and (381.23,490.55) .. (381.23,498.65);
\draw[->,line width=1.0bp,core] (246.59,244.00) .. controls (246.59,252.10) and (44.00,254.80) .. (44.00,254.80) .. controls (44.00,324.71) and (44.00,417.93) .. (44.00,487.85) .. controls (44.00,487.85) and (381.23,490.55) .. (381.23,498.65);
\draw[line width=3.5bp,panel] (399.67,558.00) .. controls (399.67,570.46) and (392.23,572.63) .. (392.23,585.09);
\draw[->,line width=1.0bp,core] (399.67,558.00) .. controls (399.67,570.46) and (392.23,572.63) .. (392.23,585.09);
\draw[line width=3.5bp,panel] (451.75,384.00) .. controls (451.75,392.10) and (741.33,394.80) .. (741.33,394.80) .. controls (741.33,448.65) and (741.33,520.44) .. (741.33,574.29) .. controls (741.33,574.29) and (407.11,576.99) .. (407.11,585.09);
\draw[->,line width=1.0bp,core] (451.75,384.00) .. controls (451.75,392.10) and (741.33,394.80) .. (741.33,394.80) .. controls (741.33,448.65) and (741.33,520.44) .. (741.33,574.29) .. controls (741.33,574.29) and (407.11,576.99) .. (407.11,585.09);
\draw[line width=3.5bp,panel] (192.48,549.50) .. controls (192.48,568.16) and (347.59,571.40) .. (347.59,590.06);
\draw[->,line width=1.0bp,core] (192.48,549.50) .. controls (192.48,568.16) and (347.59,571.40) .. (347.59,590.06);
\draw[line width=3.5bp,panel] (399.67,632.00) .. controls (399.67,650.40) and (421.46,653.60) .. (421.46,672.00);
\draw[->,line width=1.0bp,core] (399.67,632.00) .. controls (399.67,650.40) and (421.46,653.60) .. (421.46,672.00);
\draw[line width=3.5bp,panel] (244.56,549.50) .. controls (244.56,557.60) and (744.83,560.30) .. (744.83,560.30) .. controls (744.83,590.56) and (744.83,630.90) .. (744.83,661.16) .. controls (744.83,661.16) and (573.35,663.86) .. (573.35,671.96);
\draw[->,line width=1.0bp,core] (244.56,549.50) .. controls (244.56,557.60) and (744.83,560.30) .. (744.83,560.30) .. controls (744.83,590.56) and (744.83,630.90) .. (744.83,661.16) .. controls (744.83,661.16) and (573.35,663.86) .. (573.35,671.96);
\draw[line width=3.5bp,panel] (410.08,78.00) .. controls (410.08,86.10) and (748.33,88.80) .. (748.33,88.80) .. controls (748.33,210.26) and (748.33,372.20) .. (748.33,493.66) .. controls (748.33,493.66) and (454.96,496.36) .. (454.96,504.46);
\draw[->,line width=1.0bp,api,dash pattern=on 5bp off 4bp] (410.08,78.00) .. controls (410.08,86.10) and (748.33,88.80) .. (748.33,88.80) .. controls (748.33,210.26) and (748.33,372.20) .. (748.33,493.66) .. controls (748.33,493.66) and (454.96,496.36) .. (454.96,504.46);
\draw[line width=3.5bp,panel] (298.67,161.00) .. controls (298.67,169.10) and (40.50,171.80) .. (40.50,171.80) .. controls (40.50,292.77) and (40.50,454.07) .. (40.50,575.04) .. controls (40.50,575.04) and (377.35,577.74) .. (377.35,585.84);
\draw[->,line width=1.0bp,api,dash pattern=on 5bp off 4bp] (298.67,161.00) .. controls (298.67,169.10) and (40.50,171.80) .. (40.50,171.80) .. controls (40.50,292.77) and (40.50,454.07) .. (40.50,575.04) .. controls (40.50,575.04) and (377.35,577.74) .. (377.35,585.84);
\draw[line width=3.5bp,panel] (430.91,78.00) .. controls (430.91,86.10) and (751.83,88.80) .. (751.83,88.80) .. controls (751.83,234.67) and (751.83,429.17) .. (751.83,575.04) .. controls (751.83,575.04) and (421.99,577.74) .. (421.99,585.84);
\draw[->,line width=1.0bp,api,dash pattern=on 5bp off 4bp] (430.91,78.00) .. controls (430.91,86.10) and (751.83,88.80) .. (751.83,88.80) .. controls (751.83,234.67) and (751.83,429.17) .. (751.83,575.04) .. controls (751.83,575.04) and (421.99,577.74) .. (421.99,585.84);
\draw[line width=3.5bp,panel] (552.75,161.00) .. controls (552.75,169.10) and (755.33,171.80) .. (755.33,171.80) .. controls (755.33,316.67) and (755.33,509.83) .. (755.33,654.70) .. controls (755.33,654.70) and (628.64,657.40) .. (628.64,665.50);
\draw[->,line width=1.0bp,api,dash pattern=on 5bp off 4bp] (552.75,161.00) .. controls (552.75,169.10) and (755.33,171.80) .. (755.33,171.80) .. controls (755.33,316.67) and (755.33,509.83) .. (755.33,654.70) .. controls (755.33,654.70) and (628.64,657.40) .. (628.64,665.50);
\draw[line width=3.5bp,panel] (324.71,161.00) .. controls (324.71,169.10) and (37.00,171.80) .. (37.00,171.80) .. controls (37.00,198.32) and (37.00,233.68) .. (37.00,260.20) .. controls (37.00,260.20) and (347.59,262.90) .. (347.59,271.00);
\draw[->,line width=1.0bp,core] (324.71,161.00) .. controls (324.71,169.10) and (37.00,171.80) .. (37.00,171.80) .. controls (37.00,198.32) and (37.00,233.68) .. (37.00,260.20) .. controls (37.00,260.20) and (347.59,262.90) .. (347.59,271.00);
\draw[line width=3.5bp,panel] (350.75,244.00) .. controls (350.75,256.42) and (399.67,258.58) .. (399.67,271.00);
\draw[->,line width=1.0bp,core] (350.75,244.00) .. controls (350.75,256.42) and (399.67,258.58) .. (399.67,271.00);
\draw[line width=3.5bp,panel] (448.59,244.00) .. controls (448.59,256.42) and (451.75,258.58) .. (451.75,271.00);
\draw[->,line width=1.0bp,core] (448.59,244.00) .. controls (448.59,256.42) and (451.75,258.58) .. (451.75,271.00);
\draw[line width=3.5bp,panel] (451.75,314.00) .. controls (451.75,322.10) and (758.83,324.80) .. (758.83,324.80) .. controls (758.83,350.89) and (758.83,385.67) .. (758.83,411.76) .. controls (758.83,411.76) and (682.12,414.46) .. (682.12,422.56);
\draw[->,line width=1.0bp,core] (451.75,314.00) .. controls (451.75,322.10) and (758.83,324.80) .. (758.83,324.80) .. controls (758.83,350.89) and (758.83,385.67) .. (758.83,411.76) .. controls (758.83,411.76) and (682.12,414.46) .. (682.12,422.56);
\draw[line width=3.5bp,panel] (350.75,161.00) .. controls (350.75,169.10) and (762.33,171.80) .. (762.33,171.80) .. controls (762.33,243.79) and (762.33,339.77) .. (762.33,411.76) .. controls (762.33,411.76) and (577.96,414.46) .. (577.96,422.56);
\draw[->,line width=1.0bp,api,dash pattern=on 5bp off 4bp] (350.75,161.00) .. controls (350.75,169.10) and (762.33,171.80) .. (762.33,171.80) .. controls (762.33,243.79) and (762.33,339.77) .. (762.33,411.76) .. controls (762.33,411.76) and (577.96,414.46) .. (577.96,422.56);
\draw[line width=3.5bp,panel] (399.67,314.00) .. controls (399.67,326.42) and (417.03,328.58) .. (417.03,341.00);
\draw[->,line width=1.0bp,core] (399.67,314.00) .. controls (399.67,326.42) and (417.03,328.58) .. (417.03,341.00);
\draw[line width=3.5bp,panel] (451.75,78.00) .. controls (451.75,86.10) and (33.50,88.80) .. (33.50,88.80) .. controls (33.50,161.22) and (33.50,257.78) .. (33.50,330.20) .. controls (33.50,330.20) and (451.75,332.90) .. (451.75,341.00);
\draw[->,line width=1.0bp,core] (451.75,78.00) .. controls (451.75,86.10) and (33.50,88.80) .. (33.50,88.80) .. controls (33.50,161.22) and (33.50,257.78) .. (33.50,330.20) .. controls (33.50,330.20) and (451.75,332.90) .. (451.75,341.00);
\draw[line width=3.5bp,panel] (410.08,384.00) .. controls (410.08,396.42) and (197.67,398.58) .. (197.67,411.00);
\draw[->,line width=1.0bp,core] (410.08,384.00) .. controls (410.08,396.42) and (197.67,398.58) .. (197.67,411.00);
\draw[line width=3.5bp,panel] (347.59,314.00) .. controls (347.59,322.10) and (30.00,324.80) .. (30.00,324.80) .. controls (30.00,376.07) and (30.00,444.43) .. (30.00,495.70) .. controls (30.00,495.70) and (175.12,498.40) .. (175.12,506.50);
\draw[->,line width=1.0bp,core] (347.59,314.00) .. controls (347.59,322.10) and (30.00,324.80) .. (30.00,324.80) .. controls (30.00,376.07) and (30.00,444.43) .. (30.00,495.70) .. controls (30.00,495.70) and (175.12,498.40) .. (175.12,506.50);
\draw[line width=3.5bp,panel] (347.59,78.00) .. controls (347.59,86.10) and (26.50,88.80) .. (26.50,88.80) .. controls (26.50,210.87) and (26.50,373.63) .. (26.50,495.70) .. controls (26.50,495.70) and (209.84,498.40) .. (209.84,506.50);
\draw[->,line width=1.0bp,core] (347.59,78.00) .. controls (347.59,86.10) and (26.50,88.80) .. (26.50,88.80) .. controls (26.50,210.87) and (26.50,373.63) .. (26.50,495.70) .. controls (26.50,495.70) and (209.84,498.40) .. (209.84,506.50);
\draw[line width=3.5bp,panel] (368.42,384.00) .. controls (368.42,392.10) and (23.00,394.80) .. (23.00,394.80) .. controls (23.00,425.07) and (23.00,465.43) .. (23.00,495.70) .. controls (23.00,495.70) and (244.56,498.40) .. (244.56,506.50);
\draw[->,line width=1.0bp,core] (368.42,384.00) .. controls (368.42,392.10) and (23.00,394.80) .. (23.00,394.80) .. controls (23.00,425.07) and (23.00,465.43) .. (23.00,495.70) .. controls (23.00,495.70) and (244.56,498.40) .. (244.56,506.50);
\draw[line width=3.5bp,panel] (577.96,459.44) .. controls (577.96,467.54) and (765.83,470.24) .. (765.83,470.24) .. controls (765.83,502.95) and (765.83,546.55) .. (765.83,579.26) .. controls (765.83,579.26) and (451.75,581.96) .. (451.75,590.06);
\draw[->,line width=1.0bp,core] (577.96,459.44) .. controls (577.96,467.54) and (765.83,470.24) .. (765.83,470.24) .. controls (765.83,502.95) and (765.83,546.55) .. (765.83,579.26) .. controls (765.83,579.26) and (451.75,581.96) .. (451.75,590.06);
\draw[line width=3.5bp,panel] (127.99,464.54) .. controls (127.99,482.90) and (344.37,486.10) .. (344.37,504.46);
\draw[->,line width=1.0bp,core] (127.99,464.54) .. controls (127.99,482.90) and (344.37,486.10) .. (344.37,504.46);
\draw[line width=3.5bp,panel] (267.34,464.54) .. controls (267.34,472.64) and (19.50,475.34) .. (19.50,475.34) .. controls (19.50,505.73) and (19.50,546.24) .. (19.50,576.63) .. controls (19.50,576.63) and (362.47,579.33) .. (362.47,587.43);
\draw[->,line width=1.0bp,core] (267.34,464.54) .. controls (267.34,472.64) and (19.50,475.34) .. (19.50,475.34) .. controls (19.50,505.73) and (19.50,546.24) .. (19.50,576.63) .. controls (19.50,576.63) and (362.47,579.33) .. (362.47,587.43);
\draw[line width=3.5bp,panel] (554.77,556.00) .. controls (554.77,570.46) and (436.87,572.97) .. (436.87,587.43);
\draw[->,line width=1.0bp,api,dash pattern=on 5bp off 4bp] (554.77,556.00) .. controls (554.77,570.46) and (436.87,572.97) .. (436.87,587.43);
\draw[line width=3.5bp,panel] (682.12,459.44) .. controls (682.12,478.10) and (658.93,481.34) .. (658.93,500.00);
\draw[->,line width=1.0bp,api,dash pattern=on 5bp off 4bp] (682.12,459.44) .. controls (682.12,478.10) and (658.93,481.34) .. (658.93,500.00);
\draw[line width=3.5bp,panel] (552.75,244.00) .. controls (552.75,252.10) and (769.33,254.80) .. (769.33,254.80) .. controls (769.33,325.12) and (769.33,418.88) .. (769.33,489.20) .. controls (769.33,489.20) and (554.77,491.90) .. (554.77,500.00);
\draw[->,line width=1.0bp,api,dash pattern=on 5bp off 4bp] (552.75,244.00) .. controls (552.75,252.10) and (769.33,254.80) .. (769.33,254.80) .. controls (769.33,325.12) and (769.33,418.88) .. (769.33,489.20) .. controls (769.33,489.20) and (554.77,491.90) .. (554.77,500.00);
\draw[line width=3.5bp,panel] (658.93,556.00) .. controls (658.93,564.10) and (772.83,566.80) .. (772.83,566.80) .. controls (772.83,595.11) and (772.83,632.85) .. (772.83,661.16) .. controls (772.83,661.16) and (683.94,663.86) .. (683.94,671.96);
\draw[->,line width=1.0bp,api,dash pattern=on 5bp off 4bp] (658.93,556.00) .. controls (658.93,564.10) and (772.83,566.80) .. (772.83,566.80) .. controls (772.83,595.11) and (772.83,632.85) .. (772.83,661.16) .. controls (772.83,661.16) and (683.94,663.86) .. (683.94,671.96);
\filldraw[fill=setupbg,draw=setupfg,line width=1.0bp] (399.67,608.50) ellipse (84.00bp and 23.50bp);
\node[inner sep=0pt,font=\sffamily\fontsize{10.7bp}{12bp}\selectfont\bfseries,text=ink] at (399.67,615.00) {\hyperlink{node-W2P6_2}{WZ2 Proposition 6.2}};
\node[inner sep=0pt,font=\sffamily\fontsize{10.7bp}{12bp}\selectfont,text=ink] at (399.67,602.00) {\hyperlink{node-W2P6_2}{Coarse/fine extremality}};
\filldraw[fill=setupbg,draw=setupfg,line width=1.0bp] (399.67,528.00) ellipse (89.19bp and 30.00bp);
\node[inner sep=0pt,font=\sffamily\fontsize{10.7bp}{12bp}\selectfont\bfseries,text=ink] at (399.67,541.00) {\hyperlink{node-L3_3}{WZ1 Lemma 3.3 / CW}};
\node[inner sep=0pt,font=\sffamily\fontsize{10.7bp}{12bp}\selectfont,text=ink] at (399.67,528.00) {\hyperlink{node-L3_3}{Local rescaling and}};
\node[inner sep=0pt,font=\sffamily\fontsize{10.7bp}{12bp}\selectfont,text=ink] at (399.67,515.00) {\hyperlink{node-L3_3}{multiplicity control}};
\filldraw[fill=setupbg,draw=setupfg,line width=1.0bp] (421.46,695.50) ellipse (84.00bp and 23.50bp);
\node[inner sep=0pt,font=\sffamily\fontsize{10.7bp}{12bp}\selectfont\bfseries,text=ink] at (421.46,702.00) {\hyperlink{node-L3_4}{WZ1 Lemma 3.4 / CW}};
\node[inner sep=0pt,font=\sffamily\fontsize{10.7bp}{12bp}\selectfont,text=ink] at (421.46,689.00) {\hyperlink{node-L3_4}{Robust transversality}};
\filldraw[fill=setupbg,draw=setupfg,line width=1.0bp] (628.64,695.50) ellipse (89.19bp and 30.00bp);
\node[inner sep=0pt,font=\sffamily\fontsize{10.7bp}{12bp}\selectfont\bfseries,text=ink] at (628.64,708.50) {\hyperlink{node-L3_5}{WZ1 Lemma 3.5 / CW}};
\node[inner sep=0pt,font=\sffamily\fontsize{10.7bp}{12bp}\selectfont,text=ink] at (628.64,695.50) {\hyperlink{node-L3_5}{Mild anisotropic}};
\node[inner sep=0pt,font=\sffamily\fontsize{10.7bp}{12bp}\selectfont,text=ink] at (628.64,682.50) {\hyperlink{node-L3_5}{rescaling preserves}};
\filldraw[fill=setupbg,draw=setupfg,line width=1.0bp] (192.48,695.50) ellipse (110.98bp and 36.50bp);
\node[inner sep=0pt,font=\sffamily\fontsize{10.7bp}{12bp}\selectfont\bfseries,text=ink] at (192.48,715.00) {\hyperlink{node-W2Exists}{WZ2 Sec. 6.2 / WZ1 Lem.}};
\node[inner sep=0pt,font=\sffamily\fontsize{10.7bp}{12bp}\selectfont,text=ink] at (192.48,702.00) {\hyperlink{node-W2Exists}{3.1}};
\node[inner sep=0pt,font=\sffamily\fontsize{10.7bp}{12bp}\selectfont,text=ink] at (192.48,689.00) {\hyperlink{node-W2Exists}{Extremizers exist at}};
\node[inner sep=0pt,font=\sffamily\fontsize{10.7bp}{12bp}\selectfont,text=ink] at (192.48,676.00) {\hyperlink{node-W2Exists}{arbitrarily small scales}};
\filldraw[fill=setupbg,draw=setupfg,line width=1.0bp] (630.04,441.00) ellipse (84.00bp and 23.50bp);
\node[inner sep=0pt,font=\sffamily\fontsize{10.7bp}{12bp}\selectfont\bfseries,text=ink] at (630.04,447.50) {\hyperlink{node-W2L2_13}{WZ2 Lemma 2.13}};
\node[inner sep=0pt,font=\sffamily\fontsize{10.7bp}{12bp}\selectfont,text=ink] at (630.04,434.50) {\hyperlink{node-W2L2_13}{Multi-scale inheritance}};
\filldraw[fill=apibg,draw=apifg,line width=1.0bp] (197.67,441.00) ellipse (112.38bp and 30.00bp);
\node[inner sep=0pt,font=\sffamily\fontsize{10.7bp}{12bp}\selectfont\bfseries,text=ink] at (197.67,454.00) {\hyperlink{node-Ienvelope}{WZ2 Sec. 6.2: consequence}};
\node[inner sep=0pt,font=\sffamily\fontsize{10.7bp}{12bp}\selectfont,text=ink] at (197.67,441.00) {\hyperlink{node-Ienvelope}{Uniform lower envelope}};
\node[inner sep=0pt,font=\sffamily\fontsize{10.7bp}{12bp}\selectfont,text=ink] at (197.67,428.00) {\hyperlink{node-Ienvelope}{for admissible unions}};
\filldraw[fill=setupbg,draw=setupfg,line width=1.0bp,rounded corners=5bp] (108.48,506.50) rectangle (276.48,549.50);
\node[inner sep=0pt,font=\sffamily\fontsize{10.7bp}{12bp}\selectfont\bfseries,text=ink] at (192.48,534.50) {\hyperlink{node-W2D6_1}{WZ2 Definition 6.1}};
\node[inner sep=0pt,font=\sffamily\fontsize{10.7bp}{12bp}\selectfont,text=ink] at (192.48,521.50) {\hyperlink{node-W2D6_1}{CW-extremal pair}};
\filldraw[fill=setupbg,draw=setupfg,line width=1.0bp,rounded corners=5bp] (315.67,341.00) rectangle (483.67,384.00);
\node[inner sep=0pt,font=\sffamily\fontsize{10.7bp}{12bp}\selectfont\bfseries,text=ink] at (399.67,369.00) {\hyperlink{node-W2Sigma}{WZ2 Section 6.2}};
\node[inner sep=0pt,font=\sffamily\fontsize{10.7bp}{12bp}\selectfont,text=ink] at (399.67,356.00) {\hyperlink{node-W2Sigma}{Counterexample deficit}};
\filldraw[fill=setupbg,draw=setupfg,line width=1.0bp,rounded corners=5bp] (315.67,271.00) rectangle (483.67,314.00);
\node[inner sep=0pt,font=\sffamily\fontsize{10.7bp}{12bp}\selectfont\bfseries,text=ink] at (399.67,299.00) {\hyperlink{node-W2MW}{WZ2 Definition 2.12}};
\node[inner sep=0pt,font=\sffamily\fontsize{10.7bp}{12bp}\selectfont,text=ink] at (399.67,286.00) {\hyperlink{node-W2MW}{Multi-scale convex CW}};
\filldraw[fill=setupbg,draw=setupfg,line width=1.0bp,rounded corners=5bp] (416.67,188.00) rectangle (584.67,244.00);
\node[inner sep=0pt,font=\sffamily\fontsize{10.7bp}{12bp}\selectfont\bfseries,text=ink] at (500.67,229.00) {\hyperlink{node-W2CWA}{WZ2 Def. 2.5}};
\node[inner sep=0pt,font=\sffamily\fontsize{10.7bp}{12bp}\selectfont,text=ink] at (500.67,216.00) {\hyperlink{node-W2CWA}{Normalized}};
\node[inner sep=0pt,font=\sffamily\fontsize{10.7bp}{12bp}\selectfont,text=ink] at (500.67,203.00) {\hyperlink{node-W2CWA}{convex-containment count}};
\filldraw[fill=setupbg,draw=setupfg,line width=1.0bp,rounded corners=5bp] (214.67,188.00) rectangle (382.67,244.00);
\node[inner sep=0pt,font=\sffamily\fontsize{10.7bp}{12bp}\selectfont\bfseries,text=ink] at (298.67,229.00) {\hyperlink{node-W2Unit}{WZ2 Def. 2.8 / Sec. 6.1}};
\node[inner sep=0pt,font=\sffamily\fontsize{10.7bp}{12bp}\selectfont,text=ink] at (298.67,216.00) {\hyperlink{node-W2Unit}{Geometric and shaded unit}};
\node[inner sep=0pt,font=\sffamily\fontsize{10.7bp}{12bp}\selectfont,text=ink] at (298.67,203.00) {\hyperlink{node-W2Unit}{rescaling}};
\filldraw[fill=apibg,draw=apifg,line width=1.0bp,rounded corners=5bp] (214.67,105.00) rectangle (382.67,161.00);
\node[inner sep=0pt,font=\sffamily\fontsize{10.7bp}{12bp}\selectfont\bfseries,text=ink] at (298.67,146.00) {\hyperlink{node-Icover}{Cover API}};
\node[inner sep=0pt,font=\sffamily\fontsize{10.7bp}{12bp}\selectfont,text=ink] at (298.67,133.00) {\hyperlink{node-Icover}{Balanced coarse/fine}};
\node[inner sep=0pt,font=\sffamily\fontsize{10.7bp}{12bp}\selectfont,text=ink] at (298.67,120.00) {\hyperlink{node-Icover}{configurations}};
\filldraw[fill=apibg,draw=apifg,line width=1.0bp,rounded corners=5bp] (315.67,22.00) rectangle (483.67,78.00);
\node[inner sep=0pt,font=\sffamily\fontsize{10.7bp}{12bp}\selectfont\bfseries,text=ink] at (399.67,63.00) {\hyperlink{node-Imass}{Refinement API}};
\node[inner sep=0pt,font=\sffamily\fontsize{10.7bp}{12bp}\selectfont,text=ink] at (399.67,50.00) {\hyperlink{node-Imass}{Mass and constant}};
\node[inner sep=0pt,font=\sffamily\fontsize{10.7bp}{12bp}\selectfont,text=ink] at (399.67,37.00) {\hyperlink{node-Imass}{multiplicity}};
\filldraw[fill=apibg,draw=apifg,line width=1.0bp,rounded corners=5bp] (344.04,413.00) rectangle (512.04,469.00);
\node[inner sep=0pt,font=\sffamily\fontsize{10.7bp}{12bp}\selectfont\bfseries,text=ink] at (428.04,454.00) {\hyperlink{node-Iparams}{Quantifier API}};
\node[inner sep=0pt,font=\sffamily\fontsize{10.7bp}{12bp}\selectfont,text=ink] at (428.04,441.00) {\hyperlink{node-Iparams}{Small-parameter and}};
\node[inner sep=0pt,font=\sffamily\fontsize{10.7bp}{12bp}\selectfont,text=ink] at (428.04,428.00) {\hyperlink{node-Iparams}{finite-iteration}};
\filldraw[fill=apibg,draw=apifg,line width=1.0bp,rounded corners=5bp] (416.67,105.00) rectangle (584.67,161.00);
\node[inner sep=0pt,font=\sffamily\fontsize{10.7bp}{12bp}\selectfont\bfseries,text=ink] at (500.67,146.00) {\hyperlink{node-Iaffine}{Affine transport API}};
\node[inner sep=0pt,font=\sffamily\fontsize{10.7bp}{12bp}\selectfont,text=ink] at (500.67,133.00) {\hyperlink{node-Iaffine}{Jacobian, direction and}};
\node[inner sep=0pt,font=\sffamily\fontsize{10.7bp}{12bp}\selectfont,text=ink] at (500.67,120.00) {\hyperlink{node-Iaffine}{grid transport}};
\filldraw[fill=apibg,draw=apifg,line width=1.0bp,rounded corners=5bp] (522.85,500.00) rectangle (690.85,556.00);
\node[inner sep=0pt,font=\sffamily\fontsize{10.7bp}{12bp}\selectfont\bfseries,text=ink] at (606.85,541.00) {\hyperlink{node-Iwolffbalance}{WZ2 proof operation}};
\node[inner sep=0pt,font=\sffamily\fontsize{10.7bp}{12bp}\selectfont,text=ink] at (606.85,528.00) {\hyperlink{node-Iwolffbalance}{Restriction and}};
\node[inner sep=0pt,font=\sffamily\fontsize{10.7bp}{12bp}\selectfont,text=ink] at (606.85,515.00) {\hyperlink{node-Iwolffbalance}{rebuilding of cover}};
\end{tikzpicture}}\end{center}
{\footnotesize Solid: proof input. Dashed: API. Dash-dot: assembly. Long-dashed: overview path. Read from the top down; arrows point to consumers. Full prerequisites appear on the cards and in the ledger.}\par
\sectitle{Plane map and Lipschitz refinement}
\hypertarget{view-plane}{}
{\sffamily\small\color{muted}The WZ1 Section 4.1 proof is expanded with the new coarse/fine interface.}\par
\begin{center}\resizebox{493.00bp}{!}{\begin{tikzpicture}[x=1bp,y=1bp, >={Stealth[length=4bp,width=3.2bp]}]
\path[use as bounding box] (0,0) rectangle (675.50,718.00);
\fill[panel,rounded corners=5bp] (0,0) rectangle (675.50,718.00);
\draw[line width=3.5bp,panel] (520.96,156.50) .. controls (520.96,164.60) and (645.50,167.30) .. (645.50,167.30) .. controls (645.50,274.67) and (645.50,417.83) .. (645.50,525.20) .. controls (645.50,525.20) and (445.69,527.90) .. (445.69,536.00);
\draw[->,line width=1.0bp,core] (520.96,156.50) .. controls (520.96,164.60) and (645.50,167.30) .. (645.50,167.30) .. controls (645.50,274.67) and (645.50,417.83) .. (645.50,525.20) .. controls (645.50,525.20) and (445.69,527.90) .. (445.69,536.00);
\draw[line width=3.5bp,panel] (416.80,156.50) .. controls (416.80,177.23) and (492.58,180.83) .. (492.58,201.56);
\draw[->,line width=1.0bp,core] (416.80,156.50) .. controls (416.80,177.23) and (492.58,180.83) .. (492.58,201.56);
\draw[line width=3.5bp,panel] (440.50,243.50) .. controls (440.50,260.98) and (339.50,264.02) .. (339.50,281.50);
\draw[->,line width=1.0bp,core] (440.50,243.50) .. controls (440.50,260.98) and (339.50,264.02) .. (339.50,281.50);
\draw[line width=3.5bp,panel] (137.50,335.00) .. controls (137.50,352.74) and (181.23,355.82) .. (181.23,373.56);
\draw[->,line width=1.0bp,core] (137.50,335.00) .. controls (137.50,352.74) and (181.23,355.82) .. (181.23,373.56);
\draw[line width=3.5bp,panel] (287.42,323.44) .. controls (287.42,344.40) and (250.67,348.05) .. (250.67,369.01);
\draw[->,line width=1.0bp,core] (287.42,323.44) .. controls (287.42,344.40) and (250.67,348.05) .. (250.67,369.01);
\draw[line width=3.5bp,panel] (388.42,238.44) .. controls (388.42,246.54) and (33.50,249.24) .. (33.50,249.24) .. controls (33.50,305.93) and (33.50,381.51) .. (33.50,438.20) .. controls (33.50,438.20) and (238.50,440.90) .. (238.50,449.00);
\draw[->,line width=1.0bp,core] (388.42,238.44) .. controls (388.42,246.54) and (33.50,249.24) .. (33.50,249.24) .. controls (33.50,305.93) and (33.50,381.51) .. (33.50,438.20) .. controls (33.50,438.20) and (238.50,440.90) .. (238.50,449.00);
\draw[line width=3.5bp,panel] (391.58,323.44) .. controls (391.58,344.15) and (385.20,347.75) .. (385.20,368.46);
\draw[->,line width=1.0bp,core] (391.58,323.44) .. controls (391.58,344.15) and (385.20,347.75) .. (385.20,368.46);
\draw[line width=3.5bp,panel] (492.58,238.44) .. controls (492.58,246.54) and (649.00,249.24) .. (649.00,249.24) .. controls (649.00,279.83) and (649.00,320.61) .. (649.00,351.20) .. controls (649.00,351.20) and (440.50,353.90) .. (440.50,362.00);
\draw[->,line width=1.0bp,core] (492.58,238.44) .. controls (492.58,246.54) and (649.00,249.24) .. (649.00,249.24) .. controls (649.00,279.83) and (649.00,320.61) .. (649.00,351.20) .. controls (649.00,351.20) and (440.50,353.90) .. (440.50,362.00);
\draw[line width=3.5bp,panel] (440.50,422.00) .. controls (440.50,434.42) and (445.69,436.58) .. (445.69,449.00);
\draw[->,line width=1.0bp,core] (440.50,422.00) .. controls (440.50,434.42) and (445.69,436.58) .. (445.69,449.00);
\draw[line width=3.5bp,panel] (445.69,509.00) .. controls (445.69,527.38) and (266.82,530.58) .. (266.82,548.96);
\draw[->,line width=1.0bp,core] (445.69,509.00) .. controls (445.69,527.38) and (266.82,530.58) .. (266.82,548.96);
\draw[line width=3.5bp,panel] (445.69,609.00) .. controls (445.69,624.39) and (394.80,627.07) .. (394.80,642.46);
\draw[->,line width=1.0bp,core] (445.69,609.00) .. controls (445.69,624.39) and (394.80,627.07) .. (394.80,642.46);
\draw[line width=3.5bp,panel] (285.39,410.44) .. controls (285.39,418.54) and (30.00,421.24) .. (30.00,421.24) .. controls (30.00,482.62) and (30.00,564.46) .. (30.00,625.85) .. controls (30.00,625.85) and (321.07,628.55) .. (321.07,636.65);
\draw[->,line width=1.0bp,core] (285.39,410.44) .. controls (285.39,418.54) and (30.00,421.24) .. (30.00,421.24) .. controls (30.00,482.62) and (30.00,564.46) .. (30.00,625.85) .. controls (30.00,625.85) and (321.07,628.55) .. (321.07,636.65);
\draw[line width=3.5bp,panel] (293.80,502.54) .. controls (293.80,510.64) and (26.50,513.34) .. (26.50,513.34) .. controls (26.50,547.09) and (26.50,592.09) .. (26.50,625.85) .. controls (26.50,625.85) and (357.93,628.55) .. (357.93,636.65);
\draw[->,line width=1.0bp,core] (293.80,502.54) .. controls (293.80,510.64) and (26.50,513.34) .. (26.50,513.34) .. controls (26.50,547.09) and (26.50,592.09) .. (26.50,625.85) .. controls (26.50,625.85) and (357.93,628.55) .. (357.93,636.65);
\draw[line width=3.5bp,panel] (211.52,602.50) .. controls (211.52,620.88) and (284.20,624.08) .. (284.20,642.46);
\draw[->,line width=1.0bp,core] (211.52,602.50) .. controls (211.52,620.88) and (284.20,624.08) .. (284.20,642.46);
\draw[line width=3.5bp,panel] (238.50,248.00) .. controls (238.50,260.42) and (137.50,262.58) .. (137.50,275.00);
\draw[->,line width=1.0bp,core] (238.50,248.00) .. controls (238.50,260.42) and (137.50,262.58) .. (137.50,275.00);
\draw[line width=3.5bp,panel] (339.50,78.00) .. controls (339.50,86.10) and (652.50,88.80) .. (652.50,88.80) .. controls (652.50,117.87) and (652.50,156.63) .. (652.50,185.70) .. controls (652.50,185.70) and (440.50,188.40) .. (440.50,196.50);
\draw[->,line width=1.0bp,api,dash pattern=on 5bp off 4bp] (339.50,78.00) .. controls (339.50,86.10) and (652.50,88.80) .. (652.50,88.80) .. controls (652.50,117.87) and (652.50,156.63) .. (652.50,185.70) .. controls (652.50,185.70) and (440.50,188.40) .. (440.50,196.50);
\draw[line width=3.5bp,panel] (287.42,78.00) .. controls (287.42,86.10) and (23.00,88.80) .. (23.00,88.80) .. controls (23.00,170.99) and (23.00,280.57) .. (23.00,362.76) .. controls (23.00,362.76) and (285.39,365.46) .. (285.39,373.56);
\draw[->,line width=1.0bp,api,dash pattern=on 5bp off 4bp] (287.42,78.00) .. controls (287.42,86.10) and (23.00,88.80) .. (23.00,88.80) .. controls (23.00,170.99) and (23.00,280.57) .. (23.00,362.76) .. controls (23.00,362.76) and (285.39,365.46) .. (285.39,373.56);
\draw[line width=3.5bp,panel] (489.42,326.50) .. controls (489.42,334.60) and (656.00,337.30) .. (656.00,337.30) .. controls (656.00,369.51) and (656.00,412.45) .. (656.00,444.66) .. controls (656.00,444.66) and (293.80,447.36) .. (293.80,455.46);
\draw[->,line width=1.0bp,api,dash pattern=on 5bp off 4bp] (489.42,326.50) .. controls (489.42,334.60) and (656.00,337.30) .. (656.00,337.30) .. controls (656.00,369.51) and (656.00,412.45) .. (656.00,444.66) .. controls (656.00,444.66) and (293.80,447.36) .. (293.80,455.46);
\draw[line width=3.5bp,panel] (593.58,326.50) .. controls (593.58,345.80) and (495.80,349.16) .. (495.80,368.46);
\draw[->,line width=1.0bp,api,dash pattern=on 5bp off 4bp] (593.58,326.50) .. controls (593.58,345.80) and (495.80,349.16) .. (495.80,368.46);
\draw[line width=3.5bp,panel] (181.23,410.44) .. controls (181.23,431.15) and (183.20,434.75) .. (183.20,455.46);
\draw[->,line width=1.0bp,assembly,dash pattern=on 7bp off 3bp on 1bp off 3bp] (181.23,410.44) .. controls (181.23,431.15) and (183.20,434.75) .. (183.20,455.46);
\draw[line width=3.5bp,panel] (183.20,502.54) .. controls (183.20,523.89) and (156.23,527.61) .. (156.23,548.96);
\draw[->,line width=1.0bp,assembly,dash pattern=on 7bp off 3bp on 1bp off 3bp] (183.20,502.54) .. controls (183.20,523.89) and (156.23,527.61) .. (156.23,548.96);
\draw[line width=3.5bp,panel] (391.58,78.00) .. controls (391.58,94.33) and (468.88,97.17) .. (468.88,113.50);
\draw[->,line width=1.0bp,core] (391.58,78.00) .. controls (391.58,94.33) and (468.88,97.17) .. (468.88,113.50);
\draw[line width=3.5bp,panel] (308.17,158.54) .. controls (308.17,178.33) and (388.42,181.77) .. (388.42,201.56);
\draw[->,line width=1.0bp,core] (308.17,158.54) .. controls (308.17,178.33) and (388.42,181.77) .. (388.42,201.56);
\draw[line width=3.5bp,panel] (168.83,158.54) .. controls (168.83,166.64) and (19.50,169.34) .. (19.50,169.34) .. controls (19.50,226.00) and (19.50,301.55) .. (19.50,358.21) .. controls (19.50,358.21) and (215.95,360.91) .. (215.95,369.01);
\draw[->,line width=1.0bp,core] (168.83,158.54) .. controls (168.83,166.64) and (19.50,169.34) .. (19.50,169.34) .. controls (19.50,226.00) and (19.50,301.55) .. (19.50,358.21) .. controls (19.50,358.21) and (215.95,360.91) .. (215.95,369.01);
\filldraw[fill=planebg,draw=planefg,line width=1.0bp] (339.50,666.00) ellipse (89.19bp and 30.00bp);
\node[inner sep=0pt,font=\sffamily\fontsize{10.7bp}{12bp}\selectfont\bfseries,text=ink] at (339.50,679.00) {\hyperlink{node-L4_8}{WZ1 Lemma 4.8 / CW}};
\node[inner sep=0pt,font=\sffamily\fontsize{10.7bp}{12bp}\selectfont,text=ink] at (339.50,666.00) {\hyperlink{node-L4_8}{Existence of Lipschitz}};
\node[inner sep=0pt,font=\sffamily\fontsize{10.7bp}{12bp}\selectfont,text=ink] at (339.50,653.00) {\hyperlink{node-L4_8}{plane data}};
\filldraw[fill=setupbg,draw=setupfg,line width=1.0bp] (445.69,572.50) ellipse (110.98bp and 36.50bp);
\node[inner sep=0pt,font=\sffamily\fontsize{10.7bp}{12bp}\selectfont\bfseries,text=ink] at (445.69,592.00) {\hyperlink{node-W2Exists}{WZ2 Sec. 6.2 / WZ1 Lem.}};
\node[inner sep=0pt,font=\sffamily\fontsize{10.7bp}{12bp}\selectfont,text=ink] at (445.69,579.00) {\hyperlink{node-W2Exists}{3.1}};
\node[inner sep=0pt,font=\sffamily\fontsize{10.7bp}{12bp}\selectfont,text=ink] at (445.69,566.00) {\hyperlink{node-W2Exists}{Extremizers exist at}};
\node[inner sep=0pt,font=\sffamily\fontsize{10.7bp}{12bp}\selectfont,text=ink] at (445.69,553.00) {\hyperlink{node-W2Exists}{arbitrarily small scales}};
\filldraw[fill=planebg,draw=planefg,line width=1.0bp] (238.50,479.00) ellipse (89.19bp and 30.00bp);
\node[inner sep=0pt,font=\sffamily\fontsize{10.7bp}{12bp}\selectfont\bfseries,text=ink] at (238.50,492.00) {\hyperlink{node-L4_4}{WZ1 Lemma 4.4 / CW}};
\node[inner sep=0pt,font=\sffamily\fontsize{10.7bp}{12bp}\selectfont,text=ink] at (238.50,479.00) {\hyperlink{node-L4_4}{Coarsen weak planiness to}};
\node[inner sep=0pt,font=\sffamily\fontsize{10.7bp}{12bp}\selectfont,text=ink] at (238.50,466.00) {\hyperlink{node-L4_4}{a plane map}};
\filldraw[fill=planebg,draw=planefg,line width=1.0bp] (211.52,572.50) ellipse (89.19bp and 30.00bp);
\node[inner sep=0pt,font=\sffamily\fontsize{10.7bp}{12bp}\selectfont\bfseries,text=ink] at (211.52,585.50) {\hyperlink{node-L4_7}{WZ1 Lemma 4.7 / CW}};
\node[inner sep=0pt,font=\sffamily\fontsize{10.7bp}{12bp}\selectfont,text=ink] at (211.52,572.50) {\hyperlink{node-L4_7}{Iterate to a Lipschitz}};
\node[inner sep=0pt,font=\sffamily\fontsize{10.7bp}{12bp}\selectfont,text=ink] at (211.52,559.50) {\hyperlink{node-L4_7}{plane map}};
\filldraw[fill=planebg,draw=planefg,line width=1.0bp] (440.50,392.00) ellipse (89.19bp and 30.00bp);
\node[inner sep=0pt,font=\sffamily\fontsize{10.7bp}{12bp}\selectfont\bfseries,text=ink] at (440.50,405.00) {\hyperlink{node-L4_5}{WZ1 Lemma 4.5 / CW}};
\node[inner sep=0pt,font=\sffamily\fontsize{10.7bp}{12bp}\selectfont,text=ink] at (440.50,392.00) {\hyperlink{node-L4_5}{Small oscillation on a}};
\node[inner sep=0pt,font=\sffamily\fontsize{10.7bp}{12bp}\selectfont,text=ink] at (440.50,379.00) {\hyperlink{node-L4_5}{fixed cube scale}};
\filldraw[fill=planebg,draw=planefg,line width=1.0bp] (445.69,479.00) ellipse (84.00bp and 30.00bp);
\node[inner sep=0pt,font=\sffamily\fontsize{10.7bp}{12bp}\selectfont\bfseries,text=ink] at (445.69,492.00) {\hyperlink{node-C4_6}{WZ1 Cor. 4.6 / CW}};
\node[inner sep=0pt,font=\sffamily\fontsize{10.7bp}{12bp}\selectfont,text=ink] at (445.69,479.00) {\hyperlink{node-C4_6}{Small oscillation for}};
\node[inner sep=0pt,font=\sffamily\fontsize{10.7bp}{12bp}\selectfont,text=ink] at (445.69,466.00) {\hyperlink{node-C4_6}{nearby points}};
\filldraw[fill=planebg,draw=planefg,line width=1.0bp] (233.31,392.00) ellipse (84.00bp and 23.50bp);
\node[inner sep=0pt,font=\sffamily\fontsize{10.7bp}{12bp}\selectfont\bfseries,text=ink] at (233.31,398.50) {\hyperlink{node-L4_3}{WZ1 Lemma 4.3 / CW}};
\node[inner sep=0pt,font=\sffamily\fontsize{10.7bp}{12bp}\selectfont,text=ink] at (233.31,385.50) {\hyperlink{node-L4_3}{Weak planiness}};
\filldraw[fill=setupbg,draw=setupfg,line width=1.0bp] (440.50,220.00) ellipse (84.00bp and 23.50bp);
\node[inner sep=0pt,font=\sffamily\fontsize{10.7bp}{12bp}\selectfont\bfseries,text=ink] at (440.50,226.50) {\hyperlink{node-W2P6_2}{WZ2 Proposition 6.2}};
\node[inner sep=0pt,font=\sffamily\fontsize{10.7bp}{12bp}\selectfont,text=ink] at (440.50,213.50) {\hyperlink{node-W2P6_2}{Coarse/fine extremality}};
\filldraw[fill=setupbg,draw=setupfg,line width=1.0bp] (339.50,305.00) ellipse (84.00bp and 23.50bp);
\node[inner sep=0pt,font=\sffamily\fontsize{10.7bp}{12bp}\selectfont\bfseries,text=ink] at (339.50,311.50) {\hyperlink{node-L3_4}{WZ1 Lemma 3.4 / CW}};
\node[inner sep=0pt,font=\sffamily\fontsize{10.7bp}{12bp}\selectfont,text=ink] at (339.50,298.50) {\hyperlink{node-L3_4}{Robust transversality}};
\filldraw[fill=planebg,draw=planefg,line width=1.0bp] (137.50,305.00) ellipse (84.00bp and 30.00bp);
\node[inner sep=0pt,font=\sffamily\fontsize{10.7bp}{12bp}\selectfont\bfseries,text=ink] at (137.50,318.00) {\hyperlink{node-T4_2}{WZ1 Thm. 4.2}};
\node[inner sep=0pt,font=\sffamily\fontsize{10.7bp}{12bp}\selectfont,text=ink] at (137.50,305.00) {\hyperlink{node-T4_2}{Endpoint multilinear}};
\node[inner sep=0pt,font=\sffamily\fontsize{10.7bp}{12bp}\selectfont,text=ink] at (137.50,292.00) {\hyperlink{node-T4_2}{Kakeya input}};
\filldraw[fill=setupbg,draw=setupfg,line width=1.0bp,rounded corners=5bp] (384.88,113.50) rectangle (552.88,156.50);
\node[inner sep=0pt,font=\sffamily\fontsize{10.7bp}{12bp}\selectfont\bfseries,text=ink] at (468.88,141.50) {\hyperlink{node-W2D6_1}{WZ2 Definition 6.1}};
\node[inner sep=0pt,font=\sffamily\fontsize{10.7bp}{12bp}\selectfont,text=ink] at (468.88,128.50) {\hyperlink{node-W2D6_1}{CW-extremal pair}};
\filldraw[fill=planebg,draw=planefg,line width=1.0bp,rounded corners=5bp] (457.50,283.50) rectangle (625.50,326.50);
\node[inner sep=0pt,font=\sffamily\fontsize{10.7bp}{12bp}\selectfont\bfseries,text=ink] at (541.50,311.50) {\hyperlink{node-Dplane}{WZ2 Sec. 6.1: plane map}};
\node[inner sep=0pt,font=\sffamily\fontsize{10.7bp}{12bp}\selectfont,text=ink] at (541.50,298.50) {\hyperlink{node-Dplane}{Plane map}};
\filldraw[fill=externalbg,draw=externalfg,line width=1.0bp,rounded corners=5bp,dashed] (154.50,192.00) rectangle (322.50,248.00);
\node[inner sep=0pt,font=\sffamily\fontsize{10.7bp}{12bp}\selectfont\bfseries,text=ink] at (238.50,233.00) {\hyperlink{node-XBCT}{External [1] / [8]}};
\node[inner sep=0pt,font=\sffamily\fontsize{10.7bp}{12bp}\selectfont,text=ink] at (238.50,220.00) {\hyperlink{node-XBCT}{Endpoint multilinear}};
\node[inner sep=0pt,font=\sffamily\fontsize{10.7bp}{12bp}\selectfont,text=ink] at (238.50,207.00) {\hyperlink{node-XBCT}{Kakeya}};
\filldraw[fill=apibg,draw=apifg,line width=1.0bp,rounded corners=5bp] (255.50,22.00) rectangle (423.50,78.00);
\node[inner sep=0pt,font=\sffamily\fontsize{10.7bp}{12bp}\selectfont\bfseries,text=ink] at (339.50,63.00) {\hyperlink{node-Imass}{Refinement API}};
\node[inner sep=0pt,font=\sffamily\fontsize{10.7bp}{12bp}\selectfont,text=ink] at (339.50,50.00) {\hyperlink{node-Imass}{Mass and constant}};
\node[inner sep=0pt,font=\sffamily\fontsize{10.7bp}{12bp}\selectfont,text=ink] at (339.50,37.00) {\hyperlink{node-Imass}{multiplicity}};
\filldraw[fill=apibg,draw=apifg,line width=1.0bp] (238.50,135.00) ellipse (112.38bp and 30.00bp);
\node[inner sep=0pt,font=\sffamily\fontsize{10.7bp}{12bp}\selectfont\bfseries,text=ink] at (238.50,148.00) {\hyperlink{node-Ienvelope}{WZ2 Sec. 6.2: consequence}};
\node[inner sep=0pt,font=\sffamily\fontsize{10.7bp}{12bp}\selectfont,text=ink] at (238.50,135.00) {\hyperlink{node-Ienvelope}{Uniform lower envelope}};
\node[inner sep=0pt,font=\sffamily\fontsize{10.7bp}{12bp}\selectfont,text=ink] at (238.50,122.00) {\hyperlink{node-Ienvelope}{for admissible unions}};
\end{tikzpicture}}\end{center}
{\footnotesize Solid: proof input. Dashed: API. Dash-dot: assembly. Long-dashed: overview path. Read from the top down; arrows point to consumers. Full prerequisites appear on the cards and in the ledger.}\par
\sectitle{Local and global grains}
\hypertarget{view-grains}{}
{\sffamily\small\color{muted}The upper AD covering structure from WZ1 Section 4.2 supplies WZ2 Proposition 6.3.}\par
\begin{center}\resizebox{481.87bp}{!}{\begin{tikzpicture}[x=1bp,y=1bp, >={Stealth[length=4bp,width=3.2bp]}]
\path[use as bounding box] (0,0) rectangle (726.47,992.00);
\fill[panel,rounded corners=5bp] (0,0) rectangle (726.47,992.00);
\draw[line width=3.5bp,panel] (417.07,78.00) .. controls (417.07,86.10) and (685.97,88.80) .. (685.97,88.80) .. controls (685.97,310.32) and (685.97,605.68) .. (685.97,827.20) .. controls (685.97,827.20) and (581.97,829.90) .. (581.97,838.00);
\draw[->,line width=1.0bp,api,dash pattern=on 5bp off 4bp] (417.07,78.00) .. controls (417.07,86.10) and (685.97,88.80) .. (685.97,88.80) .. controls (685.97,310.32) and (685.97,605.68) .. (685.97,827.20) .. controls (685.97,827.20) and (581.97,829.90) .. (581.97,838.00);
\draw[line width=3.5bp,panel] (338.95,150.84) .. controls (338.95,163.79) and (364.99,166.05) .. (364.99,179.00);
\draw[->,line width=1.0bp,core] (338.95,150.84) .. controls (338.95,163.79) and (364.99,166.05) .. (364.99,179.00);
\draw[line width=3.5bp,panel] (417.07,220.94) .. controls (417.07,235.69) and (465.99,238.25) .. (465.99,253.00);
\draw[->,line width=1.0bp,core] (417.07,220.94) .. controls (417.07,235.69) and (465.99,238.25) .. (465.99,253.00);
\draw[line width=3.5bp,panel] (312.91,146.94) .. controls (312.91,155.04) and (44.00,157.74) .. (44.00,157.74) .. controls (44.00,205.28) and (44.00,268.66) .. (44.00,316.20) .. controls (44.00,316.20) and (263.99,318.90) .. (263.99,327.00);
\draw[->,line width=1.0bp,core] (312.91,146.94) .. controls (312.91,155.04) and (44.00,157.74) .. (44.00,157.74) .. controls (44.00,205.28) and (44.00,268.66) .. (44.00,316.20) .. controls (44.00,316.20) and (263.99,318.90) .. (263.99,327.00);
\draw[line width=3.5bp,panel] (413.91,294.94) .. controls (413.91,303.04) and (40.50,305.74) .. (40.50,305.74) .. controls (40.50,363.02) and (40.50,439.38) .. (40.50,496.66) .. controls (40.50,496.66) and (319.28,499.36) .. (319.28,507.46);
\draw[->,line width=1.0bp,core] (413.91,294.94) .. controls (413.91,303.04) and (40.50,305.74) .. (40.50,305.74) .. controls (40.50,363.02) and (40.50,439.38) .. (40.50,496.66) .. controls (40.50,496.66) and (319.28,499.36) .. (319.28,507.46);
\draw[line width=3.5bp,panel] (208.69,380.54) .. controls (208.69,388.64) and (37.00,391.34) .. (37.00,391.34) .. controls (37.00,422.94) and (37.00,465.06) .. (37.00,496.66) .. controls (37.00,496.66) and (208.69,499.36) .. (208.69,507.46);
\draw[->,line width=1.0bp,core] (208.69,380.54) .. controls (208.69,388.64) and (37.00,391.34) .. (37.00,391.34) .. controls (37.00,422.94) and (37.00,465.06) .. (37.00,496.66) .. controls (37.00,496.66) and (208.69,499.36) .. (208.69,507.46);
\draw[line width=3.5bp,panel] (309.69,467.54) .. controls (309.69,482.93) and (263.99,485.61) .. (263.99,501.00);
\draw[->,line width=1.0bp,core] (309.69,467.54) .. controls (309.69,482.93) and (263.99,485.61) .. (263.99,501.00);
\draw[line width=3.5bp,panel] (364.99,152.00) .. controls (364.99,160.10) and (689.47,162.80) .. (689.47,162.80) .. controls (689.47,287.27) and (689.47,453.24) .. (689.47,577.71) .. controls (689.47,577.71) and (347.57,580.41) .. (347.57,588.51);
\draw[->,line width=1.0bp,core] (364.99,152.00) .. controls (364.99,160.10) and (689.47,162.80) .. (689.47,162.80) .. controls (689.47,287.27) and (689.47,453.24) .. (689.47,577.71) .. controls (689.47,577.71) and (347.57,580.41) .. (347.57,588.51);
\draw[line width=3.5bp,panel] (312.91,220.94) .. controls (312.91,229.04) and (692.97,231.74) .. (692.97,231.74) .. controls (692.97,335.53) and (692.97,473.92) .. (692.97,577.71) .. controls (692.97,577.71) and (382.40,580.41) .. (382.40,588.51);
\draw[->,line width=1.0bp,core] (312.91,220.94) .. controls (312.91,229.04) and (692.97,231.74) .. (692.97,231.74) .. controls (692.97,335.53) and (692.97,473.92) .. (692.97,577.71) .. controls (692.97,577.71) and (382.40,580.41) .. (382.40,588.51);
\draw[line width=3.5bp,panel] (471.17,559.00) .. controls (471.17,574.67) and (417.23,577.39) .. (417.23,593.06);
\draw[->,line width=1.0bp,core] (471.17,559.00) .. controls (471.17,574.67) and (417.23,577.39) .. (417.23,593.06);
\draw[line width=3.5bp,panel] (364.99,635.00) .. controls (364.99,650.39) and (306.83,653.07) .. (306.83,668.46);
\draw[->,line width=1.0bp,core] (364.99,635.00) .. controls (364.99,650.39) and (306.83,653.07) .. (306.83,668.46);
\draw[line width=3.5bp,panel] (391.03,150.84) .. controls (391.03,158.94) and (33.50,161.64) .. (33.50,161.64) .. controls (33.50,310.45) and (33.50,508.86) .. (33.50,657.66) .. controls (33.50,657.66) and (423.14,660.36) .. (423.14,668.46);
\draw[->,line width=1.0bp,core] (391.03,150.84) .. controls (391.03,158.94) and (33.50,161.64) .. (33.50,161.64) .. controls (33.50,310.45) and (33.50,508.86) .. (33.50,657.66) .. controls (33.50,657.66) and (423.14,660.36) .. (423.14,668.46);
\draw[line width=3.5bp,panel] (364.99,722.00) .. controls (364.99,737.39) and (407.83,740.07) .. (407.83,755.46);
\draw[->,line width=1.0bp,core] (364.99,722.00) .. controls (364.99,737.39) and (407.83,740.07) .. (407.83,755.46);
\draw[line width=3.5bp,panel] (465.99,809.00) .. controls (465.99,824.39) and (312.02,827.07) .. (312.02,842.46);
\draw[->,line width=1.0bp,core] (465.99,809.00) .. controls (465.99,824.39) and (312.02,827.07) .. (312.02,842.46);
\draw[line width=3.5bp,panel] (208.69,554.54) .. controls (208.69,562.64) and (30.00,565.34) .. (30.00,565.34) .. controls (30.00,669.74) and (30.00,808.95) .. (30.00,913.36) .. controls (30.00,913.36) and (338.95,916.06) .. (338.95,924.16);
\draw[->,line width=1.0bp,core] (208.69,554.54) .. controls (208.69,562.64) and (30.00,565.34) .. (30.00,565.34) .. controls (30.00,669.74) and (30.00,808.95) .. (30.00,913.36) .. controls (30.00,913.36) and (338.95,916.06) .. (338.95,924.16);
\draw[line width=3.5bp,panel] (370.17,896.00) .. controls (370.17,908.95) and (391.03,911.21) .. (391.03,924.16);
\draw[->,line width=1.0bp,core] (370.17,896.00) .. controls (370.17,908.95) and (391.03,911.21) .. (391.03,924.16);
\draw[line width=3.5bp,panel] (417.07,146.94) .. controls (417.07,155.04) and (696.47,157.74) .. (696.47,157.74) .. controls (696.47,384.08) and (696.47,685.86) .. (696.47,912.20) .. controls (696.47,912.20) and (364.99,914.90) .. (364.99,923.00);
\draw[->,line width=1.0bp,core] (417.07,146.94) .. controls (417.07,155.04) and (696.47,157.74) .. (696.47,157.74) .. controls (696.47,384.08) and (696.47,685.86) .. (696.47,912.20) .. controls (696.47,912.20) and (364.99,914.90) .. (364.99,923.00);
\draw[line width=3.5bp,panel] (581.97,894.00) .. controls (581.97,909.67) and (417.07,912.39) .. (417.07,928.06);
\draw[->,line width=1.0bp,core] (581.97,894.00) .. controls (581.97,909.67) and (417.07,912.39) .. (417.07,928.06);
\draw[line width=3.5bp,panel] (518.07,294.94) .. controls (518.07,303.04) and (699.97,305.74) .. (699.97,305.74) .. controls (699.97,488.49) and (699.97,732.15) .. (699.97,914.90) .. controls (699.97,914.90) and (404.05,917.60) .. (404.05,925.70);
\draw[->,line width=1.0bp,core] (518.07,294.94) .. controls (518.07,303.04) and (699.97,305.74) .. (699.97,305.74) .. controls (699.97,488.49) and (699.97,732.15) .. (699.97,914.90) .. controls (699.97,914.90) and (404.05,917.60) .. (404.05,925.70);
\draw[line width=3.5bp,panel] (263.99,387.00) .. controls (263.99,395.10) and (26.50,397.80) .. (26.50,397.80) .. controls (26.50,552.21) and (26.50,758.08) .. (26.50,912.48) .. controls (26.50,912.48) and (351.97,915.18) .. (351.97,923.28);
\draw[->,line width=1.0bp,core] (263.99,387.00) .. controls (263.99,395.10) and (26.50,397.80) .. (26.50,397.80) .. controls (26.50,552.21) and (26.50,758.08) .. (26.50,912.48) .. controls (26.50,912.48) and (351.97,915.18) .. (351.97,923.28);
\draw[line width=3.5bp,panel] (153.19,896.00) .. controls (153.19,910.75) and (312.91,913.31) .. (312.91,928.06);
\draw[->,line width=1.0bp,core] (153.19,896.00) .. controls (153.19,910.75) and (312.91,913.31) .. (312.91,928.06);
\draw[line width=3.5bp,panel] (254.19,807.00) .. controls (254.19,815.10) and (23.00,817.80) .. (23.00,817.80) .. controls (23.00,846.93) and (23.00,885.77) .. (23.00,914.90) .. controls (23.00,914.90) and (325.93,917.60) .. (325.93,925.70);
\draw[->,line width=1.0bp,core] (254.19,807.00) .. controls (254.19,815.10) and (23.00,817.80) .. (23.00,817.80) .. controls (23.00,846.93) and (23.00,885.77) .. (23.00,914.90) .. controls (23.00,914.90) and (325.93,917.60) .. (325.93,925.70);
\draw[line width=3.5bp,panel] (312.91,78.00) .. controls (312.91,90.42) and (364.99,92.58) .. (364.99,105.00);
\draw[->,line width=1.0bp,api,dash pattern=on 5bp off 4bp] (312.91,78.00) .. controls (312.91,90.42) and (364.99,92.58) .. (364.99,105.00);
\draw[line width=3.5bp,panel] (263.99,298.00) .. controls (263.99,314.31) and (208.69,317.15) .. (208.69,333.46);
\draw[->,line width=1.0bp,api,dash pattern=on 5bp off 4bp] (263.99,298.00) .. controls (263.99,314.31) and (208.69,317.15) .. (208.69,333.46);
\draw[line width=3.5bp,panel] (419.09,385.00) .. controls (419.09,401.31) and (420.28,404.15) .. (420.28,420.46);
\draw[->,line width=1.0bp,api,dash pattern=on 5bp off 4bp] (419.09,385.00) .. controls (419.09,401.31) and (420.28,404.15) .. (420.28,420.46);
\draw[line width=3.5bp,panel] (523.25,385.00) .. controls (523.25,393.10) and (703.47,395.80) .. (703.47,395.80) .. controls (703.47,500.46) and (703.47,640.00) .. (703.47,744.66) .. controls (703.47,744.66) and (524.14,747.36) .. (524.14,755.46);
\draw[->,line width=1.0bp,api,dash pattern=on 5bp off 4bp] (523.25,385.00) .. controls (523.25,393.10) and (703.47,395.80) .. (703.47,395.80) .. controls (703.47,500.46) and (703.47,640.00) .. (703.47,744.66) .. controls (703.47,744.66) and (524.14,747.36) .. (524.14,755.46);
\draw[line width=3.5bp,panel] (471.17,385.00) .. controls (471.17,393.10) and (706.97,395.80) .. (706.97,395.80) .. controls (706.97,526.56) and (706.97,700.90) .. (706.97,831.66) .. controls (706.97,831.66) and (428.33,834.36) .. (428.33,842.46);
\draw[->,line width=1.0bp,api,dash pattern=on 5bp off 4bp] (471.17,385.00) .. controls (471.17,393.10) and (706.97,395.80) .. (706.97,395.80) .. controls (706.97,526.56) and (706.97,700.90) .. (706.97,831.66) .. controls (706.97,831.66) and (428.33,834.36) .. (428.33,842.46);
\draw[line width=3.5bp,panel] (465.99,300.00) .. controls (465.99,315.39) and (319.28,318.07) .. (319.28,333.46);
\draw[->,line width=1.0bp,assembly,dash pattern=on 7bp off 3bp on 1bp off 3bp] (465.99,300.00) .. controls (465.99,315.39) and (319.28,318.07) .. (319.28,333.46);
\draw[line width=3.5bp,panel] (319.28,380.54) .. controls (319.28,398.90) and (309.69,402.10) .. (309.69,420.46);
\draw[->,line width=1.0bp,assembly,dash pattern=on 7bp off 3bp on 1bp off 3bp] (319.28,380.54) .. controls (319.28,398.90) and (309.69,402.10) .. (309.69,420.46);
\draw[line width=3.5bp,panel] (319.28,554.54) .. controls (319.28,572.26) and (312.74,575.34) .. (312.74,593.06);
\draw[->,line width=1.0bp,assembly,dash pattern=on 7bp off 3bp on 1bp off 3bp] (319.28,554.54) .. controls (319.28,572.26) and (312.74,575.34) .. (312.74,593.06);
\draw[line width=3.5bp,panel] (420.28,467.54) .. controls (420.28,475.64) and (19.50,478.34) .. (19.50,478.34) .. controls (19.50,608.58) and (19.50,782.24) .. (19.50,912.48) .. controls (19.50,912.48) and (378.01,915.18) .. (378.01,923.28);
\draw[->,line width=1.0bp,assembly,dash pattern=on 7bp off 3bp on 1bp off 3bp] (420.28,467.54) .. controls (420.28,475.64) and (19.50,478.34) .. (19.50,478.34) .. controls (19.50,608.58) and (19.50,782.24) .. (19.50,912.48) .. controls (19.50,912.48) and (378.01,915.18) .. (378.01,923.28);
\filldraw[fill=grainsbg,draw=grainsfg,line width=1.0bp] (364.99,946.50) ellipse (84.00bp and 23.50bp);
\node[inner sep=0pt,font=\sffamily\fontsize{10.7bp}{12bp}\selectfont\bfseries,text=ink] at (364.99,953.00) {\hyperlink{node-W2P6_3}{WZ2 Proposition 6.3}};
\node[inner sep=0pt,font=\sffamily\fontsize{10.7bp}{12bp}\selectfont,text=ink] at (364.99,940.00) {\hyperlink{node-W2P6_3}{Lipschitz grains}};
\filldraw[fill=setupbg,draw=setupfg,line width=1.0bp] (153.19,866.00) ellipse (89.19bp and 30.00bp);
\node[inner sep=0pt,font=\sffamily\fontsize{10.7bp}{12bp}\selectfont\bfseries,text=ink] at (153.19,879.00) {\hyperlink{node-L3_5}{WZ1 Lemma 3.5 / CW}};
\node[inner sep=0pt,font=\sffamily\fontsize{10.7bp}{12bp}\selectfont,text=ink] at (153.19,866.00) {\hyperlink{node-L3_5}{Mild anisotropic}};
\node[inner sep=0pt,font=\sffamily\fontsize{10.7bp}{12bp}\selectfont,text=ink] at (153.19,853.00) {\hyperlink{node-L3_5}{rescaling preserves}};
\filldraw[fill=grainsbg,draw=grainsfg,line width=1.0bp] (465.99,779.00) ellipse (93.80bp and 30.00bp);
\node[inner sep=0pt,font=\sffamily\fontsize{10.7bp}{12bp}\selectfont\bfseries,text=ink] at (465.99,792.00) {\hyperlink{node-L4_11}{WZ1 Lemma 4.11 / CW}};
\node[inner sep=0pt,font=\sffamily\fontsize{10.7bp}{12bp}\selectfont,text=ink] at (465.99,779.00) {\hyperlink{node-L4_11}{AD spacing at one grain}};
\node[inner sep=0pt,font=\sffamily\fontsize{10.7bp}{12bp}\selectfont,text=ink] at (465.99,766.00) {\hyperlink{node-L4_11}{scale}};
\filldraw[fill=grainsbg,draw=grainsfg,line width=1.0bp] (370.17,866.00) ellipse (93.80bp and 30.00bp);
\node[inner sep=0pt,font=\sffamily\fontsize{10.7bp}{12bp}\selectfont\bfseries,text=ink] at (370.17,879.00) {\hyperlink{node-L4_12}{WZ1 Lemma 4.12 / CW}};
\node[inner sep=0pt,font=\sffamily\fontsize{10.7bp}{12bp}\selectfont,text=ink] at (370.17,866.00) {\hyperlink{node-L4_12}{AD spacing at every grain}};
\node[inner sep=0pt,font=\sffamily\fontsize{10.7bp}{12bp}\selectfont,text=ink] at (370.17,853.00) {\hyperlink{node-L4_12}{scale}};
\filldraw[fill=planebg,draw=planefg,line width=1.0bp] (263.99,531.00) ellipse (89.19bp and 30.00bp);
\node[inner sep=0pt,font=\sffamily\fontsize{10.7bp}{12bp}\selectfont\bfseries,text=ink] at (263.99,544.00) {\hyperlink{node-L4_8}{WZ1 Lemma 4.8 / CW}};
\node[inner sep=0pt,font=\sffamily\fontsize{10.7bp}{12bp}\selectfont,text=ink] at (263.99,531.00) {\hyperlink{node-L4_8}{Existence of Lipschitz}};
\node[inner sep=0pt,font=\sffamily\fontsize{10.7bp}{12bp}\selectfont,text=ink] at (263.99,518.00) {\hyperlink{node-L4_8}{plane data}};
\filldraw[fill=grainsbg,draw=grainsfg,line width=1.0bp] (364.99,611.50) ellipse (84.26bp and 23.50bp);
\node[inner sep=0pt,font=\sffamily\fontsize{10.7bp}{12bp}\selectfont\bfseries,text=ink] at (364.99,618.00) {\hyperlink{node-L4_9}{WZ1 Lemma 4.9 / CW}};
\node[inner sep=0pt,font=\sffamily\fontsize{10.7bp}{12bp}\selectfont,text=ink] at (364.99,605.00) {\hyperlink{node-L4_9}{Grain count in tau-cubes}};
\filldraw[fill=grainsbg,draw=grainsfg,line width=1.0bp] (364.99,692.00) ellipse (93.80bp and 30.00bp);
\node[inner sep=0pt,font=\sffamily\fontsize{10.7bp}{12bp}\selectfont\bfseries,text=ink] at (364.99,705.00) {\hyperlink{node-L4_10}{WZ1 Lemma 4.10 / CW}};
\node[inner sep=0pt,font=\sffamily\fontsize{10.7bp}{12bp}\selectfont,text=ink] at (364.99,692.00) {\hyperlink{node-L4_10}{Grain count in thin}};
\node[inner sep=0pt,font=\sffamily\fontsize{10.7bp}{12bp}\selectfont,text=ink] at (364.99,679.00) {\hyperlink{node-L4_10}{prisms}};
\filldraw[fill=planebg,draw=planefg,line width=1.0bp] (465.99,276.50) ellipse (84.00bp and 23.50bp);
\node[inner sep=0pt,font=\sffamily\fontsize{10.7bp}{12bp}\selectfont\bfseries,text=ink] at (465.99,283.00) {\hyperlink{node-L4_3}{WZ1 Lemma 4.3 / CW}};
\node[inner sep=0pt,font=\sffamily\fontsize{10.7bp}{12bp}\selectfont,text=ink] at (465.99,270.00) {\hyperlink{node-L4_3}{Weak planiness}};
\filldraw[fill=planebg,draw=planefg,line width=1.0bp] (263.99,357.00) ellipse (89.19bp and 30.00bp);
\node[inner sep=0pt,font=\sffamily\fontsize{10.7bp}{12bp}\selectfont\bfseries,text=ink] at (263.99,370.00) {\hyperlink{node-L4_4}{WZ1 Lemma 4.4 / CW}};
\node[inner sep=0pt,font=\sffamily\fontsize{10.7bp}{12bp}\selectfont,text=ink] at (263.99,357.00) {\hyperlink{node-L4_4}{Coarsen weak planiness to}};
\node[inner sep=0pt,font=\sffamily\fontsize{10.7bp}{12bp}\selectfont,text=ink] at (263.99,344.00) {\hyperlink{node-L4_4}{a plane map}};
\filldraw[fill=planebg,draw=planefg,line width=1.0bp] (364.99,444.00) ellipse (89.19bp and 30.00bp);
\node[inner sep=0pt,font=\sffamily\fontsize{10.7bp}{12bp}\selectfont\bfseries,text=ink] at (364.99,457.00) {\hyperlink{node-L4_7}{WZ1 Lemma 4.7 / CW}};
\node[inner sep=0pt,font=\sffamily\fontsize{10.7bp}{12bp}\selectfont,text=ink] at (364.99,444.00) {\hyperlink{node-L4_7}{Iterate to a Lipschitz}};
\node[inner sep=0pt,font=\sffamily\fontsize{10.7bp}{12bp}\selectfont,text=ink] at (364.99,431.00) {\hyperlink{node-L4_7}{plane map}};
\filldraw[fill=setupbg,draw=setupfg,line width=1.0bp] (364.99,128.50) ellipse (84.00bp and 23.50bp);
\node[inner sep=0pt,font=\sffamily\fontsize{10.7bp}{12bp}\selectfont\bfseries,text=ink] at (364.99,135.00) {\hyperlink{node-W2P6_2}{WZ2 Proposition 6.2}};
\node[inner sep=0pt,font=\sffamily\fontsize{10.7bp}{12bp}\selectfont,text=ink] at (364.99,122.00) {\hyperlink{node-W2P6_2}{Coarse/fine extremality}};
\filldraw[fill=setupbg,draw=setupfg,line width=1.0bp] (364.99,202.50) ellipse (84.00bp and 23.50bp);
\node[inner sep=0pt,font=\sffamily\fontsize{10.7bp}{12bp}\selectfont\bfseries,text=ink] at (364.99,209.00) {\hyperlink{node-L3_4}{WZ1 Lemma 3.4 / CW}};
\node[inner sep=0pt,font=\sffamily\fontsize{10.7bp}{12bp}\selectfont,text=ink] at (364.99,196.00) {\hyperlink{node-L3_4}{Robust transversality}};
\filldraw[fill=setupbg,draw=setupfg,line width=1.0bp,rounded corners=5bp] (497.97,838.00) rectangle (665.97,894.00);
\node[inner sep=0pt,font=\sffamily\fontsize{10.7bp}{12bp}\selectfont\bfseries,text=ink] at (581.97,879.00) {\hyperlink{node-W2Unit}{WZ2 Def. 2.8 / Sec. 6.1}};
\node[inner sep=0pt,font=\sffamily\fontsize{10.7bp}{12bp}\selectfont,text=ink] at (581.97,866.00) {\hyperlink{node-W2Unit}{Geometric and shaded unit}};
\node[inner sep=0pt,font=\sffamily\fontsize{10.7bp}{12bp}\selectfont,text=ink] at (581.97,853.00) {\hyperlink{node-W2Unit}{rescaling}};
\filldraw[fill=grainsbg,draw=grainsfg,line width=1.0bp,rounded corners=5bp] (387.17,503.00) rectangle (555.17,559.00);
\node[inner sep=0pt,font=\sffamily\fontsize{10.7bp}{12bp}\selectfont\bfseries,text=ink] at (471.17,544.00) {\hyperlink{node-DAD}{WZ2 Sec. 6.1: ADset}};
\node[inner sep=0pt,font=\sffamily\fontsize{10.7bp}{12bp}\selectfont,text=ink] at (471.17,531.00) {\hyperlink{node-DAD}{Upper AD covering}};
\node[inner sep=0pt,font=\sffamily\fontsize{10.7bp}{12bp}\selectfont,text=ink] at (471.17,518.00) {\hyperlink{node-DAD}{condition}};
\filldraw[fill=externalbg,draw=externalfg,line width=1.0bp,rounded corners=5bp,dashed] (170.19,751.00) rectangle (338.19,807.00);
\node[inner sep=0pt,font=\sffamily\fontsize{10.7bp}{12bp}\selectfont\bfseries,text=ink] at (254.19,792.00) {\hyperlink{node-XExt}{External [29] / Mathlib}};
\node[inner sep=0pt,font=\sffamily\fontsize{10.7bp}{12bp}\selectfont,text=ink] at (254.19,779.00) {\hyperlink{node-XExt}{Scalar Lipschitz}};
\node[inner sep=0pt,font=\sffamily\fontsize{10.7bp}{12bp}\selectfont,text=ink] at (254.19,766.00) {\hyperlink{node-XExt}{extension}};
\filldraw[fill=planebg,draw=planefg,line width=1.0bp,rounded corners=5bp] (179.99,255.00) rectangle (347.99,298.00);
\node[inner sep=0pt,font=\sffamily\fontsize{10.7bp}{12bp}\selectfont\bfseries,text=ink] at (263.99,283.00) {\hyperlink{node-Dplane}{WZ2 Sec. 6.1: plane map}};
\node[inner sep=0pt,font=\sffamily\fontsize{10.7bp}{12bp}\selectfont,text=ink] at (263.99,270.00) {\hyperlink{node-Dplane}{Plane map}};
\filldraw[fill=apibg,draw=apifg,line width=1.0bp,rounded corners=5bp] (280.99,22.00) rectangle (448.99,78.00);
\node[inner sep=0pt,font=\sffamily\fontsize{10.7bp}{12bp}\selectfont\bfseries,text=ink] at (364.99,63.00) {\hyperlink{node-Icover}{Cover API}};
\node[inner sep=0pt,font=\sffamily\fontsize{10.7bp}{12bp}\selectfont,text=ink] at (364.99,50.00) {\hyperlink{node-Icover}{Balanced coarse/fine}};
\node[inner sep=0pt,font=\sffamily\fontsize{10.7bp}{12bp}\selectfont,text=ink] at (364.99,37.00) {\hyperlink{node-Icover}{configurations}};
\filldraw[fill=apibg,draw=apifg,line width=1.0bp,rounded corners=5bp] (387.17,329.00) rectangle (555.17,385.00);
\node[inner sep=0pt,font=\sffamily\fontsize{10.7bp}{12bp}\selectfont\bfseries,text=ink] at (471.17,370.00) {\hyperlink{node-Iparams}{Quantifier API}};
\node[inner sep=0pt,font=\sffamily\fontsize{10.7bp}{12bp}\selectfont,text=ink] at (471.17,357.00) {\hyperlink{node-Iparams}{Small-parameter and}};
\node[inner sep=0pt,font=\sffamily\fontsize{10.7bp}{12bp}\selectfont,text=ink] at (471.17,344.00) {\hyperlink{node-Iparams}{finite-iteration}};
\end{tikzpicture}}\end{center}
{\footnotesize Solid: proof input. Dashed: API. Dash-dot: assembly. Long-dashed: overview path. Read from the top down; arrows point to consumers. Full prerequisites appear on the cards and in the ledger.}\par
\sectitle{C2 regularity of the global slope}
\hypertarget{view-regularity}{}
{\sffamily\small\color{muted}Retain WZ1 Sections 5 and 8: dot-product expansion yields a compatible trapezoid tree.}\par
\begin{center}\resizebox{493.00bp}{!}{\begin{tikzpicture}[x=1bp,y=1bp, >={Stealth[length=4bp,width=3.2bp]}]
\path[use as bounding box] (0,0) rectangle (678.87,779.00);
\fill[panel,rounded corners=5bp] (0,0) rectangle (678.87,779.00);
\draw[line width=3.5bp,panel] (341.19,78.00) .. controls (341.19,90.42) and (341.19,92.58) .. (341.19,105.00);
\draw[->,line width=1.0bp,api,dash pattern=on 5bp off 4bp] (341.19,78.00) .. controls (341.19,90.42) and (341.19,92.58) .. (341.19,105.00);
\draw[line width=3.5bp,panel] (182.92,236.44) .. controls (182.92,259.50) and (87.11,263.50) .. (87.11,286.56);
\draw[->,line width=1.0bp,core] (182.92,236.44) .. controls (182.92,259.50) and (87.11,263.50) .. (87.11,286.56);
\draw[line width=3.5bp,panel] (386.89,241.54) .. controls (386.89,262.25) and (191.27,265.85) .. (191.27,286.56);
\draw[->,line width=1.0bp,core] (386.89,241.54) .. controls (386.89,262.25) and (191.27,265.85) .. (191.27,286.56);
\draw[line width=3.5bp,panel] (543.19,335.00) .. controls (543.19,350.39) and (393.27,353.07) .. (393.27,368.46);
\draw[->,line width=1.0bp,core] (543.19,335.00) .. controls (543.19,350.39) and (393.27,353.07) .. (393.27,368.46);
\draw[line width=3.5bp,panel] (341.19,422.00) .. controls (341.19,434.42) and (442.19,436.58) .. (442.19,449.00);
\draw[->,line width=1.0bp,core] (341.19,422.00) .. controls (341.19,434.42) and (442.19,436.58) .. (442.19,449.00);
\draw[line width=3.5bp,panel] (287.08,236.44) .. controls (287.08,244.54) and (30.00,247.24) .. (30.00,247.24) .. controls (30.00,306.47) and (30.00,385.43) .. (30.00,444.66) .. controls (30.00,444.66) and (386.89,447.36) .. (386.89,455.46);
\draw[->,line width=1.0bp,core] (287.08,236.44) .. controls (287.08,244.54) and (30.00,247.24) .. (30.00,247.24) .. controls (30.00,306.47) and (30.00,385.43) .. (30.00,444.66) .. controls (30.00,444.66) and (386.89,447.36) .. (386.89,455.46);
\draw[line width=3.5bp,panel] (442.19,509.00) .. controls (442.19,524.39) and (396.48,527.07) .. (396.48,542.46);
\draw[->,line width=1.0bp,core] (442.19,509.00) .. controls (442.19,524.39) and (396.48,527.07) .. (396.48,542.46);
\draw[line width=3.5bp,panel] (341.19,596.00) .. controls (341.19,611.39) and (390.11,614.07) .. (390.11,629.46);
\draw[->,line width=1.0bp,core] (341.19,596.00) .. controls (341.19,611.39) and (390.11,614.07) .. (390.11,629.46);
\draw[line width=3.5bp,panel] (600.45,594.00) .. controls (600.45,610.31) and (494.27,613.15) .. (494.27,629.46);
\draw[->,line width=1.0bp,core] (600.45,594.00) .. controls (600.45,610.31) and (494.27,613.15) .. (494.27,629.46);
\draw[line width=3.5bp,panel] (134.00,596.00) .. controls (134.00,611.39) and (188.11,614.07) .. (188.11,629.46);
\draw[->,line width=1.0bp,core] (134.00,596.00) .. controls (134.00,611.39) and (188.11,614.07) .. (188.11,629.46);
\draw[line width=3.5bp,panel] (496.29,594.00) .. controls (496.29,610.31) and (292.27,613.15) .. (292.27,629.46);
\draw[->,line width=1.0bp,core] (496.29,594.00) .. controls (496.29,610.31) and (292.27,613.15) .. (292.27,629.46);
\draw[line width=3.5bp,panel] (87.11,323.44) .. controls (87.11,331.54) and (26.50,334.24) .. (26.50,334.24) .. controls (26.50,445.25) and (26.50,593.25) .. (26.50,704.26) .. controls (26.50,704.26) and (289.11,706.96) .. (289.11,715.06);
\draw[->,line width=1.0bp,core] (87.11,323.44) .. controls (87.11,331.54) and (26.50,334.24) .. (26.50,334.24) .. controls (26.50,445.25) and (26.50,593.25) .. (26.50,704.26) .. controls (26.50,704.26) and (289.11,706.96) .. (289.11,715.06);
\draw[line width=3.5bp,panel] (442.19,683.00) .. controls (442.19,695.65) and (358.55,697.85) .. (358.55,710.51);
\draw[->,line width=1.0bp,core] (442.19,683.00) .. controls (442.19,695.65) and (358.55,697.85) .. (358.55,710.51);
\draw[line width=3.5bp,panel] (240.19,683.00) .. controls (240.19,695.65) and (323.83,697.85) .. (323.83,710.51);
\draw[->,line width=1.0bp,core] (240.19,683.00) .. controls (240.19,695.65) and (323.83,697.85) .. (323.83,710.51);
\draw[line width=3.5bp,panel] (497.48,241.54) .. controls (497.48,249.64) and (652.37,252.34) .. (652.37,252.34) .. controls (652.37,387.92) and (652.37,568.68) .. (652.37,704.26) .. controls (652.37,704.26) and (393.27,706.96) .. (393.27,715.06);
\draw[->,line width=1.0bp,core] (497.48,241.54) .. controls (497.48,249.64) and (652.37,252.34) .. (652.37,252.34) .. controls (652.37,387.92) and (652.37,568.68) .. (652.37,704.26) .. controls (652.37,704.26) and (393.27,706.96) .. (393.27,715.06);
\draw[line width=3.5bp,panel] (289.11,161.00) .. controls (289.11,178.74) and (182.92,181.82) .. (182.92,199.56);
\draw[->,line width=1.0bp,api,dash pattern=on 5bp off 4bp] (289.11,161.00) .. controls (289.11,178.74) and (182.92,181.82) .. (182.92,199.56);
\draw[line width=3.5bp,panel] (289.11,78.00) .. controls (289.11,86.10) and (23.00,88.80) .. (23.00,88.80) .. controls (23.00,118.79) and (23.00,158.77) .. (23.00,188.76) .. controls (23.00,188.76) and (287.08,191.46) .. (287.08,199.56);
\draw[->,line width=1.0bp,api,dash pattern=on 5bp off 4bp] (289.11,78.00) .. controls (289.11,86.10) and (23.00,88.80) .. (23.00,88.80) .. controls (23.00,118.79) and (23.00,158.77) .. (23.00,188.76) .. controls (23.00,188.76) and (287.08,191.46) .. (287.08,199.56);
\draw[line width=3.5bp,panel] (341.19,333.00) .. controls (341.19,346.34) and (341.19,348.66) .. (341.19,362.00);
\draw[->,line width=1.0bp,api,dash pattern=on 5bp off 4bp] (341.19,333.00) .. controls (341.19,346.34) and (341.19,348.66) .. (341.19,362.00);
\draw[line width=3.5bp,panel] (393.27,161.00) .. controls (393.27,169.10) and (655.87,171.80) .. (655.87,171.80) .. controls (655.87,253.66) and (655.87,362.80) .. (655.87,444.66) .. controls (655.87,444.66) and (497.48,447.36) .. (497.48,455.46);
\draw[->,line width=1.0bp,api,dash pattern=on 5bp off 4bp] (393.27,161.00) .. controls (393.27,169.10) and (655.87,171.80) .. (655.87,171.80) .. controls (655.87,253.66) and (655.87,362.80) .. (655.87,444.66) .. controls (655.87,444.66) and (497.48,447.36) .. (497.48,455.46);
\draw[line width=3.5bp,panel] (393.27,78.00) .. controls (393.27,86.10) and (659.37,88.80) .. (659.37,88.80) .. controls (659.37,221.66) and (659.37,398.80) .. (659.37,531.66) .. controls (659.37,531.66) and (285.89,534.36) .. (285.89,542.46);
\draw[->,line width=1.0bp,api,dash pattern=on 5bp off 4bp] (393.27,78.00) .. controls (393.27,86.10) and (659.37,88.80) .. (659.37,88.80) .. controls (659.37,221.66) and (659.37,398.80) .. (659.37,531.66) .. controls (659.37,531.66) and (285.89,534.36) .. (285.89,542.46);
\draw[line width=3.5bp,panel] (182.92,507.00) .. controls (182.92,520.34) and (134.00,522.66) .. (134.00,536.00);
\draw[->,line width=1.0bp,api,dash pattern=on 5bp off 4bp] (182.92,507.00) .. controls (182.92,520.34) and (134.00,522.66) .. (134.00,536.00);
\draw[line width=3.5bp,panel] (287.08,507.00) .. controls (287.08,515.10) and (19.50,517.80) .. (19.50,517.80) .. controls (19.50,546.12) and (19.50,583.88) .. (19.50,612.20) .. controls (19.50,612.20) and (240.19,614.90) .. (240.19,623.00);
\draw[->,line width=1.0bp,api,dash pattern=on 5bp off 4bp] (287.08,507.00) .. controls (287.08,515.10) and (19.50,517.80) .. (19.50,517.80) .. controls (19.50,546.12) and (19.50,583.88) .. (19.50,612.20) .. controls (19.50,612.20) and (240.19,614.90) .. (240.19,623.00);
\draw[line width=3.5bp,panel] (191.27,323.44) .. controls (191.27,344.15) and (289.11,347.75) .. (289.11,368.46);
\draw[->,line width=1.0bp,assembly,dash pattern=on 7bp off 3bp on 1bp off 3bp] (191.27,323.44) .. controls (191.27,344.15) and (289.11,347.75) .. (289.11,368.46);
\filldraw[fill=regularitybg,draw=regularityfg,line width=1.0bp] (341.19,733.50) ellipse (84.00bp and 23.50bp);
\node[inner sep=0pt,font=\sffamily\fontsize{10.7bp}{12bp}\selectfont\bfseries,text=ink] at (341.19,740.00) {\hyperlink{node-W2P6_4}{WZ2 Proposition 6.4}};
\node[inner sep=0pt,font=\sffamily\fontsize{10.7bp}{12bp}\selectfont,text=ink] at (341.19,727.00) {\hyperlink{node-W2P6_4}{C2 grains}};
\filldraw[fill=regularitybg,draw=regularityfg,line width=1.0bp] (134.00,566.00) ellipse (84.00bp and 30.00bp);
\node[inner sep=0pt,font=\sffamily\fontsize{10.7bp}{12bp}\selectfont\bfseries,text=ink] at (134.00,579.00) {\hyperlink{node-L5_8}{WZ1 Lemma 5.8}};
\node[inner sep=0pt,font=\sffamily\fontsize{10.7bp}{12bp}\selectfont,text=ink] at (134.00,566.00) {\hyperlink{node-L5_8}{Smooth extension of an}};
\node[inner sep=0pt,font=\sffamily\fontsize{10.7bp}{12bp}\selectfont,text=ink] at (134.00,553.00) {\hyperlink{node-L5_8}{affine segment}};
\filldraw[fill=regularitybg,draw=regularityfg,line width=1.0bp] (240.19,653.00) ellipse (84.00bp and 30.00bp);
\node[inner sep=0pt,font=\sffamily\fontsize{10.7bp}{12bp}\selectfont\bfseries,text=ink] at (240.19,666.00) {\hyperlink{node-P5_7}{WZ1 Prop. 5.7}};
\node[inner sep=0pt,font=\sffamily\fontsize{10.7bp}{12bp}\selectfont,text=ink] at (240.19,653.00) {\hyperlink{node-P5_7}{A trapezoid tree produces}};
\node[inner sep=0pt,font=\sffamily\fontsize{10.7bp}{12bp}\selectfont,text=ink] at (240.19,640.00) {\hyperlink{node-P5_7}{a C2 approximation}};
\filldraw[fill=regularitybg,draw=regularityfg,line width=1.0bp] (442.19,653.00) ellipse (84.00bp and 30.00bp);
\node[inner sep=0pt,font=\sffamily\fontsize{10.7bp}{12bp}\selectfont\bfseries,text=ink] at (442.19,666.00) {\hyperlink{node-C5_6}{WZ1 Cor. 5.6 / CW}};
\node[inner sep=0pt,font=\sffamily\fontsize{10.7bp}{12bp}\selectfont,text=ink] at (442.19,653.00) {\hyperlink{node-C5_6}{A nested trapezoid tree}};
\node[inner sep=0pt,font=\sffamily\fontsize{10.7bp}{12bp}\selectfont,text=ink] at (442.19,640.00) {\hyperlink{node-C5_6}{for the slope graph}};
\filldraw[fill=regularitybg,draw=regularityfg,line width=1.0bp] (341.19,392.00) ellipse (84.00bp and 30.00bp);
\node[inner sep=0pt,font=\sffamily\fontsize{10.7bp}{12bp}\selectfont\bfseries,text=ink] at (341.19,405.00) {\hyperlink{node-L5_3}{WZ1 Lemma 5.3}};
\node[inner sep=0pt,font=\sffamily\fontsize{10.7bp}{12bp}\selectfont,text=ink] at (341.19,392.00) {\hyperlink{node-L5_3}{Local linearity from the}};
\node[inner sep=0pt,font=\sffamily\fontsize{10.7bp}{12bp}\selectfont,text=ink] at (341.19,379.00) {\hyperlink{node-L5_3}{two grain structures}};
\filldraw[fill=regularitybg,draw=regularityfg,line width=1.0bp] (341.19,566.00) ellipse (89.19bp and 30.00bp);
\node[inner sep=0pt,font=\sffamily\fontsize{10.7bp}{12bp}\selectfont\bfseries,text=ink] at (341.19,579.00) {\hyperlink{node-L5_5}{WZ1 Lemma 5.5 / CW}};
\node[inner sep=0pt,font=\sffamily\fontsize{10.7bp}{12bp}\selectfont,text=ink] at (341.19,566.00) {\hyperlink{node-L5_5}{Compatible approximations}};
\node[inner sep=0pt,font=\sffamily\fontsize{10.7bp}{12bp}\selectfont,text=ink] at (341.19,553.00) {\hyperlink{node-L5_5}{at many scales}};
\filldraw[fill=regularitybg,draw=regularityfg,line width=1.0bp] (442.19,479.00) ellipse (89.19bp and 30.00bp);
\node[inner sep=0pt,font=\sffamily\fontsize{10.7bp}{12bp}\selectfont\bfseries,text=ink] at (442.19,492.00) {\hyperlink{node-L5_4}{WZ1 Lemma 5.4 / CW}};
\node[inner sep=0pt,font=\sffamily\fontsize{10.7bp}{12bp}\selectfont,text=ink] at (442.19,479.00) {\hyperlink{node-L5_4}{Piecewise linear}};
\node[inner sep=0pt,font=\sffamily\fontsize{10.7bp}{12bp}\selectfont,text=ink] at (442.19,466.00) {\hyperlink{node-L5_4}{approximation at one}};
\filldraw[fill=regularitybg,draw=regularityfg,line width=1.0bp] (543.19,305.00) ellipse (84.00bp and 30.00bp);
\node[inner sep=0pt,font=\sffamily\fontsize{10.7bp}{12bp}\selectfont\bfseries,text=ink] at (543.19,318.00) {\hyperlink{node-T5_2}{WZ1 Thm. 5.2}};
\node[inner sep=0pt,font=\sffamily\fontsize{10.7bp}{12bp}\selectfont,text=ink] at (543.19,305.00) {\hyperlink{node-T5_2}{Dot-product expansion or}};
\node[inner sep=0pt,font=\sffamily\fontsize{10.7bp}{12bp}\selectfont,text=ink] at (543.19,292.00) {\hyperlink{node-T5_2}{orthogonal lines}};
\filldraw[fill=setupbg,draw=setupfg,line width=1.0bp] (235.00,218.00) ellipse (84.00bp and 23.50bp);
\node[inner sep=0pt,font=\sffamily\fontsize{10.7bp}{12bp}\selectfont\bfseries,text=ink] at (235.00,224.50) {\hyperlink{node-W2P6_2}{WZ2 Proposition 6.2}};
\node[inner sep=0pt,font=\sffamily\fontsize{10.7bp}{12bp}\selectfont,text=ink] at (235.00,211.50) {\hyperlink{node-W2P6_2}{Coarse/fine extremality}};
\filldraw[fill=grainsbg,draw=grainsfg,line width=1.0bp] (139.19,305.00) ellipse (84.00bp and 23.50bp);
\node[inner sep=0pt,font=\sffamily\fontsize{10.7bp}{12bp}\selectfont\bfseries,text=ink] at (139.19,311.50) {\hyperlink{node-W2P6_3}{WZ2 Proposition 6.3}};
\node[inner sep=0pt,font=\sffamily\fontsize{10.7bp}{12bp}\selectfont,text=ink] at (139.19,298.50) {\hyperlink{node-W2P6_3}{Lipschitz grains}};
\filldraw[fill=setupbg,draw=setupfg,line width=1.0bp] (442.19,218.00) ellipse (89.19bp and 30.00bp);
\node[inner sep=0pt,font=\sffamily\fontsize{10.7bp}{12bp}\selectfont\bfseries,text=ink] at (442.19,231.00) {\hyperlink{node-L3_5}{WZ1 Lemma 3.5 / CW}};
\node[inner sep=0pt,font=\sffamily\fontsize{10.7bp}{12bp}\selectfont,text=ink] at (442.19,218.00) {\hyperlink{node-L3_5}{Mild anisotropic}};
\node[inner sep=0pt,font=\sffamily\fontsize{10.7bp}{12bp}\selectfont,text=ink] at (442.19,205.00) {\hyperlink{node-L3_5}{rescaling preserves}};
\filldraw[fill=apibg,draw=apifg,line width=1.0bp,rounded corners=5bp] (151.00,451.00) rectangle (319.00,507.00);
\node[inner sep=0pt,font=\sffamily\fontsize{10.7bp}{12bp}\selectfont\bfseries,text=ink] at (235.00,492.00) {\hyperlink{node-Icalc}{Calculus API}};
\node[inner sep=0pt,font=\sffamily\fontsize{10.7bp}{12bp}\selectfont,text=ink] at (235.00,479.00) {\hyperlink{node-Icalc}{C2 bounds and compactly}};
\node[inner sep=0pt,font=\sffamily\fontsize{10.7bp}{12bp}\selectfont,text=ink] at (235.00,466.00) {\hyperlink{node-Icalc}{supported patches}};
\filldraw[fill=regularitybg,draw=regularityfg,line width=1.0bp,rounded corners=5bp] (257.19,277.00) rectangle (425.19,333.00);
\node[inner sep=0pt,font=\sffamily\fontsize{10.7bp}{12bp}\selectfont\bfseries,text=ink] at (341.19,318.00) {\hyperlink{node-Iadgap}{AD exponent gap}};
\node[inner sep=0pt,font=\sffamily\fontsize{10.7bp}{12bp}\selectfont,text=ink] at (341.19,305.00) {\hyperlink{node-Iadgap}{Why the expansion}};
\node[inner sep=0pt,font=\sffamily\fontsize{10.7bp}{12bp}\selectfont,text=ink] at (341.19,292.00) {\hyperlink{node-Iadgap}{alternative is impossible}};
\filldraw[fill=regularitybg,draw=regularityfg,line width=1.0bp,rounded corners=5bp] (464.37,538.00) rectangle (632.37,594.00);
\node[inner sep=0pt,font=\sffamily\fontsize{10.7bp}{12bp}\selectfont\bfseries,text=ink] at (548.37,579.00) {\hyperlink{node-Dtrap}{WZ1 Def. 5.1}};
\node[inner sep=0pt,font=\sffamily\fontsize{10.7bp}{12bp}\selectfont,text=ink] at (548.37,566.00) {\hyperlink{node-Dtrap}{Vertical trapezoids and}};
\node[inner sep=0pt,font=\sffamily\fontsize{10.7bp}{12bp}\selectfont,text=ink] at (548.37,553.00) {\hyperlink{node-Dtrap}{parents}};
\filldraw[fill=apibg,draw=apifg,line width=1.0bp,rounded corners=5bp] (257.19,22.00) rectangle (425.19,78.00);
\node[inner sep=0pt,font=\sffamily\fontsize{10.7bp}{12bp}\selectfont\bfseries,text=ink] at (341.19,63.00) {\hyperlink{node-Imass}{Refinement API}};
\node[inner sep=0pt,font=\sffamily\fontsize{10.7bp}{12bp}\selectfont,text=ink] at (341.19,50.00) {\hyperlink{node-Imass}{Mass and constant}};
\node[inner sep=0pt,font=\sffamily\fontsize{10.7bp}{12bp}\selectfont,text=ink] at (341.19,37.00) {\hyperlink{node-Imass}{multiplicity}};
\filldraw[fill=apibg,draw=apifg,line width=1.0bp,rounded corners=5bp] (257.19,105.00) rectangle (425.19,161.00);
\node[inner sep=0pt,font=\sffamily\fontsize{10.7bp}{12bp}\selectfont\bfseries,text=ink] at (341.19,146.00) {\hyperlink{node-Icover}{Cover API}};
\node[inner sep=0pt,font=\sffamily\fontsize{10.7bp}{12bp}\selectfont,text=ink] at (341.19,133.00) {\hyperlink{node-Icover}{Balanced coarse/fine}};
\node[inner sep=0pt,font=\sffamily\fontsize{10.7bp}{12bp}\selectfont,text=ink] at (341.19,120.00) {\hyperlink{node-Icover}{configurations}};
\end{tikzpicture}}\end{center}
{\footnotesize Solid: proof input. Dashed: API. Dash-dot: assembly. Long-dashed: overview path. Read from the top down; arrows point to consumers. Full prerequisites appear on the cards and in the ledger.}\par
\sectitle{Nonsingular slope: the new slab argument}
\hypertarget{view-slope}{}
{\sffamily\small\color{muted}WZ2 Lemma 6.6 replaces only the direction-separated endpoint of WZ1 Lemma 6.3.}\par
\begin{center}\resizebox{493.00bp}{!}{\begin{tikzpicture}[x=1bp,y=1bp, >={Stealth[length=4bp,width=3.2bp]}]
\path[use as bounding box] (0,0) rectangle (731.67,749.00);
\fill[panel,rounded corners=5bp] (0,0) rectangle (731.67,749.00);
\draw[line width=3.5bp,panel] (447.43,407.00) .. controls (447.43,415.10) and (698.17,417.80) .. (698.17,417.80) .. controls (698.17,468.32) and (698.17,535.68) .. (698.17,586.20) .. controls (698.17,586.20) and (466.84,588.90) .. (466.84,597.00);
\draw[->,line width=1.0bp,api,dash pattern=on 5bp off 4bp] (447.43,407.00) .. controls (447.43,415.10) and (698.17,417.80) .. (698.17,417.80) .. controls (698.17,468.32) and (698.17,535.68) .. (698.17,586.20) .. controls (698.17,586.20) and (466.84,588.90) .. (466.84,597.00);
\draw[line width=3.5bp,panel] (313.76,65.00) .. controls (313.76,73.10) and (33.50,75.80) .. (33.50,75.80) .. controls (33.50,180.62) and (33.50,320.38) .. (33.50,425.20) .. controls (33.50,425.20) and (264.84,427.90) .. (264.84,436.00);
\draw[->,line width=1.0bp,core] (313.76,65.00) .. controls (313.76,73.10) and (33.50,75.80) .. (33.50,75.80) .. controls (33.50,180.62) and (33.50,320.38) .. (33.50,425.20) .. controls (33.50,425.20) and (264.84,427.90) .. (264.84,436.00);
\draw[line width=3.5bp,panel] (417.92,65.00) .. controls (417.92,77.42) and (365.84,79.58) .. (365.84,92.00);
\draw[->,line width=1.0bp,core] (417.92,65.00) .. controls (417.92,77.42) and (365.84,79.58) .. (365.84,92.00);
\draw[line width=3.5bp,panel] (310.54,145.54) .. controls (310.54,163.00) and (264.84,166.04) .. (264.84,183.50);
\draw[->,line width=1.0bp,core] (310.54,145.54) .. controls (310.54,163.00) and (264.84,166.04) .. (264.84,183.50);
\draw[line width=3.5bp,panel] (518.92,235.00) .. controls (518.92,247.42) and (392.17,249.58) .. (392.17,262.00);
\draw[->,line width=1.0bp,core] (518.92,235.00) .. controls (518.92,247.42) and (392.17,249.58) .. (392.17,262.00);
\draw[line width=3.5bp,panel] (444.25,315.54) .. controls (444.25,323.64) and (701.67,326.34) .. (701.67,326.34) .. controls (701.67,356.34) and (701.67,396.35) .. (701.67,426.36) .. controls (701.67,426.36) and (492.88,429.06) .. (492.88,437.16);
\draw[->,line width=1.0bp,core] (444.25,315.54) .. controls (444.25,323.64) and (701.67,326.34) .. (701.67,326.34) .. controls (701.67,356.34) and (701.67,396.35) .. (701.67,426.36) .. controls (701.67,426.36) and (492.88,429.06) .. (492.88,437.16);
\draw[line width=3.5bp,panel] (282.20,229.99) .. controls (282.20,238.09) and (30.00,240.79) .. (30.00,240.79) .. controls (30.00,320.25) and (30.00,426.20) .. (30.00,505.66) .. controls (30.00,505.66) and (310.54,508.36) .. (310.54,516.46);
\draw[->,line width=1.0bp,core] (282.20,229.99) .. controls (282.20,238.09) and (30.00,240.79) .. (30.00,240.79) .. controls (30.00,320.25) and (30.00,426.20) .. (30.00,505.66) .. controls (30.00,505.66) and (310.54,508.36) .. (310.54,516.46);
\draw[line width=3.5bp,panel] (264.84,483.00) .. controls (264.84,495.42) and (365.84,497.58) .. (365.84,510.00);
\draw[->,line width=1.0bp,core] (264.84,483.00) .. controls (264.84,495.42) and (365.84,497.58) .. (365.84,510.00);
\draw[line width=3.5bp,panel] (466.84,483.00) .. controls (466.84,498.39) and (421.13,501.07) .. (421.13,516.46);
\draw[->,line width=1.0bp,core] (466.84,483.00) .. controls (466.84,498.39) and (421.13,501.07) .. (421.13,516.46);
\draw[line width=3.5bp,panel] (365.84,570.00) .. controls (365.84,584.49) and (264.84,587.01) .. (264.84,601.50);
\draw[->,line width=1.0bp,core] (365.84,570.00) .. controls (365.84,584.49) and (264.84,587.01) .. (264.84,601.50);
\draw[line width=3.5bp,panel] (421.13,145.54) .. controls (421.13,153.64) and (26.50,156.34) .. (26.50,156.34) .. controls (26.50,288.17) and (26.50,463.93) .. (26.50,595.76) .. controls (26.50,595.76) and (316.92,598.46) .. (316.92,606.56);
\draw[->,line width=1.0bp,core] (421.13,145.54) .. controls (421.13,153.64) and (26.50,156.34) .. (26.50,156.34) .. controls (26.50,288.17) and (26.50,463.93) .. (26.50,595.76) .. controls (26.50,595.76) and (316.92,598.46) .. (316.92,606.56);
\draw[line width=3.5bp,panel] (264.84,648.50) .. controls (264.84,665.32) and (313.76,668.24) .. (313.76,685.06);
\draw[->,line width=1.0bp,core] (264.84,648.50) .. controls (264.84,665.32) and (313.76,668.24) .. (313.76,685.06);
\draw[line width=3.5bp,panel] (466.84,653.00) .. controls (466.84,665.65) and (348.48,667.85) .. (348.48,680.51);
\draw[->,line width=1.0bp,core] (466.84,653.00) .. controls (466.84,665.65) and (348.48,667.85) .. (348.48,680.51);
\draw[line width=3.5bp,panel] (551.59,407.00) .. controls (551.59,422.67) and (518.92,425.39) .. (518.92,441.06);
\draw[->,line width=1.0bp,api,dash pattern=on 5bp off 4bp] (551.59,407.00) .. controls (551.59,422.67) and (518.92,425.39) .. (518.92,441.06);
\draw[line width=3.5bp,panel] (316.92,225.44) .. controls (316.92,233.54) and (705.17,236.24) .. (705.17,236.24) .. controls (705.17,293.27) and (705.17,369.32) .. (705.17,426.36) .. controls (705.17,426.36) and (440.80,429.06) .. (440.80,437.16);
\draw[->,line width=1.0bp,assembly,dash pattern=on 7bp off 3bp on 1bp off 3bp] (316.92,225.44) .. controls (316.92,233.54) and (705.17,236.24) .. (705.17,236.24) .. controls (705.17,293.27) and (705.17,369.32) .. (705.17,426.36) .. controls (705.17,426.36) and (440.80,429.06) .. (440.80,437.16);
\draw[line width=3.5bp,panel] (212.76,225.44) .. controls (212.76,245.23) and (95.43,248.67) .. (95.43,268.46);
\draw[->,line width=1.0bp,assembly,dash pattern=on 7bp off 3bp on 1bp off 3bp] (212.76,225.44) .. controls (212.76,245.23) and (95.43,248.67) .. (95.43,268.46);
\draw[line width=3.5bp,panel] (414.76,235.00) .. controls (414.76,250.39) and (232.25,253.07) .. (232.25,268.46);
\draw[->,line width=1.0bp,core] (414.76,235.00) .. controls (414.76,250.39) and (232.25,253.07) .. (232.25,268.46);
\draw[line width=3.5bp,panel] (95.43,315.54) .. controls (95.43,333.90) and (192.50,337.10) .. (192.50,355.46);
\draw[->,line width=1.0bp,core] (95.43,315.54) .. controls (95.43,333.90) and (192.50,337.10) .. (192.50,355.46);
\draw[line width=3.5bp,panel] (340.09,315.54) .. controls (340.09,330.93) and (264.84,333.61) .. (264.84,349.00);
\draw[->,line width=1.0bp,core] (340.09,315.54) .. controls (340.09,330.93) and (264.84,333.61) .. (264.84,349.00);
\draw[line width=3.5bp,panel] (542.09,320.00) .. controls (542.09,336.31) and (337.17,339.15) .. (337.17,355.46);
\draw[->,line width=1.0bp,core] (542.09,320.00) .. controls (542.09,336.31) and (337.17,339.15) .. (337.17,355.46);
\draw[line width=3.5bp,panel] (264.84,409.00) .. controls (264.84,421.42) and (466.84,423.58) .. (466.84,436.00);
\draw[->,line width=1.0bp,core] (264.84,409.00) .. controls (264.84,421.42) and (466.84,423.58) .. (466.84,436.00);
\draw[line width=3.5bp,panel] (232.25,315.54) .. controls (232.25,323.64) and (23.00,326.34) .. (23.00,326.34) .. controls (23.00,357.52) and (23.00,399.08) .. (23.00,430.26) .. controls (23.00,430.26) and (414.76,432.96) .. (414.76,441.06);
\draw[->,line width=1.0bp,core] (232.25,315.54) .. controls (232.25,323.64) and (23.00,326.34) .. (23.00,326.34) .. controls (23.00,357.52) and (23.00,399.08) .. (23.00,430.26) .. controls (23.00,430.26) and (414.76,432.96) .. (414.76,441.06);
\draw[line width=3.5bp,panel] (247.48,229.99) .. controls (247.48,238.09) and (19.50,240.79) .. (19.50,240.79) .. controls (19.50,347.28) and (19.50,489.27) .. (19.50,595.76) .. controls (19.50,595.76) and (212.76,598.46) .. (212.76,606.56);
\draw[->,line width=1.0bp,assembly,dash pattern=on 7bp off 3bp on 1bp off 3bp] (247.48,229.99) .. controls (247.48,238.09) and (19.50,240.79) .. (19.50,240.79) .. controls (19.50,347.28) and (19.50,489.27) .. (19.50,595.76) .. controls (19.50,595.76) and (212.76,598.46) .. (212.76,606.56);
\draw[line width=3.5bp,panel] (646.25,320.00) .. controls (646.25,328.10) and (708.67,330.80) .. (708.67,330.80) .. controls (708.67,433.84) and (708.67,571.22) .. (708.67,674.26) .. controls (708.67,674.26) and (417.92,676.96) .. (417.92,685.06);
\draw[->,line width=1.0bp,api,dash pattern=on 5bp off 4bp] (646.25,320.00) .. controls (646.25,328.10) and (708.67,330.80) .. (708.67,330.80) .. controls (708.67,433.84) and (708.67,571.22) .. (708.67,674.26) .. controls (708.67,674.26) and (417.92,676.96) .. (417.92,685.06);
\draw[line width=3.5bp,panel] (466.84,235.00) .. controls (466.84,243.10) and (712.17,245.80) .. (712.17,245.80) .. controls (712.17,372.97) and (712.17,542.54) .. (712.17,669.71) .. controls (712.17,669.71) and (383.20,672.41) .. (383.20,680.51);
\draw[->,line width=1.0bp,core] (466.84,235.00) .. controls (466.84,243.10) and (712.17,245.80) .. (712.17,245.80) .. controls (712.17,372.97) and (712.17,542.54) .. (712.17,669.71) .. controls (712.17,669.71) and (383.20,672.41) .. (383.20,680.51);
\filldraw[fill=slopebg,draw=slopefg,line width=1.0bp] (365.84,703.50) ellipse (84.00bp and 23.50bp);
\node[inner sep=0pt,font=\sffamily\fontsize{10.7bp}{12bp}\selectfont\bfseries,text=ink] at (365.84,710.00) {\hyperlink{node-W2L6_7}{WZ2 Lemma 6.7}};
\node[inner sep=0pt,font=\sffamily\fontsize{10.7bp}{12bp}\selectfont,text=ink] at (365.84,697.00) {\hyperlink{node-W2L6_7}{Small twisted image}};
\filldraw[fill=slopebg,draw=slopefg,line width=1.0bp] (264.84,625.00) ellipse (84.00bp and 23.50bp);
\node[inner sep=0pt,font=\sffamily\fontsize{10.7bp}{12bp}\selectfont\bfseries,text=ink] at (264.84,631.50) {\hyperlink{node-W2P6_5}{WZ2 Prop. 6.5 (1'')}};
\node[inner sep=0pt,font=\sffamily\fontsize{10.7bp}{12bp}\selectfont,text=ink] at (264.84,618.50) {\hyperlink{node-W2P6_5}{Nonsingular C2 slope}};
\filldraw[fill=slopebg,draw=slopefg,line width=1.0bp] (365.84,540.00) ellipse (89.19bp and 30.00bp);
\node[inner sep=0pt,font=\sffamily\fontsize{10.7bp}{12bp}\selectfont\bfseries,text=ink] at (365.84,553.00) {\hyperlink{node-L6_4}{WZ1 Lemma 6.4 / CW}};
\node[inner sep=0pt,font=\sffamily\fontsize{10.7bp}{12bp}\selectfont,text=ink] at (365.84,540.00) {\hyperlink{node-L6_4}{A thick slab with}};
\node[inner sep=0pt,font=\sffamily\fontsize{10.7bp}{12bp}\selectfont,text=ink] at (365.84,527.00) {\hyperlink{node-L6_4}{comparable nonzero}};
\filldraw[fill=slopebg,draw=slopefg,line width=1.0bp] (466.84,459.50) ellipse (84.00bp and 23.50bp);
\node[inner sep=0pt,font=\sffamily\fontsize{10.7bp}{12bp}\selectfont\bfseries,text=ink] at (466.84,466.00) {\hyperlink{node-W2L6_6}{WZ2 Lemma 6.6}};
\node[inner sep=0pt,font=\sffamily\fontsize{10.7bp}{12bp}\selectfont,text=ink] at (466.84,453.00) {\hyperlink{node-W2L6_6}{Popular slab is nonflat}};
\filldraw[fill=slopebg,draw=slopefg,line width=1.0bp] (264.84,379.00) ellipse (116.67bp and 30.00bp);
\node[inner sep=0pt,font=\sffamily\fontsize{10.7bp}{12bp}\selectfont\bfseries,text=ink] at (264.84,392.00) {\hyperlink{node-W2SlabCount}{WZ2 Lemma 6.6: new Step 4}};
\node[inner sep=0pt,font=\sffamily\fontsize{10.7bp}{12bp}\selectfont,text=ink] at (264.84,379.00) {\hyperlink{node-W2SlabCount}{Convex Wolff contradicts}};
\node[inner sep=0pt,font=\sffamily\fontsize{10.7bp}{12bp}\selectfont,text=ink] at (264.84,366.00) {\hyperlink{node-W2SlabCount}{a popular thin slab}};
\filldraw[fill=slopebg,draw=slopefg,line width=1.0bp] (163.84,292.00) ellipse (110.34bp and 30.00bp);
\node[inner sep=0pt,font=\sffamily\fontsize{10.7bp}{12bp}\selectfont\bfseries,text=ink] at (163.84,305.00) {\hyperlink{node-W1SlabPrep}{WZ1 Lemma 6.3: Steps 1-3}};
\node[inner sep=0pt,font=\sffamily\fontsize{10.7bp}{12bp}\selectfont,text=ink] at (163.84,292.00) {\hyperlink{node-W1SlabPrep}{Popular grains in a thin}};
\node[inner sep=0pt,font=\sffamily\fontsize{10.7bp}{12bp}\selectfont,text=ink] at (163.84,279.00) {\hyperlink{node-W1SlabPrep}{horizontal slab}};
\filldraw[fill=slopebg,draw=slopefg,line width=1.0bp] (392.17,292.00) ellipse (84.00bp and 30.00bp);
\node[inner sep=0pt,font=\sffamily\fontsize{10.7bp}{12bp}\selectfont\bfseries,text=ink] at (392.17,305.00) {\hyperlink{node-L6_2}{WZ1 Lemma 6.2}};
\node[inner sep=0pt,font=\sffamily\fontsize{10.7bp}{12bp}\selectfont,text=ink] at (392.17,292.00) {\hyperlink{node-L6_2}{A thin projection forces}};
\node[inner sep=0pt,font=\sffamily\fontsize{10.7bp}{12bp}\selectfont,text=ink] at (392.17,279.00) {\hyperlink{node-L6_2}{near orthogonality}};
\filldraw[fill=regularitybg,draw=regularityfg,line width=1.0bp] (264.84,207.00) ellipse (84.00bp and 23.50bp);
\node[inner sep=0pt,font=\sffamily\fontsize{10.7bp}{12bp}\selectfont\bfseries,text=ink] at (264.84,213.50) {\hyperlink{node-W2P6_4}{WZ2 Proposition 6.4}};
\node[inner sep=0pt,font=\sffamily\fontsize{10.7bp}{12bp}\selectfont,text=ink] at (264.84,200.50) {\hyperlink{node-W2P6_4}{C2 grains}};
\filldraw[fill=setupbg,draw=setupfg,line width=1.0bp] (264.84,459.50) ellipse (84.00bp and 23.50bp);
\node[inner sep=0pt,font=\sffamily\fontsize{10.7bp}{12bp}\selectfont\bfseries,text=ink] at (264.84,466.00) {\hyperlink{node-W2P6_2}{WZ2 Proposition 6.2}};
\node[inner sep=0pt,font=\sffamily\fontsize{10.7bp}{12bp}\selectfont,text=ink] at (264.84,453.00) {\hyperlink{node-W2P6_2}{Coarse/fine extremality}};
\filldraw[fill=setupbg,draw=setupfg,line width=1.0bp] (365.84,122.00) ellipse (89.19bp and 30.00bp);
\node[inner sep=0pt,font=\sffamily\fontsize{10.7bp}{12bp}\selectfont\bfseries,text=ink] at (365.84,135.00) {\hyperlink{node-L3_5}{WZ1 Lemma 3.5 / CW}};
\node[inner sep=0pt,font=\sffamily\fontsize{10.7bp}{12bp}\selectfont,text=ink] at (365.84,122.00) {\hyperlink{node-L3_5}{Mild anisotropic}};
\node[inner sep=0pt,font=\sffamily\fontsize{10.7bp}{12bp}\selectfont,text=ink] at (365.84,109.00) {\hyperlink{node-L3_5}{rescaling preserves}};
\filldraw[fill=setupbg,draw=setupfg,line width=1.0bp,rounded corners=5bp] (510.17,264.00) rectangle (678.17,320.00);
\node[inner sep=0pt,font=\sffamily\fontsize{10.7bp}{12bp}\selectfont\bfseries,text=ink] at (594.17,305.00) {\hyperlink{node-W2CWA}{WZ2 Def. 2.5}};
\node[inner sep=0pt,font=\sffamily\fontsize{10.7bp}{12bp}\selectfont,text=ink] at (594.17,292.00) {\hyperlink{node-W2CWA}{Normalized}};
\node[inner sep=0pt,font=\sffamily\fontsize{10.7bp}{12bp}\selectfont,text=ink] at (594.17,279.00) {\hyperlink{node-W2CWA}{convex-containment count}};
\filldraw[fill=setupbg,draw=setupfg,line width=1.0bp,rounded corners=5bp] (281.84,22.00) rectangle (449.84,65.00);
\node[inner sep=0pt,font=\sffamily\fontsize{10.7bp}{12bp}\selectfont\bfseries,text=ink] at (365.84,50.00) {\hyperlink{node-W2D6_1}{WZ2 Definition 6.1}};
\node[inner sep=0pt,font=\sffamily\fontsize{10.7bp}{12bp}\selectfont,text=ink] at (365.84,37.00) {\hyperlink{node-W2D6_1}{CW-extremal pair}};
\filldraw[fill=slopebg,draw=slopefg,line width=1.0bp,rounded corners=5bp] (382.84,597.00) rectangle (550.84,653.00);
\node[inner sep=0pt,font=\sffamily\fontsize{10.7bp}{12bp}\selectfont\bfseries,text=ink] at (466.84,638.00) {\hyperlink{node-Dtwist}{WZ1 Def. 6.1 / WZ2 Sec.}};
\node[inner sep=0pt,font=\sffamily\fontsize{10.7bp}{12bp}\selectfont,text=ink] at (466.84,625.00) {\hyperlink{node-Dtwist}{6.5}};
\node[inner sep=0pt,font=\sffamily\fontsize{10.7bp}{12bp}\selectfont,text=ink] at (466.84,612.00) {\hyperlink{node-Dtwist}{Twisted projection}};
\filldraw[fill=apibg,draw=apifg,line width=1.0bp,rounded corners=5bp] (415.51,351.00) rectangle (583.51,407.00);
\node[inner sep=0pt,font=\sffamily\fontsize{10.7bp}{12bp}\selectfont\bfseries,text=ink] at (499.51,392.00) {\hyperlink{node-Icalc}{Calculus API}};
\node[inner sep=0pt,font=\sffamily\fontsize{10.7bp}{12bp}\selectfont,text=ink] at (499.51,379.00) {\hyperlink{node-Icalc}{C2 bounds and compactly}};
\node[inner sep=0pt,font=\sffamily\fontsize{10.7bp}{12bp}\selectfont,text=ink] at (499.51,366.00) {\hyperlink{node-Icalc}{supported patches}};
\filldraw[fill=grainsbg,draw=grainsfg,line width=1.0bp,rounded corners=5bp] (382.84,179.00) rectangle (550.84,235.00);
\node[inner sep=0pt,font=\sffamily\fontsize{10.7bp}{12bp}\selectfont\bfseries,text=ink] at (466.84,220.00) {\hyperlink{node-DAD}{WZ2 Sec. 6.1: ADset}};
\node[inner sep=0pt,font=\sffamily\fontsize{10.7bp}{12bp}\selectfont,text=ink] at (466.84,207.00) {\hyperlink{node-DAD}{Upper AD covering}};
\node[inner sep=0pt,font=\sffamily\fontsize{10.7bp}{12bp}\selectfont,text=ink] at (466.84,194.00) {\hyperlink{node-DAD}{condition}};
\end{tikzpicture}}\end{center}
{\footnotesize Solid: proof input. Dashed: API. Dash-dot: assembly. Long-dashed: overview path. Read from the top down; arrows point to consumers. Full prerequisites appear on the cards and in the ledger.}\par
\sectitle{Twisted projection and scale iteration}
\hypertarget{view-projection}{}
{\sffamily\small\color{muted}Section 7 proves a lower bound using only Tube Wolff; the two area bounds close the contradiction.}\par
\begin{center}\resizebox{493.00bp}{!}{\begin{tikzpicture}[x=1bp,y=1bp, >={Stealth[length=4bp,width=3.2bp]}]
\path[use as bounding box] (0,0) rectangle (732.75,753.00);
\fill[panel,rounded corners=5bp] (0,0) rectangle (732.75,753.00);
\draw[line width=3.5bp,panel] (312.54,165.00) .. controls (312.54,173.10) and (30.00,175.80) .. (30.00,175.80) .. controls (30.00,282.14) and (30.00,423.92) .. (30.00,530.26) .. controls (30.00,530.26) and (315.70,532.96) .. (315.70,541.06);
\draw[->,line width=1.0bp,core] (312.54,165.00) .. controls (312.54,173.10) and (30.00,175.80) .. (30.00,175.80) .. controls (30.00,282.14) and (30.00,423.92) .. (30.00,530.26) .. controls (30.00,530.26) and (315.70,532.96) .. (315.70,541.06);
\draw[line width=3.5bp,panel] (263.62,583.00) .. controls (263.62,595.42) and (364.62,597.58) .. (364.62,610.00);
\draw[->,line width=1.0bp,core] (263.62,583.00) .. controls (263.62,595.42) and (364.62,597.58) .. (364.62,610.00);
\draw[line width=3.5bp,panel] (211.54,80.00) .. controls (211.54,88.10) and (26.50,90.80) .. (26.50,90.80) .. controls (26.50,222.64) and (26.50,398.42) .. (26.50,530.26) .. controls (26.50,530.26) and (211.54,532.96) .. (211.54,541.06);
\draw[->,line width=1.0bp,api,dash pattern=on 5bp off 4bp] (211.54,80.00) .. controls (211.54,88.10) and (26.50,90.80) .. (26.50,90.80) .. controls (26.50,222.64) and (26.50,398.42) .. (26.50,530.26) .. controls (26.50,530.26) and (211.54,532.96) .. (211.54,541.06);
\draw[line width=3.5bp,panel] (315.70,80.00) .. controls (315.70,88.10) and (699.25,90.80) .. (699.25,90.80) .. controls (699.25,120.47) and (699.25,160.03) .. (699.25,189.70) .. controls (699.25,189.70) and (595.25,192.40) .. (595.25,200.50);
\draw[->,line width=1.0bp,core] (315.70,80.00) .. controls (315.70,88.10) and (699.25,90.80) .. (699.25,90.80) .. controls (699.25,120.47) and (699.25,160.03) .. (699.25,189.70) .. controls (699.25,189.70) and (595.25,192.40) .. (595.25,200.50);
\draw[line width=3.5bp,panel] (413.54,75.54) .. controls (413.54,83.64) and (702.75,86.34) .. (702.75,86.34) .. controls (702.75,114.80) and (702.75,152.74) .. (702.75,181.20) .. controls (702.75,181.20) and (393.25,183.90) .. (393.25,192.00);
\draw[->,line width=1.0bp,core] (413.54,75.54) .. controls (413.54,83.64) and (702.75,86.34) .. (702.75,86.34) .. controls (702.75,114.80) and (702.75,152.74) .. (702.75,181.20) .. controls (702.75,181.20) and (393.25,183.90) .. (393.25,192.00);
\draw[line width=3.5bp,panel] (416.70,165.00) .. controls (416.70,181.04) and (341.17,183.83) .. (341.17,199.86);
\draw[->,line width=1.0bp,api,dash pattern=on 5bp off 4bp] (416.70,165.00) .. controls (416.70,181.04) and (341.17,183.83) .. (341.17,199.86);
\draw[line width=3.5bp,panel] (566.62,165.00) .. controls (566.62,181.04) and (445.33,183.83) .. (445.33,199.86);
\draw[->,line width=1.0bp,api,dash pattern=on 5bp off 4bp] (566.62,165.00) .. controls (566.62,181.04) and (445.33,183.83) .. (445.33,199.86);
\draw[line width=3.5bp,panel] (514.54,165.00) .. controls (514.54,181.04) and (232.45,183.83) .. (232.45,199.86);
\draw[->,line width=1.0bp,api,dash pattern=on 5bp off 4bp] (514.54,165.00) .. controls (514.54,181.04) and (232.45,183.83) .. (232.45,199.86);
\draw[line width=3.5bp,panel] (110.54,165.00) .. controls (110.54,181.04) and (92.80,183.83) .. (92.80,199.86);
\draw[->,line width=1.0bp,api,dash pattern=on 5bp off 4bp] (110.54,165.00) .. controls (110.54,181.04) and (92.80,183.83) .. (92.80,199.86);
\draw[line width=3.5bp,panel] (162.62,265.00) .. controls (162.62,281.82) and (211.54,284.74) .. (211.54,301.56);
\draw[->,line width=1.0bp,core] (162.62,265.00) .. controls (162.62,281.82) and (211.54,284.74) .. (211.54,301.56);
\draw[line width=3.5bp,panel] (393.25,265.00) .. controls (393.25,279.49) and (263.62,282.01) .. (263.62,296.50);
\draw[->,line width=1.0bp,core] (393.25,265.00) .. controls (393.25,279.49) and (263.62,282.01) .. (263.62,296.50);
\draw[line width=3.5bp,panel] (543.17,256.50) .. controls (543.17,277.23) and (315.70,280.83) .. (315.70,301.56);
\draw[->,line width=1.0bp,api,dash pattern=on 5bp off 4bp] (543.17,256.50) .. controls (543.17,277.23) and (315.70,280.83) .. (315.70,301.56);
\draw[line width=3.5bp,panel] (517.70,75.54) .. controls (517.70,83.64) and (706.25,86.34) .. (706.25,86.34) .. controls (706.25,144.80) and (706.25,222.74) .. (706.25,281.20) .. controls (706.25,281.20) and (413.54,283.90) .. (413.54,292.00);
\draw[->,line width=1.0bp,core] (517.70,75.54) .. controls (517.70,83.64) and (706.25,86.34) .. (706.25,86.34) .. controls (706.25,144.80) and (706.25,222.74) .. (706.25,281.20) .. controls (706.25,281.20) and (413.54,283.90) .. (413.54,292.00);
\draw[line width=3.5bp,panel] (595.25,256.50) .. controls (595.25,272.83) and (517.70,275.67) .. (517.70,292.00);
\draw[->,line width=1.0bp,core] (595.25,256.50) .. controls (595.25,272.83) and (517.70,275.67) .. (517.70,292.00);
\draw[line width=3.5bp,panel] (618.70,165.00) .. controls (618.70,173.10) and (709.75,175.80) .. (709.75,175.80) .. controls (709.75,207.42) and (709.75,249.58) .. (709.75,281.20) .. controls (709.75,281.20) and (465.62,283.90) .. (465.62,292.00);
\draw[->,line width=1.0bp,api,dash pattern=on 5bp off 4bp] (618.70,165.00) .. controls (618.70,173.10) and (709.75,175.80) .. (709.75,175.80) .. controls (709.75,207.42) and (709.75,249.58) .. (709.75,281.20) .. controls (709.75,281.20) and (465.62,283.90) .. (465.62,292.00);
\draw[line width=3.5bp,panel] (263.62,343.50) .. controls (263.62,357.99) and (364.62,360.51) .. (364.62,375.00);
\draw[->,line width=1.0bp,core] (263.62,343.50) .. controls (263.62,357.99) and (364.62,360.51) .. (364.62,375.00);
\draw[line width=3.5bp,panel] (465.62,348.00) .. controls (465.62,363.39) and (426.78,366.07) .. (426.78,381.46);
\draw[->,line width=1.0bp,core] (465.62,348.00) .. controls (465.62,363.39) and (426.78,366.07) .. (426.78,381.46);
\draw[line width=3.5bp,panel] (162.62,165.00) .. controls (162.62,173.10) and (23.00,175.80) .. (23.00,175.80) .. controls (23.00,234.26) and (23.00,312.20) .. (23.00,370.66) .. controls (23.00,370.66) and (302.46,373.36) .. (302.46,381.46);
\draw[->,line width=1.0bp,api,dash pattern=on 5bp off 4bp] (162.62,165.00) .. controls (162.62,173.10) and (23.00,175.80) .. (23.00,175.80) .. controls (23.00,234.26) and (23.00,312.20) .. (23.00,370.66) .. controls (23.00,370.66) and (302.46,373.36) .. (302.46,381.46);
\draw[line width=3.5bp,panel] (364.62,435.00) .. controls (364.62,447.42) and (364.62,449.58) .. (364.62,462.00);
\draw[->,line width=1.0bp,core] (364.62,435.00) .. controls (364.62,447.42) and (364.62,449.58) .. (364.62,462.00);
\draw[line width=3.5bp,panel] (364.62,509.00) .. controls (364.62,523.75) and (413.54,526.31) .. (413.54,541.06);
\draw[->,line width=1.0bp,core] (364.62,509.00) .. controls (364.62,523.75) and (413.54,526.31) .. (413.54,541.06);
\draw[line width=3.5bp,panel] (647.33,256.50) .. controls (647.33,264.60) and (713.25,267.30) .. (713.25,267.30) .. controls (713.25,346.19) and (713.25,451.37) .. (713.25,530.26) .. controls (713.25,530.26) and (517.70,532.96) .. (517.70,541.06);
\draw[->,line width=1.0bp,core] (647.33,256.50) .. controls (647.33,264.60) and (713.25,267.30) .. (713.25,267.30) .. controls (713.25,346.19) and (713.25,451.37) .. (713.25,530.26) .. controls (713.25,530.26) and (517.70,532.96) .. (517.70,541.06);
\draw[line width=3.5bp,panel] (465.62,583.00) .. controls (465.62,597.75) and (416.70,600.31) .. (416.70,615.06);
\draw[->,line width=1.0bp,core] (465.62,583.00) .. controls (465.62,597.75) and (416.70,600.31) .. (416.70,615.06);
\draw[line width=3.5bp,panel] (214.70,165.00) .. controls (214.70,173.10) and (19.50,175.80) .. (19.50,175.80) .. controls (19.50,304.34) and (19.50,475.72) .. (19.50,604.26) .. controls (19.50,604.26) and (312.54,606.96) .. (312.54,615.06);
\draw[->,line width=1.0bp,api,dash pattern=on 5bp off 4bp] (214.70,165.00) .. controls (214.70,173.10) and (19.50,175.80) .. (19.50,175.80) .. controls (19.50,304.34) and (19.50,475.72) .. (19.50,604.26) .. controls (19.50,604.26) and (312.54,606.96) .. (312.54,615.06);
\draw[line width=3.5bp,panel] (364.62,657.00) .. controls (364.62,669.42) and (364.62,671.58) .. (364.62,684.00);
\draw[->,line width=1.0bp,core] (364.62,657.00) .. controls (364.62,669.42) and (364.62,671.58) .. (364.62,684.00);
\filldraw[fill=setupbg,draw=setupfg,line width=1.0bp] (364.62,707.50) ellipse (84.00bp and 23.50bp);
\node[inner sep=0pt,font=\sffamily\fontsize{10.7bp}{12bp}\selectfont\bfseries,text=ink] at (364.62,714.00) {\hyperlink{node-W2T5_2}{WZ2 Theorem 5.2}};
\node[inner sep=0pt,font=\sffamily\fontsize{10.7bp}{12bp}\selectfont,text=ink] at (364.62,701.00) {\hyperlink{node-W2T5_2}{Multi-scale CW union}};
\filldraw[fill=setupbg,draw=setupfg,line width=1.0bp] (364.62,633.50) ellipse (84.00bp and 23.50bp);
\node[inner sep=0pt,font=\sffamily\fontsize{10.7bp}{12bp}\selectfont\bfseries,text=ink] at (364.62,640.00) {\hyperlink{node-W2Zero}{WZ2 Contradiction}};
\node[inner sep=0pt,font=\sffamily\fontsize{10.7bp}{12bp}\selectfont,text=ink] at (364.62,627.00) {\hyperlink{node-W2Zero}{sigma\_CW = 0}};
\filldraw[fill=cinematicbg,draw=cinematicfg,line width=1.0bp] (465.62,559.50) ellipse (84.00bp and 23.50bp);
\node[inner sep=0pt,font=\sffamily\fontsize{10.7bp}{12bp}\selectfont\bfseries,text=ink] at (465.62,566.00) {\hyperlink{node-W2P6_8}{WZ2 Proposition 6.8}};
\node[inner sep=0pt,font=\sffamily\fontsize{10.7bp}{12bp}\selectfont,text=ink] at (465.62,553.00) {\hyperlink{node-W2P6_8}{Convex CW projection}};
\filldraw[fill=cinematicbg,draw=cinematicfg,line width=1.0bp] (364.62,485.50) ellipse (84.00bp and 23.50bp);
\node[inner sep=0pt,font=\sffamily\fontsize{10.7bp}{12bp}\selectfont\bfseries,text=ink] at (364.62,492.00) {\hyperlink{node-W2P7_1}{WZ2 Proposition 7.1}};
\node[inner sep=0pt,font=\sffamily\fontsize{10.7bp}{12bp}\selectfont,text=ink] at (364.62,479.00) {\hyperlink{node-W2P7_1}{Tube-Wolff projection}};
\filldraw[fill=slopebg,draw=slopefg,line width=1.0bp] (263.62,559.50) ellipse (84.00bp and 23.50bp);
\node[inner sep=0pt,font=\sffamily\fontsize{10.7bp}{12bp}\selectfont\bfseries,text=ink] at (263.62,566.00) {\hyperlink{node-W2L6_7}{WZ2 Lemma 6.7}};
\node[inner sep=0pt,font=\sffamily\fontsize{10.7bp}{12bp}\selectfont,text=ink] at (263.62,553.00) {\hyperlink{node-W2L6_7}{Small twisted image}};
\filldraw[fill=cinematicbg,draw=cinematicfg,line width=1.0bp] (364.62,405.00) ellipse (100.26bp and 30.00bp);
\node[inner sep=0pt,font=\sffamily\fontsize{10.7bp}{12bp}\selectfont\bfseries,text=ink] at (364.62,418.00) {\hyperlink{node-W2Iterate}{WZ2 Prop. 7.1: iteration}};
\node[inner sep=0pt,font=\sffamily\fontsize{10.7bp}{12bp}\selectfont,text=ink] at (364.62,405.00) {\hyperlink{node-W2Iterate}{Telescope the ratios up}};
\node[inner sep=0pt,font=\sffamily\fontsize{10.7bp}{12bp}\selectfont,text=ink] at (364.62,392.00) {\hyperlink{node-W2Iterate}{to a coarse scale}};
\filldraw[fill=cinematicbg,draw=cinematicfg,line width=1.0bp] (263.62,320.00) ellipse (84.00bp and 23.50bp);
\node[inner sep=0pt,font=\sffamily\fontsize{10.7bp}{12bp}\selectfont\bfseries,text=ink] at (263.62,326.50) {\hyperlink{node-W2L7_12}{WZ2 Lemma 7.12}};
\node[inner sep=0pt,font=\sffamily\fontsize{10.7bp}{12bp}\selectfont,text=ink] at (263.62,313.50) {\hyperlink{node-W2L7_12}{One-step area ratio}};
\filldraw[fill=cinematicbg,draw=cinematicfg,line width=1.0bp,rounded corners=5bp] (381.62,292.00) rectangle (549.62,348.00);
\node[inner sep=0pt,font=\sffamily\fontsize{10.7bp}{12bp}\selectfont\bfseries,text=ink] at (465.62,333.00) {\hyperlink{node-IcoarseTW}{WZ2 proof operation}};
\node[inner sep=0pt,font=\sffamily\fontsize{10.7bp}{12bp}\selectfont,text=ink] at (465.62,320.00) {\hyperlink{node-IcoarseTW}{Uniform Tube-Wolff}};
\node[inner sep=0pt,font=\sffamily\fontsize{10.7bp}{12bp}\selectfont,text=ink] at (465.62,307.00) {\hyperlink{node-IcoarseTW}{coarsening}};
\filldraw[fill=cinematicbg,draw=cinematicfg,line width=1.0bp] (393.25,228.50) ellipse (84.00bp and 36.50bp);
\node[inner sep=0pt,font=\sffamily\fontsize{10.7bp}{12bp}\selectfont\bfseries,text=ink] at (393.25,248.00) {\hyperlink{node-W2ImageUniform}{WZ2 Lemma 7.12:}};
\node[inner sep=0pt,font=\sffamily\fontsize{10.7bp}{12bp}\selectfont,text=ink] at (393.25,235.00) {\hyperlink{node-W2ImageUniform}{preparation}};
\node[inner sep=0pt,font=\sffamily\fontsize{10.7bp}{12bp}\selectfont,text=ink] at (393.25,222.00) {\hyperlink{node-W2ImageUniform}{Uniformize the projected}};
\node[inner sep=0pt,font=\sffamily\fontsize{10.7bp}{12bp}\selectfont,text=ink] at (393.25,209.00) {\hyperlink{node-W2ImageUniform}{incidence set}};
\filldraw[fill=cinematicbg,draw=cinematicfg,line width=1.0bp] (162.62,228.50) ellipse (112.62bp and 36.50bp);
\node[inner sep=0pt,font=\sffamily\fontsize{10.7bp}{12bp}\selectfont\bfseries,text=ink] at (162.62,248.00) {\hyperlink{node-W2AreaCompare}{WZ2 Lemma 7.12: Holder}};
\node[inner sep=0pt,font=\sffamily\fontsize{10.7bp}{12bp}\selectfont,text=ink] at (162.62,235.00) {\hyperlink{node-W2AreaCompare}{and translations}};
\node[inner sep=0pt,font=\sffamily\fontsize{10.7bp}{12bp}\selectfont,text=ink] at (162.62,222.00) {\hyperlink{node-W2AreaCompare}{From cinematic norm to a}};
\node[inner sep=0pt,font=\sffamily\fontsize{10.7bp}{12bp}\selectfont,text=ink] at (162.62,209.00) {\hyperlink{node-W2AreaCompare}{local area ratio}};
\filldraw[fill=cinematicbg,draw=cinematicfg,line width=1.0bp] (465.62,52.00) ellipse (84.00bp and 30.00bp);
\node[inner sep=0pt,font=\sffamily\fontsize{10.7bp}{12bp}\selectfont\bfseries,text=ink] at (465.62,65.00) {\hyperlink{node-W2L7_9}{WZ2 Lemma 7.9}};
\node[inner sep=0pt,font=\sffamily\fontsize{10.7bp}{12bp}\selectfont,text=ink] at (465.62,52.00) {\hyperlink{node-W2L7_9}{Sparse-scale uniform}};
\node[inner sep=0pt,font=\sffamily\fontsize{10.7bp}{12bp}\selectfont,text=ink] at (465.62,39.00) {\hyperlink{node-W2L7_9}{refinement}};
\filldraw[fill=cinematicbg,draw=cinematicfg,line width=1.0bp,rounded corners=5bp] (511.25,200.50) rectangle (679.25,256.50);
\node[inner sep=0pt,font=\sffamily\fontsize{10.7bp}{12bp}\selectfont\bfseries,text=ink] at (595.25,241.50) {\hyperlink{node-W2TW}{WZ2 Sec. 7: Tube Wolff}};
\node[inner sep=0pt,font=\sffamily\fontsize{10.7bp}{12bp}\selectfont,text=ink] at (595.25,228.50) {\hyperlink{node-W2TW}{The weaker radius-squared}};
\node[inner sep=0pt,font=\sffamily\fontsize{10.7bp}{12bp}\selectfont,text=ink] at (595.25,215.50) {\hyperlink{node-W2TW}{tube count}};
\filldraw[fill=setupbg,draw=setupfg,line width=1.0bp,rounded corners=5bp] (179.62,24.00) rectangle (347.62,80.00);
\node[inner sep=0pt,font=\sffamily\fontsize{10.7bp}{12bp}\selectfont\bfseries,text=ink] at (263.62,65.00) {\hyperlink{node-W2CWA}{WZ2 Def. 2.5}};
\node[inner sep=0pt,font=\sffamily\fontsize{10.7bp}{12bp}\selectfont,text=ink] at (263.62,52.00) {\hyperlink{node-W2CWA}{Normalized}};
\node[inner sep=0pt,font=\sffamily\fontsize{10.7bp}{12bp}\selectfont,text=ink] at (263.62,39.00) {\hyperlink{node-W2CWA}{convex-containment count}};
\filldraw[fill=apibg,draw=apifg,line width=1.0bp,rounded corners=5bp] (78.62,109.00) rectangle (246.62,165.00);
\node[inner sep=0pt,font=\sffamily\fontsize{10.7bp}{12bp}\selectfont\bfseries,text=ink] at (162.62,150.00) {\hyperlink{node-Iparams}{Quantifier API}};
\node[inner sep=0pt,font=\sffamily\fontsize{10.7bp}{12bp}\selectfont,text=ink] at (162.62,137.00) {\hyperlink{node-Iparams}{Small-parameter and}};
\node[inner sep=0pt,font=\sffamily\fontsize{10.7bp}{12bp}\selectfont,text=ink] at (162.62,124.00) {\hyperlink{node-Iparams}{finite-iteration}};
\filldraw[fill=apibg,draw=apifg,line width=1.0bp,rounded corners=5bp] (482.62,109.00) rectangle (650.62,165.00);
\node[inner sep=0pt,font=\sffamily\fontsize{10.7bp}{12bp}\selectfont\bfseries,text=ink] at (566.62,150.00) {\hyperlink{node-Imass}{Refinement API}};
\node[inner sep=0pt,font=\sffamily\fontsize{10.7bp}{12bp}\selectfont,text=ink] at (566.62,137.00) {\hyperlink{node-Imass}{Mass and constant}};
\node[inner sep=0pt,font=\sffamily\fontsize{10.7bp}{12bp}\selectfont,text=ink] at (566.62,124.00) {\hyperlink{node-Imass}{multiplicity}};
\filldraw[fill=slopebg,draw=slopefg,line width=1.0bp,rounded corners=5bp] (280.62,109.00) rectangle (448.62,165.00);
\node[inner sep=0pt,font=\sffamily\fontsize{10.7bp}{12bp}\selectfont\bfseries,text=ink] at (364.62,150.00) {\hyperlink{node-Dtwist}{WZ1 Def. 6.1 / WZ2 Sec.}};
\node[inner sep=0pt,font=\sffamily\fontsize{10.7bp}{12bp}\selectfont,text=ink] at (364.62,137.00) {\hyperlink{node-Dtwist}{6.5}};
\node[inner sep=0pt,font=\sffamily\fontsize{10.7bp}{12bp}\selectfont,text=ink] at (364.62,124.00) {\hyperlink{node-Dtwist}{Twisted projection}};
\end{tikzpicture}}\end{center}
{\footnotesize Solid: proof input. Dashed: API. Dash-dot: assembly. Long-dashed: overview path. Read from the top down; arrows point to consumers. Full prerequisites appear on the cards and in the ledger.}\par
\sectitle{Parameter extraction and cinematic comparison}
\hypertarget{view-cinematic}{}
{\sffamily\small\color{muted}Frostman in R4 is converted to an absolute one-dimensional count before applying the external maximal theorem.}\par
\begin{center}\resizebox{493.00bp}{!}{\begin{tikzpicture}[x=1bp,y=1bp, >={Stealth[length=4bp,width=3.2bp]}]
\path[use as bounding box] (0,0) rectangle (805.83,866.00);
\fill[panel,rounded corners=5bp] (0,0) rectangle (805.83,866.00);
\draw[line width=3.5bp,panel] (145.43,267.50) .. controls (145.43,275.60) and (51.00,278.30) .. (51.00,278.30) .. controls (51.00,376.09) and (51.00,506.47) .. (51.00,604.26) .. controls (51.00,604.26) and (106.30,606.96) .. (106.30,615.06);
\draw[->,line width=1.0bp,core] (145.43,267.50) .. controls (145.43,275.60) and (51.00,278.30) .. (51.00,278.30) .. controls (51.00,376.09) and (51.00,506.47) .. (51.00,604.26) .. controls (51.00,604.26) and (106.30,606.96) .. (106.30,615.06);
\draw[line width=3.5bp,panel] (199.16,470.00) .. controls (199.16,478.10) and (47.50,480.80) .. (47.50,480.80) .. controls (47.50,517.84) and (47.50,567.22) .. (47.50,604.26) .. controls (47.50,604.26) and (210.46,606.96) .. (210.46,615.06);
\draw[->,line width=1.0bp,core] (199.16,470.00) .. controls (199.16,478.10) and (47.50,480.80) .. (47.50,480.80) .. controls (47.50,517.84) and (47.50,567.22) .. (47.50,604.26) .. controls (47.50,604.26) and (210.46,606.96) .. (210.46,615.06);
\draw[line width=3.5bp,panel] (249.59,267.50) .. controls (249.59,275.60) and (44.00,278.30) .. (44.00,278.30) .. controls (44.00,313.22) and (44.00,359.78) .. (44.00,394.70) .. controls (44.00,394.70) and (401.16,397.40) .. (401.16,405.50);
\draw[->,line width=1.0bp,core] (249.59,267.50) .. controls (249.59,275.60) and (44.00,278.30) .. (44.00,278.30) .. controls (44.00,313.22) and (44.00,359.78) .. (44.00,394.70) .. controls (44.00,394.70) and (401.16,397.40) .. (401.16,405.50);
\draw[line width=3.5bp,panel] (147.08,176.00) .. controls (147.08,184.10) and (40.50,186.80) .. (40.50,186.80) .. controls (40.50,220.37) and (40.50,265.13) .. (40.50,298.70) .. controls (40.50,298.70) and (398.40,301.40) .. (398.40,309.50);
\draw[->,line width=1.0bp,core] (147.08,176.00) .. controls (147.08,184.10) and (40.50,186.80) .. (40.50,186.80) .. controls (40.50,220.37) and (40.50,265.13) .. (40.50,298.70) .. controls (40.50,298.70) and (398.40,301.40) .. (398.40,309.50);
\draw[line width=3.5bp,panel] (450.48,365.50) .. controls (450.48,373.60) and (751.33,376.30) .. (751.33,376.30) .. controls (751.33,413.16) and (751.33,462.30) .. (751.33,499.16) .. controls (751.33,499.16) and (595.25,501.86) .. (595.25,509.96);
\draw[->,line width=1.0bp,core] (450.48,365.50) .. controls (450.48,373.60) and (751.33,376.30) .. (751.33,376.30) .. controls (751.33,413.16) and (751.33,462.30) .. (751.33,499.16) .. controls (751.33,499.16) and (595.25,501.86) .. (595.25,509.96);
\draw[line width=3.5bp,panel] (655.24,176.00) .. controls (655.24,184.10) and (754.83,186.80) .. (754.83,186.80) .. controls (754.83,278.57) and (754.83,400.93) .. (754.83,492.70) .. controls (754.83,492.70) and (647.33,495.40) .. (647.33,503.50);
\draw[->,line width=1.0bp,api,dash pattern=on 5bp off 4bp] (655.24,176.00) .. controls (655.24,184.10) and (754.83,186.80) .. (754.83,186.80) .. controls (754.83,278.57) and (754.83,400.93) .. (754.83,492.70) .. controls (754.83,492.70) and (647.33,495.40) .. (647.33,503.50);
\draw[line width=3.5bp,panel] (603.16,372.00) .. controls (603.16,380.10) and (758.33,382.80) .. (758.33,382.80) .. controls (758.33,417.71) and (758.33,464.25) .. (758.33,499.16) .. controls (758.33,499.16) and (699.41,501.86) .. (699.41,509.96);
\draw[->,line width=1.0bp,core] (603.16,372.00) .. controls (603.16,380.10) and (758.33,382.80) .. (758.33,382.80) .. controls (758.33,417.71) and (758.33,464.25) .. (758.33,499.16) .. controls (758.33,499.16) and (699.41,501.86) .. (699.41,509.96);
\draw[line width=3.5bp,panel] (450.08,84.50) .. controls (450.08,102.88) and (453.24,106.08) .. (453.24,124.46);
\draw[->,line width=1.0bp,api,dash pattern=on 5bp off 4bp] (450.08,84.50) .. controls (450.08,102.88) and (453.24,106.08) .. (453.24,124.46);
\draw[line width=3.5bp,panel] (300.16,91.00) .. controls (300.16,106.39) and (349.08,109.07) .. (349.08,124.46);
\draw[->,line width=1.0bp,core] (300.16,91.00) .. controls (300.16,106.39) and (349.08,109.07) .. (349.08,124.46);
\draw[line width=3.5bp,panel] (554.24,84.50) .. controls (554.24,92.60) and (761.83,95.30) .. (761.83,95.30) .. controls (761.83,214.71) and (761.83,373.93) .. (761.83,493.35) .. controls (761.83,493.35) and (462.69,496.05) .. (462.69,504.15);
\draw[->,line width=1.0bp,core] (554.24,84.50) .. controls (554.24,92.60) and (761.83,95.30) .. (761.83,95.30) .. controls (761.83,214.71) and (761.83,373.93) .. (761.83,493.35) .. controls (761.83,493.35) and (462.69,496.05) .. (462.69,504.15);
\draw[line width=3.5bp,panel] (603.99,274.00) .. controls (603.99,282.10) and (765.33,284.80) .. (765.33,284.80) .. controls (765.33,349.11) and (765.33,434.85) .. (765.33,499.16) .. controls (765.33,499.16) and (497.41,501.86) .. (497.41,509.96);
\draw[->,line width=1.0bp,core] (603.99,274.00) .. controls (603.99,282.10) and (765.33,284.80) .. (765.33,284.80) .. controls (765.33,349.11) and (765.33,434.85) .. (765.33,499.16) .. controls (765.33,499.16) and (497.41,501.86) .. (497.41,509.96);
\draw[line width=3.5bp,panel] (400.34,274.00) .. controls (400.34,282.10) and (768.83,284.80) .. (768.83,284.80) .. controls (768.83,347.36) and (768.83,430.78) .. (768.83,493.35) .. controls (768.83,493.35) and (427.97,496.05) .. (427.97,504.15);
\draw[->,line width=1.0bp,core] (400.34,274.00) .. controls (400.34,282.10) and (768.83,284.80) .. (768.83,284.80) .. controls (768.83,347.36) and (768.83,430.78) .. (768.83,493.35) .. controls (768.83,493.35) and (427.97,496.05) .. (427.97,504.15);
\draw[line width=3.5bp,panel] (398.40,365.50) .. controls (398.40,373.60) and (772.33,376.30) .. (772.33,376.30) .. controls (772.33,413.16) and (772.33,462.30) .. (772.33,499.16) .. controls (772.33,499.16) and (393.25,501.86) .. (393.25,509.96);
\draw[->,line width=1.0bp,api,dash pattern=on 5bp off 4bp] (398.40,365.50) .. controls (398.40,373.60) and (772.33,376.30) .. (772.33,376.30) .. controls (772.33,413.16) and (772.33,462.30) .. (772.33,499.16) .. controls (772.33,499.16) and (393.25,501.86) .. (393.25,509.96);
\draw[line width=3.5bp,panel] (551.08,176.00) .. controls (551.08,184.10) and (775.83,186.80) .. (775.83,186.80) .. controls (775.83,249.17) and (775.83,332.33) .. (775.83,394.70) .. controls (775.83,394.70) and (603.16,397.40) .. (603.16,405.50);
\draw[->,line width=1.0bp,core] (551.08,176.00) .. controls (551.08,184.10) and (775.83,186.80) .. (775.83,186.80) .. controls (775.83,249.17) and (775.83,332.33) .. (775.83,394.70) .. controls (775.83,394.70) and (603.16,397.40) .. (603.16,405.50);
\draw[line width=3.5bp,panel] (349.08,171.54) .. controls (349.08,179.64) and (37.00,182.34) .. (37.00,182.34) .. controls (37.00,214.70) and (37.00,257.84) .. (37.00,290.20) .. controls (37.00,290.20) and (196.40,292.90) .. (196.40,301.00);
\draw[->,line width=1.0bp,core] (349.08,171.54) .. controls (349.08,179.64) and (37.00,182.34) .. (37.00,182.34) .. controls (37.00,214.70) and (37.00,257.84) .. (37.00,290.20) .. controls (37.00,290.20) and (196.40,292.90) .. (196.40,301.00);
\draw[line width=3.5bp,panel] (346.32,365.50) .. controls (346.32,373.60) and (33.50,376.30) .. (33.50,376.30) .. controls (33.50,409.27) and (33.50,453.23) .. (33.50,486.20) .. controls (33.50,486.20) and (199.16,488.90) .. (199.16,497.00);
\draw[->,line width=1.0bp,core] (346.32,365.50) .. controls (346.32,373.60) and (33.50,376.30) .. (33.50,376.30) .. controls (33.50,409.27) and (33.50,453.23) .. (33.50,486.20) .. controls (33.50,486.20) and (199.16,488.90) .. (199.16,497.00);
\draw[line width=3.5bp,panel] (453.24,171.54) .. controls (453.24,179.64) and (30.00,182.34) .. (30.00,182.34) .. controls (30.00,275.86) and (30.00,400.54) .. (30.00,494.06) .. controls (30.00,494.06) and (278.62,496.76) .. (278.62,504.86);
\draw[->,line width=1.0bp,core] (453.24,171.54) .. controls (453.24,179.64) and (30.00,182.34) .. (30.00,182.34) .. controls (30.00,275.86) and (30.00,400.54) .. (30.00,494.06) .. controls (30.00,494.06) and (278.62,496.76) .. (278.62,504.86);
\draw[line width=3.5bp,panel] (144.32,366.14) .. controls (144.32,374.24) and (26.50,376.94) .. (26.50,376.94) .. controls (26.50,412.08) and (26.50,458.92) .. (26.50,494.06) .. controls (26.50,494.06) and (119.70,496.76) .. (119.70,504.86);
\draw[->,line width=1.0bp,assembly,dash pattern=on 7bp off 3bp on 1bp off 3bp] (144.32,366.14) .. controls (144.32,374.24) and (26.50,376.94) .. (26.50,376.94) .. controls (26.50,412.08) and (26.50,458.92) .. (26.50,494.06) .. controls (26.50,494.06) and (119.70,496.76) .. (119.70,504.86);
\draw[line width=3.5bp,panel] (199.16,570.00) .. controls (199.16,586.04) and (525.79,588.83) .. (525.79,604.86);
\draw[->,line width=1.0bp,assembly,dash pattern=on 7bp off 3bp on 1bp off 3bp] (199.16,570.00) .. controls (199.16,586.04) and (525.79,588.83) .. (525.79,604.86);
\draw[line width=3.5bp,panel] (445.33,563.50) .. controls (445.33,579.04) and (587.69,581.74) .. (587.69,597.28);
\draw[->,line width=1.0bp,core] (445.33,563.50) .. controls (445.33,579.04) and (587.69,581.74) .. (587.69,597.28);
\draw[line width=3.5bp,panel] (647.33,563.50) .. controls (647.33,582.53) and (680.53,585.84) .. (680.53,604.86);
\draw[->,line width=1.0bp,core] (647.33,563.50) .. controls (647.33,582.53) and (680.53,585.84) .. (680.53,604.86);
\draw[line width=3.5bp,panel] (603.16,465.50) .. controls (603.16,473.60) and (779.33,476.30) .. (779.33,476.30) .. controls (779.33,509.35) and (779.33,553.43) .. (779.33,586.48) .. controls (779.33,586.48) and (618.64,589.18) .. (618.64,597.28);
\draw[->,line width=1.0bp,core] (603.16,465.50) .. controls (603.16,473.60) and (779.33,476.30) .. (779.33,476.30) .. controls (779.33,509.35) and (779.33,553.43) .. (779.33,586.48) .. controls (779.33,586.48) and (618.64,589.18) .. (618.64,597.28);
\draw[line width=3.5bp,panel] (401.16,465.50) .. controls (401.16,473.60) and (782.83,476.30) .. (782.83,476.30) .. controls (782.83,510.06) and (782.83,555.06) .. (782.83,588.82) .. controls (782.83,588.82) and (556.74,591.52) .. (556.74,599.62);
\draw[->,line width=1.0bp,core] (401.16,465.50) .. controls (401.16,473.60) and (782.83,476.30) .. (782.83,476.30) .. controls (782.83,510.06) and (782.83,555.06) .. (782.83,588.82) .. controls (782.83,588.82) and (556.74,591.52) .. (556.74,599.62);
\draw[line width=3.5bp,panel] (603.16,176.00) .. controls (603.16,184.10) and (786.33,186.80) .. (786.33,186.80) .. controls (786.33,307.41) and (786.33,468.21) .. (786.33,588.82) .. controls (786.33,588.82) and (649.58,591.52) .. (649.58,599.62);
\draw[->,line width=1.0bp,api,dash pattern=on 5bp off 4bp] (603.16,176.00) .. controls (603.16,184.10) and (786.33,186.80) .. (786.33,186.80) .. controls (786.33,307.41) and (786.33,468.21) .. (786.33,588.82) .. controls (786.33,588.82) and (649.58,591.52) .. (649.58,599.62);
\draw[line width=3.5bp,panel] (603.16,670.00) .. controls (603.16,686.04) and (470.99,688.83) .. (470.99,704.86);
\draw[->,line width=1.0bp,assembly,dash pattern=on 7bp off 3bp on 1bp off 3bp] (603.16,670.00) .. controls (603.16,686.04) and (470.99,688.83) .. (470.99,704.86);
\draw[line width=3.5bp,panel] (158.38,657.00) .. controls (158.38,679.02) and (331.34,682.85) .. (331.34,704.86);
\draw[->,line width=1.0bp,core] (158.38,657.00) .. controls (158.38,679.02) and (331.34,682.85) .. (331.34,704.86);
\draw[line width=3.5bp,panel] (360.38,661.50) .. controls (360.38,677.83) and (401.16,680.67) .. (401.16,697.00);
\draw[->,line width=1.0bp,api,dash pattern=on 5bp off 4bp] (360.38,661.50) .. controls (360.38,677.83) and (401.16,680.67) .. (401.16,697.00);
\draw[line width=3.5bp,panel] (401.16,770.00) .. controls (401.16,784.75) and (453.24,787.31) .. (453.24,802.06);
\draw[->,line width=1.0bp,core] (401.16,770.00) .. controls (401.16,784.75) and (453.24,787.31) .. (453.24,802.06);
\draw[line width=3.5bp,panel] (248.48,366.14) .. controls (248.48,374.24) and (23.00,376.94) .. (23.00,376.94) .. controls (23.00,501.24) and (23.00,666.96) .. (23.00,791.26) .. controls (23.00,791.26) and (349.08,793.96) .. (349.08,802.06);
\draw[->,line width=1.0bp,core] (248.48,366.14) .. controls (248.48,374.24) and (23.00,376.94) .. (23.00,376.94) .. controls (23.00,501.24) and (23.00,666.96) .. (23.00,791.26) .. controls (23.00,791.26) and (349.08,793.96) .. (349.08,802.06);
\draw[line width=3.5bp,panel] (251.24,176.00) .. controls (251.24,184.10) and (19.50,186.80) .. (19.50,186.80) .. controls (19.50,366.62) and (19.50,606.38) .. (19.50,786.20) .. controls (19.50,786.20) and (401.16,788.90) .. (401.16,797.00);
\draw[->,line width=1.0bp,api,dash pattern=on 5bp off 4bp] (251.24,176.00) .. controls (251.24,184.10) and (19.50,186.80) .. (19.50,186.80) .. controls (19.50,366.62) and (19.50,606.38) .. (19.50,786.20) .. controls (19.50,786.20) and (401.16,788.90) .. (401.16,797.00);
\filldraw[fill=cinematicbg,draw=cinematicfg,line width=1.0bp] (401.16,820.50) ellipse (84.00bp and 23.50bp);
\node[inner sep=0pt,font=\sffamily\fontsize{10.7bp}{12bp}\selectfont\bfseries,text=ink] at (401.16,827.00) {\hyperlink{node-W2L7_12}{WZ2 Lemma 7.12}};
\node[inner sep=0pt,font=\sffamily\fontsize{10.7bp}{12bp}\selectfont,text=ink] at (401.16,814.00) {\hyperlink{node-W2L7_12}{One-step area ratio}};
\filldraw[fill=cinematicbg,draw=cinematicfg,line width=1.0bp] (401.16,733.50) ellipse (112.62bp and 36.50bp);
\node[inner sep=0pt,font=\sffamily\fontsize{10.7bp}{12bp}\selectfont\bfseries,text=ink] at (401.16,753.00) {\hyperlink{node-W2AreaCompare}{WZ2 Lemma 7.12: Holder}};
\node[inner sep=0pt,font=\sffamily\fontsize{10.7bp}{12bp}\selectfont,text=ink] at (401.16,740.00) {\hyperlink{node-W2AreaCompare}{and translations}};
\node[inner sep=0pt,font=\sffamily\fontsize{10.7bp}{12bp}\selectfont,text=ink] at (401.16,727.00) {\hyperlink{node-W2AreaCompare}{From cinematic norm to a}};
\node[inner sep=0pt,font=\sffamily\fontsize{10.7bp}{12bp}\selectfont,text=ink] at (401.16,714.00) {\hyperlink{node-W2AreaCompare}{local area ratio}};
\filldraw[fill=cinematicbg,draw=cinematicfg,line width=1.0bp] (603.16,633.50) ellipse (124.79bp and 36.50bp);
\node[inner sep=0pt,font=\sffamily\fontsize{10.7bp}{12bp}\selectfont\bfseries,text=ink] at (603.16,653.00) {\hyperlink{node-W2JetCount}{WZ2 Lemma 7.12: cinematic}};
\node[inner sep=0pt,font=\sffamily\fontsize{10.7bp}{12bp}\selectfont,text=ink] at (603.16,640.00) {\hyperlink{node-W2JetCount}{assembly}};
\node[inner sep=0pt,font=\sffamily\fontsize{10.7bp}{12bp}\selectfont,text=ink] at (603.16,627.00) {\hyperlink{node-W2JetCount}{Convert Katz-Tao triples}};
\node[inner sep=0pt,font=\sffamily\fontsize{10.7bp}{12bp}\selectfont,text=ink] at (603.16,614.00) {\hyperlink{node-W2JetCount}{to separated curves}};
\filldraw[fill=cinematicbg,draw=cinematicfg,line width=1.0bp] (199.16,533.50) ellipse (128.16bp and 36.50bp);
\node[inner sep=0pt,font=\sffamily\fontsize{10.7bp}{12bp}\selectfont\bfseries,text=ink] at (199.16,553.00) {\hyperlink{node-W2Slice}{WZ2 Lemma 7.12: parameter}};
\node[inner sep=0pt,font=\sffamily\fontsize{10.7bp}{12bp}\selectfont,text=ink] at (199.16,540.00) {\hyperlink{node-W2Slice}{slicing}};
\node[inner sep=0pt,font=\sffamily\fontsize{10.7bp}{12bp}\selectfont,text=ink] at (199.16,527.00) {\hyperlink{node-W2Slice}{A rich fixed-c slice of}};
\node[inner sep=0pt,font=\sffamily\fontsize{10.7bp}{12bp}\selectfont,text=ink] at (199.16,514.00) {\hyperlink{node-W2Slice}{R4 parameters}};
\filldraw[fill=cinematicbg,draw=cinematicfg,line width=1.0bp] (196.40,337.50) ellipse (84.00bp and 36.50bp);
\node[inner sep=0pt,font=\sffamily\fontsize{10.7bp}{12bp}\selectfont\bfseries,text=ink] at (196.40,357.00) {\hyperlink{node-W2ImageUniform}{WZ2 Lemma 7.12:}};
\node[inner sep=0pt,font=\sffamily\fontsize{10.7bp}{12bp}\selectfont,text=ink] at (196.40,344.00) {\hyperlink{node-W2ImageUniform}{preparation}};
\node[inner sep=0pt,font=\sffamily\fontsize{10.7bp}{12bp}\selectfont,text=ink] at (196.40,331.00) {\hyperlink{node-W2ImageUniform}{Uniformize the projected}};
\node[inner sep=0pt,font=\sffamily\fontsize{10.7bp}{12bp}\selectfont,text=ink] at (196.40,318.00) {\hyperlink{node-W2ImageUniform}{incidence set}};
\filldraw[fill=cinematicbg,draw=cinematicfg,line width=1.0bp] (603.16,435.50) ellipse (84.00bp and 30.00bp);
\node[inner sep=0pt,font=\sffamily\fontsize{10.7bp}{12bp}\selectfont\bfseries,text=ink] at (603.16,448.50) {\hyperlink{node-W2L7_11}{WZ2 Lemma 7.11}};
\node[inner sep=0pt,font=\sffamily\fontsize{10.7bp}{12bp}\selectfont,text=ink] at (603.16,435.50) {\hyperlink{node-W2L7_11}{Replicate a local}};
\node[inner sep=0pt,font=\sffamily\fontsize{10.7bp}{12bp}\selectfont,text=ink] at (603.16,422.50) {\hyperlink{node-W2L7_11}{Katz-Tao cluster}};
\filldraw[fill=cinematicbg,draw=cinematicfg,line width=1.0bp] (647.33,533.50) ellipse (84.00bp and 30.00bp);
\node[inner sep=0pt,font=\sffamily\fontsize{10.7bp}{12bp}\selectfont\bfseries,text=ink] at (647.33,546.50) {\hyperlink{node-W2L7_6}{WZ2 Lemma 7.6}};
\node[inner sep=0pt,font=\sffamily\fontsize{10.7bp}{12bp}\selectfont,text=ink] at (647.33,533.50) {\hyperlink{node-W2L7_6}{Frostman sets contain}};
\node[inner sep=0pt,font=\sffamily\fontsize{10.7bp}{12bp}\selectfont,text=ink] at (647.33,520.50) {\hyperlink{node-W2L7_6}{large Katz-Tao subsets}};
\filldraw[fill=cinematicbg,draw=cinematicfg,line width=1.0bp] (445.33,533.50) ellipse (84.00bp and 30.00bp);
\node[inner sep=0pt,font=\sffamily\fontsize{10.7bp}{12bp}\selectfont\bfseries,text=ink] at (445.33,546.50) {\hyperlink{node-W2L7_10}{WZ2 Lemma 7.10}};
\node[inner sep=0pt,font=\sffamily\fontsize{10.7bp}{12bp}\selectfont,text=ink] at (445.33,533.50) {\hyperlink{node-W2L7_10}{Zoom to Frostman fibres}};
\node[inner sep=0pt,font=\sffamily\fontsize{10.7bp}{12bp}\selectfont,text=ink] at (445.33,520.50) {\hyperlink{node-W2L7_10}{of a uniform set}};
\filldraw[fill=cinematicbg,draw=cinematicfg,line width=1.0bp] (401.16,148.00) ellipse (84.00bp and 30.00bp);
\node[inner sep=0pt,font=\sffamily\fontsize{10.7bp}{12bp}\selectfont\bfseries,text=ink] at (401.16,161.00) {\hyperlink{node-W2L7_9}{WZ2 Lemma 7.9}};
\node[inner sep=0pt,font=\sffamily\fontsize{10.7bp}{12bp}\selectfont,text=ink] at (401.16,148.00) {\hyperlink{node-W2L7_9}{Sparse-scale uniform}};
\node[inner sep=0pt,font=\sffamily\fontsize{10.7bp}{12bp}\selectfont,text=ink] at (401.16,135.00) {\hyperlink{node-W2L7_9}{refinement}};
\filldraw[fill=cinematicbg,draw=cinematicfg,line width=1.0bp] (158.38,633.50) ellipse (84.00bp and 23.50bp);
\node[inner sep=0pt,font=\sffamily\fontsize{10.7bp}{12bp}\selectfont\bfseries,text=ink] at (158.38,640.00) {\hyperlink{node-T7_2}{WZ1 7.2 = WZ2 7.13}};
\node[inner sep=0pt,font=\sffamily\fontsize{10.7bp}{12bp}\selectfont,text=ink] at (158.38,627.00) {\hyperlink{node-T7_2}{Cinematic maximal}};
\filldraw[fill=cinematicbg,draw=cinematicfg,line width=1.0bp] (401.16,435.50) ellipse (84.00bp and 30.00bp);
\node[inner sep=0pt,font=\sffamily\fontsize{10.7bp}{12bp}\selectfont\bfseries,text=ink] at (401.16,448.50) {\hyperlink{node-L7_3}{WZ1 Lemma 7.3}};
\node[inner sep=0pt,font=\sffamily\fontsize{10.7bp}{12bp}\selectfont,text=ink] at (401.16,435.50) {\hyperlink{node-L7_3}{The twisted-line family}};
\node[inner sep=0pt,font=\sffamily\fontsize{10.7bp}{12bp}\selectfont,text=ink] at (401.16,422.50) {\hyperlink{node-L7_3}{is cinematic}};
\filldraw[fill=cinematicbg,draw=cinematicfg,line width=1.0bp,rounded corners=5bp] (519.16,120.00) rectangle (687.16,176.00);
\node[inner sep=0pt,font=\sffamily\fontsize{10.7bp}{12bp}\selectfont\bfseries,text=ink] at (603.16,161.00) {\hyperlink{node-W2DKT}{WZ2 Def. 7.2}};
\node[inner sep=0pt,font=\sffamily\fontsize{10.7bp}{12bp}\selectfont,text=ink] at (603.16,148.00) {\hyperlink{node-W2DKT}{Absolute Katz-Tao entropy}};
\node[inner sep=0pt,font=\sffamily\fontsize{10.7bp}{12bp}\selectfont,text=ink] at (603.16,135.00) {\hyperlink{node-W2DKT}{bound}};
\filldraw[fill=cinematicbg,draw=cinematicfg,line width=1.0bp,rounded corners=5bp] (314.40,309.50) rectangle (482.40,365.50);
\node[inner sep=0pt,font=\sffamily\fontsize{10.7bp}{12bp}\selectfont\bfseries,text=ink] at (398.40,350.50) {\hyperlink{node-W2DFrost}{WZ2 Def. 7.3}};
\node[inner sep=0pt,font=\sffamily\fontsize{10.7bp}{12bp}\selectfont,text=ink] at (398.40,337.50) {\hyperlink{node-W2DFrost}{Normalized Frostman}};
\node[inner sep=0pt,font=\sffamily\fontsize{10.7bp}{12bp}\selectfont,text=ink] at (398.40,324.50) {\hyperlink{node-W2DFrost}{entropy}};
\filldraw[fill=cinematicbg,draw=cinematicfg,line width=1.0bp,rounded corners=5bp] (418.16,28.50) rectangle (586.16,84.50);
\node[inner sep=0pt,font=\sffamily\fontsize{10.7bp}{12bp}\selectfont\bfseries,text=ink] at (502.16,69.50) {\hyperlink{node-W2DUnif}{WZ2 Def. 7.8}};
\node[inner sep=0pt,font=\sffamily\fontsize{10.7bp}{12bp}\selectfont,text=ink] at (502.16,56.50) {\hyperlink{node-W2DUnif}{Uniform occupied-cell}};
\node[inner sep=0pt,font=\sffamily\fontsize{10.7bp}{12bp}\selectfont,text=ink] at (502.16,43.50) {\hyperlink{node-W2DUnif}{counts}};
\filldraw[fill=cinematicbg,draw=cinematicfg,line width=1.0bp,rounded corners=5bp] (115.16,120.00) rectangle (283.16,176.00);
\node[inner sep=0pt,font=\sffamily\fontsize{10.7bp}{12bp}\selectfont\bfseries,text=ink] at (199.16,161.00) {\hyperlink{node-W2TW}{WZ2 Sec. 7: Tube Wolff}};
\node[inner sep=0pt,font=\sffamily\fontsize{10.7bp}{12bp}\selectfont,text=ink] at (199.16,148.00) {\hyperlink{node-W2TW}{The weaker radius-squared}};
\node[inner sep=0pt,font=\sffamily\fontsize{10.7bp}{12bp}\selectfont,text=ink] at (199.16,135.00) {\hyperlink{node-W2TW}{tube count}};
\filldraw[fill=cinematicbg,draw=cinematicfg,line width=1.0bp,rounded corners=5bp] (113.51,211.50) rectangle (281.51,267.50);
\node[inner sep=0pt,font=\sffamily\fontsize{10.7bp}{12bp}\selectfont\bfseries,text=ink] at (197.51,252.50) {\hyperlink{node-Dcinematic}{WZ1 Def. 7.1}};
\node[inner sep=0pt,font=\sffamily\fontsize{10.7bp}{12bp}\selectfont,text=ink] at (197.51,239.50) {\hyperlink{node-Dcinematic}{Cinematic families of C2}};
\node[inner sep=0pt,font=\sffamily\fontsize{10.7bp}{12bp}\selectfont,text=ink] at (197.51,226.50) {\hyperlink{node-Dcinematic}{functions}};
\filldraw[fill=externalbg,draw=externalfg,line width=1.0bp,rounded corners=5bp,dashed] (516.40,303.00) rectangle (689.92,372.00);
\node[inner sep=0pt,font=\sffamily\fontsize{10.7bp}{12bp}\selectfont\bfseries,text=ink] at (603.16,357.00) {\hyperlink{node-XFO}{External: Fässler-Orponen}};
\node[inner sep=0pt,font=\sffamily\fontsize{10.7bp}{12bp}\selectfont,text=ink] at (603.16,344.00) {\hyperlink{node-XFO}{Lem. 3.13}};
\node[inner sep=0pt,font=\sffamily\fontsize{10.7bp}{12bp}\selectfont,text=ink] at (603.16,331.00) {\hyperlink{node-XFO}{Content-to-Katz-Tao}};
\node[inner sep=0pt,font=\sffamily\fontsize{10.7bp}{12bp}\selectfont,text=ink] at (603.16,318.00) {\hyperlink{node-XFO}{extraction}};
\filldraw[fill=externalbg,draw=externalfg,line width=1.0bp,rounded corners=5bp,dashed] (216.16,22.00) rectangle (384.16,91.00);
\node[inner sep=0pt,font=\sffamily\fontsize{10.7bp}{12bp}\selectfont\bfseries,text=ink] at (300.16,76.00) {\hyperlink{node-XOSuniform}{External:}};
\node[inner sep=0pt,font=\sffamily\fontsize{10.7bp}{12bp}\selectfont,text=ink] at (300.16,63.00) {\hyperlink{node-XOSuniform}{Orponen-Shmerkin Lem.}};
\node[inner sep=0pt,font=\sffamily\fontsize{10.7bp}{12bp}\selectfont,text=ink] at (300.16,50.00) {\hyperlink{node-XOSuniform}{Uniform refinement black}};
\node[inner sep=0pt,font=\sffamily\fontsize{10.7bp}{12bp}\selectfont,text=ink] at (300.16,37.00) {\hyperlink{node-XOSuniform}{box}};
\filldraw[fill=externalbg,draw=externalfg,line width=1.0bp,rounded corners=5bp,dashed] (519.16,205.00) rectangle (688.82,274.00);
\node[inner sep=0pt,font=\sffamily\fontsize{10.7bp}{12bp}\selectfont\bfseries,text=ink] at (603.99,259.00) {\hyperlink{node-XSWdecomp}{External: Shmerkin-Wang}};
\node[inner sep=0pt,font=\sffamily\fontsize{10.7bp}{12bp}\selectfont,text=ink] at (603.99,246.00) {\hyperlink{node-XSWdecomp}{Lem. 5.21}};
\node[inner sep=0pt,font=\sffamily\fontsize{10.7bp}{12bp}\selectfont,text=ink] at (603.99,233.00) {\hyperlink{node-XSWdecomp}{Decompose the}};
\node[inner sep=0pt,font=\sffamily\fontsize{10.7bp}{12bp}\selectfont,text=ink] at (603.99,220.00) {\hyperlink{node-XSWdecomp}{scale-entropy profile}};
\filldraw[fill=externalbg,draw=externalfg,line width=1.0bp,rounded corners=5bp,dashed] (315.51,205.00) rectangle (485.16,274.00);
\node[inner sep=0pt,font=\sffamily\fontsize{10.7bp}{12bp}\selectfont\bfseries,text=ink] at (400.34,259.00) {\hyperlink{node-XSWmerge}{External: Shmerkin-Wang}};
\node[inner sep=0pt,font=\sffamily\fontsize{10.7bp}{12bp}\selectfont,text=ink] at (400.34,246.00) {\hyperlink{node-XSWmerge}{Lem. 5.20}};
\node[inner sep=0pt,font=\sffamily\fontsize{10.7bp}{12bp}\selectfont,text=ink] at (400.34,233.00) {\hyperlink{node-XSWmerge}{Merge adjacent entropy}};
\node[inner sep=0pt,font=\sffamily\fontsize{10.7bp}{12bp}\selectfont,text=ink] at (400.34,220.00) {\hyperlink{node-XSWmerge}{intervals}};
\filldraw[fill=externalbg,draw=externalfg,line width=1.0bp,rounded corners=5bp,dashed] (115.16,401.00) rectangle (283.16,470.00);
\node[inner sep=0pt,font=\sffamily\fontsize{10.7bp}{12bp}\selectfont\bfseries,text=ink] at (199.16,455.00) {\hyperlink{node-XPYZ}{External:}};
\node[inner sep=0pt,font=\sffamily\fontsize{10.7bp}{12bp}\selectfont,text=ink] at (199.16,442.00) {\hyperlink{node-XPYZ}{Pramanik-Yang-Zahl}};
\node[inner sep=0pt,font=\sffamily\fontsize{10.7bp}{12bp}\selectfont,text=ink] at (199.16,429.00) {\hyperlink{node-XPYZ}{Cinematic L3/2 maximal}};
\node[inner sep=0pt,font=\sffamily\fontsize{10.7bp}{12bp}\selectfont,text=ink] at (199.16,416.00) {\hyperlink{node-XPYZ}{theorem}};
\filldraw[fill=apibg,draw=apifg,line width=1.0bp,rounded corners=5bp] (276.38,605.50) rectangle (444.38,661.50);
\node[inner sep=0pt,font=\sffamily\fontsize{10.7bp}{12bp}\selectfont\bfseries,text=ink] at (360.38,646.50) {\hyperlink{node-Iparams}{Quantifier API}};
\node[inner sep=0pt,font=\sffamily\fontsize{10.7bp}{12bp}\selectfont,text=ink] at (360.38,633.50) {\hyperlink{node-Iparams}{Small-parameter and}};
\node[inner sep=0pt,font=\sffamily\fontsize{10.7bp}{12bp}\selectfont,text=ink] at (360.38,620.50) {\hyperlink{node-Iparams}{finite-iteration}};
\end{tikzpicture}}\end{center}
{\footnotesize Solid: proof input. Dashed: API. Dash-dot: assembly. Long-dashed: overview path. Read from the top down; arrows point to consumers. Full prerequisites appear on the cards and in the ledger.}\par
\sectitle{WZ1 radial bootstrap and dot products}
\hypertarget{view-bootstrap}{}
{\sffamily\small\color{muted}This internal WZ1 branch feeds C2 regularity, not any Hausdorff-dimension deduction.}\par
\begin{center}\resizebox{493.00bp}{!}{\begin{tikzpicture}[x=1bp,y=1bp, >={Stealth[length=4bp,width=3.2bp]}]
\path[use as bounding box] (0,0) rectangle (690.24,531.00);
\fill[panel,rounded corners=5bp] (0,0) rectangle (690.24,531.00);
\draw[line width=3.5bp,panel] (556.28,78.00) .. controls (556.28,93.41) and (555.87,96.09) .. (555.87,111.50);
\draw[->,line width=1.0bp,api,dash pattern=on 5bp off 4bp] (556.28,78.00) .. controls (556.28,93.41) and (555.87,96.09) .. (555.87,111.50);
\draw[line width=3.5bp,panel] (99.38,78.00) .. controls (99.38,92.49) and (353.87,95.01) .. (353.87,109.50);
\draw[->,line width=1.0bp,core] (99.38,78.00) .. controls (99.38,92.49) and (353.87,95.01) .. (353.87,109.50);
\draw[line width=3.5bp,panel] (353.87,156.50) .. controls (353.87,164.60) and (44.00,167.30) .. (44.00,167.30) .. controls (44.00,198.31) and (44.00,239.65) .. (44.00,270.66) .. controls (44.00,270.66) and (405.95,273.36) .. (405.95,281.46);
\draw[->,line width=1.0bp,core] (353.87,156.50) .. controls (353.87,164.60) and (44.00,167.30) .. (44.00,167.30) .. controls (44.00,198.31) and (44.00,239.65) .. (44.00,270.66) .. controls (44.00,270.66) and (405.95,273.36) .. (405.95,281.46);
\draw[line width=3.5bp,panel] (151.87,161.00) .. controls (151.87,169.10) and (40.50,171.80) .. (40.50,171.80) .. controls (40.50,201.46) and (40.50,241.00) .. (40.50,270.66) .. controls (40.50,270.66) and (301.79,273.36) .. (301.79,281.46);
\draw[->,line width=1.0bp,core] (151.87,161.00) .. controls (151.87,169.10) and (40.50,171.80) .. (40.50,171.80) .. controls (40.50,201.46) and (40.50,241.00) .. (40.50,270.66) .. controls (40.50,270.66) and (301.79,273.36) .. (301.79,281.46);
\draw[line width=3.5bp,panel] (608.36,78.00) .. controls (608.36,86.10) and (663.74,88.80) .. (663.74,88.80) .. controls (663.74,117.26) and (663.74,155.20) .. (663.74,183.66) .. controls (663.74,183.66) and (610.35,186.36) .. (610.35,194.46);
\draw[->,line width=1.0bp,core] (608.36,78.00) .. controls (608.36,86.10) and (663.74,88.80) .. (663.74,88.80) .. controls (663.74,117.26) and (663.74,155.20) .. (663.74,183.66) .. controls (663.74,183.66) and (610.35,186.36) .. (610.35,194.46);
\draw[line width=3.5bp,panel] (607.95,154.50) .. controls (607.95,172.88) and (501.39,176.08) .. (501.39,194.46);
\draw[->,line width=1.0bp,core] (607.95,154.50) .. controls (607.95,172.88) and (501.39,176.08) .. (501.39,194.46);
\draw[line width=3.5bp,panel] (353.87,335.00) .. controls (353.87,348.10) and (379.91,350.38) .. (379.91,363.48);
\draw[->,line width=1.0bp,core] (353.87,335.00) .. controls (353.87,348.10) and (379.91,350.38) .. (379.91,363.48);
\draw[line width=3.5bp,panel] (503.79,154.50) .. controls (503.79,162.60) and (667.24,165.30) .. (667.24,165.30) .. controls (667.24,223.01) and (667.24,299.95) .. (667.24,357.66) .. controls (667.24,357.66) and (405.95,360.36) .. (405.95,368.46);
\draw[->,line width=1.0bp,core] (503.79,154.50) .. controls (503.79,162.60) and (667.24,165.30) .. (667.24,165.30) .. controls (667.24,223.01) and (667.24,299.95) .. (667.24,357.66) .. controls (667.24,357.66) and (405.95,360.36) .. (405.95,368.46);
\draw[line width=3.5bp,panel] (203.54,78.00) .. controls (203.54,86.10) and (37.00,88.80) .. (37.00,88.80) .. controls (37.00,169.46) and (37.00,277.00) .. (37.00,357.66) .. controls (37.00,357.66) and (301.79,360.36) .. (301.79,368.46);
\draw[->,line width=1.0bp,core] (203.54,78.00) .. controls (203.54,86.10) and (37.00,88.80) .. (37.00,88.80) .. controls (37.00,169.46) and (37.00,277.00) .. (37.00,357.66) .. controls (37.00,357.66) and (301.79,360.36) .. (301.79,368.46);
\draw[line width=3.5bp,panel] (203.95,161.00) .. controls (203.95,169.10) and (33.50,171.80) .. (33.50,171.80) .. controls (33.50,226.06) and (33.50,298.41) .. (33.50,352.68) .. controls (33.50,352.68) and (327.83,355.38) .. (327.83,363.48);
\draw[->,line width=1.0bp,core] (203.95,161.00) .. controls (203.95,169.10) and (33.50,171.80) .. (33.50,171.80) .. controls (33.50,226.06) and (33.50,298.41) .. (33.50,352.68) .. controls (33.50,352.68) and (327.83,355.38) .. (327.83,363.48);
\draw[line width=3.5bp,panel] (151.87,326.50) .. controls (151.87,342.83) and (353.87,345.67) .. (353.87,362.00);
\draw[->,line width=1.0bp,core] (151.87,326.50) .. controls (151.87,342.83) and (353.87,345.67) .. (353.87,362.00);
\draw[line width=3.5bp,panel] (353.87,78.00) .. controls (353.87,86.10) and (30.00,88.80) .. (30.00,88.80) .. controls (30.00,117.26) and (30.00,155.20) .. (30.00,183.66) .. controls (30.00,183.66) and (297.92,186.36) .. (297.92,194.46);
\draw[->,line width=1.0bp,core] (353.87,78.00) .. controls (353.87,86.10) and (30.00,88.80) .. (30.00,88.80) .. controls (30.00,117.26) and (30.00,155.20) .. (30.00,183.66) .. controls (30.00,183.66) and (297.92,186.36) .. (297.92,194.46);
\draw[line width=3.5bp,panel] (504.20,78.00) .. controls (504.20,86.10) and (670.74,88.80) .. (670.74,88.80) .. controls (670.74,117.26) and (670.74,155.20) .. (670.74,183.66) .. controls (670.74,183.66) and (402.08,186.36) .. (402.08,194.46);
\draw[->,line width=1.0bp,core] (504.20,78.00) .. controls (504.20,86.10) and (670.74,88.80) .. (670.74,88.80) .. controls (670.74,117.26) and (670.74,155.20) .. (670.74,183.66) .. controls (670.74,183.66) and (402.08,186.36) .. (402.08,194.46);
\draw[line width=3.5bp,panel] (350.00,248.00) .. controls (350.00,263.39) and (503.79,266.07) .. (503.79,281.46);
\draw[->,line width=1.0bp,core] (350.00,248.00) .. controls (350.00,263.39) and (503.79,266.07) .. (503.79,281.46);
\draw[line width=3.5bp,panel] (555.87,248.00) .. controls (555.87,263.39) and (607.95,266.07) .. (607.95,281.46);
\draw[->,line width=1.0bp,core] (555.87,248.00) .. controls (555.87,263.39) and (607.95,266.07) .. (607.95,281.46);
\draw[line width=3.5bp,panel] (555.87,335.00) .. controls (555.87,350.39) and (607.95,353.07) .. (607.95,368.46);
\draw[->,line width=1.0bp,core] (555.87,335.00) .. controls (555.87,350.39) and (607.95,353.07) .. (607.95,368.46);
\draw[line width=3.5bp,panel] (99.79,161.00) .. controls (99.79,169.10) and (26.50,171.80) .. (26.50,171.80) .. controls (26.50,225.62) and (26.50,297.38) .. (26.50,351.20) .. controls (26.50,351.20) and (151.87,353.90) .. (151.87,362.00);
\draw[->,line width=1.0bp,core] (99.79,161.00) .. controls (99.79,169.10) and (26.50,171.80) .. (26.50,171.80) .. controls (26.50,225.62) and (26.50,297.38) .. (26.50,351.20) .. controls (26.50,351.20) and (151.87,353.90) .. (151.87,362.00);
\draw[line width=3.5bp,panel] (151.87,422.00) .. controls (151.87,434.72) and (336.51,436.93) .. (336.51,449.65);
\draw[->,line width=1.0bp,core] (151.87,422.00) .. controls (151.87,434.72) and (336.51,436.93) .. (336.51,449.65);
\draw[line width=3.5bp,panel] (555.87,422.00) .. controls (555.87,437.39) and (405.95,440.07) .. (405.95,455.46);
\draw[->,line width=1.0bp,core] (555.87,422.00) .. controls (555.87,437.39) and (405.95,440.07) .. (405.95,455.46);
\draw[line width=3.5bp,panel] (353.87,422.00) .. controls (353.87,434.72) and (371.23,436.93) .. (371.23,449.65);
\draw[->,line width=1.0bp,core] (353.87,422.00) .. controls (353.87,434.72) and (371.23,436.93) .. (371.23,449.65);
\draw[line width=3.5bp,panel] (200.08,246.00) .. controls (200.08,254.10) and (23.00,256.80) .. (23.00,256.80) .. controls (23.00,287.06) and (23.00,327.40) .. (23.00,357.66) .. controls (23.00,357.66) and (503.79,360.36) .. (503.79,368.46);
\draw[->,line width=1.0bp,api,dash pattern=on 5bp off 4bp] (200.08,246.00) .. controls (200.08,254.10) and (23.00,256.80) .. (23.00,256.80) .. controls (23.00,287.06) and (23.00,327.40) .. (23.00,357.66) .. controls (23.00,357.66) and (503.79,360.36) .. (503.79,368.46);
\draw[line width=3.5bp,panel] (95.92,246.00) .. controls (95.92,254.10) and (19.50,256.80) .. (19.50,256.80) .. controls (19.50,313.16) and (19.50,388.30) .. (19.50,444.66) .. controls (19.50,444.66) and (301.79,447.36) .. (301.79,455.46);
\draw[->,line width=1.0bp,api,dash pattern=on 5bp off 4bp] (95.92,246.00) .. controls (95.92,254.10) and (19.50,256.80) .. (19.50,256.80) .. controls (19.50,313.16) and (19.50,388.30) .. (19.50,444.66) .. controls (19.50,444.66) and (301.79,447.36) .. (301.79,455.46);
\filldraw[fill=radialbg,draw=radialfg,line width=1.0bp] (353.87,479.00) ellipse (84.00bp and 30.00bp);
\node[inner sep=0pt,font=\sffamily\fontsize{10.7bp}{12bp}\selectfont\bfseries,text=ink] at (353.87,492.00) {\hyperlink{node-P8_5}{WZ1 Prop. 8.5}};
\node[inner sep=0pt,font=\sffamily\fontsize{10.7bp}{12bp}\selectfont,text=ink] at (353.87,479.00) {\hyperlink{node-P8_5}{Strip nonconcentration}};
\node[inner sep=0pt,font=\sffamily\fontsize{10.7bp}{12bp}\selectfont,text=ink] at (353.87,466.00) {\hyperlink{node-P8_5}{forces expansion}};
\filldraw[fill=radialbg,draw=radialfg,line width=1.0bp] (151.87,392.00) ellipse (84.00bp and 30.00bp);
\node[inner sep=0pt,font=\sffamily\fontsize{10.7bp}{12bp}\selectfont\bfseries,text=ink] at (151.87,405.00) {\hyperlink{node-L8_8}{WZ1 Lemma 8.8}};
\node[inner sep=0pt,font=\sffamily\fontsize{10.7bp}{12bp}\selectfont,text=ink] at (151.87,392.00) {\hyperlink{node-L8_8}{Obtain a fixed}};
\node[inner sep=0pt,font=\sffamily\fontsize{10.7bp}{12bp}\selectfont,text=ink] at (151.87,379.00) {\hyperlink{node-L8_8}{1/4-thin-tube seed}};
\filldraw[fill=radialbg,draw=radialfg,line width=1.0bp] (353.87,392.00) ellipse (84.00bp and 30.00bp);
\node[inner sep=0pt,font=\sffamily\fontsize{10.7bp}{12bp}\selectfont\bfseries,text=ink] at (353.87,405.00) {\hyperlink{node-L8_4}{WZ1 Lemma 8.4}};
\node[inner sep=0pt,font=\sffamily\fontsize{10.7bp}{12bp}\selectfont,text=ink] at (353.87,392.00) {\hyperlink{node-L8_4}{Thin tubes force a large}};
\node[inner sep=0pt,font=\sffamily\fontsize{10.7bp}{12bp}\selectfont,text=ink] at (353.87,379.00) {\hyperlink{node-L8_4}{dot-product image}};
\filldraw[fill=radialbg,draw=radialfg,line width=1.0bp] (555.87,392.00) ellipse (84.00bp and 30.00bp);
\node[inner sep=0pt,font=\sffamily\fontsize{10.7bp}{12bp}\selectfont\bfseries,text=ink] at (555.87,405.00) {\hyperlink{node-C8_7}{WZ1 Cor. 8.7}};
\node[inner sep=0pt,font=\sffamily\fontsize{10.7bp}{12bp}\selectfont,text=ink] at (555.87,392.00) {\hyperlink{node-C8_7}{Iterate the bootstrap up}};
\node[inner sep=0pt,font=\sffamily\fontsize{10.7bp}{12bp}\selectfont,text=ink] at (555.87,379.00) {\hyperlink{node-C8_7}{to exponent 1}};
\filldraw[fill=radialbg,draw=radialfg,line width=1.0bp] (353.87,133.00) ellipse (84.00bp and 23.50bp);
\node[inner sep=0pt,font=\sffamily\fontsize{10.7bp}{12bp}\selectfont\bfseries,text=ink] at (353.87,139.50) {\hyperlink{node-L8_1}{WZ1 Lemma 8.1}};
\node[inner sep=0pt,font=\sffamily\fontsize{10.7bp}{12bp}\selectfont,text=ink] at (353.87,126.50) {\hyperlink{node-L8_1}{Hypergraph refinement}};
\filldraw[fill=radialbg,draw=radialfg,line width=1.0bp] (353.87,305.00) ellipse (84.00bp and 30.00bp);
\node[inner sep=0pt,font=\sffamily\fontsize{10.7bp}{12bp}\selectfont\bfseries,text=ink] at (353.87,318.00) {\hyperlink{node-T8_2}{WZ1 Thm. 8.2}};
\node[inner sep=0pt,font=\sffamily\fontsize{10.7bp}{12bp}\selectfont,text=ink] at (353.87,305.00) {\hyperlink{node-T8_2}{A discretized Kaufman}};
\node[inner sep=0pt,font=\sffamily\fontsize{10.7bp}{12bp}\selectfont,text=ink] at (353.87,292.00) {\hyperlink{node-T8_2}{estimate with dense}};
\filldraw[fill=radialbg,draw=radialfg,line width=1.0bp] (555.87,305.00) ellipse (84.00bp and 30.00bp);
\node[inner sep=0pt,font=\sffamily\fontsize{10.7bp}{12bp}\selectfont\bfseries,text=ink] at (555.87,318.00) {\hyperlink{node-P8_6p}{WZ1 Prop. 8.6'}};
\node[inner sep=0pt,font=\sffamily\fontsize{10.7bp}{12bp}\selectfont,text=ink] at (555.87,305.00) {\hyperlink{node-P8_6p}{Radial bootstrap for grid}};
\node[inner sep=0pt,font=\sffamily\fontsize{10.7bp}{12bp}\selectfont,text=ink] at (555.87,292.00) {\hyperlink{node-P8_6p}{sets}};
\filldraw[fill=radialbg,draw=radialfg,line width=1.0bp] (555.87,218.00) ellipse (87.87bp and 30.00bp);
\node[inner sep=0pt,font=\sffamily\fontsize{10.7bp}{12bp}\selectfont\bfseries,text=ink] at (555.87,231.00) {\hyperlink{node-L8_3}{WZ1 Lemma 8.3}};
\node[inner sep=0pt,font=\sffamily\fontsize{10.7bp}{12bp}\selectfont,text=ink] at (555.87,218.00) {\hyperlink{node-L8_3}{Pass from measure thin}};
\node[inner sep=0pt,font=\sffamily\fontsize{10.7bp}{12bp}\selectfont,text=ink] at (555.87,205.00) {\hyperlink{node-L8_3}{tubes to grid thin tubes}};
\filldraw[fill=radialbg,draw=radialfg,line width=1.0bp] (350.00,218.00) ellipse (84.00bp and 30.00bp);
\node[inner sep=0pt,font=\sffamily\fontsize{10.7bp}{12bp}\selectfont\bfseries,text=ink] at (350.00,231.00) {\hyperlink{node-P8_6}{WZ1 Prop. 8.6}};
\node[inner sep=0pt,font=\sffamily\fontsize{10.7bp}{12bp}\selectfont,text=ink] at (350.00,218.00) {\hyperlink{node-P8_6}{Radial bootstrap for}};
\node[inner sep=0pt,font=\sffamily\fontsize{10.7bp}{12bp}\selectfont,text=ink] at (350.00,205.00) {\hyperlink{node-P8_6}{measures}};
\filldraw[fill=radialbg,draw=radialfg,line width=1.0bp,rounded corners=5bp] (67.87,105.00) rectangle (235.87,161.00);
\node[inner sep=0pt,font=\sffamily\fontsize{10.7bp}{12bp}\selectfont\bfseries,text=ink] at (151.87,146.00) {\hyperlink{node-DFrost}{WZ1 Def. 8.1}};
\node[inner sep=0pt,font=\sffamily\fontsize{10.7bp}{12bp}\selectfont,text=ink] at (151.87,133.00) {\hyperlink{node-DFrost}{Normalized discretized}};
\node[inner sep=0pt,font=\sffamily\fontsize{10.7bp}{12bp}\selectfont,text=ink] at (151.87,120.00) {\hyperlink{node-DFrost}{Frostman sets}};
\filldraw[fill=radialbg,draw=radialfg,line width=1.0bp,rounded corners=5bp] (67.87,283.50) rectangle (235.87,326.50);
\node[inner sep=0pt,font=\sffamily\fontsize{10.7bp}{12bp}\selectfont\bfseries,text=ink] at (151.87,311.50) {\hyperlink{node-Dsep}{WZ1 Def. 8.6}};
\node[inner sep=0pt,font=\sffamily\fontsize{10.7bp}{12bp}\selectfont,text=ink] at (151.87,298.50) {\hyperlink{node-Dsep}{Standard separation}};
\filldraw[fill=apibg,draw=apifg,line width=1.0bp,rounded corners=5bp] (64.00,190.00) rectangle (232.00,246.00);
\node[inner sep=0pt,font=\sffamily\fontsize{10.7bp}{12bp}\selectfont\bfseries,text=ink] at (148.00,231.00) {\hyperlink{node-Iparams}{Quantifier API}};
\node[inner sep=0pt,font=\sffamily\fontsize{10.7bp}{12bp}\selectfont,text=ink] at (148.00,218.00) {\hyperlink{node-Iparams}{Small-parameter and}};
\node[inner sep=0pt,font=\sffamily\fontsize{10.7bp}{12bp}\selectfont,text=ink] at (148.00,205.00) {\hyperlink{node-Iparams}{finite-iteration}};
\filldraw[fill=radialbg,draw=radialfg,line width=1.0bp,rounded corners=5bp] (67.46,22.00) rectangle (235.46,78.00);
\node[inner sep=0pt,font=\sffamily\fontsize{10.7bp}{12bp}\selectfont\bfseries,text=ink] at (151.46,63.00) {\hyperlink{node-Dhyp}{WZ1 Def.s 8.2-8.4}};
\node[inner sep=0pt,font=\sffamily\fontsize{10.7bp}{12bp}\selectfont,text=ink] at (151.46,50.00) {\hyperlink{node-Dhyp}{Hypergraph fibres and}};
\node[inner sep=0pt,font=\sffamily\fontsize{10.7bp}{12bp}\selectfont,text=ink] at (151.46,37.00) {\hyperlink{node-Dhyp}{uniform density}};
\filldraw[fill=radialbg,draw=radialfg,line width=1.0bp,rounded corners=5bp] (471.87,111.50) rectangle (639.87,154.50);
\node[inner sep=0pt,font=\sffamily\fontsize{10.7bp}{12bp}\selectfont\bfseries,text=ink] at (555.87,139.50) {\hyperlink{node-DthinD}{WZ1 Def. 8.5'}};
\node[inner sep=0pt,font=\sffamily\fontsize{10.7bp}{12bp}\selectfont,text=ink] at (555.87,126.50) {\hyperlink{node-DthinD}{Discretized thin tubes}};
\filldraw[fill=radialbg,draw=radialfg,line width=1.0bp,rounded corners=5bp] (472.28,22.00) rectangle (640.28,78.00);
\node[inner sep=0pt,font=\sffamily\fontsize{10.7bp}{12bp}\selectfont\bfseries,text=ink] at (556.28,63.00) {\hyperlink{node-Dthin}{WZ1 Def. 8.5}};
\node[inner sep=0pt,font=\sffamily\fontsize{10.7bp}{12bp}\selectfont,text=ink] at (556.28,50.00) {\hyperlink{node-Dthin}{Thin tubes for}};
\node[inner sep=0pt,font=\sffamily\fontsize{10.7bp}{12bp}\selectfont,text=ink] at (556.28,37.00) {\hyperlink{node-Dthin}{probability measures}};
\filldraw[fill=externalbg,draw=externalfg,line width=1.0bp,rounded corners=5bp,dashed] (269.46,22.00) rectangle (438.28,78.00);
\node[inner sep=0pt,font=\sffamily\fontsize{10.7bp}{12bp}\selectfont\bfseries,text=ink] at (353.87,63.00) {\hyperlink{node-XOSW}{External [35, Lemma 2.9]}};
\node[inner sep=0pt,font=\sffamily\fontsize{10.7bp}{12bp}\selectfont,text=ink] at (353.87,50.00) {\hyperlink{node-XOSW}{Radial thin-tube}};
\node[inner sep=0pt,font=\sffamily\fontsize{10.7bp}{12bp}\selectfont,text=ink] at (353.87,37.00) {\hyperlink{node-XOSW}{bootstrap}};
\end{tikzpicture}}\end{center}
{\footnotesize Solid: proof input. Dashed: API. Dash-dot: assembly. Long-dashed: overview path. Read from the top down; arrows point to consumers. Full prerequisites appear on the cards and in the ledger.}\par
\sectitle{WZ1 strip reduction to the projection dichotomy}
\hypertarget{view-reductions}{}
{\sffamily\small\color{muted}Expand the WZ1 Theorem 5.2 interface to the same granularity as the reference atlas.}\par
\begin{center}\resizebox{493.00bp}{!}{\begin{tikzpicture}[x=1bp,y=1bp, >={Stealth[length=4bp,width=3.2bp]}]
\path[use as bounding box] (0,0) rectangle (696.50,609.00);
\fill[panel,rounded corners=5bp] (0,0) rectangle (696.50,609.00);
\draw[line width=3.5bp,panel] (203.58,150.00) .. controls (203.58,158.10) and (47.50,160.80) .. (47.50,160.80) .. controls (47.50,241.92) and (47.50,350.08) .. (47.50,431.20) .. controls (47.50,431.20) and (246.28,433.90) .. (246.28,442.00);
\draw[->,line width=1.0bp,api,dash pattern=on 5bp off 4bp] (203.58,150.00) .. controls (203.58,158.10) and (47.50,160.80) .. (47.50,160.80) .. controls (47.50,241.92) and (47.50,350.08) .. (47.50,431.20) .. controls (47.50,431.20) and (246.28,433.90) .. (246.28,442.00);
\draw[line width=3.5bp,panel] (555.50,78.00) .. controls (555.50,90.42) and (555.50,92.58) .. (555.50,105.00);
\draw[->,line width=1.0bp,core] (555.50,78.00) .. controls (555.50,90.42) and (555.50,92.58) .. (555.50,105.00);
\draw[line width=3.5bp,panel] (572.86,151.49) .. controls (572.86,167.12) and (607.58,169.84) .. (607.58,185.46);
\draw[->,line width=1.0bp,core] (572.86,151.49) .. controls (572.86,167.12) and (607.58,169.84) .. (607.58,185.46);
\draw[line width=3.5bp,panel] (405.58,78.00) .. controls (405.58,86.10) and (659.50,88.80) .. (659.50,88.80) .. controls (659.50,114.56) and (659.50,148.90) .. (659.50,174.66) .. controls (659.50,174.66) and (503.42,177.36) .. (503.42,185.46);
\draw[->,line width=1.0bp,core] (405.58,78.00) .. controls (405.58,86.10) and (659.50,88.80) .. (659.50,88.80) .. controls (659.50,114.56) and (659.50,148.90) .. (659.50,174.66) .. controls (659.50,174.66) and (503.42,177.36) .. (503.42,185.46);
\draw[line width=3.5bp,panel] (301.42,78.00) .. controls (301.42,86.10) and (44.00,88.80) .. (44.00,88.80) .. controls (44.00,114.56) and (44.00,148.90) .. (44.00,174.66) .. controls (44.00,174.66) and (203.58,177.36) .. (203.58,185.46);
\draw[->,line width=1.0bp,core] (301.42,78.00) .. controls (301.42,86.10) and (44.00,88.80) .. (44.00,88.80) .. controls (44.00,114.56) and (44.00,148.90) .. (44.00,174.66) .. controls (44.00,174.66) and (203.58,177.36) .. (203.58,185.46);
\draw[line width=3.5bp,panel] (336.14,78.00) .. controls (336.14,86.10) and (663.00,88.80) .. (663.00,88.80) .. controls (663.00,112.62) and (663.00,144.38) .. (663.00,168.20) .. controls (663.00,168.20) and (353.50,170.90) .. (353.50,179.00);
\draw[->,line width=1.0bp,core] (336.14,78.00) .. controls (336.14,86.10) and (663.00,88.80) .. (663.00,88.80) .. controls (663.00,112.62) and (663.00,144.38) .. (663.00,168.20) .. controls (663.00,168.20) and (353.50,170.90) .. (353.50,179.00);
\draw[line width=3.5bp,panel] (607.58,146.94) .. controls (607.58,155.04) and (666.50,157.74) .. (666.50,157.74) .. controls (666.50,187.62) and (666.50,227.47) .. (666.50,257.35) .. controls (666.50,257.35) and (586.75,260.05) .. (586.75,268.15);
\draw[->,line width=1.0bp,core] (607.58,146.94) .. controls (607.58,155.04) and (666.50,157.74) .. (666.50,157.74) .. controls (666.50,187.62) and (666.50,227.47) .. (666.50,257.35) .. controls (666.50,257.35) and (586.75,260.05) .. (586.75,268.15);
\draw[line width=3.5bp,panel] (203.58,232.54) .. controls (203.58,250.90) and (503.42,254.10) .. (503.42,272.46);
\draw[->,line width=1.0bp,core] (203.58,232.54) .. controls (203.58,250.90) and (503.42,254.10) .. (503.42,272.46);
\draw[line width=3.5bp,panel] (353.50,239.00) .. controls (353.50,251.53) and (545.08,253.71) .. (545.08,266.23);
\draw[->,line width=1.0bp,core] (353.50,239.00) .. controls (353.50,251.53) and (545.08,253.71) .. (545.08,266.23);
\draw[line width=3.5bp,panel] (555.50,239.00) .. controls (555.50,254.39) and (607.58,257.07) .. (607.58,272.46);
\draw[->,line width=1.0bp,core] (555.50,239.00) .. controls (555.50,254.39) and (607.58,257.07) .. (607.58,272.46);
\draw[line width=3.5bp,panel] (353.50,150.00) .. controls (353.50,158.10) and (670.00,160.80) .. (670.00,160.80) .. controls (670.00,189.19) and (670.00,227.04) .. (670.00,255.43) .. controls (670.00,255.43) and (565.92,258.13) .. (565.92,266.23);
\draw[->,line width=1.0bp,core] (353.50,150.00) .. controls (353.50,158.10) and (670.00,160.80) .. (670.00,160.80) .. controls (670.00,189.19) and (670.00,227.04) .. (670.00,255.43) .. controls (670.00,255.43) and (565.92,258.13) .. (565.92,266.23);
\draw[line width=3.5bp,panel] (151.50,326.00) .. controls (151.50,341.39) and (296.01,344.07) .. (296.01,359.46);
\draw[->,line width=1.0bp,core] (151.50,326.00) .. controls (151.50,341.39) and (296.01,344.07) .. (296.01,359.46);
\draw[line width=3.5bp,panel] (555.50,326.00) .. controls (555.50,339.41) and (388.00,341.74) .. (388.00,355.15);
\draw[->,line width=1.0bp,core] (555.50,326.00) .. controls (555.50,339.41) and (388.00,341.74) .. (388.00,355.15);
\draw[line width=3.5bp,panel] (99.42,232.54) .. controls (99.42,240.64) and (40.50,243.34) .. (40.50,243.34) .. controls (40.50,273.64) and (40.50,314.05) .. (40.50,344.35) .. controls (40.50,344.35) and (319.00,347.05) .. (319.00,355.15);
\draw[->,line width=1.0bp,core] (99.42,232.54) .. controls (99.42,240.64) and (40.50,243.34) .. (40.50,243.34) .. controls (40.50,273.64) and (40.50,314.05) .. (40.50,344.35) .. controls (40.50,344.35) and (319.00,347.05) .. (319.00,355.15);
\draw[line width=3.5bp,panel] (503.42,146.94) .. controls (503.42,155.04) and (673.50,157.74) .. (673.50,157.74) .. controls (673.50,215.02) and (673.50,291.38) .. (673.50,348.66) .. controls (673.50,348.66) and (410.99,351.36) .. (410.99,359.46);
\draw[->,line width=1.0bp,core] (503.42,146.94) .. controls (503.42,155.04) and (673.50,157.74) .. (673.50,157.74) .. controls (673.50,215.02) and (673.50,291.38) .. (673.50,348.66) .. controls (673.50,348.66) and (410.99,351.36) .. (410.99,359.46);
\draw[line width=3.5bp,panel] (353.50,326.00) .. controls (353.50,338.53) and (365.00,340.71) .. (365.00,353.23);
\draw[->,line width=1.0bp,core] (353.50,326.00) .. controls (353.50,338.53) and (365.00,340.71) .. (365.00,353.23);
\draw[line width=3.5bp,panel] (353.50,413.00) .. controls (353.50,425.42) and (454.50,427.58) .. (454.50,440.00);
\draw[->,line width=1.0bp,core] (353.50,413.00) .. controls (353.50,425.42) and (454.50,427.58) .. (454.50,440.00);
\draw[line width=3.5bp,panel] (538.14,151.49) .. controls (538.14,159.59) and (677.00,162.29) .. (677.00,162.29) .. controls (677.00,244.30) and (677.00,353.65) .. (677.00,435.66) .. controls (677.00,435.66) and (510.44,438.36) .. (510.44,446.46);
\draw[->,line width=1.0bp,core] (538.14,151.49) .. controls (538.14,159.59) and (677.00,162.29) .. (677.00,162.29) .. controls (677.00,244.30) and (677.00,353.65) .. (677.00,435.66) .. controls (677.00,435.66) and (510.44,438.36) .. (510.44,446.46);
\draw[line width=3.5bp,panel] (454.50,500.00) .. controls (454.50,515.39) and (405.58,518.07) .. (405.58,533.46);
\draw[->,line width=1.0bp,core] (454.50,500.00) .. controls (454.50,515.39) and (405.58,518.07) .. (405.58,533.46);
\draw[line width=3.5bp,panel] (370.86,78.00) .. controls (370.86,86.10) and (37.00,88.80) .. (37.00,88.80) .. controls (37.00,217.02) and (37.00,387.98) .. (37.00,516.20) .. controls (37.00,516.20) and (353.50,518.90) .. (353.50,527.00);
\draw[->,line width=1.0bp,core] (370.86,78.00) .. controls (370.86,86.10) and (37.00,88.80) .. (37.00,88.80) .. controls (37.00,217.02) and (37.00,387.98) .. (37.00,516.20) .. controls (37.00,516.20) and (353.50,518.90) .. (353.50,527.00);
\draw[line width=3.5bp,panel] (246.28,498.00) .. controls (246.28,514.31) and (301.42,517.15) .. (301.42,533.46);
\draw[->,line width=1.0bp,core] (246.28,498.00) .. controls (246.28,514.31) and (301.42,517.15) .. (301.42,533.46);
\draw[line width=3.5bp,panel] (99.42,150.00) .. controls (99.42,158.10) and (33.50,160.80) .. (33.50,160.80) .. controls (33.50,189.12) and (33.50,226.88) .. (33.50,255.20) .. controls (33.50,255.20) and (151.50,257.90) .. (151.50,266.00);
\draw[->,line width=1.0bp,api,dash pattern=on 5bp off 4bp] (99.42,150.00) .. controls (99.42,158.10) and (33.50,160.80) .. (33.50,160.80) .. controls (33.50,189.12) and (33.50,226.88) .. (33.50,255.20) .. controls (33.50,255.20) and (151.50,257.90) .. (151.50,266.00);
\draw[line width=3.5bp,panel] (99.42,78.00) .. controls (99.42,86.10) and (30.00,88.80) .. (30.00,88.80) .. controls (30.00,114.56) and (30.00,148.90) .. (30.00,174.66) .. controls (30.00,174.66) and (99.42,177.36) .. (99.42,185.46);
\draw[->,line width=1.0bp,api,dash pattern=on 5bp off 4bp] (99.42,78.00) .. controls (99.42,86.10) and (30.00,88.80) .. (30.00,88.80) .. controls (30.00,114.56) and (30.00,148.90) .. (30.00,174.66) .. controls (30.00,174.66) and (99.42,177.36) .. (99.42,185.46);
\draw[line width=3.5bp,panel] (203.58,78.00) .. controls (203.58,86.10) and (26.50,88.80) .. (26.50,88.80) .. controls (26.50,139.37) and (26.50,206.79) .. (26.50,257.35) .. controls (26.50,257.35) and (524.25,260.05) .. (524.25,268.15);
\draw[->,line width=1.0bp,api,dash pattern=on 5bp off 4bp] (203.58,78.00) .. controls (203.58,86.10) and (26.50,88.80) .. (26.50,88.80) .. controls (26.50,139.37) and (26.50,206.79) .. (26.50,257.35) .. controls (26.50,257.35) and (524.25,260.05) .. (524.25,268.15);
\draw[line width=3.5bp,panel] (134.14,78.00) .. controls (134.14,86.10) and (23.00,88.80) .. (23.00,88.80) .. controls (23.00,164.89) and (23.00,266.34) .. (23.00,342.43) .. controls (23.00,342.43) and (342.00,345.13) .. (342.00,353.23);
\draw[->,line width=1.0bp,api,dash pattern=on 5bp off 4bp] (134.14,78.00) .. controls (134.14,86.10) and (23.00,88.80) .. (23.00,88.80) .. controls (23.00,164.89) and (23.00,266.34) .. (23.00,342.43) .. controls (23.00,342.43) and (342.00,345.13) .. (342.00,353.23);
\draw[line width=3.5bp,panel] (168.86,78.00) .. controls (168.86,86.10) and (19.50,88.80) .. (19.50,88.80) .. controls (19.50,192.86) and (19.50,331.60) .. (19.50,435.66) .. controls (19.50,435.66) and (398.56,438.36) .. (398.56,446.46);
\draw[->,line width=1.0bp,api,dash pattern=on 5bp off 4bp] (168.86,78.00) .. controls (168.86,86.10) and (19.50,88.80) .. (19.50,88.80) .. controls (19.50,192.86) and (19.50,331.60) .. (19.50,435.66) .. controls (19.50,435.66) and (398.56,438.36) .. (398.56,446.46);
\filldraw[fill=regularitybg,draw=regularityfg,line width=1.0bp] (353.50,557.00) ellipse (84.00bp and 30.00bp);
\node[inner sep=0pt,font=\sffamily\fontsize{10.7bp}{12bp}\selectfont\bfseries,text=ink] at (353.50,570.00) {\hyperlink{node-T5_2}{WZ1 Thm. 5.2}};
\node[inner sep=0pt,font=\sffamily\fontsize{10.7bp}{12bp}\selectfont,text=ink] at (353.50,557.00) {\hyperlink{node-T5_2}{Dot-product expansion or}};
\node[inner sep=0pt,font=\sffamily\fontsize{10.7bp}{12bp}\selectfont,text=ink] at (353.50,544.00) {\hyperlink{node-T5_2}{orthogonal lines}};
\filldraw[fill=radialbg,draw=radialfg,line width=1.0bp] (151.50,296.00) ellipse (84.00bp and 30.00bp);
\node[inner sep=0pt,font=\sffamily\fontsize{10.7bp}{12bp}\selectfont\bfseries,text=ink] at (151.50,309.00) {\hyperlink{node-L8_10}{WZ1 Lemma 8.10}};
\node[inner sep=0pt,font=\sffamily\fontsize{10.7bp}{12bp}\selectfont,text=ink] at (151.50,296.00) {\hyperlink{node-L8_10}{Two-ends localization to}};
\node[inner sep=0pt,font=\sffamily\fontsize{10.7bp}{12bp}\selectfont,text=ink] at (151.50,283.00) {\hyperlink{node-L8_10}{a strip}};
\filldraw[fill=radialbg,draw=radialfg,line width=1.0bp] (353.50,383.00) ellipse (92.73bp and 30.00bp);
\node[inner sep=0pt,font=\sffamily\fontsize{10.7bp}{12bp}\selectfont\bfseries,text=ink] at (353.50,396.00) {\hyperlink{node-P8_9}{WZ1 Prop. 8.9}};
\node[inner sep=0pt,font=\sffamily\fontsize{10.7bp}{12bp}\selectfont,text=ink] at (353.50,383.00) {\hyperlink{node-P8_9}{Separated, uniform}};
\node[inner sep=0pt,font=\sffamily\fontsize{10.7bp}{12bp}\selectfont,text=ink] at (353.50,370.00) {\hyperlink{node-P8_9}{version of the dichotomy}};
\filldraw[fill=radialbg,draw=radialfg,line width=1.0bp] (454.50,470.00) ellipse (90.22bp and 30.00bp);
\node[inner sep=0pt,font=\sffamily\fontsize{10.7bp}{12bp}\selectfont\bfseries,text=ink] at (454.50,483.00) {\hyperlink{node-T5_2p}{WZ1 Thm. 5.2'}};
\node[inner sep=0pt,font=\sffamily\fontsize{10.7bp}{12bp}\selectfont,text=ink] at (454.50,470.00) {\hyperlink{node-T5_2p}{Frostman form, without}};
\node[inner sep=0pt,font=\sffamily\fontsize{10.7bp}{12bp}\selectfont,text=ink] at (454.50,457.00) {\hyperlink{node-T5_2p}{separation or uniformity}};
\filldraw[fill=radialbg,draw=radialfg,line width=1.0bp] (151.50,209.00) ellipse (84.00bp and 30.00bp);
\node[inner sep=0pt,font=\sffamily\fontsize{10.7bp}{12bp}\selectfont\bfseries,text=ink] at (151.50,222.00) {\hyperlink{node-L8_11}{WZ1 Lemma 8.11}};
\node[inner sep=0pt,font=\sffamily\fontsize{10.7bp}{12bp}\selectfont,text=ink] at (151.50,209.00) {\hyperlink{node-L8_11}{Restore Frostman after}};
\node[inner sep=0pt,font=\sffamily\fontsize{10.7bp}{12bp}\selectfont,text=ink] at (151.50,196.00) {\hyperlink{node-L8_11}{anisotropic rescaling}};
\filldraw[fill=radialbg,draw=radialfg,line width=1.0bp] (555.50,296.00) ellipse (84.00bp and 30.00bp);
\node[inner sep=0pt,font=\sffamily\fontsize{10.7bp}{12bp}\selectfont\bfseries,text=ink] at (555.50,309.00) {\hyperlink{node-L8_13}{WZ1 Lemma 8.13}};
\node[inner sep=0pt,font=\sffamily\fontsize{10.7bp}{12bp}\selectfont,text=ink] at (555.50,296.00) {\hyperlink{node-L8_13}{Comparable orthogonal}};
\node[inner sep=0pt,font=\sffamily\fontsize{10.7bp}{12bp}\selectfont,text=ink] at (555.50,283.00) {\hyperlink{node-L8_13}{strips or expansion}};
\filldraw[fill=radialbg,draw=radialfg,line width=1.0bp] (353.50,296.00) ellipse (84.00bp and 30.00bp);
\node[inner sep=0pt,font=\sffamily\fontsize{10.7bp}{12bp}\selectfont\bfseries,text=ink] at (353.50,309.00) {\hyperlink{node-P8_5}{WZ1 Prop. 8.5}};
\node[inner sep=0pt,font=\sffamily\fontsize{10.7bp}{12bp}\selectfont,text=ink] at (353.50,296.00) {\hyperlink{node-P8_5}{Strip nonconcentration}};
\node[inner sep=0pt,font=\sffamily\fontsize{10.7bp}{12bp}\selectfont,text=ink] at (353.50,283.00) {\hyperlink{node-P8_5}{forces expansion}};
\filldraw[fill=radialbg,draw=radialfg,line width=1.0bp] (555.50,128.50) ellipse (84.00bp and 23.50bp);
\node[inner sep=0pt,font=\sffamily\fontsize{10.7bp}{12bp}\selectfont\bfseries,text=ink] at (555.50,135.00) {\hyperlink{node-L8_1}{WZ1 Lemma 8.1}};
\node[inner sep=0pt,font=\sffamily\fontsize{10.7bp}{12bp}\selectfont,text=ink] at (555.50,122.00) {\hyperlink{node-L8_1}{Hypergraph refinement}};
\filldraw[fill=radialbg,draw=radialfg,line width=1.0bp] (555.50,209.00) ellipse (84.00bp and 30.00bp);
\node[inner sep=0pt,font=\sffamily\fontsize{10.7bp}{12bp}\selectfont\bfseries,text=ink] at (555.50,222.00) {\hyperlink{node-T8_2}{WZ1 Thm. 8.2}};
\node[inner sep=0pt,font=\sffamily\fontsize{10.7bp}{12bp}\selectfont,text=ink] at (555.50,209.00) {\hyperlink{node-T8_2}{A discretized Kaufman}};
\node[inner sep=0pt,font=\sffamily\fontsize{10.7bp}{12bp}\selectfont,text=ink] at (555.50,196.00) {\hyperlink{node-T8_2}{estimate with dense}};
\filldraw[fill=radialbg,draw=radialfg,line width=1.0bp] (353.50,209.00) ellipse (84.00bp and 30.00bp);
\node[inner sep=0pt,font=\sffamily\fontsize{10.7bp}{12bp}\selectfont\bfseries,text=ink] at (353.50,222.00) {\hyperlink{node-L8_12}{WZ1 Lemma 8.12}};
\node[inner sep=0pt,font=\sffamily\fontsize{10.7bp}{12bp}\selectfont,text=ink] at (353.50,209.00) {\hyperlink{node-L8_12}{Frostman coarsening after}};
\node[inner sep=0pt,font=\sffamily\fontsize{10.7bp}{12bp}\selectfont,text=ink] at (353.50,196.00) {\hyperlink{node-L8_12}{pigeonholing}};
\filldraw[fill=apibg,draw=apifg,line width=1.0bp,rounded corners=5bp] (67.50,107.00) rectangle (235.50,150.00);
\node[inner sep=0pt,font=\sffamily\fontsize{10.7bp}{12bp}\selectfont\bfseries,text=ink] at (151.50,135.00) {\hyperlink{node-Ientropy}{Entropy API}};
\node[inner sep=0pt,font=\sffamily\fontsize{10.7bp}{12bp}\selectfont,text=ink] at (151.50,122.00) {\hyperlink{node-Ientropy}{Finite covering numbers}};
\filldraw[fill=radialbg,draw=radialfg,line width=1.0bp,rounded corners=5bp] (269.50,107.00) rectangle (437.50,150.00);
\node[inner sep=0pt,font=\sffamily\fontsize{10.7bp}{12bp}\selectfont\bfseries,text=ink] at (353.50,135.00) {\hyperlink{node-Dsep}{WZ1 Def. 8.6}};
\node[inner sep=0pt,font=\sffamily\fontsize{10.7bp}{12bp}\selectfont,text=ink] at (353.50,122.00) {\hyperlink{node-Dsep}{Standard separation}};
\filldraw[fill=grainsbg,draw=grainsfg,line width=1.0bp,rounded corners=5bp] (162.28,442.00) rectangle (330.28,498.00);
\node[inner sep=0pt,font=\sffamily\fontsize{10.7bp}{12bp}\selectfont\bfseries,text=ink] at (246.28,483.00) {\hyperlink{node-DAD}{WZ2 Sec. 6.1: ADset}};
\node[inner sep=0pt,font=\sffamily\fontsize{10.7bp}{12bp}\selectfont,text=ink] at (246.28,470.00) {\hyperlink{node-DAD}{Upper AD covering}};
\node[inner sep=0pt,font=\sffamily\fontsize{10.7bp}{12bp}\selectfont,text=ink] at (246.28,457.00) {\hyperlink{node-DAD}{condition}};
\filldraw[fill=apibg,draw=apifg,line width=1.0bp,rounded corners=5bp] (67.50,22.00) rectangle (235.50,78.00);
\node[inner sep=0pt,font=\sffamily\fontsize{10.7bp}{12bp}\selectfont\bfseries,text=ink] at (151.50,63.00) {\hyperlink{node-Iaffine}{Affine transport API}};
\node[inner sep=0pt,font=\sffamily\fontsize{10.7bp}{12bp}\selectfont,text=ink] at (151.50,50.00) {\hyperlink{node-Iaffine}{Jacobian, direction and}};
\node[inner sep=0pt,font=\sffamily\fontsize{10.7bp}{12bp}\selectfont,text=ink] at (151.50,37.00) {\hyperlink{node-Iaffine}{grid transport}};
\filldraw[fill=radialbg,draw=radialfg,line width=1.0bp,rounded corners=5bp] (471.50,22.00) rectangle (639.50,78.00);
\node[inner sep=0pt,font=\sffamily\fontsize{10.7bp}{12bp}\selectfont\bfseries,text=ink] at (555.50,63.00) {\hyperlink{node-Dhyp}{WZ1 Def.s 8.2-8.4}};
\node[inner sep=0pt,font=\sffamily\fontsize{10.7bp}{12bp}\selectfont,text=ink] at (555.50,50.00) {\hyperlink{node-Dhyp}{Hypergraph fibres and}};
\node[inner sep=0pt,font=\sffamily\fontsize{10.7bp}{12bp}\selectfont,text=ink] at (555.50,37.00) {\hyperlink{node-Dhyp}{uniform density}};
\filldraw[fill=radialbg,draw=radialfg,line width=1.0bp,rounded corners=5bp] (269.50,22.00) rectangle (437.50,78.00);
\node[inner sep=0pt,font=\sffamily\fontsize{10.7bp}{12bp}\selectfont\bfseries,text=ink] at (353.50,63.00) {\hyperlink{node-DFrost}{WZ1 Def. 8.1}};
\node[inner sep=0pt,font=\sffamily\fontsize{10.7bp}{12bp}\selectfont,text=ink] at (353.50,50.00) {\hyperlink{node-DFrost}{Normalized discretized}};
\node[inner sep=0pt,font=\sffamily\fontsize{10.7bp}{12bp}\selectfont,text=ink] at (353.50,37.00) {\hyperlink{node-DFrost}{Frostman sets}};
\end{tikzpicture}}\end{center}
{\footnotesize Solid: proof input. Dashed: API. Dash-dot: assembly. Long-dashed: overview path. Read from the top down; arrows point to consumers. Full prerequisites appear on the cards and in the ledger.}\par
\sectitle{Node catalogue: Wolff setup}
\begin{minipage}{\linewidth}
\hypertarget{node-L3_3}{}
\begin{tcolorbox}[colback=white,colframe=setupfg,colbacktitle=setupbg,coltitle=ink,boxrule=0.45pt,arc=1.6mm,left=3mm,right=3mm,top=2mm,bottom=2mm,title={\sffamily\bfseries WZ1 Lemma 3.3 / CW \hfill \texttt{L3\_3}}]
{\sffamily\small\color{muted}Local rescaling and multiplicity control | WZ1 | WZ1 Lemma 3.3, pp. 13-14; normalized CW form in the proof of WZ2 Proposition 6.2}\par
\textbf{Mathematical statement.} For an admissible unit-rescaled child family at mesh $s=\delta/\rho$, apply the WZ2 deficit lower-envelope estimate to its constant-multiplicity refinement.
Since its total shading mass is at most $O(s^2\#\mathcal T_{\rm child})$ and its union has volume at least $s^{\sigma+\varepsilon}$, its multiplicity obeys
\[\mu_{\rm fine}\lesssim s^{2-\sigma-\varepsilon}\#\mathcal T_{\rm child}.\]
The rescaled family must satisfy the multi-scale convex Wolff condition and retain the stipulated density; Lemma 2.13 and cover rebuilding supply these hypotheses.
\smallskip\textbf{Lean statement sketch.}
\begin{Verbatim}[fontsize=\small,breaklines=true,breakanywhere=true,tabsize=2]
theorem child_multiplicity_CW (he : 0 < eps) :
  ∃ eta > 0, ∃ s0 > 0, ∀ s ∈ Set.Ioc 0 s0,
  ∀ C : Config 3 s, MWAdmissible eta C → ConstantMultiplicity C →
  ∀ p, (multiplicity C p : ℝ) ≤
    K * s^(2-sigmaCW-eps) * (C.tubes.card : ℝ)
\end{Verbatim}
{\footnotesize\textbf{Mathlib support.} A, F, G, M, Q.\par
\textbf{Incoming edges.} \nodecard{W2Sigma} (core); \nodecard{W2Unit} (core); \nodecard{Imass} (api); \nodecard{Ienvelope} (core)\par}
{\footnotesize\textbf{Interface note.} This is the transferred proof interface, not the literal WZ1 formula (rho/delta)\textasciicircum{}(sigma\_n+epsilon). WZ2 may have many more than delta\textasciicircum{}(-2) tubes, so the normalized tube count must remain visible.}
\end{tcolorbox}\end{minipage}\par\medskip
\begin{minipage}{\linewidth}
\hypertarget{node-L3_4}{}
\begin{tcolorbox}[colback=white,colframe=setupfg,colbacktitle=setupbg,coltitle=ink,boxrule=0.45pt,arc=1.6mm,left=3mm,right=3mm,top=2mm,bottom=2mm,title={\sffamily\bfseries WZ1 Lemma 3.4 / CW \hfill \texttt{L3\_4}}]
{\sffamily\small\color{muted}Robust transversality | WZ1 | WZ1 Lemma 3.4, pp. 17-18; WZ2 Proposition 6.2 supplies its CW input}\par
\textbf{Mathematical statement.} For every small $\theta>0$, a sufficiently CW-extremal configuration admits a refinement with robust angular transversality.
In normalized counting form, at every shaded point $p$ and unit vector $v$,
\[\#\{T\in\mathcal T(p):\angle(v,\operatorname{dir}T)<\delta^\theta\}
\lesssim\delta^{2-\sigma+\theta\sigma/2}\#\mathcal T.\]
Use a partitioning $\delta^\theta$-cover and the fine multiplicity bound from Proposition 6.2; the participating tubes occupy $O(1)$ parents.
\smallskip\textbf{Lean statement sketch.}
\begin{Verbatim}[fontsize=\small,breaklines=true,breakanywhere=true,tabsize=2]
theorem robust_transversality_CW (ht : 0 < theta) :
  ∃ eta > 0, ∃ delta0 > 0, ∀ delta ∈ Set.Ioc 0 delta0,
  CWExtremal sigmaCW eta C →
  ∃ C', Refinement C' C ∧ ∀ p v, ‖v‖ = 1 →
    (capIncidences C' p v (delta^theta) : ℝ) ≤
      K * delta^(2-sigmaCW+theta*sigmaCW/2) * C'.tubes.card
\end{Verbatim}
{\footnotesize\textbf{Mathlib support.} F, G, M, Q.\par
\textbf{Incoming edges.} \nodecard{W2P6_2} (core)\par}
{\footnotesize\textbf{Interface note.} Constants and the small margin are absorbed by the eta/theta hierarchy. This angular-incidence conclusion does not say that the full tube family has separated directions.}
\end{tcolorbox}\end{minipage}\par\medskip
\begin{minipage}{\linewidth}
\hypertarget{node-L3_5}{}
\begin{tcolorbox}[colback=white,colframe=setupfg,colbacktitle=setupbg,coltitle=ink,boxrule=0.45pt,arc=1.6mm,left=3mm,right=3mm,top=2mm,bottom=2mm,title={\sffamily\bfseries WZ1 Lemma 3.5 / CW \hfill \texttt{L3\_5}}]
{\sffamily\small\color{muted}Mild anisotropic rescaling preserves extremality | WZ1 | WZ1 Lemma 3.5, p. 18; CW transfer used by WZ2 Sections 6.3-6.4}\par
\textbf{Mathematical statement.} Restrict a CW-extremal configuration to a prism retaining a stipulated small-power fraction of its shading mass.
Transport by a translation and diagonal dilation with all singular values between $\delta^\eta$ and $\delta^{-\eta}$.
After re-cubulation, removal of bounded multiplicities and rebuilding the cover witnesses, the transported family is $\varepsilon$-CW-extremal at a comparable power scale $\delta\le s\le\delta^{1-O(\eta)}$.
The density and convex-containment errors incur $\delta^{-O(\eta)}$ losses.
\smallskip\textbf{Lean statement sketch.}
\begin{Verbatim}[fontsize=\small,breaklines=true,breakanywhere=true,tabsize=2]
theorem mild_rescale_CW (he : 0 < eps) :
  ∃ eta > 0, ∃ delta0 > 0, ∀ delta ∈ Set.Ioc 0 delta0,
  MildCWTransportHyp eta C Q phi →
  ∃ s D, PowerComparableScale delta eta s ∧
    RecubulatedImage C Q phi D ∧ CWExtremal sigmaCW eps D
\end{Verbatim}
{\footnotesize\textbf{Mathlib support.} A, F, G, M, Q.\par
\textbf{Incoming edges.} \nodecard{W2D6_1} (core); \nodecard{Iaffine} (api); \nodecard{Iwolffbalance} (api)\par}
{\footnotesize\textbf{Interface note.} The old absolute mass threshold and automatic sticky-subfamily observation are replaced by density relative to total tube volume and explicit multi-scale cover reconstruction. This is the transport interface implicit in the author-endorsed WZ2 transfer, not a new independent black-box theorem.}
\end{tcolorbox}\end{minipage}\par\medskip
\begin{minipage}{\linewidth}
\hypertarget{node-W2CWA}{}
\begin{tcolorbox}[colback=white,colframe=setupfg,colbacktitle=setupbg,coltitle=ink,boxrule=0.45pt,arc=1.6mm,left=3mm,right=3mm,top=2mm,bottom=2mm,title={\sffamily\bfseries WZ2 Definition 2.5 \hfill \texttt{W2CWA}}]
{\sffamily\small\color{muted}Normalized convex-containment count | WZ2 | WZ2 PDF p. 8}\par
\textbf{Mathematical statement.} A finite multi-family $\mathcal S$ satisfies the convex Wolff axiom with error $A$ if
\[\#\{S\in\mathcal S:S\subset W\}\le A|W|\#\mathcal S\quad(W\text{ convex}).\]
For a tube family this implies $\#\mathcal T\gtrsim A^{-1}\delta^{-2}$ by testing $W=T$.
No upper bound $\#\mathcal T\lesssim\delta^{-2}$ is imposed.
\smallskip\textbf{Lean statement sketch.}
\begin{Verbatim}[fontsize=\small,breaklines=true,breakanywhere=true,tabsize=2]
def ConvexWolff (A : ℝ) (S : Finset (Set R3)) : Prop :=
  ∀ W, Convex ℝ W →
    ((S.filter (fun U => U ⊆ W)).card : ℝ) ≤ A * vol W * S.card
\end{Verbatim}
{\footnotesize\textbf{Mathlib support.} F, G, M, Q.\par
\textbf{Incoming edges.} \nodecard{Igeom} (api)\par}
{\footnotesize\textbf{Interface note.} The normalization is by the whole family cardinality. It differs from the traditional cardinality-normalized Wolff axiom. Tube volumes and the bounded ambient chart are kept explicit.}
\end{tcolorbox}\end{minipage}\par\medskip
\begin{minipage}{\linewidth}
\hypertarget{node-W2Unit}{}
\begin{tcolorbox}[colback=white,colframe=setupfg,colbacktitle=setupbg,coltitle=ink,boxrule=0.45pt,arc=1.6mm,left=3mm,right=3mm,top=2mm,bottom=2mm,title={\sffamily\bfseries WZ2 Definition 2.8 / Sec. 6.1 \hfill \texttt{W2Unit}}]
{\sffamily\small\color{muted}Geometric and shaded unit rescaling | WZ2 | WZ2 PDF pp. 8, 38; WZ1 Definition 3.2}\par
\textbf{Mathematical statement.} For a convex parent $U$, map its circumscribed John ellipsoid to the unit ball by $\phi_U$.
For structural tube arguments use the compatible shaded rescaling from WZ2 Section 6.1: rotate the parent axis to $e_3$, dilate transverse coordinates by $\rho^{-1}$, and re-cubulate at $\delta/\rho$.
The two choices differ by bounded normalizations; in dimension three the volume Jacobian is comparable to $\rho^{-2}$.
\smallskip\textbf{Lean statement sketch.}
\begin{Verbatim}[fontsize=\small,breaklines=true,breakanywhere=true,tabsize=2]
noncomputable def geometricUnitRescale (U S : Set R3) : Set R3 := phi U '' S
noncomputable def unitRescale (C : Config 3 delta) (P : Tube rho) :
    Config 3 (delta/rho) := recubulate (delta/rho) (transport (phi P) C)
-- Volume, containment and chart comparisons are explicit transport obligations.
\end{Verbatim}
{\footnotesize\textbf{Mathlib support.} F, G, M, Q.\par
\textbf{Incoming edges.} \nodecard{Iaffine} (api); \nodecard{Icover} (api)\par}

\end{tcolorbox}\end{minipage}\par\medskip
\begin{minipage}{\linewidth}
\hypertarget{node-W2MW}{}
\begin{tcolorbox}[colback=white,colframe=setupfg,colbacktitle=setupbg,coltitle=ink,boxrule=0.45pt,arc=1.6mm,left=3mm,right=3mm,top=2mm,bottom=2mm,title={\sffamily\bfseries WZ2 Definition 2.12 \hfill \texttt{W2MW}}]
{\sffamily\small\color{muted}Convex Wolff axioms at every scale | WZ2 | WZ2 PDF p. 9}\par
\textbf{Mathematical statement.} The tubes are essentially distinct. For every $\rho_0\in[\delta,1]$, there are $\rho\in[\rho_0,A\rho_0)$ and $\rho$-tubes $\mathcal T_\rho$ such that:
(i) $\mathcal T_\rho$ is an $A$-uniform partitioning cover of $\mathcal T$;
(ii) for each parent, the unit-rescaled child family satisfies the convex Wolff axiom with error $A$.
This is the root hypothesis; neither a direction-separated cover nor self-similar child counts is stipulated.
\smallskip\textbf{Lean statement sketch.}
\begin{Verbatim}[fontsize=\small,breaklines=true,breakanywhere=true,tabsize=2]
def MultiscaleConvexWolff (A : ℝ) (C : Config 3 delta) : Prop :=
  EssentiallyDistinct C ∧ ∀ rho0 ∈ Set.Icc delta 1,
  ∃ rho ∈ Set.Ico rho0 (A*rho0), ∃ P : Finset (Tube rho),
    UniformPartitioningCover A P C.tubes ∧
    ∀ U ∈ P, ConvexWolff A (unitRescaledChildren C U)
\end{Verbatim}
{\footnotesize\textbf{Mathlib support.} F, G, M, Q.\par
\textbf{Incoming edges.} \nodecard{Igeom} (api); \nodecard{Icover} (core); \nodecard{W2Unit} (core); \nodecard{W2CWA} (core)\par}

\end{tcolorbox}\end{minipage}\par\medskip
\begin{minipage}{\linewidth}
\hypertarget{node-W2L2_13}{}
\begin{tcolorbox}[colback=white,colframe=setupfg,colbacktitle=setupbg,coltitle=ink,boxrule=0.45pt,arc=1.6mm,left=3mm,right=3mm,top=2mm,bottom=2mm,title={\sffamily\bfseries WZ2 Lemma 2.13 \hfill \texttt{W2L2\_13}}]
{\sffamily\small\color{muted}Coarse and fine multi-scale inheritance | WZ2 | WZ2 PDF pp. 9-10}\par
\textbf{Mathematical statement.} If $\mathcal T$ satisfies the multi-scale convex Wolff axiom with error $A$, then for every $\rho_0\in[\delta,1]$ one can choose $\rho\in[\rho_0,A\rho_0]$ and a $\rho$-cover whose coarse family and every unit-rescaled child family satisfy the same multi-scale condition with error $O(A)$.
Partitioning forces the child-incidence clusters at successive scales to form a nested tree.
\smallskip\textbf{Lean statement sketch.}
\begin{Verbatim}[fontsize=\small,breaklines=true,breakanywhere=true,tabsize=2]
theorem mw_coarse_fine_inheritance (h : MultiscaleConvexWolff A C) :
  ∀ rho0 ∈ Set.Icc delta 1, ∃ rho P,
    rho0 ≤ rho ∧ rho ≤ A*rho0 ∧ CoarseFineMW (K*A) C P
\end{Verbatim}
{\footnotesize\textbf{Mathlib support.} F, G, M, Q.\par
\textbf{Incoming edges.} \nodecard{W2MW} (core); \nodecard{Icover} (api)\par}

\end{tcolorbox}\end{minipage}\par\medskip
\begin{minipage}{\linewidth}
\hypertarget{node-W2Sigma}{}
\begin{tcolorbox}[colback=white,colframe=setupfg,colbacktitle=setupbg,coltitle=ink,boxrule=0.45pt,arc=1.6mm,left=3mm,right=3mm,top=2mm,bottom=2mm,title={\sffamily\bfseries WZ2 Section 6.2 \hfill \texttt{W2Sigma}}]
{\sffamily\small\color{muted}The counterexample deficit sigma\_CW | WZ2 | WZ2 PDF p. 38}\par
\textbf{Mathematical statement.} Let $\mathsf{Bad}(s)$ mean: for every $\eta,\delta_0>0$, some $0<\delta\le\delta_0$ has multi-scale convex Wolff error $\delta^{-\eta}$, shading density $\delta^\eta$, and $|E_C|\le\delta^s$.
Define $\sigma=\sigma_{\rm CW}=\sup(\{0\}\cup\{s>0:\mathsf{Bad}(s)\})$.
Failure of Theorem 5.2 gives $\sigma>0$. This is a new deficit for a different admissible class, not WZ1's $\sigma_3$.
\smallskip\textbf{Lean statement sketch.}
\begin{Verbatim}[fontsize=\small,breaklines=true,breakanywhere=true,tabsize=2]
def BadExponent (s : ℝ) : Prop :=
  ∀ eta > 0, ∀ delta0 > 0, ∃ delta ∈ Set.Ioc 0 delta0,
  ∃ C : Config 3 delta, MWAdmissible eta C ∧ vol (union C) ≤ delta^s
noncomputable def sigmaCW : ℝ := sSup (insert 0 {s | 0 < s ∧ BadExponent s})
\end{Verbatim}
{\footnotesize\textbf{Mathlib support.} F, G, M, Q.\par
\textbf{Incoming edges.} \nodecard{Ientropy} (api); \nodecard{Imass} (api); \nodecard{Icover} (api); \nodecard{W2MW} (core); \nodecard{Imass} (core)\par}

\end{tcolorbox}\end{minipage}\par\medskip
\begin{minipage}{\linewidth}
\hypertarget{node-W2D6_1}{}
\begin{tcolorbox}[colback=white,colframe=setupfg,colbacktitle=setupbg,coltitle=ink,boxrule=0.45pt,arc=1.6mm,left=3mm,right=3mm,top=2mm,bottom=2mm,title={\sffamily\bfseries WZ2 Definition 6.1 \hfill \texttt{W2D6\_1}}]
{\sffamily\small\color{muted}Extremal configurations under convex Wolff | WZ2 | WZ2 PDF p. 39}\par
\textbf{Mathematical statement.} The pair $C=(\mathcal T,Y)_\delta$ is $\varepsilon$-extremal if
$(1)$ it has multi-scale convex Wolff error $\delta^{-\varepsilon}$;
$(2)$ its shading is $\delta^\varepsilon$-dense;
$(3)$ $|E_C|\le\delta^{\sigma-\varepsilon}$.
Density is relative to $\sum_T|T|$. All transferred WZ1 configuration nodes use this definition.
\smallskip\textbf{Lean statement sketch.}
\begin{Verbatim}[fontsize=\small,breaklines=true,breakanywhere=true,tabsize=2]
def CWExtremal (sigma eps : ℝ) (C : Config 3 delta) : Prop :=
  MultiscaleConvexWolff (delta^(-eps)) C ∧
  Dense (delta^eps) C ∧ vol (union C) ≤ delta^(sigma-eps)
\end{Verbatim}
{\footnotesize\textbf{Mathlib support.} F, G, M, Q.\par
\textbf{Incoming edges.} \nodecard{W2Sigma} (api); \nodecard{W2MW} (core); \nodecard{Imass} (core); \nodecard{W2Sigma} (core)\par}

\end{tcolorbox}\end{minipage}\par\medskip
\begin{minipage}{\linewidth}
\hypertarget{node-W2Exists}{}
\begin{tcolorbox}[colback=white,colframe=setupfg,colbacktitle=setupbg,coltitle=ink,boxrule=0.45pt,arc=1.6mm,left=3mm,right=3mm,top=2mm,bottom=2mm,title={\sffamily\bfseries WZ2 Section 6.2 / WZ1 Lem. 3.1 \hfill \texttt{W2Exists}}]
{\sffamily\small\color{muted}Extremizers exist at arbitrarily small scales | WZ2 | WZ2 paragraph after Definition 6.1, PDF p. 39}\par
\textbf{Mathematical statement.} In the contradiction regime $\sigma>0$, for every $\varepsilon,\delta_0>0$ there is an $\varepsilon$-extremal pair at some $0<\delta\le\delta_0$.
Choose a bad exponent greater than $\sigma-\varepsilon$ and use the defining arbitrarily-small-scale witness.
The scale is existential, not prescribed.
\smallskip\textbf{Lean statement sketch.}
\begin{Verbatim}[fontsize=\small,breaklines=true,breakanywhere=true,tabsize=2]
theorem CW_extremizers_exist (hs : 0 < sigmaCW) (he : 0 < eps) (hd : 0 < delta0) :
  ∃ delta ∈ Set.Ioc 0 delta0, ∃ C : Config 3 delta, CWExtremal sigmaCW eps C
\end{Verbatim}
{\footnotesize\textbf{Mathlib support.} F, G, M, Q.\par
\textbf{Incoming edges.} \nodecard{W2Sigma} (core); \nodecard{W2D6_1} (core)\par}

\end{tcolorbox}\end{minipage}\par\medskip
\begin{minipage}{\linewidth}
\hypertarget{node-W2P6_2}{}
\begin{tcolorbox}[colback=white,colframe=setupfg,colbacktitle=setupbg,coltitle=ink,boxrule=0.45pt,arc=1.6mm,left=3mm,right=3mm,top=2mm,bottom=2mm,title={\sffamily\bfseries WZ2 Proposition 6.2 \hfill \texttt{W2P6\_2}}]
{\sffamily\small\color{muted}Balanced extremal coarse and fine families | WZ2 | WZ2 PDF pp. 39-40; replaces WZ1 Proposition 3.2}\par
\textbf{Mathematical statement.} For every $\varepsilon>0$ choose $\eta,\delta_0>0$.
If $C$ is $\eta$-extremal and $\delta^{1-\varepsilon}\le\rho\le\delta^\varepsilon$, there are a refinement $C'$ and a balanced $\rho$-cover $D$ such that $D$ and each unit-rescaled child family are $\varepsilon$-extremal.
For every point $p$,
\[\begin{split}\#\mathcal T_D(p)&\le\rho^{2-\sigma-\varepsilon}\#\mathcal T_D,\\
\#\mathcal T_{C'}[U](p)&\le(\delta/\rho)^{2-\sigma-\varepsilon}\#\mathcal T_{C'}[U].\end{split}\]
The fine/coarse product comparison also supplies matching lower multiplicity estimates with a reserved small loss.
\smallskip\textbf{Lean statement sketch.}
\begin{Verbatim}[fontsize=\small,breaklines=true,breakanywhere=true,tabsize=2]
theorem CW_self_similarity (he : 0 < eps) :
  ∃ eta > 0, ∃ delta0 > 0, ∀ delta ∈ Set.Ioc 0 delta0,
  CWExtremal sigmaCW eta C → ∀ rho ∈ Set.Icc (delta^(1-eps)) (delta^eps),
  ∃ C' D, Refinement C' C ∧ BalancedCover D C' ∧
    CWExtremal sigmaCW eps D ∧ EveryChildCWExtremal sigmaCW eps C' D ∧
    NormalizedCoarseFineMultiplicity sigmaCW eps C' D
\end{Verbatim}
{\footnotesize\textbf{Mathlib support.} F, G, M, Q.\par
\textbf{Incoming edges.} \nodecard{L3_3} (core); \nodecard{W2Sigma} (core); \nodecard{W2D6_1} (core); \nodecard{Icover} (api); \nodecard{Imass} (api); \nodecard{W2L2_13} (core); \nodecard{Ienvelope} (core); \nodecard{Iwolffbalance} (api)\par}

\end{tcolorbox}\end{minipage}\par\medskip
\begin{minipage}{\linewidth}
\hypertarget{node-W2Zero}{}
\begin{tcolorbox}[colback=white,colframe=setupfg,colbacktitle=setupbg,coltitle=ink,boxrule=0.45pt,arc=1.6mm,left=3mm,right=3mm,top=2mm,bottom=2mm,title={\sffamily\bfseries WZ2 Section 6.5: contradiction \hfill \texttt{W2Zero}}]
{\sffamily\small\color{muted}The multi-scale convex Wolff deficit is zero | WZ2 | WZ2 PDF pp. 43-44}\par
\textbf{Mathematical statement.} If $\sigma>0$, apply the lower bound with exponent $\sigma/4$.
Choose the small-image construction with loss $\kappa<\sigma/4$ and $\kappa$ smaller than the lower-bound admissibility threshold.
At arbitrarily small mesh the same pair obeys
$\delta^{\sigma/4}\le |X_f(C)|\lesssim\delta^{\sigma-\kappa}$,
which is impossible. Therefore $\sigma_{\rm CW}=0$.
\smallskip\textbf{Lean statement sketch.}
\begin{Verbatim}[fontsize=\small,breaklines=true,breakanywhere=true,tabsize=2]
theorem sigmaCW_eq_zero : sigmaCW = 0 := by
  -- Assume positive deficit; combine W2L6_7 and W2P6_8 with compatible losses.
  -- The small-image exponent is strictly larger than sigmaCW / 4.
  ...
\end{Verbatim}
{\footnotesize\textbf{Mathlib support.} F, G, M, Q.\par
\textbf{Incoming edges.} \nodecard{W2L6_7} (core); \nodecard{W2Sigma} (api); \nodecard{W2P6_8} (core); \nodecard{Iparams} (api)\par}

\end{tcolorbox}\end{minipage}\par\medskip
\begin{minipage}{\linewidth}
\hypertarget{node-W2T5_2}{}
\begin{tcolorbox}[colback=white,colframe=setupfg,colbacktitle=setupbg,coltitle=ink,boxrule=0.45pt,arc=1.6mm,left=3mm,right=3mm,top=2mm,bottom=2mm,title={\sffamily\bfseries WZ2 Theorem 5.2 \hfill \texttt{W2T5\_2}}]
{\sffamily\small\color{muted}Discrete union estimate for multi-scale convex Wolff tubes | WZ2 | WZ2 PDF p. 35; proof in Sections 6-7}\par
\textbf{Mathematical statement.} For every $\varepsilon>0$ there exist $\eta,\delta_0>0$ such that for every $0<\delta\le\delta_0$:
if $\mathcal T$ satisfies the convex Wolff axioms at every scale with error $\delta^{-\eta}$ and $Y$ is $\delta^\eta$-dense, then
\[\left|\bigcup_{T\in\mathcal T}Y(T)\right|\ge\delta^\varepsilon.\]
This is the sole root theorem. No Hausdorff/Assouad-dimension deduction is part of the graph.
\smallskip\textbf{Lean statement sketch.}
\begin{Verbatim}[fontsize=\small,breaklines=true,breakanywhere=true,tabsize=2]
theorem multiscale_convex_Wolff_union (he : 0 < eps) :
  ∃ eta > 0, ∃ delta0 > 0, ∀ delta ∈ Set.Ioc 0 delta0,
  ∀ C : Config 3 delta,
  MultiscaleConvexWolff (delta^(-eta)) C → Dense (delta^eta) C →
    delta^eps ≤ vol (union C)
\end{Verbatim}
{\footnotesize\textbf{Mathlib support.} F, G, M, Q.\par
\textbf{Incoming edges.} \nodecard{W2Zero} (core); \nodecard{W2Sigma} (core); \nodecard{Igeom} (api)\par}

\end{tcolorbox}\end{minipage}\par\medskip
\sectitle{Node catalogue: Plane map}
\begin{minipage}{\linewidth}
\hypertarget{node-Dplane}{}
\begin{tcolorbox}[colback=white,colframe=planefg,colbacktitle=planebg,coltitle=ink,boxrule=0.45pt,arc=1.6mm,left=3mm,right=3mm,top=2mm,bottom=2mm,title={\sffamily\bfseries WZ2 Section 6.1: plane map \hfill \texttt{Dplane}}]
{\sffamily\small\color{muted}Plane map | WZ2 | WZ2 Section 6.1, PDF p. 38; WZ1 Definition 4.1}\par
\textbf{Mathematical statement.} A plane map $V:E_C\to S^2$ satisfies $|\operatorname{dir}(T)\cdot V(x)|\le\delta$ for every shaded incidence $x\in Y(T)$.
The normal determines the plane $V(x)^\perp$. For Lipschitz statements, use a consistently oriented normal on the selected chart, or retain a projective-plane metric.
\smallskip\textbf{Lean statement sketch.}
\begin{Verbatim}[fontsize=\small,breaklines=true,breakanywhere=true,tabsize=2]
def PlaneMap (C : Config 3 delta) (V : Rn 3 → Rn 3) : Prop :=
  (∀ x ∈ union C, ‖V x‖ = 1) ∧
  ∀ T ∈ C.tubes, ∀ x ∈ C.shade T,
    |inner ℝ (direction T) (V x)| ≤ delta
\end{Verbatim}
{\footnotesize\textbf{Mathlib support.} G, L.\par
\textbf{Incoming edges.} \nodecard{Igeom} (api)\par}
{\footnotesize\textbf{Interface note.} Changing the sign of a normal does not change its plane. An implementation should state explicitly which orientation convention is used when taking |V(p)-V(q)|.}
\end{tcolorbox}\end{minipage}\par\medskip
\begin{minipage}{\linewidth}
\hypertarget{node-T4_2}{}
\begin{tcolorbox}[colback=white,colframe=planefg,colbacktitle=planebg,coltitle=ink,boxrule=0.45pt,arc=1.6mm,left=3mm,right=3mm,top=2mm,bottom=2mm,title={\sffamily\bfseries WZ1 Theorem 4.2 \hfill \texttt{T4\_2}}]
{\sffamily\small\color{muted}Endpoint multilinear Kakeya input | WZ1 | WZ1: (4.2), p. 19; [1], version in [8]}\par
\textbf{Mathematical statement.} For a finite family of $\delta$-tubes in $\mathbb R^3$ with unit directions $v_T$,
\[
\int_{\mathbb R^3}\left(\sum_{T_1,T_2,T_3}
1_{T_1}1_{T_2}1_{T_3}\,|\det(v_{T_1},v_{T_2},v_{T_3})|\right)^{1/2}
\lesssim(\delta^2\#\mathcal T)^{3/2}.
\]
The implicit constant is absolute in the normalized tube geometry.
\smallskip\textbf{Lean statement sketch.}
\begin{Verbatim}[fontsize=\small,breaklines=true,breakanywhere=true,tabsize=2]
theorem multilinear_kakeya
  (hT : NormalizedTubes delta Ts) :
  ∫⁻ x, ENNReal.ofReal (Real.sqrt (trilinearDensity Ts x)) ∂volume
    ≤ ENNReal.ofReal (C_MLK*(delta^2*(Ts.card:ℝ))^(3/2:ℝ))
\end{Verbatim}
{\footnotesize\textbf{Mathlib support.} A, F, G, M, Q.\par
\textbf{Incoming edges.} \nodecard{XBCT} (core)\par}
{\footnotesize\textbf{Interface note.} This theorem is an imported analytic black box. The determinant is the volume of the parallelepiped, despite the source's descriptive word 'area'.}
\end{tcolorbox}\end{minipage}\par\medskip
\begin{minipage}{\linewidth}
\hypertarget{node-L4_3}{}
\begin{tcolorbox}[colback=white,colframe=planefg,colbacktitle=planebg,coltitle=ink,boxrule=0.45pt,arc=1.6mm,left=3mm,right=3mm,top=2mm,bottom=2mm,title={\sffamily\bfseries WZ1 Lemma 4.3 / CW \hfill \texttt{L4\_3}}]
{\sffamily\small\color{muted}Weak planiness | WZ1 | WZ1: (4.3), pp. 19-20; CW configuration transfer: WZ2 Sections 6.2-6.3}\par
\textbf{Mathematical statement.} There are $\eta,\delta_0>0$ such that every $\eta$-CW-extremal configuration in $\mathbb R^3$ at $0<\delta\le\delta_0$ has a refinement $C'$ and a unit-normal map $V$ with
$|V(x)\cdot\operatorname{dir}(T)|\le\delta^{\sigma/2}$ for all shaded incidences of $C'$.
Here $\sigma=\sigma_{\rm CW}$ is the WZ2 deficit.
\smallskip\textbf{Lean statement sketch.}
\begin{Verbatim}[fontsize=\small,breaklines=true,breakanywhere=true,tabsize=2]
theorem weak_planiness :
  ∃ eta > 0, ∃ delta0 > 0, ∀ delta ∈ Set.Ioc 0 delta0,
  CWExtremal sigmaCW eta C →
  ∃ C' V, Refinement C' C ∧ UnitNormalsOn V (union C') ∧
    ∀ T ∈ C'.tubes, ∀ x ∈ C'.shade T,
      |inner ℝ (V x) (direction T)| ≤ delta^(sigma/2)
\end{Verbatim}
{\footnotesize\textbf{Mathlib support.} F, G, M, Q.\par
\textbf{Incoming edges.} \nodecard{T4_2} (core); \nodecard{L3_4} (core); \nodecard{Imass} (api); \nodecard{Ienvelope} (core)\par}
{\footnotesize\textbf{Interface note.}  Ported configuration statement: use WZ2 Definition 6.1 and Proposition 6.2, not WZ1 admissibility or an absolute tube-count normalization. WZ2 Sections 6.2-6.4 authorize this transfer; the numerical count adjustments are displayed on the setup cards.}
\end{tcolorbox}\end{minipage}\par\medskip
\begin{minipage}{\linewidth}
\hypertarget{node-L4_4}{}
\begin{tcolorbox}[colback=white,colframe=planefg,colbacktitle=planebg,coltitle=ink,boxrule=0.45pt,arc=1.6mm,left=3mm,right=3mm,top=2mm,bottom=2mm,title={\sffamily\bfseries WZ1 Lemma 4.4 / CW \hfill \texttt{L4\_4}}]
{\sffamily\small\color{muted}Coarsen weak planiness to a plane map | WZ1 | WZ1: (4.10), pp. 20-21; CW configuration transfer: WZ2 Sections 6.2-6.3}\par
\textbf{Mathematical statement.} For each $\varepsilon>0$, sufficiently extremal $C$ and $\rho\in[\delta^{1-\varepsilon},\delta^\varepsilon]$ with $|V(x)\cdot\operatorname{dir}T|\le\rho/2$ admit a refinement $C'$ and an $\varepsilon$-extremal $\rho$-cover $D$ with plane map $\widetilde V$.
The map $\widetilde V$ is constant on $\rho$-cubes and equals some original value $V(p_Q)$ in each occupied cube.
\smallskip\textbf{Lean statement sketch.}
\begin{Verbatim}[fontsize=\small,breaklines=true,breakanywhere=true,tabsize=2]
theorem coarsen_plane_map (he : 0 < eps) :
  ∃ eta > 0, ∃ delta0 > 0, ∀ delta ∈ Set.Ioc 0 delta0,
  CoarsenPlaneHyp sigma eta eps rho C V →
  ∃ C' D W, Refinement C' C ∧ Covers D C' ∧
    CWExtremal sigmaCW eps D ∧ PlaneMap D W ∧
    CubeConstant rho W ∧ CubewiseConsistent V W D
\end{Verbatim}
{\footnotesize\textbf{Mathlib support.} F, G, M, Q.\par
\textbf{Incoming edges.} \nodecard{W2P6_2} (core); \nodecard{Dplane} (api); \nodecard{L4_3} (assembly)\par}
{\footnotesize\textbf{Interface note.}  Ported configuration statement: use WZ2 Definition 6.1 and Proposition 6.2, not WZ1 admissibility or an absolute tube-count normalization. WZ2 Sections 6.2-6.4 authorize this transfer; the numerical count adjustments are displayed on the setup cards.}
\end{tcolorbox}\end{minipage}\par\medskip
\begin{minipage}{\linewidth}
\hypertarget{node-L4_5}{}
\begin{tcolorbox}[colback=white,colframe=planefg,colbacktitle=planebg,coltitle=ink,boxrule=0.45pt,arc=1.6mm,left=3mm,right=3mm,top=2mm,bottom=2mm,title={\sffamily\bfseries WZ1 Lemma 4.5 / CW \hfill \texttt{L4\_5}}]
{\sffamily\small\color{muted}Small oscillation on a fixed cube scale | WZ1 | WZ1: pp. 21-22; CW configuration transfer: WZ2 Sections 6.2-6.3}\par
\textbf{Mathematical statement.} For every $\varepsilon>0$, sufficiently extremal $C$ with plane map $V$ and middle scale $\rho$ has an $\varepsilon$-CW-extremal subcollection $C'$ such that
$|V(p)-V(q)|\le\rho$ whenever $p,q\in E_{C'}$ lie in a common $\rho$-cube.
\smallskip\textbf{Lean statement sketch.}
\begin{Verbatim}[fontsize=\small,breaklines=true,breakanywhere=true,tabsize=2]
theorem plane_oscillation_on_cubes (he : 0 < eps) :
  ∃ eta > 0, ∃ delta0 > 0, ∀ delta ∈ Set.Ioc 0 delta0,
  PlaneOscillationHyp sigma eta eps rho C V →
  ∃ C', Subcollection C' C ∧ CWExtremal sigmaCW eps C' ∧
    ∀ p q ∈ union C', SameCube rho p q → ‖V p-V q‖ ≤ rho
\end{Verbatim}
{\footnotesize\textbf{Mathlib support.} F, G, L, M.\par
\textbf{Incoming edges.} \nodecard{L3_4} (core); \nodecard{W2P6_2} (core); \nodecard{Dplane} (api)\par}
{\footnotesize\textbf{Interface note.}  Ported configuration statement: use WZ2 Definition 6.1 and Proposition 6.2, not WZ1 admissibility or an absolute tube-count normalization. WZ2 Sections 6.2-6.4 authorize this transfer; the numerical count adjustments are displayed on the setup cards.}
\end{tcolorbox}\end{minipage}\par\medskip
\begin{minipage}{\linewidth}
\hypertarget{node-C4_6}{}
\begin{tcolorbox}[colback=white,colframe=planefg,colbacktitle=planebg,coltitle=ink,boxrule=0.45pt,arc=1.6mm,left=3mm,right=3mm,top=2mm,bottom=2mm,title={\sffamily\bfseries WZ1 Corollary 4.6 / CW \hfill \texttt{C4\_6}}]
{\sffamily\small\color{muted}Small oscillation for nearby points | WZ1 | WZ1: p. 22; CW configuration transfer: WZ2 Sections 6.2-6.3}\par
\textbf{Mathematical statement.} Under the same hypotheses, an $\varepsilon$-CW-extremal subcollection satisfies
$|p-q|\le\rho\Longrightarrow |V(p)-V(q)|\le2\rho$ on its union.
The proof discards a constant fraction of coarse cubes to enforce separation.
\smallskip\textbf{Lean statement sketch.}
\begin{Verbatim}[fontsize=\small,breaklines=true,breakanywhere=true,tabsize=2]
theorem plane_oscillation_nearby (he : 0 < eps) :
  ∃ eta > 0, ∃ delta0 > 0, ∀ delta ∈ Set.Ioc 0 delta0,
  PlaneOscillationHyp sigma eta eps rho C V →
  ∃ C', Subcollection C' C ∧ CWExtremal sigmaCW eps C' ∧
    ∀ p q ∈ union C', dist p q ≤ rho → dist (V p) (V q) ≤ 2*rho
\end{Verbatim}
{\footnotesize\textbf{Mathlib support.} F, G, L, M.\par
\textbf{Incoming edges.} \nodecard{L4_5} (core)\par}
{\footnotesize\textbf{Interface note.}  Ported configuration statement: use WZ2 Definition 6.1 and Proposition 6.2, not WZ1 admissibility or an absolute tube-count normalization. WZ2 Sections 6.2-6.4 authorize this transfer; the numerical count adjustments are displayed on the setup cards.}
\end{tcolorbox}\end{minipage}\par\medskip
\begin{minipage}{\linewidth}
\hypertarget{node-L4_7}{}
\begin{tcolorbox}[colback=white,colframe=planefg,colbacktitle=planebg,coltitle=ink,boxrule=0.45pt,arc=1.6mm,left=3mm,right=3mm,top=2mm,bottom=2mm,title={\sffamily\bfseries WZ1 Lemma 4.7 / CW \hfill \texttt{L4\_7}}]
{\sffamily\small\color{muted}Iterate to a Lipschitz plane map | WZ1 | WZ1: pp. 22-23; CW configuration transfer: WZ2 Sections 6.2-6.3}\par
\textbf{Mathematical statement.} A sufficiently CW-extremal configuration whose plane map $V$ is constant on $\delta$-cubes has an $\varepsilon$-CW-extremal subcollection on which $V$ is $\delta^{-\varepsilon}$-Lipschitz.
This iterates Corollary 4.6 at finitely many power scales.
\smallskip\textbf{Lean statement sketch.}
\begin{Verbatim}[fontsize=\small,breaklines=true,breakanywhere=true,tabsize=2]
theorem lipschitz_plane_refinement (he : 0 < eps) :
  ∃ eta > 0, ∃ delta0 > 0, ∀ delta ∈ Set.Ioc 0 delta0,
  CWExtremal sigmaCW eta C → PlaneMap C V → CubeConstant delta V →
  ∃ C', Subcollection C' C ∧ CWExtremal sigmaCW eps C' ∧
    LipschitzOnWith (Real.toNNReal (delta^(-eps))) V (union C')
\end{Verbatim}
{\footnotesize\textbf{Mathlib support.} F, G, L, M, Q.\par
\textbf{Incoming edges.} \nodecard{C4_6} (core); \nodecard{Iparams} (api); \nodecard{L4_4} (assembly)\par}
{\footnotesize\textbf{Interface note.}  Ported configuration statement: use WZ2 Definition 6.1 and Proposition 6.2, not WZ1 admissibility or an absolute tube-count normalization. WZ2 Sections 6.2-6.4 authorize this transfer; the numerical count adjustments are displayed on the setup cards.}
\end{tcolorbox}\end{minipage}\par\medskip
\begin{minipage}{\linewidth}
\hypertarget{node-L4_8}{}
\begin{tcolorbox}[colback=white,colframe=planefg,colbacktitle=planebg,coltitle=ink,boxrule=0.45pt,arc=1.6mm,left=3mm,right=3mm,top=2mm,bottom=2mm,title={\sffamily\bfseries WZ1 Lemma 4.8 / CW \hfill \texttt{L4\_8}}]
{\sffamily\small\color{muted}Existence of Lipschitz plane data | WZ1 | WZ1: p. 23; CW configuration transfer: WZ2 Sections 6.2-6.3}\par
\textbf{Mathematical statement.} For every $\varepsilon,\delta_0>0$, some $0<\delta\le\delta_0$ admits an $\varepsilon$-CW-extremal configuration with a $\delta^{-\varepsilon}$-Lipschitz plane map.
Assemble extremizer existence, weak planiness, coarsening and Lipschitz refinement, using WZ2 Proposition 6.2 throughout.
\smallskip\textbf{Lean statement sketch.}
\begin{Verbatim}[fontsize=\small,breaklines=true,breakanywhere=true,tabsize=2]
theorem exists_lipschitz_plane_data (he : 0 < eps) (hd : 0 < delta0) :
  ∃ delta ∈ Set.Ioc 0 delta0, ∃ C : Config 3 delta, ∃ V,
    CWExtremal sigmaCW eps C ∧ PlaneMap C V ∧
    LipschitzOnWith (Real.toNNReal (delta^(-eps))) V (union C)
\end{Verbatim}
{\footnotesize\textbf{Mathlib support.} F, G, L, M, Q.\par
\textbf{Incoming edges.} \nodecard{W2Exists} (core); \nodecard{L4_3} (core); \nodecard{L4_4} (core); \nodecard{L4_7} (core)\par}
{\footnotesize\textbf{Interface note.}  Ported configuration statement: use WZ2 Definition 6.1 and Proposition 6.2, not WZ1 admissibility or an absolute tube-count normalization. WZ2 Sections 6.2-6.4 authorize this transfer; the numerical count adjustments are displayed on the setup cards.}
\end{tcolorbox}\end{minipage}\par\medskip
\sectitle{Node catalogue: Grains}
\begin{minipage}{\linewidth}
\hypertarget{node-DAD}{}
\begin{tcolorbox}[colback=white,colframe=grainsfg,colbacktitle=grainsbg,coltitle=ink,boxrule=0.45pt,arc=1.6mm,left=3mm,right=3mm,top=2mm,bottom=2mm,title={\sffamily\bfseries WZ2 Section 6.1: ADset \hfill \texttt{DAD}}]
{\sffamily\small\color{muted}Upper AD covering condition | WZ2 | WZ2 Section 6.1, PDF p. 38; WZ1 Definition 4.2}\par
\textbf{Mathematical statement.} For $E\subset\mathbb R$, the condition $(\delta,\alpha,A)_1$-ADset means
$\mathcal E_\rho(E\cap I)\le A(r/\rho)^\alpha$ for every $\rho\ge\delta$ and interval $I$ of length $r\ge\rho$.
Only an upper covering bound is required. It passes to arbitrary subsets.
The higher-dimensional ball version from WZ1 is used within its Section 8 projection argument.
\smallskip\textbf{Lean statement sketch.}
\begin{Verbatim}[fontsize=\small,breaklines=true,breakanywhere=true,tabsize=2]
def ADSet (delta alpha C : ℝ) (E : Set (Rn n)) : Prop :=
  ∀ x rho r, delta ≤ rho → rho ≤ r →
    (coverNum rho (E ∩ Metric.ball x r) : ℝ) ≤ C*(r/rho)^alpha
theorem ADSet.mono (h : ADSet delta alpha C E) (hS : S ⊆ E) :
  ADSet delta alpha C S
\end{Verbatim}
{\footnotesize\textbf{Mathlib support.} F, G, Q.\par
\textbf{Incoming edges.} \nodecard{Ientropy} (api)\par}
{\footnotesize\textbf{Interface note.} This is the convention recalled in WZ2, not two-sided classical AD regularity. Keep it distinct from normalized Frostman sets and absolute Katz-Tao sets.}
\end{tcolorbox}\end{minipage}\par\medskip
\begin{minipage}{\linewidth}
\hypertarget{node-L4_9}{}
\begin{tcolorbox}[colback=white,colframe=grainsfg,colbacktitle=grainsbg,coltitle=ink,boxrule=0.45pt,arc=1.6mm,left=3mm,right=3mm,top=2mm,bottom=2mm,title={\sffamily\bfseries WZ1 Lemma 4.9 / CW \hfill \texttt{L4\_9}}]
{\sffamily\small\color{muted}Grain count in tau-cubes | WZ1 | WZ1: (4.15), pp. 23-25; CW configuration transfer: WZ2 Sections 6.2-6.3}\par
\textbf{Mathematical statement.} For sufficiently extremal $C$ with a $\delta^{-\eta}$-Lipschitz plane map, $\rho\in[\delta^{1-\varepsilon},\delta^\varepsilon]$ and $\tau\in[\rho^{1-\varepsilon},\rho^{1/2}]$, an $\varepsilon$-CW-extremal subcollection $C'$ satisfies
$\mathcal E_\rho(V(Q)\cdot(E_{C'}\cap Q))\le\rho^{-\varepsilon}(\tau/\rho)^{1-\sigma}$
for every $\tau$-cube $Q$. Here $V(Q)$ is evaluated at its centre using the chosen extension.
\smallskip\textbf{Lean statement sketch.}
\begin{Verbatim}[fontsize=\small,breaklines=true,breakanywhere=true,tabsize=2]
theorem grains_in_cubes (he : 0 < eps) :
  ∃ eta > 0, ∃ delta0 > 0, ∀ delta ∈ Set.Ioc 0 delta0,
  GrainCountHyp sigma eta eps rho tau C V →
  ∃ C', Subcollection C' C ∧ CWExtremal sigmaCW eps C' ∧
    ∀ Q, GridCube tau Q →
      coverNumR rho (normalImage V Q (union C' ∩ Q)) ≤
        rho^(-eps)*(tau/rho)^(1-sigma)
\end{Verbatim}
{\footnotesize\textbf{Mathlib support.} F, G, L, M, P, Q.\par
\textbf{Incoming edges.} \nodecard{W2P6_2} (core); \nodecard{L3_4} (core); \nodecard{DAD} (core); \nodecard{Imass} (api); \nodecard{L4_8} (assembly)\par}
{\footnotesize\textbf{Interface note.} The proof's Cordoba-type L2 estimate is elementary tube-intersection geometry and Cauchy-Schwarz. It is not a separate imported Cordoba theorem. Ported configuration statement: use WZ2 Definition 6.1 and Proposition 6.2, not WZ1 admissibility or an absolute tube-count normalization. WZ2 Sections 6.2-6.4 authorize this transfer; the numerical count adjustments are displayed on the setup cards.}
\end{tcolorbox}\end{minipage}\par\medskip
\begin{minipage}{\linewidth}
\hypertarget{node-L4_10}{}
\begin{tcolorbox}[colback=white,colframe=grainsfg,colbacktitle=grainsbg,coltitle=ink,boxrule=0.45pt,arc=1.6mm,left=3mm,right=3mm,top=2mm,bottom=2mm,title={\sffamily\bfseries WZ1 Lemma 4.10 / CW \hfill \texttt{L4\_10}}]
{\sffamily\small\color{muted}Grain count in thin prisms | WZ1 | WZ1: (4.20), pp. 26-28; CW configuration transfer: WZ2 Sections 6.2-6.3}\par
\textbf{Mathematical statement.} Under the same scale and plane-map hypotheses, an $\varepsilon$-CW-extremal subcollection satisfies:
for every $\rho^{1/2}$-cube $Q$, put $F_Q=V(Q)\cdot(E_{C'}\cap Q)$.
For each interval $J$ of length $\tau$,
$\mathcal E_\rho(J\cap F_Q)\le\rho^{-\varepsilon}(\tau/\rho)^{1-\sigma}$.
This is a localized interval bound, stronger than only counting the entire projection.
\smallskip\textbf{Lean statement sketch.}
\begin{Verbatim}[fontsize=\small,breaklines=true,breakanywhere=true,tabsize=2]
theorem grains_in_prisms (he : 0 < eps) :
  ∃ eta > 0, ∃ delta0 > 0, ∀ delta ∈ Set.Ioc 0 delta0,
  GrainCountHyp sigma eta eps rho tau C V →
  ∃ C', Subcollection C' C ∧ CWExtremal sigmaCW eps C' ∧
    ∀ Q J, GridCube (Real.sqrt rho) Q → IntervalLength J tau →
      coverNumR rho (J ∩ normalImage V Q (union C' ∩ Q)) ≤
        rho^(-eps)*(tau/rho)^(1-sigma)
\end{Verbatim}
{\footnotesize\textbf{Mathlib support.} F, G, L, M, Q.\par
\textbf{Incoming edges.} \nodecard{L4_9} (core); \nodecard{W2P6_2} (core)\par}
{\footnotesize\textbf{Interface note.}  Ported configuration statement: use WZ2 Definition 6.1 and Proposition 6.2, not WZ1 admissibility or an absolute tube-count normalization. WZ2 Sections 6.2-6.4 authorize this transfer; the numerical count adjustments are displayed on the setup cards.}
\end{tcolorbox}\end{minipage}\par\medskip
\begin{minipage}{\linewidth}
\hypertarget{node-L4_11}{}
\begin{tcolorbox}[colback=white,colframe=grainsfg,colbacktitle=grainsbg,coltitle=ink,boxrule=0.45pt,arc=1.6mm,left=3mm,right=3mm,top=2mm,bottom=2mm,title={\sffamily\bfseries WZ1 Lemma 4.11 / CW \hfill \texttt{L4\_11}}]
{\sffamily\small\color{muted}AD spacing at one grain scale | WZ1 | WZ1: pp. 28-29; CW configuration transfer: WZ2 Sections 6.2-6.3}\par
\textbf{Mathematical statement.} For sufficiently extremal $C$ with a Lipschitz plane map and any middle scale $\rho$, an $\varepsilon$-CW-extremal subcollection satisfies:
$V(Q)\cdot(E_{C'}\cap Q)$ is a $(\rho,1-\sigma,\delta^{-\varepsilon})_1$-ADset for every $\rho^{1/2}$-cube $Q$.
Iterate Lemma 4.10 over the interval-length parameter $\tau$.
\smallskip\textbf{Lean statement sketch.}
\begin{Verbatim}[fontsize=\small,breaklines=true,breakanywhere=true,tabsize=2]
theorem grains_AD_one_scale (he : 0 < eps) :
  ∃ eta > 0, ∃ delta0 > 0, ∀ delta ∈ Set.Ioc 0 delta0,
  OneScaleGrainHyp sigma eta eps rho C V →
  ∃ C', Subcollection C' C ∧ CWExtremal sigmaCW eps C' ∧
    ∀ Q, GridCube (Real.sqrt rho) Q →
      ADSet rho (1-sigma) (delta^(-eps))
        (normalImage V Q (union C' ∩ Q))
\end{Verbatim}
{\footnotesize\textbf{Mathlib support.} F, G, L, Q.\par
\textbf{Incoming edges.} \nodecard{L4_10} (core); \nodecard{Iparams} (api)\par}
{\footnotesize\textbf{Interface note.}  Ported configuration statement: use WZ2 Definition 6.1 and Proposition 6.2, not WZ1 admissibility or an absolute tube-count normalization. WZ2 Sections 6.2-6.4 authorize this transfer; the numerical count adjustments are displayed on the setup cards.}
\end{tcolorbox}\end{minipage}\par\medskip
\begin{minipage}{\linewidth}
\hypertarget{node-L4_12}{}
\begin{tcolorbox}[colback=white,colframe=grainsfg,colbacktitle=grainsbg,coltitle=ink,boxrule=0.45pt,arc=1.6mm,left=3mm,right=3mm,top=2mm,bottom=2mm,title={\sffamily\bfseries WZ1 Lemma 4.12 / CW \hfill \texttt{L4\_12}}]
{\sffamily\small\color{muted}AD spacing at every grain scale | WZ1 | WZ1: (4.29), pp. 29-30; CW configuration transfer: WZ2 Sections 6.2-6.3}\par
\textbf{Mathematical statement.} For sufficiently extremal $C$ with a $\delta^{-\eta}$-Lipschitz plane map, an $\varepsilon$-CW-extremal subcollection satisfies:
$V(p)\cdot(E_{C'}\cap B(p,\rho^{1/2}))$ is a $(\rho,1-\sigma,\delta^{-\varepsilon})_1$-ADset for all $\rho\in[\delta,1]$ and all centres $p$.
Iterate Lemma 4.11 over $\rho$, then compare intervening scales.
\smallskip\textbf{Lean statement sketch.}
\begin{Verbatim}[fontsize=\small,breaklines=true,breakanywhere=true,tabsize=2]
theorem grains_AD_all_scales (he : 0 < eps) :
  ∃ eta > 0, ∃ delta0 > 0, ∀ delta ∈ Set.Ioc 0 delta0,
  LipschitzPlaneHyp sigma eta C V →
  ∃ C', Subcollection C' C ∧ CWExtremal sigmaCW eps C' ∧
    ∀ rho ∈ Set.Icc delta 1, ∀ p,
      ADSet rho (1-sigma) (delta^(-eps))
        (dotImage (V p) (union C' ∩ Metric.ball p (Real.sqrt rho)))
\end{Verbatim}
{\footnotesize\textbf{Mathlib support.} F, G, L, Q.\par
\textbf{Incoming edges.} \nodecard{L4_11} (core); \nodecard{Iparams} (api)\par}
{\footnotesize\textbf{Interface note.} Plane values outside the union need an explicit extension or a harmless centre-selection convention. The original exposition suppresses this domain bookkeeping. Ported configuration statement: use WZ2 Definition 6.1 and Proposition 6.2, not WZ1 admissibility or an absolute tube-count normalization. WZ2 Sections 6.2-6.4 authorize this transfer; the numerical count adjustments are displayed on the setup cards.}
\end{tcolorbox}\end{minipage}\par\medskip
\begin{minipage}{\linewidth}
\hypertarget{node-W2P6_3}{}
\begin{tcolorbox}[colback=white,colframe=grainsfg,colbacktitle=grainsbg,coltitle=ink,boxrule=0.45pt,arc=1.6mm,left=3mm,right=3mm,top=2mm,bottom=2mm,title={\sffamily\bfseries WZ2 Proposition 6.3 \hfill \texttt{W2P6\_3}}]
{\sffamily\small\color{muted}Global and local grains with Lipschitz data | WZ2 | WZ2 PDF p. 40; transferred WZ1 Proposition 4.1}\par
\textbf{Mathematical statement.} For every $\varepsilon,\delta_0>0$, some $\varepsilon$-extremal pair at $0<\delta\le\delta_0$ has:
(1) a $1$-Lipschitz $f:[-1,1]\to\mathbb R$ with each fixed-$z$ image $\{x+f(z)y:(x,y,z)\in E_C\}$ an $(\delta,1-\sigma,\delta^{-\varepsilon})_1$-ADset;
(2) a $1$-Lipschitz plane map $V$ with $V(p)\cdot(E_C\cap B(p,\sqrt\rho))$ an $(\rho,1-\sigma,\delta^{-\varepsilon})_1$-ADset for all $p$ and $\rho\in[\delta,1]$.
\smallskip\textbf{Lean statement sketch.}
\begin{Verbatim}[fontsize=\small,breaklines=true,breakanywhere=true,tabsize=2]
theorem exists_Lipschitz_grains_CW (he : 0 < eps) (hd : 0 < delta0) :
  ∃ delta ∈ Set.Ioc 0 delta0, ∃ C f V,
    CWExtremal sigmaCW eps C ∧ LipschitzGrainsCW sigmaCW eps C f V
\end{Verbatim}
{\footnotesize\textbf{Mathlib support.} F, G, M, Q.\par
\textbf{Incoming edges.} \nodecard{L4_8} (core); \nodecard{L4_12} (core); \nodecard{W2P6_2} (core); \nodecard{W2Unit} (core); \nodecard{L4_3} (core); \nodecard{L4_4} (core); \nodecard{L3_5} (core); \nodecard{XExt} (core); \nodecard{L4_7} (assembly)\par}
{\footnotesize\textbf{Interface note.} The entire WZ1 Section 4 proof is expanded in the plane and grain views. Proposition 6.2 replaces Proposition 3.2; numerical multiplicities are normalized by cardinality.}
\end{tcolorbox}\end{minipage}\par\medskip
\sectitle{Node catalogue: C2 regularity}
\begin{minipage}{\linewidth}
\hypertarget{node-Iadgap}{}
\begin{tcolorbox}[colback=white,colframe=regularityfg,colbacktitle=regularitybg,coltitle=ink,boxrule=0.45pt,arc=1.6mm,left=3mm,right=3mm,top=2mm,bottom=2mm,title={\sffamily\bfseries AD exponent gap \hfill \texttt{Iadgap}}]
{\sffamily\small\color{muted}Why the expansion alternative is impossible | API | WZ1: Proof of Lemma 5.3, p. 37; discussion p. 5}\par
\textbf{Mathematical statement.} If $A$ is an $(r_0,1-\sigma,C)$-ADset, then for $\rho\ge r_0$ and $R=|I|/\rho$,
$\mathcal E_\rho(A\cap I)\le C R^{1-\sigma}$.
An expansion lower bound $\mathcal E_\rho(A\cap I)\ge R^{1-\varepsilon}$ forces $R^{\sigma-\varepsilon}\le C$.
Choose $\varepsilon<\sigma$ and make the permitted interval-to-scale ratio exceed this threshold.
\smallskip\textbf{Lean statement sketch.}
\begin{Verbatim}[fontsize=\small,breaklines=true,breakanywhere=true,tabsize=2]
theorem ad_expansion_incompatible
  (hAD : ADSet r0 (1-sigma) C A)
  (hexpand : (len I / rho)^(1-eps) ≤ coverNumR rho (A ∩ I))
  (hscale : r0 ≤ rho) :
  (len I / rho)^(sigma-eps) ≤ C
\end{Verbatim}
{\footnotesize\textbf{Mathlib support.} F, Q.\par
\textbf{Incoming edges.} \nodecard{DAD} (api)\par}
{\footnotesize\textbf{Interface note.} The direct proof is an exponent comparison. The introduction attributes a related observation to [12, Proposition 6.13]; that citation is not invoked as a separate black box in the written proof of Lemma 5.3.}
\end{tcolorbox}\end{minipage}\par\medskip
\begin{minipage}{\linewidth}
\hypertarget{node-T5_2}{}
\begin{tcolorbox}[colback=white,colframe=regularityfg,colbacktitle=regularitybg,coltitle=ink,boxrule=0.45pt,arc=1.6mm,left=3mm,right=3mm,top=2mm,bottom=2mm,title={\sffamily\bfseries WZ1 Theorem 5.2 \hfill \texttt{T5\_2}}]
{\sffamily\small\color{muted}Dot-product expansion or orthogonal lines | WZ1 | WZ1: (5.1)-(5.2), p. 32; proof Section 8}\par
\textbf{Mathematical statement.} For every $\varepsilon>0$, choose $\eta,\delta_0>0$.
For $(\delta,1,\delta^{-\eta})$-ADsets $F,G\subset[0,1]^2$, either
(A) an affine line $\ell$ and the orthogonal line $\ell^\perp$ through the origin have
$\mathcal E_\delta(F\cap N_\delta\ell),\mathcal E_\delta(G\cap N_\delta\ell^\perp)\ge\delta^{\varepsilon-1}$; or
(B) for every $H\subset F\times G\times G$ with $\mathcal E_\delta(H)\ge\delta^{\eta-3}$, some $\rho\ge\delta$ and interval $I$, $|I|\ge\delta^{-\eta}\rho$, satisfy
$\mathcal E_\rho(I\cap\{a\cdot(b_1-b_2):(a,b_1,b_2)\in H\})\ge(|I|/\rho)^{1-\varepsilon}$.
\smallskip\textbf{Lean statement sketch.}
\begin{Verbatim}[fontsize=\small,breaklines=true,breakanywhere=true,tabsize=2]
theorem dot_product_AD_dichotomy (he : 0 < eps) :
  ∃ eta > 0, ∃ delta0 > 0, ∀ delta ∈ Set.Ioc 0 delta0,
  ADPairHyp delta eta F G →
    OrthogonalLineConcentration delta eps F G G ∨
    ∀ H, HyperSubset H F G G →
      delta^(eta-3) ≤ coverNumR delta H →
      ExpansionConclusion delta eta eps (dotDiffImage H)
\end{Verbatim}
{\footnotesize\textbf{Mathlib support.} F, G, Q.\par
\textbf{Incoming edges.} \nodecard{T5_2p} (core); \nodecard{DFrost} (core); \nodecard{DAD} (core)\par}
{\footnotesize\textbf{Interface note.} The second alternative is universal in H. It is not the weaker assertion that a specially selected H expands.}
\end{tcolorbox}\end{minipage}\par\medskip
\begin{minipage}{\linewidth}
\hypertarget{node-L5_3}{}
\begin{tcolorbox}[colback=white,colframe=regularityfg,colbacktitle=regularitybg,coltitle=ink,boxrule=0.45pt,arc=1.6mm,left=3mm,right=3mm,top=2mm,bottom=2mm,title={\sffamily\bfseries WZ1 Lemma 5.3 \hfill \texttt{L5\_3}}]
{\sffamily\small\color{muted}Local linearity from the two grain structures | WZ1 | WZ1: (5.3)-(5.4), pp. 33-38; CW configuration transfer: WZ2 Sections 6.2-6.3}\par
\textbf{Mathematical statement.} For suitable $\eta,\rho_0$, let $I$ have length $\rho^{1/2}$, $Z\subset I$, and let $f:I\to[0,1]$ and $V:E\to S^2$ be $1$-Lipschitz.
Assume $E\subset[-1,1]^2\times Z$ is a $\rho$-cube union lying in a $\rho^{1/2}$-neighbourhood of a line with direction $(-f(z_*),1,0)$, and $|E|\ge\rho^{1+\sigma/2+\eta}$.
Assume $|V\cdot e_3|\le1/2$, the global slice images are $(\rho,1-\sigma,\rho^{-\eta})_1$-ADsets, and
$|V(p)\cdot(E\cap B(p,\rho^{1/2}))|\le\rho^{(1+\sigma)/2-\eta}$.
Then an affine function $L:I\to[0,1]$ and $Z_{\rm lin}\subset Z$ satisfy
$|Z_{\rm lin}|\ge\rho^{1/2+\varepsilon}$ and $|f(z)-L(z)|\le\rho$ on $Z_{\rm lin}$.
\smallskip\textbf{Lean statement sketch.}
\begin{Verbatim}[fontsize=\small,breaklines=true,breakanywhere=true,tabsize=2]
theorem local_linearity (he : 0 < eps) :
  ∃ eta > 0, ∃ rho0 > 0, ∀ rho ∈ Set.Ioc 0 rho0,
  LocalLinearityHyp sigma eta rho I Z f E V →
  ∃ Zlin ⊆ Z, ∃ L : ℝ → ℝ, AffineOn I L ∧
    MapsTo L I (Set.Icc 0 1) ∧
    rho^(1/2+eps) ≤ vol Zlin ∧
    ∀ z ∈ Zlin, |f z-L z| ≤ rho
\end{Verbatim}
{\footnotesize\textbf{Mathlib support.} F, G, L, M, Q.\par
\textbf{Incoming edges.} \nodecard{T5_2} (core); \nodecard{Iadgap} (api); \nodecard{W2P6_3} (assembly)\par}
{\footnotesize\textbf{Interface note.} LocalLinearityHyp is exactly the preceding hypothesis list. The proof encodes closed grain paths as dot products, applies Theorem 5.2, and rules out expansion by the AD exponent gap.}
\end{tcolorbox}\end{minipage}\par\medskip
\begin{minipage}{\linewidth}
\hypertarget{node-L5_4}{}
\begin{tcolorbox}[colback=white,colframe=regularityfg,colbacktitle=regularitybg,coltitle=ink,boxrule=0.45pt,arc=1.6mm,left=3mm,right=3mm,top=2mm,bottom=2mm,title={\sffamily\bfseries WZ1 Lemma 5.4 / CW \hfill \texttt{L5\_4}}]
{\sffamily\small\color{muted}Piecewise linear approximation at one scale | WZ1 | WZ1: pp. 38-39; CW configuration transfer: WZ2 Sections 6.2-6.3}\par
\textbf{Mathematical statement.} For a sufficiently CW-extremal configuration with the two grain properties of Proposition 4.1 and middle scale $\rho$, there are an $\varepsilon$-CW-extremal subcollection $C'$ and a $2$-Lipschitz function $L$.
The function $L$ is affine on each interval $[k\rho^{1/2},(k+1)\rho^{1/2}]$ and satisfies
$|f(z)-L(z)|\le\rho$ for $z\in\pi_z(E_{C'})$.
\smallskip\textbf{Lean statement sketch.}
\begin{Verbatim}[fontsize=\small,breaklines=true,breakanywhere=true,tabsize=2]
theorem piecewise_linear_one_scale (he : 0 < eps) :
  ∃ eta > 0, ∃ delta0 > 0, ∀ delta ∈ Set.Ioc 0 delta0,
  LipschitzGrainsHyp sigma eta eps rho C f V →
  ∃ C' L, Subcollection C' C ∧ CWExtremal sigmaCW eps C' ∧
    LipschitzOnWith 2 L (Set.Icc (-1) 1) ∧
    PiecewiseAffine (Real.sqrt rho) L ∧
    ∀ z ∈ zImage (union C'), |f z-L z| ≤ rho
\end{Verbatim}
{\footnotesize\textbf{Mathlib support.} F, G, L, M, Q.\par
\textbf{Incoming edges.} \nodecard{L5_3} (core); \nodecard{W2P6_2} (core); \nodecard{Icover} (api)\par}
{\footnotesize\textbf{Interface note.}  Ported configuration statement: use WZ2 Definition 6.1 and Proposition 6.2, not WZ1 admissibility or an absolute tube-count normalization. WZ2 Sections 6.2-6.4 authorize this transfer; the numerical count adjustments are displayed on the setup cards.}
\end{tcolorbox}\end{minipage}\par\medskip
\begin{minipage}{\linewidth}
\hypertarget{node-L5_5}{}
\begin{tcolorbox}[colback=white,colframe=regularityfg,colbacktitle=regularitybg,coltitle=ink,boxrule=0.45pt,arc=1.6mm,left=3mm,right=3mm,top=2mm,bottom=2mm,title={\sffamily\bfseries WZ1 Lemma 5.5 / CW \hfill \texttt{L5\_5}}]
{\sffamily\small\color{muted}Compatible approximations at many scales | WZ1 | WZ1: pp. 39-40; CW configuration transfer: WZ2 Sections 6.2-6.3}\par
\textbf{Mathematical statement.} For $\varepsilon>0$ and integer $N\ge1$, let $\rho_j=\delta^{j/N}$.
A sufficiently CW-extremal configuration with Proposition 4.1 data has an $\varepsilon$-CW-extremal subcollection and functions $L_j$, each $2$-Lipschitz and affine on intervals of length $\rho_j^{1/2}$, with $|f-L_j|\le\rho_j$ on $D=\pi_z(E_{C'})$.
The subsequent refinement on p.~40 also ensures each occupied interval has $|I\cap D|\ge4^{-N}\delta^\varepsilon |I|$, and occupied intervals are non-adjacent at each scale.
\smallskip\textbf{Lean statement sketch.}
\begin{Verbatim}[fontsize=\small,breaklines=true,breakanywhere=true,tabsize=2]
theorem piecewise_linear_many_scales (he : 0 < eps) (hN : 1 ≤ N) :
  ∃ eta > 0, ∃ delta0 > 0, ∀ delta ∈ Set.Ioc 0 delta0,
  LipschitzGrainsCW sigmaCW eta C f V →
  ∃ C' Ls, Subcollection C' C ∧ CWExtremal sigmaCW eps C' ∧
    ∀ j ∈ Finset.Icc 1 N,
      LinearApproximationAt (delta^((j:ℝ)/N)) f Ls[j] (zImage (union C'))
\end{Verbatim}
{\footnotesize\textbf{Mathlib support.} F, L, M, Q.\par
\textbf{Incoming edges.} \nodecard{L5_4} (core); \nodecard{Imass} (api); \nodecard{Iparams} (api)\par}
{\footnotesize\textbf{Interface note.} The popularity and separation conclusions are the explicit unnumbered refinement after Lemma 5.5; Corollary 5.6 uses this strengthened version. The approximation is only on D, not on the whole interval. Ported configuration statement: use WZ2 Definition 6.1 and Proposition 6.2, not WZ1 admissibility or an absolute tube-count normalization. WZ2 Sections 6.2-6.4 authorize this transfer; the numerical count adjustments are displayed on the setup cards.}
\end{tcolorbox}\end{minipage}\par\medskip
\begin{minipage}{\linewidth}
\hypertarget{node-Dtrap}{}
\begin{tcolorbox}[colback=white,colframe=regularityfg,colbacktitle=regularitybg,coltitle=ink,boxrule=0.45pt,arc=1.6mm,left=3mm,right=3mm,top=2mm,bottom=2mm,title={\sffamily\bfseries WZ1 Definition 5.1 \hfill \texttt{Dtrap}}]
{\sffamily\small\color{muted}Vertical trapezoids and parents | WZ1 | WZ1: p. 40}\par
\textbf{Mathematical statement.} A vertical trapezoid is $R=\{(z,t):z\in J,\ |t-L(z)|\le s\}$, with affine centre line $L$.
Its length is $|J|$, height is $s$, and slope is the slope of $L$.
Across scales, a parent relation means geometric containment of the child trapezoid in the unique parent.
\smallskip\textbf{Lean statement sketch.}
\begin{Verbatim}[fontsize=\small,breaklines=true,breakanywhere=true,tabsize=2]
structure VerticalTrapezoid where
  interval : Set ℝ
  slope intercept height : ℝ
def VerticalTrapezoid.carrier (R : VerticalTrapezoid) :=
  {p : ℝ × ℝ | p.1 ∈ R.interval ∧
    |p.2-(R.slope*p.1+R.intercept)| ≤ R.height}
def Parent (R S : VerticalTrapezoid) :=
  R.carrier ⊆ S.carrier
\end{Verbatim}
{\footnotesize\textbf{Mathlib support.} F, G.\par
\textbf{Incoming edges.} \nodecard{Ientropy} (api)\par}

\end{tcolorbox}\end{minipage}\par\medskip
\begin{minipage}{\linewidth}
\hypertarget{node-C5_6}{}
\begin{tcolorbox}[colback=white,colframe=regularityfg,colbacktitle=regularitybg,coltitle=ink,boxrule=0.45pt,arc=1.6mm,left=3mm,right=3mm,top=2mm,bottom=2mm,title={\sffamily\bfseries WZ1 Corollary 5.6 / CW \hfill \texttt{C5\_6}}]
{\sffamily\small\color{muted}A nested trapezoid tree for the slope graph | WZ1 | WZ1: pp. 40-41; CW configuration transfer: WZ2 Sections 6.2-6.3}\par
\textbf{Mathematical statement.} With the hypotheses of Lemma 5.5, a suitable $\varepsilon$-CW-extremal subcollection has trapezoid families $\mathcal R_j$ with $\rho_j=\delta^{j/N}$:
slopes lie in $[-2,2]$; height is $\rho_j$; lengths lie in $[\rho_j^{1/2+\varepsilon},\rho_j^{1/2}]$; horizontal projections are $\rho_j^{1/2}$-separated; every child has a unique containing parent; and $\operatorname{graph}(f|_D)\subset\bigcup_{R\in\mathcal R_N}R$.
\smallskip\textbf{Lean statement sketch.}
\begin{Verbatim}[fontsize=\small,breaklines=true,breakanywhere=true,tabsize=2]
theorem slope_trapezoid_tree (he : 0 < eps) (hN : 1 ≤ N) :
  ∃ eta > 0, ∃ delta0 > 0, ∀ delta ∈ Set.Ioc 0 delta0,
  LipschitzGrainsCW sigmaCW eta C f V →
  ∃ C' Rs, Subcollection C' C ∧ CWExtremal sigmaCW eps C' ∧
    TrapezoidTree delta eps N Rs ∧
    graphOn f (zImage (union C')) ⊆ lastLevelUnion Rs
\end{Verbatim}
{\footnotesize\textbf{Mathlib support.} F, G, L, Q.\par
\textbf{Incoming edges.} \nodecard{L5_5} (core); \nodecard{Dtrap} (core)\par}
{\footnotesize\textbf{Interface note.}  Ported configuration statement: use WZ2 Definition 6.1 and Proposition 6.2, not WZ1 admissibility or an absolute tube-count normalization. WZ2 Sections 6.2-6.4 authorize this transfer; the numerical count adjustments are displayed on the setup cards.}
\end{tcolorbox}\end{minipage}\par\medskip
\begin{minipage}{\linewidth}
\hypertarget{node-L5_8}{}
\begin{tcolorbox}[colback=white,colframe=regularityfg,colbacktitle=regularitybg,coltitle=ink,boxrule=0.45pt,arc=1.6mm,left=3mm,right=3mm,top=2mm,bottom=2mm,title={\sffamily\bfseries WZ1 Lemma 5.8 \hfill \texttt{L5\_8}}]
{\sffamily\small\color{muted}Smooth extension of an affine segment | WZ1 | WZ1: pp. 41-42; CW configuration transfer: WZ2 Sections 6.2-6.3}\par
\textbf{Mathematical statement.} For $x_1<x_2$, $a>0$ and the segment from $(x_1,y_1)$ to $(x_2,y_2)$, there is a $C^2$ function $G:\mathbb R\to\mathbb R$ whose graph contains that segment and which vanishes outside $[x_1-a,x_2+a]$.
With $h=|y_1|+|y_2|$ and $s=(y_2-y_1)/(x_2-x_1)$,
$|G'|\lesssim h/a+|s|$ and $|G''|\lesssim h/a^2+|s|/a$.
\smallskip\textbf{Lean statement sketch.}
\begin{Verbatim}[fontsize=\small,breaklines=true,breakanywhere=true,tabsize=2]
theorem smooth_affine_patch (hx : x1 < x2) (ha : 0 < a) :
  ∃ G : ℝ → ℝ, ContDiff ℝ 2 G ∧
    EqOn G (fun x => s*(x-x1)+y1) (Set.Icc x1 x2) ∧
    support G ⊆ Set.Icc (x1-a) (x2+a) ∧
    ∀ x, |deriv G x| ≤ C_patch*(h/a+|s|) ∧
      |deriv (deriv G) x| ≤ C_patch*(h/a^2+|s|/a)
\end{Verbatim}
{\footnotesize\textbf{Mathlib support.} C, G, Q.\par
\textbf{Incoming edges.} \nodecard{Icalc} (api)\par}

\end{tcolorbox}\end{minipage}\par\medskip
\begin{minipage}{\linewidth}
\hypertarget{node-P5_7}{}
\begin{tcolorbox}[colback=white,colframe=regularityfg,colbacktitle=regularitybg,coltitle=ink,boxrule=0.45pt,arc=1.6mm,left=3mm,right=3mm,top=2mm,bottom=2mm,title={\sffamily\bfseries WZ1 Proposition 5.7 \hfill \texttt{P5\_7}}]
{\sffamily\small\color{muted}A trapezoid tree produces a C2 approximation | WZ1 | WZ1: pp. 41-43}\par
\textbf{Mathematical statement.} Let $N\ge2$ and $\eta>0$. Suppose the trapezoid families satisfy the same bounds and nesting as Corollary 5.6, with $\eta$ in place of $\varepsilon$, and $\operatorname{graph}(f|_D)$ lies in the last-level union.
Then some $C^2$ function $g$ satisfies $|f-g|\le\delta$ on $D$,
$|g''|\lesssim\delta^{-1/N-2\eta}$, and for $x\in\pi_z(R)$ with $R\in\mathcal R_j$,
$|g'(x)-\operatorname{slope}(R)|\lesssim\rho_j/\rho_{j+1}^{1/2+\eta}$.
\smallskip\textbf{Lean statement sketch.}
\begin{Verbatim}[fontsize=\small,breaklines=true,breakanywhere=true,tabsize=2]
theorem C2_from_trapezoid_tree
  (hN : 2 ≤ N) (htree : TrapezoidTree delta eta N Rs)
  (hgraph : graphOn f D ⊆ lastLevelUnion Rs) :
  ∃ g, C2Near g (Set.Icc (-1) 1) ∧
    (∀ x ∈ D, |f x-g x| ≤ delta) ∧
    (∀ x ∈ Set.Icc 0 1, |deriv (deriv g) x| ≤ C_tree*delta^(-1/(N:ℝ)-2*eta)) ∧
    TreeDerivativeApproximation delta eta N Rs g
\end{Verbatim}
{\footnotesize\textbf{Mathlib support.} C, F, G, Q.\par
\textbf{Incoming edges.} \nodecard{L5_8} (core); \nodecard{Dtrap} (core); \nodecard{Icalc} (api)\par}
{\footnotesize\textbf{Interface note.} The source labels the proof 'Proof of Lemma 5.7', but the numbered statement is Proposition 5.7. The derivative comparison uses j<N; at the last level the tail of the patch sum is zero. The sketch preserves the printed [0,1] domain of the displayed second-derivative bound.}
\end{tcolorbox}\end{minipage}\par\medskip
\begin{minipage}{\linewidth}
\hypertarget{node-W2P6_4}{}
\begin{tcolorbox}[colback=white,colframe=regularityfg,colbacktitle=regularitybg,coltitle=ink,boxrule=0.45pt,arc=1.6mm,left=3mm,right=3mm,top=2mm,bottom=2mm,title={\sffamily\bfseries WZ2 Proposition 6.4 \hfill \texttt{W2P6\_4}}]
{\sffamily\small\color{muted}Global grains with a C2 slope | WZ2 | WZ2 PDF p. 40; transferred WZ1 Proposition 5.1}\par
\textbf{Mathematical statement.} For every $\varepsilon,\delta_0>0$, an $\varepsilon$-extremal pair at some $0<\delta\le\delta_0$ has global grains with $\|f\|_{C^2([-1,1])}\le1$ and retains the same local grains with a $1$-Lipschitz plane map.
The WZ1 Section 5 construction uses its Section 8 dot-product projection theorem; that entire internal branch remains in this atlas.
\smallskip\textbf{Lean statement sketch.}
\begin{Verbatim}[fontsize=\small,breaklines=true,breakanywhere=true,tabsize=2]
theorem exists_C2_grains_CW (he : 0 < eps) (hd : 0 < delta0) :
  ∃ delta ∈ Set.Ioc 0 delta0, ∃ C f V,
    CWExtremal sigmaCW eps C ∧ C2GrainsCW sigmaCW eps C f V
\end{Verbatim}
{\footnotesize\textbf{Mathlib support.} F, G, M, Q.\par
\textbf{Incoming edges.} \nodecard{W2P6_3} (core); \nodecard{C5_6} (core); \nodecard{P5_7} (core); \nodecard{L3_5} (core)\par}

\end{tcolorbox}\end{minipage}\par\medskip
\sectitle{Node catalogue: Nonsingular slope}
\begin{minipage}{\linewidth}
\hypertarget{node-L6_2}{}
\begin{tcolorbox}[colback=white,colframe=slopefg,colbacktitle=slopebg,coltitle=ink,boxrule=0.45pt,arc=1.6mm,left=3mm,right=3mm,top=2mm,bottom=2mm,title={\sffamily\bfseries WZ1 Lemma 6.2 \hfill \texttt{L6\_2}}]
{\sffamily\small\color{muted}A thin projection forces near orthogonality | WZ1 | WZ1: (6.1)-(6.4), p. 44}\par
\textbf{Mathematical statement.} Let $0\le\alpha<1$, $\delta\le\rho\le1$, and let $U$ be a length-$\rho^{1/2}$ segment of a $\delta$-tube $T$.
If $F\subset U$, $|F|\ge\delta^\varepsilon|U|$, and $v\cdot F$ is a $(\rho,\alpha,\delta^{-\varepsilon})_1$-ADset for a unit $v$, then
$|v\cdot\operatorname{dir}T|\lesssim\delta^{-2\varepsilon/(1-\alpha)}\rho^{1/2}$.
\smallskip\textbf{Lean statement sketch.}
\begin{Verbatim}[fontsize=\small,breaklines=true,breakanywhere=true,tabsize=2]
theorem projection_forces_orthogonality
  (halpha : 0 ≤ alpha ∧ alpha < 1)
  (hF : F ⊆ U) (hm : delta^eps*vol U ≤ vol F)
  (hAD : ADSet rho alpha (delta^(-eps)) (dotImage v F)) :
  |inner ℝ v (direction T)| ≤
    C_ortho*delta^(-2*eps/(1-alpha))*Real.sqrt rho
\end{Verbatim}
{\footnotesize\textbf{Mathlib support.} F, G, M, Q.\par
\textbf{Incoming edges.} \nodecard{DAD} (core)\par}

\end{tcolorbox}\end{minipage}\par\medskip
\begin{minipage}{\linewidth}
\hypertarget{node-L6_4}{}
\begin{tcolorbox}[colback=white,colframe=slopefg,colbacktitle=slopebg,coltitle=ink,boxrule=0.45pt,arc=1.6mm,left=3mm,right=3mm,top=2mm,bottom=2mm,title={\sffamily\bfseries WZ1 Lemma 6.4 / CW \hfill \texttt{L6\_4}}]
{\sffamily\small\color{muted}A thick slab with comparable nonzero derivative | WZ1 | WZ1 Lemma 6.4, pp. 47-48; WZ2 proof of Proposition 6.5 uses Lemma 6.6 instead of WZ1 Lemma 6.3}\par
\textbf{Mathematical statement.} For every $\varepsilon,\delta_0>0$, some $0<\delta\le\delta_0$ admits a configuration satisfying Proposition 5.1 and with all shadings in $[-1,1]^2\times[0,\delta^\varepsilon]$.
There is $m\ge\delta^\varepsilon$ such that $m\le|f'(z)|\le2m$ on $[0,\delta^\varepsilon]$.
\smallskip\textbf{Lean statement sketch.}
\begin{Verbatim}[fontsize=\small,breaklines=true,breakanywhere=true,tabsize=2]
theorem exists_comparable_slope_slab (he : 0 < eps) (hd : 0 < delta0) :
  ∃ delta ∈ Set.Ioc 0 delta0, ∃ C : Config 3 delta, ∃ f V m,
    C2GrainsCW sigmaCW eps C f V ∧
    ShadingsInSlab C (Set.Icc 0 (delta^eps)) ∧
    delta^eps ≤ m ∧
    ∀ z ∈ Set.Icc 0 (delta^eps), m ≤ |deriv f z| ∧ |deriv f z| ≤ 2*m
\end{Verbatim}
{\footnotesize\textbf{Mathlib support.} C, F, G, M, Q.\par
\textbf{Incoming edges.} \nodecard{W2P6_4} (core); \nodecard{W2P6_2} (core); \nodecard{W2L6_6} (core)\par}
{\footnotesize\textbf{Interface note.} Remove the initial direction-separated subcollection selection in the WZ1 proof. The popular slab is selected using the normalized coarse counts from WZ2 Proposition 6.2, and nonflatness is supplied by WZ2 Lemma 6.6. No direction selection node survives.}
\end{tcolorbox}\end{minipage}\par\medskip
\begin{minipage}{\linewidth}
\hypertarget{node-Dtwist}{}
\begin{tcolorbox}[colback=white,colframe=slopefg,colbacktitle=slopebg,coltitle=ink,boxrule=0.45pt,arc=1.6mm,left=3mm,right=3mm,top=2mm,bottom=2mm,title={\sffamily\bfseries WZ1 Def. 6.1 / WZ2 Sec. 6.5 \hfill \texttt{Dtwist}}]
{\sffamily\small\color{muted}Twisted projection | WZ1 | WZ1 p. 48; WZ2 PDF p. 43}\par
\textbf{Mathematical statement.} For a slope function $f$, set $\pi_f(x,y,z)=(x+yf(z),z)$.
For a line $(a,b,0)+t(c,d,1)$, its image has graph
$t\mapsto a+ct+(b+dt)f(t)$.
Global AD slice bounds and Fubini give $|\pi_f(E_C)|\lesssim\delta^{\sigma-\varepsilon}$ for the configuration from Proposition 6.1, after the usual exponent allowance.
\smallskip\textbf{Lean statement sketch.}
\begin{Verbatim}[fontsize=\small,breaklines=true,breakanywhere=true,tabsize=2]
def twistedProjection (f : ℝ → ℝ) (p : Rn 3) : Rn 2 :=
  ![p 0 + p 1*f (p 2), p 2]
def twistedLine (f : ℝ → ℝ) (a b c d : ℝ) (t : ℝ) :=
  a+c*t+(b+d*t)*f t
theorem twisted_union_volume_bound
  (h : GlobalGrains sigma (delta^(-eps)) C f) :
  vol (twistedProjection f '' union C) ≤ C_slice*delta^(sigma-eps)
\end{Verbatim}
{\footnotesize\textbf{Mathlib support.} C, G, M, Q.\par
\textbf{Incoming edges.} \nodecard{Igeom} (api); \nodecard{Icalc} (api)\par}
{\footnotesize\textbf{Interface note.} The coordinate literal is schematic for EuclideanSpace; a real implementation wraps the vector literal in the appropriate WithLp constructor.}
\end{tcolorbox}\end{minipage}\par\medskip
\begin{minipage}{\linewidth}
\hypertarget{node-W1SlabPrep}{}
\begin{tcolorbox}[colback=white,colframe=slopefg,colbacktitle=slopebg,coltitle=ink,boxrule=0.45pt,arc=1.6mm,left=3mm,right=3mm,top=2mm,bottom=2mm,title={\sffamily\bfseries WZ1 Lemma 6.3: Steps 1-3 \hfill \texttt{W1SlabPrep}}]
{\sffamily\small\color{muted}Popular grains in a thin horizontal slab | WZ1 | WZ1 pp. 44-46; summarized by WZ2 PDF pp. 41-42}\par
\textbf{Mathematical statement.} Assume $|J|=\sqrt\rho$ and that $f'$ is small at a point of the popular slab.
After normalizing $f(0)=0$, bounded $f''$ gives $|f|\lesssim\rho$ on $J$.
Refine to a set $F_2$ with dense global grains, select a popular $y$-slice, and cover its $(x,z)$-projection by $\lesssim\delta^{-O(\eta)}\rho^{-1/2+\sigma/2}$ squares $S$ of side $\sqrt\rho$.
For $P_S=S\times\mathbb R_y$ the projection $(F_2\cap P_S)\cdot\widetilde V(p_S)$ is a $(\rho,1-\sigma,\delta^{-O(\eta)})_1$-ADset, where $\widetilde V=(V_x,0,V_z)$ has norm bounded below.
\smallskip\textbf{Lean statement sketch.}
\begin{Verbatim}[fontsize=\small,breaklines=true,breakanywhere=true,tabsize=2]
theorem popular_slab_preparation_CW (h : FlatPopularSlabHyp sigmaCW eta C f V J) :
  ∃ F2 squares, SlabPrepData sigmaCW eta rho C F2 squares V
\end{Verbatim}
{\footnotesize\textbf{Mathlib support.} F, G, M, Q.\par
\textbf{Incoming edges.} \nodecard{W2P6_4} (assembly); \nodecard{DAD} (core); \nodecard{Imass} (api)\par}
{\footnotesize\textbf{Interface note.} A grouped unnumbered construction inside the old lemma, not a second invocation of its direction-separated conclusion. The local and global grain hypotheses are outputs from WZ2 Proposition 6.4.}
\end{tcolorbox}\end{minipage}\par\medskip
\begin{minipage}{\linewidth}
\hypertarget{node-W2SlabCount}{}
\begin{tcolorbox}[colback=white,colframe=slopefg,colbacktitle=slopebg,coltitle=ink,boxrule=0.45pt,arc=1.6mm,left=3mm,right=3mm,top=2mm,bottom=2mm,title={\sffamily\bfseries WZ2 Lemma 6.6: new Step 4 \hfill \texttt{W2SlabCount}}]
{\sffamily\small\color{muted}Convex Wolff contradicts a popular thin slab | WZ2 | WZ2 equations (6.8)-(6.10), PDF pp. 42-43}\par
\textbf{Mathematical statement.} The popular-square argument gives a family $\mathcal T_S$ with
$\#\mathcal T_S\gtrsim\delta^{O(\eta)}\rho^{(1-\sigma)/2}\#\mathcal T$.
WZ1 Lemma 6.2 makes its directions nearly orthogonal to $\widetilde V(p_S)$; since all tubes meet $P_S$, their full tubes lie in one slab $\Sigma_S$ of thickness $\delta^{-O(\eta)}\sqrt\rho$.
Convex Wolff gives $\#\mathcal T_S\lesssim\delta^{-O(\eta)}\sqrt\rho\#\mathcal T$.
Thus $\rho^{-\sigma/2}\lesssim\delta^{-O(\eta)}$, impossible when $\sqrt\rho=|J|\le\delta^\varepsilon$ and $\eta\ll\varepsilon\sigma$.
\smallskip\textbf{Lean statement sketch.}
\begin{Verbatim}[fontsize=\small,breaklines=true,breakanywhere=true,tabsize=2]
theorem popular_slab_CW_contradiction
  (hprep : SlabPrepData sigmaCW eta rho C F2 squares V)
  (hCW : ConvexWolff (delta^(-eta)) C.tubeSets)
  (hsmall : SlopeLossHierarchy sigmaCW eps eta delta rho) : False
\end{Verbatim}
{\footnotesize\textbf{Mathlib support.} F, G, M, Q.\par
\textbf{Incoming edges.} \nodecard{W1SlabPrep} (core); \nodecard{L6_2} (core); \nodecard{W2CWA} (core)\par}
{\footnotesize\textbf{Interface note.} This replaces the parallel-direction pigeonhole contradiction in WZ1. It uses a normalized convex slab count, so the common factor \#T cancels; no separated directions or upper tube cardinality is used.}
\end{tcolorbox}\end{minipage}\par\medskip
\begin{minipage}{\linewidth}
\hypertarget{node-W2L6_6}{}
\begin{tcolorbox}[colback=white,colframe=slopefg,colbacktitle=slopebg,coltitle=ink,boxrule=0.45pt,arc=1.6mm,left=3mm,right=3mm,top=2mm,bottom=2mm,title={\sffamily\bfseries WZ2 Lemma 6.6 \hfill \texttt{W2L6\_6}}]
{\sffamily\small\color{muted}The slope cannot be flat on a popular slab | WZ2 | WZ2 PDF pp. 41-43; replaces WZ1 Lemma 6.3}\par
\textbf{Mathematical statement.} Let $C,f,V$ satisfy Proposition 6.4 with loss $\eta$. Let $\delta^{1/2}\le|J|\le\delta^\varepsilon$.
Assume $\sum_T|Y(T)\cap(\mathbb R^2\times J)|\ge\delta^\eta|J|$ and that the shaded union in this slab is covered by at most $\delta^{-\eta}|J|^{-2+\sigma}$ cubes of side $|J|$.
For sufficiently small $\eta,\delta_0$ one has $|f'(z)|\ge|J|$ for every $z\in J$.
\smallskip\textbf{Lean statement sketch.}
\begin{Verbatim}[fontsize=\small,breaklines=true,breakanywhere=true,tabsize=2]
theorem nonflat_popular_slab_CW (he : 0 < eps) :
  ∃ eta > 0, ∃ delta0 > 0, ∀ delta ∈ Set.Ioc 0 delta0,
  C2GrainsCW sigmaCW eta C f V → PopularSlabHyp sigmaCW eta eps C J →
    ∀ z ∈ J, intervalLength J ≤ |deriv f z|
\end{Verbatim}
{\footnotesize\textbf{Mathlib support.} F, G, M, Q.\par
\textbf{Incoming edges.} \nodecard{L6_2} (core); \nodecard{Icalc} (api); \nodecard{W2P6_4} (assembly); \nodecard{W2SlabCount} (core); \nodecard{W1SlabPrep} (core)\par}

\end{tcolorbox}\end{minipage}\par\medskip
\begin{minipage}{\linewidth}
\hypertarget{node-W2P6_5}{}
\begin{tcolorbox}[colback=white,colframe=slopefg,colbacktitle=slopebg,coltitle=ink,boxrule=0.45pt,arc=1.6mm,left=3mm,right=3mm,top=2mm,bottom=2mm,title={\sffamily\bfseries WZ2 Proposition 6.5, part 1 double-prime \hfill \texttt{W2P6\_5}}]
{\sffamily\small\color{muted}Nonsingular C2 global grains | WZ2 | WZ2 PDF pp. 41, 43; assembly from WZ1 Proposition 6.1}\par
\textbf{Mathematical statement.} For every $\varepsilon,\delta_0>0$, an $\varepsilon$-extremal pair at some $0<\delta\le\delta_0$ has global grains for a $C^2$ slope satisfying
\[1\le|f'(z)|\le2,\qquad |f''(z)|\le1/100\quad(-1\le z\le1).\]
The local grains and $1$-Lipschitz plane map are retained.
Use the transferred comparable-slope slab lemma, WZ2 Lemma 6.6 and mild anisotropic normalization.
\smallskip\textbf{Lean statement sketch.}
\begin{Verbatim}[fontsize=\small,breaklines=true,breakanywhere=true,tabsize=2]
def NonsingularSlope (f : ℝ → ℝ) (I : Set ℝ) : Prop :=
  C2Near f I ∧ ∀ z ∈ I,
    1 ≤ |deriv f z| ∧ |deriv f z| ≤ 2 ∧ |deriv (deriv f) z| ≤ 1/100
theorem exists_nonsingular_grains_CW (he : 0 < eps) (hd : 0 < delta0) :
  ∃ delta ∈ Set.Ioc 0 delta0, ∃ C f V,
    CWExtremal sigmaCW eps C ∧ NonsingularSlope f (Set.Icc (-1) 1) ∧
    GlobalGrains sigmaCW (delta^(-eps)) C f ∧ LocalPlaneDataCW sigmaCW eps C V
\end{Verbatim}
{\footnotesize\textbf{Mathlib support.} F, G, M, Q.\par
\textbf{Incoming edges.} \nodecard{L6_4} (core); \nodecard{L3_5} (core); \nodecard{Iaffine} (api); \nodecard{W2P6_4} (assembly)\par}

\end{tcolorbox}\end{minipage}\par\medskip
\begin{minipage}{\linewidth}
\hypertarget{node-W2L6_7}{}
\begin{tcolorbox}[colback=white,colframe=slopefg,colbacktitle=slopebg,coltitle=ink,boxrule=0.45pt,arc=1.6mm,left=3mm,right=3mm,top=2mm,bottom=2mm,title={\sffamily\bfseries WZ2 Lemma 6.7: small-image interface \hfill \texttt{W2L6\_7}}]
{\sffamily\small\color{muted}Extremality produces a small twisted image | WZ2 | WZ2 PDF p. 43, equation (6.11); WZ1 Corollary 6.5 ancestry}\par
\textbf{Mathematical statement.} For every small loss $\varepsilon>0$ and cutoff $\delta_0>0$, obtain a convex-Wolff family with error $\delta^{-\varepsilon}$, shading density $\delta^\varepsilon$, and a nonsingular slope such that
\[|X_f(C)|\lesssim\delta^{\sigma-\varepsilon},\qquad
X_f(C)=\bigcup_T\pi_f(Y(T)).\]
Indeed, the global AD slice bound gives $\mathcal E_\delta(X_f(C)\cap\{z=t\})\le\delta^{-\varepsilon}\delta^{-(1-\sigma)}$; multiply by $O(\delta)$ and integrate in $t$.
\smallskip\textbf{Lean statement sketch.}
\begin{Verbatim}[fontsize=\small,breaklines=true,breakanywhere=true,tabsize=2]
theorem small_twisted_image_CW (he : 0 < eps) (hd : 0 < delta0) :
  ∃ delta ∈ Set.Ioc 0 delta0, ∃ C f,
    ConvexWolff (delta^(-eps)) C.tubeSets ∧ Dense (delta^eps) C ∧
    NonsingularSlope f (Set.Icc (-1) 1) ∧
    vol (twistedImage C f) ≤ K * delta^(sigmaCW-eps)
\end{Verbatim}
{\footnotesize\textbf{Mathlib support.} F, G, M, Q.\par
\textbf{Incoming edges.} \nodecard{W2P6_5} (core); \nodecard{Dtwist} (core); \nodecard{W2CWA} (api); \nodecard{W2MW} (api); \nodecard{DAD} (core)\par}
{\footnotesize\textbf{Interface note.} Equation (6.11) in the supplied WZ2 PDF prints sigma+epsilon. The direct consequence of the displayed AD bound is sigma-epsilon (with harmless constant losses). This weaker, explicitly derived statement is sufficient for the contradiction and is the one used in all three roadmap formats.}
\end{tcolorbox}\end{minipage}\par\medskip
\sectitle{Node catalogue: Twisted projection}
\begin{minipage}{\linewidth}
\hypertarget{node-Dcinematic}{}
\begin{tcolorbox}[colback=white,colframe=cinematicfg,colbacktitle=cinematicbg,coltitle=ink,boxrule=0.45pt,arc=1.6mm,left=3mm,right=3mm,top=2mm,bottom=2mm,title={\sffamily\bfseries WZ1 Definition 7.1 \hfill \texttt{Dcinematic}}]
{\sffamily\small\color{muted}Cinematic families of C2 functions | WZ1 | WZ1: (7.2), p. 49; [36, Definition 1.6]}\par
\textbf{Mathematical statement.} A family $\mathcal F\subset C^2(I)$ is cinematic with constants $K,D$ if its diameter in the $C^2$ norm is at most $K$, it is a doubling metric space with constant $D$, and for every $g,h\in\mathcal F$,
$\inf_{t\in I}(|g-h|+|g'-h'|+|g''-h''|)(t)\ge K^{-1}\|g-h\|_{C^2(I)}$.
\smallskip\textbf{Lean statement sketch.}
\begin{Verbatim}[fontsize=\small,breaklines=true,breakanywhere=true,tabsize=2]
def CinematicFamily (I : Set ℝ) (K D : ℝ)
  (Fs : Set (C2Function I)) : Prop :=
  DiameterBound K Fs ∧ DoublingBound D Fs ∧
  ∀ g ∈ Fs, ∀ h ∈ Fs, ∀ t ∈ I,
    c2Dist g h / K ≤
      |g t-h t|+|g.deriv t-h.deriv t|+|g.deriv2 t-h.deriv2 t|
\end{Verbatim}
{\footnotesize\textbf{Mathlib support.} C, G, L.\par
\textbf{Incoming edges.} \nodecard{Icalc} (api)\par}
{\footnotesize\textbf{Interface note.} C2Function, the jet metric, and DoublingBound are proposed project definitions. The inequality may be expressed pointwise instead of using an infimum.}
\end{tcolorbox}\end{minipage}\par\medskip
\begin{minipage}{\linewidth}
\hypertarget{node-T7_2}{}
\begin{tcolorbox}[colback=white,colframe=cinematicfg,colbacktitle=cinematicbg,coltitle=ink,boxrule=0.45pt,arc=1.6mm,left=3mm,right=3mm,top=2mm,bottom=2mm,title={\sffamily\bfseries WZ1 Thm. 7.2 = WZ2 Thm. 7.13 \hfill \texttt{T7\_2}}]
{\sffamily\small\color{muted}Cinematic maximal estimate | WZ1/WZ2 | WZ1 p. 49; WZ2 PDF pp. 47-48; PYZ Theorem 1.7}\par
\textbf{Mathematical statement.} For a cinematic family on a compact interval and every $\varepsilon>0$, sufficiently small $\delta$ has the following property.
A $\delta$-separated finite subfamily $F$ satisfying
$\#(F\cap B)\le\delta^{-\varepsilon}(r/\delta)$ for every $C^2$-ball of diameter $r$
obeys
$\|\sum_{g\in F}1_{N_\delta(\operatorname{graph}g)}\|_{L^{3/2}}\le\delta^{-\varepsilon}$.
The small-scale threshold depends on the cinematic and doubling constants.
\smallskip\textbf{Lean statement sketch.}
\begin{Verbatim}[fontsize=\small,breaklines=true,breakanywhere=true,tabsize=2]
theorem cinematic_maximal (hcin : CinematicFamily I K D Fs) (he : 0 < eps) :
  ∃ delta0 > 0, ∀ delta ∈ Set.Ioc 0 delta0,
  ∀ F : Finset (C2Function I),
  FiniteCinematicFrostmanHyp delta eps Fs F →
    eLpNorm (curveNeighborhoodSum delta F) (3/2) volume ≤
      ENNReal.ofReal (delta^(-eps))
\end{Verbatim}
{\footnotesize\textbf{Mathlib support.} C, F, G, M, P, Q.\par
\textbf{Incoming edges.} \nodecard{Dcinematic} (core); \nodecard{XPYZ} (core)\par}
{\footnotesize\textbf{Interface note.} An imported cinematic maximal theorem. This interface is used only after extracting an absolute one-dimensional nonconcentration bound in C2 space; a two-dimensional normalized Frostman bound on R4 parameters is not a substitute.}
\end{tcolorbox}\end{minipage}\par\medskip
\begin{minipage}{\linewidth}
\hypertarget{node-L7_3}{}
\begin{tcolorbox}[colback=white,colframe=cinematicfg,colbacktitle=cinematicbg,coltitle=ink,boxrule=0.45pt,arc=1.6mm,left=3mm,right=3mm,top=2mm,bottom=2mm,title={\sffamily\bfseries WZ1 Lemma 7.3 \hfill \texttt{L7\_3}}]
{\sffamily\small\color{muted}The twisted-line family is cinematic | WZ1 | WZ1 Lemma 7.3, pp. 49-50; invoked in WZ2 proof of Lemma 7.12, PDF p. 49}\par
\textbf{Mathematical statement.} If $f(0)=0$, $1\le|f'|\le2$ and $|f''|\le1/100$ on $[0,1]$, then
$M_f=\{g_{a,b,d}:g_{a,b,d}(t)=a+bf(t)+dtf(t),\ a,b,d\in[-1,1]\}$
is cinematic with absolute constants.
The proof compares the $C^2$ distance to the distance between parameter triples and bounds the size of the 2-jet from below at every point.
\smallskip\textbf{Lean statement sketch.}
\begin{Verbatim}[fontsize=\small,breaklines=true,breakanywhere=true,tabsize=2]
def cinematicModel (f : ℝ → ℝ) (a b d t : ℝ) :=
  a+b*f t+d*t*f t
theorem twisted_model_cinematic
  (h0 : f 0 = 0) (hf : NonsingularSlope f (Set.Icc 0 1)) :
  ∃ K D, AbsoluteCinematicConstants K D ∧
    CinematicFamily (Set.Icc 0 1) K D (modelFamily f)
\end{Verbatim}
{\footnotesize\textbf{Mathlib support.} C, F, G, L.\par
\textbf{Incoming edges.} \nodecard{Dcinematic} (core); \nodecard{Icalc} (api)\par}
{\footnotesize\textbf{Interface note.} Normalize the bounded t-interval and subtract f(0) if needed. A fixed c contributes the common target shear (u,t) -> (u-ct,t), or equivalently the common summand ct in every curve. Area and all pairwise jet distances are preserved.}
\end{tcolorbox}\end{minipage}\par\medskip
\begin{minipage}{\linewidth}
\hypertarget{node-W2TW}{}
\begin{tcolorbox}[colback=white,colframe=cinematicfg,colbacktitle=cinematicbg,coltitle=ink,boxrule=0.45pt,arc=1.6mm,left=3mm,right=3mm,top=2mm,bottom=2mm,title={\sffamily\bfseries WZ2 Section 7: Tube Wolff \hfill \texttt{W2TW}}]
{\sffamily\small\color{muted}The weaker radius-squared tube count | WZ2 | WZ2 PDF p. 44}\par
\textbf{Mathematical statement.} The Tube Wolff axiom with error $A$ requires
$\#\mathcal T[U_r]\le Ar^2\#\mathcal T$ for every $r$-tube $U_r$, $\delta\le r\le1$.
Convex Wolff implies this condition by testing a convex tube of volume $O(r^2)$.
Section 7 needs only this weaker condition; it does not assume multi-scale convex Wolff.
\smallskip\textbf{Lean statement sketch.}
\begin{Verbatim}[fontsize=\small,breaklines=true,breakanywhere=true,tabsize=2]
def TubeWolff (A : ℝ) (C : Config 3 delta) : Prop :=
  ∀ r ∈ Set.Icc delta 1, ∀ U : Tube r,
    (children C U).card ≤ A*r^2*(C.tubes.card : ℝ)
theorem convexWolff_implies_tubeWolff : ConvexWolff A C.tubeSets → TubeWolff (K*A) C
\end{Verbatim}
{\footnotesize\textbf{Mathlib support.} F, G, M, Q.\par
\textbf{Incoming edges.} \nodecard{W2CWA} (core)\par}

\end{tcolorbox}\end{minipage}\par\medskip
\begin{minipage}{\linewidth}
\hypertarget{node-W2DFrost}{}
\begin{tcolorbox}[colback=white,colframe=cinematicfg,colbacktitle=cinematicbg,coltitle=ink,boxrule=0.45pt,arc=1.6mm,left=3mm,right=3mm,top=2mm,bottom=2mm,title={\sffamily\bfseries WZ2 Definition 7.3 \hfill \texttt{W2DFrost}}]
{\sffamily\small\color{muted}Normalized Frostman entropy | WZ2 | WZ2 PDF p. 44}\par
\textbf{Mathematical statement.} A bounded nonempty $A\subset\mathbb R^n$ is $(\delta,s,C)$-Frostman if
\[\mathcal E_\delta(A\cap B(x,r))\le Cr^s\mathcal E_\delta(A),\qquad\delta\le r\le1.\]
In a bounded line chart, Tube Wolff makes the parameter set
$\mathcal A=\{(a,b,c,d):\ell(t)=(a+ct,b+dt,t)\}$ a normalized $(\delta,2,O(C))$-Frostman set in $\mathbb R^4$.
\smallskip\textbf{Lean statement sketch.}
\begin{Verbatim}[fontsize=\small,breaklines=true,breakanywhere=true,tabsize=2]
def EntropyFrostman (delta s C : ℝ) (A : Set (Rn n)) : Prop :=
  BoundedNonempty A ∧ ∀ x r, delta ≤ r → r ≤ 1 →
    coverNumR delta (A ∩ Metric.ball x r) ≤ C*r^s*coverNumR delta A
\end{Verbatim}
{\footnotesize\textbf{Mathlib support.} F, G, M, Q.\par
\textbf{Incoming edges.} \nodecard{Ientropy} (api); \nodecard{W2TW} (core)\par}

\end{tcolorbox}\end{minipage}\par\medskip
\begin{minipage}{\linewidth}
\hypertarget{node-W2DKT}{}
\begin{tcolorbox}[colback=white,colframe=cinematicfg,colbacktitle=cinematicbg,coltitle=ink,boxrule=0.45pt,arc=1.6mm,left=3mm,right=3mm,top=2mm,bottom=2mm,title={\sffamily\bfseries WZ2 Definition 7.2 \hfill \texttt{W2DKT}}]
{\sffamily\small\color{muted}Absolute Katz-Tao entropy bound | WZ2 | WZ2 PDF p. 44}\par
\textbf{Mathematical statement.} A bounded nonempty set is $(\delta,s,C)$-Katz-Tao if
\[\mathcal E_\delta(A\cap B(x,r))\le C(r/\delta)^s,\qquad\delta\le r\le1.\]
This is an absolute count, not a normalized fraction of $\mathcal E_\delta(A)$.
For $s\le1$ it implies the one-dimensional absolute ball condition required by the cinematic maximal theorem, after parameter/jet metric comparison.
\smallskip\textbf{Lean statement sketch.}
\begin{Verbatim}[fontsize=\small,breaklines=true,breakanywhere=true,tabsize=2]
def KatzTao (delta s C : ℝ) (A : Set (Rn n)) : Prop :=
  BoundedNonempty A ∧ ∀ x r, delta ≤ r → r ≤ 1 →
    coverNumR delta (A ∩ Metric.ball x r) ≤ C*(r/delta)^s
\end{Verbatim}
{\footnotesize\textbf{Mathlib support.} F, G, M, Q.\par
\textbf{Incoming edges.} \nodecard{Ientropy} (api)\par}
{\footnotesize\textbf{Interface note.} The supplied PDF displays (delta/r)\textasciicircum{}s in Definition 7.2. The counting arguments, Lemma 7.11 and the standard Katz-Tao convention require (r/delta)\textasciicircum{}s. The atlas displays that corrected ratio and records the source discrepancy here.}
\end{tcolorbox}\end{minipage}\par\medskip
\begin{minipage}{\linewidth}
\hypertarget{node-W2DUnif}{}
\begin{tcolorbox}[colback=white,colframe=cinematicfg,colbacktitle=cinematicbg,coltitle=ink,boxrule=0.45pt,arc=1.6mm,left=3mm,right=3mm,top=2mm,bottom=2mm,title={\sffamily\bfseries WZ2 Definition 7.8 \hfill \texttt{W2DUnif}}]
{\sffamily\small\color{muted}Uniform occupied-cell counts | WZ2 | WZ2 PDF p. 45}\par
\textbf{Mathematical statement.} At the prescribed sparse dyadic scales $r=2^{-j/\eta}$ between $\delta$ and $1$, the counts $\mathcal E_\delta(A\cap Q)$ in occupied $r$-cells are comparable within the fixed factor $100$.
Use a consistent grid convention (or bounded-overlap ball version).
This uniformity concerns a set's nested occupied cells, not the hypergraph uniform density used in WZ1 Section 8.
\smallskip\textbf{Lean statement sketch.}
\begin{Verbatim}[fontsize=\small,breaklines=true,breakanywhere=true,tabsize=2]
def EtaUniform (eta delta : ℝ) (A : Set (Rn n)) : Prop :=
  ∀ r ∈ sparseScales eta delta, UniformOccupiedCounts 100 delta r A
\end{Verbatim}
{\footnotesize\textbf{Mathlib support.} F, G, M, Q.\par
\textbf{Incoming edges.} \nodecard{Ientropy} (api)\par}

\end{tcolorbox}\end{minipage}\par\medskip
\begin{minipage}{\linewidth}
\hypertarget{node-W2L7_6}{}
\begin{tcolorbox}[colback=white,colframe=cinematicfg,colbacktitle=cinematicbg,coltitle=ink,boxrule=0.45pt,arc=1.6mm,left=3mm,right=3mm,top=2mm,bottom=2mm,title={\sffamily\bfseries WZ2 Lemma 7.6 \hfill \texttt{W2L7\_6}}]
{\sffamily\small\color{muted}Frostman sets contain large Katz-Tao subsets | WZ2 | WZ2 PDF pp. 44-45}\par
\textbf{Mathematical statement.} A $(\delta,s,C)$-Frostman set has a $(\delta,s,100)$-Katz-Tao subset $A'$ with
$\mathcal E_\delta(A')\gtrsim_{n,s}\delta^{-s}|\log\delta|^{-1}C^{-1}$.
Apply the Hausdorff-content discretization after thickening to the mesh and, if necessary, pull the selected net back to $A$ with bounded losses.
\smallskip\textbf{Lean statement sketch.}
\begin{Verbatim}[fontsize=\small,breaklines=true,breakanywhere=true,tabsize=2]
theorem Frostman_extracts_KatzTao (h : EntropyFrostman delta s C A) :
  ∃ A' ⊆ A, KatzTao delta s 100 A' ∧
    K_ns * delta^(-s) / (|Real.log delta|*C) ≤ coverNumR delta A' 
\end{Verbatim}
{\footnotesize\textbf{Mathlib support.} F, G, M, Q.\par
\textbf{Incoming edges.} \nodecard{W2DFrost} (core); \nodecard{W2DKT} (api); \nodecard{XFO} (core)\par}
{\footnotesize\textbf{Interface note.} For a finite parameter net use delta-content, or thicken by delta before applying the ordinary content lemma. A finite set has zero ordinary positive-dimensional Hausdorff content; do not apply that continuous hypothesis directly to a finite net.}
\end{tcolorbox}\end{minipage}\par\medskip
\begin{minipage}{\linewidth}
\hypertarget{node-W2L7_9}{}
\begin{tcolorbox}[colback=white,colframe=cinematicfg,colbacktitle=cinematicbg,coltitle=ink,boxrule=0.45pt,arc=1.6mm,left=3mm,right=3mm,top=2mm,bottom=2mm,title={\sffamily\bfseries WZ2 Lemma 7.9 \hfill \texttt{W2L7\_9}}]
{\sffamily\small\color{muted}Sparse-scale uniform refinement | WZ2 | WZ2 PDF p. 46; WZ2 [11, Lemma 2.15]}\par
\textbf{Mathematical statement.} If $\delta=2^{-T/\eta}$ and $A\subset[0,1]^n$, there is an $\eta$-uniform subset $A'$ with
$\mathcal E_\delta(A')\ge(2T)^{-1/\eta}\mathcal E_\delta(A)$.
For fixed $\eta$ this polylogarithmic loss is smaller than any reserved $\delta$-power when $\delta$ is sufficiently small.
\smallskip\textbf{Lean statement sketch.}
\begin{Verbatim}[fontsize=\small,breaklines=true,breakanywhere=true,tabsize=2]
theorem uniform_refinement (hmesh : delta = (2:ℝ)^(-(T:ℝ)/eta)) :
  ∃ A' ⊆ A, EtaUniform eta delta A' ∧
    (2*T)^(-1/eta) * coverNumR delta A ≤ coverNumR delta A' 
\end{Verbatim}
{\footnotesize\textbf{Mathlib support.} F, G, M, Q.\par
\textbf{Incoming edges.} \nodecard{W2DUnif} (api); \nodecard{XOSuniform} (core)\par}

\end{tcolorbox}\end{minipage}\par\medskip
\begin{minipage}{\linewidth}
\hypertarget{node-W2L7_10}{}
\begin{tcolorbox}[colback=white,colframe=cinematicfg,colbacktitle=cinematicbg,coltitle=ink,boxrule=0.45pt,arc=1.6mm,left=3mm,right=3mm,top=2mm,bottom=2mm,title={\sffamily\bfseries WZ2 Lemma 7.10 \hfill \texttt{W2L7\_10}}]
{\sffamily\small\color{muted}Zoom to Frostman fibres of a uniform set | WZ2 | WZ2 PDF pp. 46-47}\par
\textbf{Mathematical statement.} For every $\varepsilon>0$ choose $\eta>0$.
If $A\subset[0,1]^n$ is $\eta$-uniform and $\mathcal E_\delta(A)\ge\delta^{-s}$, there is $q\in(\delta^{1-\varepsilon},1]$ such that every occupied $q$-cube $Q$ has
$S_Q(A\cap Q)$ a $(\delta/q,s,(\delta/q)^{-4n\varepsilon})$-Frostman set.
The proof merges the scale profile until its exponents increase, then selects the first interval at or above $s-\varepsilon$ and absorbs that small exponent deficit into the error.
\smallskip\textbf{Lean statement sketch.}
\begin{Verbatim}[fontsize=\small,breaklines=true,breakanywhere=true,tabsize=2]
theorem uniform_zoom_Frostman (he : 0 < eps) :
  ∃ eta > 0, ∀ delta A s, SparseMesh eta delta → EtaUniform eta delta A →
  delta^(-s) ≤ coverNumR delta A → ∃ q ∈ Set.Ioc (delta^(1-eps)) 1,
    ∀ Q ∈ occupiedCubes q A,
      EntropyFrostman (delta/q) s ((delta/q)^(-4*n*eps)) (rescale Q (A ∩ Q))
\end{Verbatim}
{\footnotesize\textbf{Mathlib support.} F, G, M, Q.\par
\textbf{Incoming edges.} \nodecard{W2DUnif} (core); \nodecard{XSWdecomp} (core); \nodecard{XSWmerge} (core); \nodecard{W2DFrost} (api)\par}
{\footnotesize\textbf{Interface note.} Ambient dimension n replaces the inconsistent d/n notation in the source. The statements quoted from Shmerkin-Wang retain their complete quantitative contracts; only the local proof assembly is expanded.}
\end{tcolorbox}\end{minipage}\par\medskip
\begin{minipage}{\linewidth}
\hypertarget{node-W2L7_11}{}
\begin{tcolorbox}[colback=white,colframe=cinematicfg,colbacktitle=cinematicbg,coltitle=ink,boxrule=0.45pt,arc=1.6mm,left=3mm,right=3mm,top=2mm,bottom=2mm,title={\sffamily\bfseries WZ2 Lemma 7.11 \hfill \texttt{W2L7\_11}}]
{\sffamily\small\color{muted}Replicate a local Katz-Tao cluster | WZ2 | WZ2 PDF p. 47; density-budget version of the same proof}\par
\textbf{Mathematical statement.} Let $A_0$ lie in a $q$-cube and satisfy a $(\delta,s,C)$-Katz-Tao bound.
Place $\asymp q^{-s}$ translated copies on a grid of spacing $\asymp q^{s/n}$.
The union $A^*$ is $(\delta,s,O(C))$-Katz-Tao and
$\mathcal E_\delta(A^*)\asymp q^{-s}\mathcal E_\delta(A_0)$.
If $\mathcal E_\delta(A_0)\ge\kappa(q/\delta)^s$, then $\mathcal E_\delta(A^*)\gtrsim\kappa\delta^{-s}$.
The printed lemma is the saturated case $\kappa\asymp1$.
\smallskip\textbf{Lean statement sketch.}
\begin{Verbatim}[fontsize=\small,breaklines=true,breakanywhere=true,tabsize=2]
theorem replicate_KatzTao (hKT : KatzTao delta s C A0) (hbox : InCube q A0) :
  ∃ Astar translations, KatzTao delta s (K*C) Astar ∧
    ReplicatedBy translations A0 Astar ∧
    Comparable (coverNumR delta Astar) (q^(-s)*coverNumR delta A0)
\end{Verbatim}
{\footnotesize\textbf{Mathlib support.} F, G, M, Q.\par
\textbf{Incoming edges.} \nodecard{W2DKT} (core)\par}
{\footnotesize\textbf{Interface note.} The extraction in Lemma 7.12 is only nearly saturated. Keeping kappa visible avoids using the saturated statement without its hypothesis. The upper count proof is unchanged: use one cluster below the grid spacing and the grid count above it.}
\end{tcolorbox}\end{minipage}\par\medskip
\begin{minipage}{\linewidth}
\hypertarget{node-W2ImageUniform}{}
\begin{tcolorbox}[colback=white,colframe=cinematicfg,colbacktitle=cinematicbg,coltitle=ink,boxrule=0.45pt,arc=1.6mm,left=3mm,right=3mm,top=2mm,bottom=2mm,title={\sffamily\bfseries WZ2 Lemma 7.12: preparation \hfill \texttt{W2ImageUniform}}]
{\sffamily\small\color{muted}Uniformize the projected incidence set | WZ2 | WZ2 proof of Lemma 7.12, PDF p. 48}\par
\textbf{Mathematical statement.} Pigeonhole the fibre incidence multiplicity of $\pi_f$ to a dyadic level $X_\lambda$.
Apply uniform refinement to obtain an $\eta$-uniform target set $X\subset X_\lambda$ and pull it back to a sub-shading $Y'$.
The constant fibre-incidence level makes projected-area retention comparable to retained shading mass; after selecting popular tubes the shading loss is a controlled small power.
\smallskip\textbf{Lean statement sketch.}
\begin{Verbatim}[fontsize=\small,breaklines=true,breakanywhere=true,tabsize=2]
theorem uniform_twisted_incidence_refinement (h : TwistedProjectionHyp eta C f) :
  ∃ C' X, Subcollection C' C ∧ EtaUniform eta delta X ∧
    X = twistedImage C' f ∧ ControlledDensityRetention eta C' C
\end{Verbatim}
{\footnotesize\textbf{Mathlib support.} F, G, M, Q.\par
\textbf{Incoming edges.} \nodecard{W2L7_9} (core); \nodecard{Dtwist} (api); \nodecard{Imass} (api)\par}

\end{tcolorbox}\end{minipage}\par\medskip
\begin{minipage}{\linewidth}
\hypertarget{node-W2Slice}{}
\begin{tcolorbox}[colback=white,colframe=cinematicfg,colbacktitle=cinematicbg,coltitle=ink,boxrule=0.45pt,arc=1.6mm,left=3mm,right=3mm,top=2mm,bottom=2mm,title={\sffamily\bfseries WZ2 Lemma 7.12: parameter slicing \hfill \texttt{W2Slice}}]
{\sffamily\small\color{muted}A rich fixed-c slice of R4 parameters | WZ2 | WZ2 Remark 7.4 and proof of Lemma 7.12, PDF pp. 44, 48}\par
\textbf{Mathematical statement.} Tube Wolff gives a normalized two-dimensional Frostman bound for $\mathcal A\subset\mathbb R^4$ and hence $\mathcal E_\delta(\mathcal A)\gtrsim\delta^{-2+O(\eta)}$.
After the density and fibre refinements, some $\delta$-slab of the $c$-coordinate contains a uniformly shaded triple set $A_c\subset\mathbb R^3_{a,b,d}$ with
$\mathcal E_\delta(A_c)\gtrsim\delta^{-1+O(\eta)}$.
Freeze $c$ with an $O(\delta)$ target error and then apply a sparse-scale uniform refinement to this slice.
\smallskip\textbf{Lean statement sketch.}
\begin{Verbatim}[fontsize=\small,breaklines=true,breakanywhere=true,tabsize=2]
theorem rich_fixed_c_slice (hTW : TubeWolff (delta^(-eta)) C)
  (hd : UniformlyDense (delta^eta) C) :
  ∃ c A, FixedCParameterSlice delta c C A ∧
    delta^(-1+K*eta) ≤ coverNumR delta A
\end{Verbatim}
{\footnotesize\textbf{Mathlib support.} F, G, M, Q.\par
\textbf{Incoming edges.} \nodecard{W2DFrost} (core); \nodecard{W2L7_9} (core); \nodecard{W2ImageUniform} (assembly)\par}
{\footnotesize\textbf{Interface note.} This is a finite delta-slab pigeonhole, not an exact continuum fibre theorem. Exact c-fibres can be empty for a finite net. Shade transport and target thickening absorb the O(delta) freezing error.}
\end{tcolorbox}\end{minipage}\par\medskip
\begin{minipage}{\linewidth}
\hypertarget{node-W2JetCount}{}
\begin{tcolorbox}[colback=white,colframe=cinematicfg,colbacktitle=cinematicbg,coltitle=ink,boxrule=0.45pt,arc=1.6mm,left=3mm,right=3mm,top=2mm,bottom=2mm,title={\sffamily\bfseries WZ2 Lemma 7.12: cinematic assembly \hfill \texttt{W2JetCount}}]
{\sffamily\small\color{muted}Convert Katz-Tao triples to separated curves | WZ2 | WZ2 PDF pp. 48-49; WZ1 Lemma 7.3}\par
\textbf{Mathematical statement.} Zoom the uniform triple set to a $q$-cube, extract Katz-Tao triples at $h=\delta/q$, and replicate them back at mesh $\delta$.
For fixed $c$ the functions $g(t)=a+bf(t)+dtf(t)+ct$ are cinematic; their parameter metric and $C^2$ metric are comparable.
A $(\delta,1-\tau,O(1))$-Katz-Tao parameter set therefore gives a $\delta$-separated curve family obeying
$\#(G\cap B_r)\lesssim(r/\delta)$.
All cardinality and shading deficits from extraction remain explicit small powers of $h$.
\smallskip\textbf{Lean statement sketch.}
\begin{Verbatim}[fontsize=\small,breaklines=true,breakanywhere=true,tabsize=2]
theorem replicated_curves_cinematic (hf : NonsingularSlope f I)
  (hKT : KatzTao delta (1-tau) K Astar) :
  ∃ G, CurveFamilyOfTriples f c Astar G ∧
    CinematicFamily I Kcin Dcin G ∧ AbsoluteCurveBallCount delta K' G
\end{Verbatim}
{\footnotesize\textbf{Mathlib support.} F, G, M, Q.\par
\textbf{Incoming edges.} \nodecard{W2Slice} (assembly); \nodecard{W2L7_10} (core); \nodecard{W2L7_6} (core); \nodecard{W2L7_11} (core); \nodecard{L7_3} (core); \nodecard{W2DKT} (api)\par}

\end{tcolorbox}\end{minipage}\par\medskip
\begin{minipage}{\linewidth}
\hypertarget{node-W2AreaCompare}{}
\begin{tcolorbox}[colback=white,colframe=cinematicfg,colbacktitle=cinematicbg,coltitle=ink,boxrule=0.45pt,arc=1.6mm,left=3mm,right=3mm,top=2mm,bottom=2mm,title={\sffamily\bfseries WZ2 Lemma 7.12: Holder and translations \hfill \texttt{W2AreaCompare}}]
{\sffamily\small\color{muted}From cinematic norm to a local area ratio | WZ2 | WZ2 equations (7.4)-(7.6), PDF pp. 49-50}\par
\textbf{Mathematical statement.} For $H=\sum_{g\in G}1_{N_{O(\delta)}(\operatorname{graph}g)}$ and the replicated shaded image $Z$, Holder gives
$M_{\rm proj}\le |Z|^{1/3}\|H\|_{3/2}$.
Here $M_{\rm proj}=\sum_T|\pi_f(Y'(T))|$ is projected area mass; a uniformly dense tube contributes $\gtrsim\lambda\delta$.
Translated parameter clusters have identical slice-image area: at time $t$ their displacement is $a_0+b_0f(t)+d_0tf(t)$.
Divide by the number of replicas to obtain an area lower bound for one cluster, relative to its $O(q)$-neighbourhood.
\smallskip\textbf{Lean statement sketch.}
\begin{Verbatim}[fontsize=\small,breaklines=true,breakanywhere=true,tabsize=2]
theorem cinematic_cluster_area_ratio
  (h : ReplicatedCinematicComparisonData tau eta delta q C f) :
  vol (twistedImage C f) ≥ (delta/q)^(K*tau) *
    vol (neighborhood q (twistedImage C f))
\end{Verbatim}
{\footnotesize\textbf{Mathlib support.} F, G, M, Q.\par
\textbf{Incoming edges.} \nodecard{W2JetCount} (assembly); \nodecard{T7_2} (core); \nodecard{Imass} (api); \nodecard{Iparams} (api)\par}
{\footnotesize\textbf{Interface note.} Use a separately chosen cinematic exponent: h=delta/q lies between delta and delta\textasciicircum{}tau, so delta\textasciicircum{}(-tau\textasciicircum{}2) is at most h\textasciicircum{}(-tau). The fixed maximal-theorem constants are absorbed by another reserved margin. This prevents a delta-loss from being confused with an h-loss.}
\end{tcolorbox}\end{minipage}\par\medskip
\begin{minipage}{\linewidth}
\hypertarget{node-W2L7_12}{}
\begin{tcolorbox}[colback=white,colframe=cinematicfg,colbacktitle=cinematicbg,coltitle=ink,boxrule=0.45pt,arc=1.6mm,left=3mm,right=3mm,top=2mm,bottom=2mm,title={\sffamily\bfseries WZ2 Lemma 7.12 \hfill \texttt{W2L7\_12}}]
{\sffamily\small\color{muted}One-step fine-to-coarse projected-area comparison | WZ2 | WZ2 PDF pp. 47-50, equation (7.2)}\par
\textbf{Mathematical statement.} For every $\varepsilon>0$ choose $\eta,\delta_0>0$.
For a Tube-Wolff family with error $\delta^{-\eta}$, $\delta^\eta$-dense shading and nonsingular slope, there are $r\in(\delta^{1-\varepsilon^2},1)$ and a sub-shading $Y'$ that is $\delta^{4\eta}$-dense such that
\[\left|\bigcup_T\pi_f(Y'(T))\right|\ge(\delta/r)^\varepsilon
\left|N_r\!\left(\bigcup_T\pi_f(Y'(T))\right)\right|.\]
The exact mesh can be rounded to the sparse scale set with bounded losses.
\smallskip\textbf{Lean statement sketch.}
\begin{Verbatim}[fontsize=\small,breaklines=true,breakanywhere=true,tabsize=2]
theorem one_step_twisted_area (he : 0 < eps) :
  ∃ eta > 0, ∃ delta0 > 0, ∀ delta ∈ Set.Ioc 0 delta0,
  TwistedProjectionHyp eta C f → ∃ r ∈ Set.Ioo (delta^(1-eps^2)) 1,
  ∃ C', Subcollection C' C ∧ Dense (delta^(4*eta)) C' ∧
    vol (twistedImage C' f) ≥ (delta/r)^eps *
      vol (neighborhood r (twistedImage C' f))
\end{Verbatim}
{\footnotesize\textbf{Mathlib support.} F, G, M, Q.\par
\textbf{Incoming edges.} \nodecard{W2AreaCompare} (core); \nodecard{W2ImageUniform} (core); \nodecard{W2TW} (api)\par}
{\footnotesize\textbf{Interface note.} The source states this quantitative lemma and proves it by the sliced/replicated cinematic construction. The cards separate the overloaded source scales, keep the near-saturation deficit and reserve independent small-loss budgets.}
\end{tcolorbox}\end{minipage}\par\medskip
\begin{minipage}{\linewidth}
\hypertarget{node-IcoarseTW}{}
\begin{tcolorbox}[colback=white,colframe=cinematicfg,colbacktitle=cinematicbg,coltitle=ink,boxrule=0.45pt,arc=1.6mm,left=3mm,right=3mm,top=2mm,bottom=2mm,title={\sffamily\bfseries WZ2 proof operation \hfill \texttt{IcoarseTW}}]
{\sffamily\small\color{muted}Uniform Tube-Wolff coarsening | WZ2 | WZ2 proof of Proposition 7.1, PDF pp. 50-51}\par
\textbf{Mathematical statement.} For a uniform parameter family, choose occupied $r$-cells as parent tubes and comparable numbers of children per cell.
The normalized radius-squared Tube-Wolff count passes to this coarse family with controlled losses.
Define the coarse shading by $Y_r(U)=N_r(\bigcup_{T\in\mathcal T[U]}Y'(T))$ and retain its density and compatible target-neighbourhood bounds.
This is the coarsening operation used at every step of the projection iteration.
\smallskip\textbf{Lean statement sketch.}
\begin{Verbatim}[fontsize=\small,breaklines=true,breakanywhere=true,tabsize=2]
structure TubeWolffCoarseningWitness (C : Config 3 delta) (D : Config 3 r) where
  balanced_children : ComparableChildCounts C D
  tube_count : TubeWolff A' D
  shading_density : Dense lambda' D
  target_compatibility : TargetNeighborhoodCompatibility C D f
\end{Verbatim}
{\footnotesize\textbf{Mathlib support.} F, G, M, Q.\par
\textbf{Incoming edges.} \nodecard{W2L7_9} (core); \nodecard{W2TW} (core); \nodecard{Imass} (api)\par}
{\footnotesize\textbf{Interface note.} Uniformity is applied in line-parameter space. It is required to convert normalized fine counts into normalized coarse counts; thickening without controlling child counts is insufficient.}
\end{tcolorbox}\end{minipage}\par\medskip
\begin{minipage}{\linewidth}
\hypertarget{node-W2Iterate}{}
\begin{tcolorbox}[colback=white,colframe=cinematicfg,colbacktitle=cinematicbg,coltitle=ink,boxrule=0.45pt,arc=1.6mm,left=3mm,right=3mm,top=2mm,bottom=2mm,title={\sffamily\bfseries WZ2 Proposition 7.1: iteration \hfill \texttt{W2Iterate}}]
{\sffamily\small\color{muted}Telescope the ratios up to a coarse scale | WZ2 | WZ2 PDF pp. 50-51}\par
\textbf{Mathematical statement.} Let $r_0=\delta$ and apply the one-step comparison successively to uniform coarse families.
Take $r_{j+1}>r_j^{1-\theta}$ with $\theta>0$ fixed by the final error, until $r_N\ge\delta^\theta$.
One may use
$N\le 1+\lceil\log\theta/\log(1-\theta)\rceil$, a constant depending only on $\theta$.
The ratios telescope to $(\delta/r_N)^\varepsilon$; the terminal nonempty coarse curve neighbourhood has area $\gtrsim r_N$.
Choose the initial density/error budget for these finitely many calls, and rename the final exponent.
\smallskip\textbf{Lean statement sketch.}
\begin{Verbatim}[fontsize=\small,breaklines=true,breakanywhere=true,tabsize=2]
theorem iterate_twisted_area (hstep : OneStepTwistedAreaContract)
  (hcoarse : UniformTubeWolffCoarseningContract) :
  TwistedProjectionLowerBoundContract
\end{Verbatim}
{\footnotesize\textbf{Mathlib support.} F, G, M, Q.\par
\textbf{Incoming edges.} \nodecard{W2L7_12} (core); \nodecard{IcoarseTW} (core); \nodecard{Iparams} (api)\par}
{\footnotesize\textbf{Interface note.} The source writes N <= epsilon\textasciicircum{}(-2) and uses rho\_2 in its last telescoping line. The atlas uses a valid finite iteration bound and r\_N as the terminal scale; neither literal numerical slip is needed for the theorem.}
\end{tcolorbox}\end{minipage}\par\medskip
\begin{minipage}{\linewidth}
\hypertarget{node-W2P7_1}{}
\begin{tcolorbox}[colback=white,colframe=cinematicfg,colbacktitle=cinematicbg,coltitle=ink,boxrule=0.45pt,arc=1.6mm,left=3mm,right=3mm,top=2mm,bottom=2mm,title={\sffamily\bfseries WZ2 Proposition 7.1 \hfill \texttt{W2P7\_1}}]
{\sffamily\small\color{muted}Twisted-image lower bound under Tube Wolff | WZ2 | WZ2 PDF pp. 44, 50-51}\par
\textbf{Mathematical statement.} For every $\varepsilon>0$ choose $\eta,\delta_0>0$.
For $0<\delta\le\delta_0$, Tube-Wolff error $\delta^{-\eta}$, $\delta^\eta$-dense shading and a nonsingular slope imply
$|\bigcup_T\pi_f(Y(T))|\ge\delta^\varepsilon$.
This is the stronger Section 7 statement that proves Proposition 6.8.
\smallskip\textbf{Lean statement sketch.}
\begin{Verbatim}[fontsize=\small,breaklines=true,breakanywhere=true,tabsize=2]
theorem twisted_projection_TubeWolff (he : 0 < eps) :
  ∃ eta > 0, ∃ delta0 > 0, ∀ delta ∈ Set.Ioc 0 delta0,
  TubeWolff (delta^(-eta)) C → Dense (delta^eta) C →
  NonsingularSlope f (Set.Icc (-1) 1) → delta^eps ≤ vol (twistedImage C f)
\end{Verbatim}
{\footnotesize\textbf{Mathlib support.} F, G, M, Q.\par
\textbf{Incoming edges.} \nodecard{W2Iterate} (core)\par}

\end{tcolorbox}\end{minipage}\par\medskip
\begin{minipage}{\linewidth}
\hypertarget{node-W2P6_8}{}
\begin{tcolorbox}[colback=white,colframe=cinematicfg,colbacktitle=cinematicbg,coltitle=ink,boxrule=0.45pt,arc=1.6mm,left=3mm,right=3mm,top=2mm,bottom=2mm,title={\sffamily\bfseries WZ2 Proposition 6.8 \hfill \texttt{W2P6\_8}}]
{\sffamily\small\color{muted}Twisted-image lower bound under convex Wolff | WZ2 | WZ2 PDF pp. 43-44; proved by Proposition 7.1}\par
\textbf{Mathematical statement.} For every $\varepsilon>0$, convex-Wolff error $\delta^{-\eta}$ and $\delta^\eta$-dense shading imply
$|\bigcup_T\pi_f(Y(T))|\ge\delta^\varepsilon$ for any nonsingular $C^2$ slope and all sufficiently small $\delta$.
The required $\eta,\delta_0$ depend on $\varepsilon$.
Use convex Wolff $\Rightarrow$ Tube Wolff, then Proposition 7.1.
\smallskip\textbf{Lean statement sketch.}
\begin{Verbatim}[fontsize=\small,breaklines=true,breakanywhere=true,tabsize=2]
theorem twisted_projection_ConvexWolff (he : 0 < eps) :
  ∃ eta > 0, ∃ delta0 > 0, ∀ delta ∈ Set.Ioc 0 delta0,
  ConvexWolff (delta^(-eta)) C.tubeSets → Dense (delta^eta) C →
  NonsingularSlope f (Set.Icc (-1) 1) → delta^eps ≤ vol (twistedImage C f)
\end{Verbatim}
{\footnotesize\textbf{Mathlib support.} F, G, M, Q.\par
\textbf{Incoming edges.} \nodecard{W2P7_1} (core); \nodecard{W2TW} (core)\par}

\end{tcolorbox}\end{minipage}\par\medskip
\sectitle{Node catalogue: Radial bootstrap}
\begin{minipage}{\linewidth}
\hypertarget{node-DFrost}{}
\begin{tcolorbox}[colback=white,colframe=radialfg,colbacktitle=radialbg,coltitle=ink,boxrule=0.45pt,arc=1.6mm,left=3mm,right=3mm,top=2mm,bottom=2mm,title={\sffamily\bfseries WZ1 Definition 8.1 \hfill \texttt{DFrost}}]
{\sffamily\small\color{muted}Normalized discretized Frostman sets | WZ1 | WZ1: (8.1), p. 51}\par
\textbf{Mathematical statement.} A $(\delta,\alpha,C)_n$-Frostman set is a union of $\delta$-cubes with
$|E\cap B(x,r)|\le Cr^\alpha |E|$ for every $r\ge\delta$.
The analogue on $S^1$ uses unions of $\delta$-arcs. Unlike Definition 4.2, this condition is normalized by $|E|$, controls balls at one discretization scale, and need not pass to subsets with the same constant.
An ADset of volume at least $c\delta^{n-\alpha}$ is Frostman with constant $C/c$.
\smallskip\textbf{Lean statement sketch.}
\begin{Verbatim}[fontsize=\small,breaklines=true,breakanywhere=true,tabsize=2]
def FrostmanSet (delta alpha C : ℝ) (E : Set (Rn n)) : Prop :=
  GridUnion delta E ∧ ∀ x r, delta ≤ r →
    vol (E ∩ Metric.ball x r) ≤ C*r^alpha*vol E
theorem AD_to_Frostman
  (hAD : ADSet delta alpha C E)
  (hm : c*delta^((n:ℝ)-alpha) ≤ vol E) :
  FrostmanSet delta alpha (C/c) E
\end{Verbatim}
{\footnotesize\textbf{Mathlib support.} F, G, M, Q.\par
\textbf{Incoming edges.} \nodecard{Igeom} (api)\par}
{\footnotesize\textbf{Interface note.} Whenever probability measures are normalized by |E|, add 0<|E| and finite measurability hypotheses. Empty sets cause vacuous alternatives rather than legitimate normalized measures.}
\end{tcolorbox}\end{minipage}\par\medskip
\begin{minipage}{\linewidth}
\hypertarget{node-Dhyp}{}
\begin{tcolorbox}[colback=white,colframe=radialfg,colbacktitle=radialbg,coltitle=ink,boxrule=0.45pt,arc=1.6mm,left=3mm,right=3mm,top=2mm,bottom=2mm,title={\sffamily\bfseries WZ1 Definitions 8.2-8.4 \hfill \texttt{Dhyp}}]
{\sffamily\small\color{muted}Hypergraph fibres and uniform density | WZ1 | WZ1: pp. 52-53; (8.4)}\par
\textbf{Mathematical statement.} For $H\subset A_1\times\cdots\times A_k$, a coordinate subset $I$ and an edge $h$,
$n_H[h_I]=\{e\in H:e_i=h_i\ {\rm for}\ i\in I\}$.
The hypergraph is uniformly $c$-dense when
$\#n_H[h_I]\ge c\prod_{i\notin I}\#A_i$
for every $h\in H$ and every $I\subset[k]$.
For Section 8, vertices are occupied $\delta$-cells, as specified in Section 8.0.1.
\smallskip\textbf{Lean statement sketch.}
\begin{Verbatim}[fontsize=\small,breaklines=true,breakanywhere=true,tabsize=2]
def fiber (H : Finset (∀ i : Fin k, A i))
  (I : Finset (Fin k)) (h : ∀ i, A i) :=
  H.filter (fun e => ∀ i ∈ I, e i = h i)
def UniformDense (c : ℝ) (H : Finset (∀ i, A i)) : Prop :=
  ∀ h ∈ H, ∀ I : Finset (Fin k),
    c * ∏ i ∈ Iᶜ, (Fintype.card (A i) : ℝ) ≤ (fiber H I h).card
\end{Verbatim}
{\footnotesize\textbf{Mathlib support.} F, M, Q.\par
\textbf{Incoming edges.} \nodecard{Imass} (api)\par}
{\footnotesize\textbf{Interface note.} The product is over the unfixed coordinates i not in I. Surjective projection onto a chosen vertex class is needed for the induced-subgraph density observation in Remark 8.2.}
\end{tcolorbox}\end{minipage}\par\medskip
\begin{minipage}{\linewidth}
\hypertarget{node-L8_1}{}
\begin{tcolorbox}[colback=white,colframe=radialfg,colbacktitle=radialbg,coltitle=ink,boxrule=0.45pt,arc=1.6mm,left=3mm,right=3mm,top=2mm,bottom=2mm,title={\sffamily\bfseries WZ1 Lemma 8.1 \hfill \texttt{L8\_1}}]
{\sffamily\small\color{muted}Hypergraph refinement | WZ1 | WZ1: p. 53}\par
\textbf{Mathematical statement.} For a finite $k$-partite hypergraph $H$ of density $d=\#H/\prod_i\#A_i$ and $0<\varepsilon<1$, there is $H'\subset H$ with
$\#H'\ge(1-\varepsilon)\#H$ that is uniformly $(2^{-k}\varepsilon d)$-dense.
The source proves the general statement by pruning all under-popular coordinate fibres.
\smallskip\textbf{Lean statement sketch.}
\begin{Verbatim}[fontsize=\small,breaklines=true,breakanywhere=true,tabsize=2]
theorem uniform_dense_refinement
  (he : 0 < eps ∧ eps < 1) :
  ∃ H' ⊆ H, (1-eps)*(H.card:ℝ) ≤ (H'.card:ℝ) ∧
    UniformDense ((2:ℝ)^(-(k:ℤ))*eps*hypergraphDensity H) H'
\end{Verbatim}
{\footnotesize\textbf{Mathlib support.} F, Q.\par
\textbf{Incoming edges.} \nodecard{Dhyp} (core)\par}
{\footnotesize\textbf{Interface note.} The k=2, epsilon=1/2 antecedent is attributed to [14]. The general lemma is proved inside this paper and is not an external axiom.}
\end{tcolorbox}\end{minipage}\par\medskip
\begin{minipage}{\linewidth}
\hypertarget{node-T8_2}{}
\begin{tcolorbox}[colback=white,colframe=radialfg,colbacktitle=radialbg,coltitle=ink,boxrule=0.45pt,arc=1.6mm,left=3mm,right=3mm,top=2mm,bottom=2mm,title={\sffamily\bfseries WZ1 Theorem 8.2 \hfill \texttt{T8\_2}}]
{\sffamily\small\color{muted}A discretized Kaufman estimate with dense fibres | WZ1 | WZ1: p. 54}\par
\textbf{Mathematical statement.} For $0<\alpha\le\beta\le1$ and $\varepsilon>0$, let $F\subset\mathbb R^2$ and $\Lambda\subset S^1$ be $(\delta,\alpha,C)$- and $(\delta,\beta,C)$-Frostman sets.
If $H\subset\Lambda\times F$ has density $d=|H|/(|\Lambda||F|)$, there is $\Lambda'\subset\Lambda$ with $|\Lambda'|\gtrsim d|\Lambda|$ and, for every $\theta\in\Lambda'$,
$|\theta\cdot F_\theta|\gtrsim C^{-2}d^3\delta^{1-\alpha+\varepsilon}$, where $F_\theta=\{x:(\theta,x)\in H\}$.
\smallskip\textbf{Lean statement sketch.}
\begin{Verbatim}[fontsize=\small,breaklines=true,breakanywhere=true,tabsize=2]
theorem discretized_kaufman
  (hab : 0 < alpha ∧ alpha ≤ beta ∧ beta ≤ 1) (he : 0 < eps) :
  ∃ delta0 > 0, ∀ delta ∈ Set.Ioc 0 delta0,
  KaufmanHyp delta alpha beta C F Lambda H →
  ∃ Lambda' ⊆ Lambda,
    vol Lambda' ≥ c_K*density H*vol Lambda ∧
    ∀ theta ∈ Lambda',
      vol (dotImage theta (section H theta)) ≥
        c_K*C^(-2:ℤ)*(density H)^3*delta^(1-alpha+eps)
\end{Verbatim}
{\footnotesize\textbf{Mathlib support.} F, G, M, Q.\par
\textbf{Incoming edges.} \nodecard{L8_1} (core); \nodecard{DFrost} (core)\par}
{\footnotesize\textbf{Interface note.} This version is re-proved on pp. 54-55 using energy estimates and Lemma 8.1. [24] records historical provenance rather than an unproved input to this proof.}
\end{tcolorbox}\end{minipage}\par\medskip
\begin{minipage}{\linewidth}
\hypertarget{node-Dthin}{}
\begin{tcolorbox}[colback=white,colframe=radialfg,colbacktitle=radialbg,coltitle=ink,boxrule=0.45pt,arc=1.6mm,left=3mm,right=3mm,top=2mm,bottom=2mm,title={\sffamily\bfseries WZ1 Definition 8.5 \hfill \texttt{Dthin}}]
{\sffamily\small\color{muted}Thin tubes for probability measures | WZ1 | WZ1: (8.6), p. 55}\par
\textbf{Mathematical statement.} Probability measures $(\nu_1,\nu_2)$ have $(\beta,K,1-c)$-thin tubes when some Borel
$E\subset\operatorname{supp}\nu_1\times\operatorname{supp}\nu_2$ has $(\nu_1\times\nu_2)(E)\ge1-c$ and
$\nu_2(\{b_2\in N_r(\ell):(b_1,b_2)\in E\})\le Kr^\beta$
for all $b_1\in\operatorname{supp}\nu_1$, all lines $\ell$ through $b_1$, and all $r>0$.
The condition is directional: interchanging the measures is a separate hypothesis.
\smallskip\textbf{Lean statement sketch.}
\begin{Verbatim}[fontsize=\small,breaklines=true,breakanywhere=true,tabsize=2]
def ThinTubes (beta K c : ℝ) (nu1 nu2 : Measure (Rn 2)) : Prop :=
  IsProbabilityMeasure nu1 ∧ IsProbabilityMeasure nu2 ∧
  ∃ E, MeasurableSet E ∧ E ⊆ supportProduct nu1 nu2 ∧
    ENNReal.ofReal (1-c) ≤ (nu1.prod nu2) E ∧
    ∀ b1 ∈ measureSupport nu1, ∀ ell, b1 ∈ ell →
      ∀ r > 0, nu2 (section E b1 ∩ neighborhood r ell) ≤
        ENNReal.ofReal (K*r^beta)
\end{Verbatim}
{\footnotesize\textbf{Mathlib support.} G, M, Q.\par
\textbf{Incoming edges.} \nodecard{Igeom} (api)\par}

\end{tcolorbox}\end{minipage}\par\medskip
\begin{minipage}{\linewidth}
\hypertarget{node-DthinD}{}
\begin{tcolorbox}[colback=white,colframe=radialfg,colbacktitle=radialbg,coltitle=ink,boxrule=0.45pt,arc=1.6mm,left=3mm,right=3mm,top=2mm,bottom=2mm,title={\sffamily\bfseries WZ1 Definition 8.5' \hfill \texttt{DthinD}}]
{\sffamily\small\color{muted}Discretized thin tubes | WZ1 | WZ1: (8.7), p. 55}\par
\textbf{Mathematical statement.} For unions $G_1,G_2$ of $\delta$-squares, require a product-grid set $E\subset G_1\times G_2$ with $|E|\ge(1-c)|G_1||G_2|$ and
$|\{b_2\in G_2\cap N_r(\ell):(b_1,b_2)\in E\}|\le Kr^\beta|G_2|$
for every $b_1\in G_1$, every line through $b_1$, and every $r\ge\delta$.
We write this as $\operatorname{Thin}_\delta(\beta,K,c;G_1,G_2)$, where the final argument c denotes discarded mass.
\smallskip\textbf{Lean statement sketch.}
\begin{Verbatim}[fontsize=\small,breaklines=true,breakanywhere=true,tabsize=2]
def DiscreteThinTubes (delta beta K c : ℝ) (G1 G2 : Set (Rn 2)) : Prop :=
  ∃ E, ProductGridUnion delta E ∧ E ⊆ G1 ×ˢ G2 ∧
    (1-c)*vol G1*vol G2 ≤ vol E ∧
    ∀ b1 ∈ G1, ∀ ell, b1 ∈ ell → ∀ r, delta ≤ r →
      vol (section E b1 ∩ G2 ∩ neighborhood r ell) ≤ K*r^beta*vol G2
\end{Verbatim}
{\footnotesize\textbf{Mathlib support.} F, G, M, Q.\par
\textbf{Incoming edges.} \nodecard{Dthin} (api)\par}

\end{tcolorbox}\end{minipage}\par\medskip
\begin{minipage}{\linewidth}
\hypertarget{node-L8_3}{}
\begin{tcolorbox}[colback=white,colframe=radialfg,colbacktitle=radialbg,coltitle=ink,boxrule=0.45pt,arc=1.6mm,left=3mm,right=3mm,top=2mm,bottom=2mm,title={\sffamily\bfseries WZ1 Lemma 8.3 \hfill \texttt{L8\_3}}]
{\sffamily\small\color{muted}Pass from measure thin tubes to grid thin tubes | WZ1 | WZ1: pp. 55-56}\par
\textbf{Mathematical statement.} Let $\nu_i$ be normalized Lebesgue measures on the positive-volume grid sets $G_i$.
If $(\nu_1,\nu_2)$ has $(\beta,K,1-c)$-thin tubes, the paper states that $(G_1,G_2)$ has $\delta$-discretized $(\beta,K',1-c)$-thin tubes, with $K'=4^\beta K/c$.
This quantitative statement is recorded as printed; its use requires $c>0$.
\smallskip\textbf{Lean statement sketch.}
\begin{Verbatim}[fontsize=\small,breaklines=true,breakanywhere=true,tabsize=2]
theorem discretize_thin_tubes
  (hc : 0 < c ∧ c < 1)
  (h : ThinTubes beta K c (normalizedVolume G1) (normalizedVolume G2)) :
  DiscreteThinTubes delta beta ((4:ℝ)^beta*K/c) c G1 G2
\end{Verbatim}
{\footnotesize\textbf{Mathlib support.} F, G, M, Q.\par
\textbf{Incoming edges.} \nodecard{Dthin} (core); \nodecard{DthinD} (core)\par}
{\footnotesize\textbf{Interface note.} This atlas is a statement blueprint, not a verification of the measure-to-grid proof. Any quantitative implementation must check the bad-mass bookkeeping in that proof before adopting the printed constants.}
\end{tcolorbox}\end{minipage}\par\medskip
\begin{minipage}{\linewidth}
\hypertarget{node-Dsep}{}
\begin{tcolorbox}[colback=white,colframe=radialfg,colbacktitle=radialbg,coltitle=ink,boxrule=0.45pt,arc=1.6mm,left=3mm,right=3mm,top=2mm,bottom=2mm,title={\sffamily\bfseries WZ1 Definition 8.6 \hfill \texttt{Dsep}}]
{\sffamily\small\color{muted}Standard separation | WZ1 | WZ1: (8.11), p. 56}\par
\textbf{Mathematical statement.} For $F,G_1,G_2\subset[0,1]^2$, standard separation means
$\operatorname{diam}F,\operatorname{diam}G_1,\operatorname{diam}G_2\le1/10$,
$\operatorname{dist}(G_1,G_2)\ge1/2$, and $\operatorname{dist}(F,0)\ge1/2$.
The final reduction removes these restrictions by localization, common translation of $G_1,G_2$, and controlled dilations.
\smallskip\textbf{Lean statement sketch.}
\begin{Verbatim}[fontsize=\small,breaklines=true,breakanywhere=true,tabsize=2]
def StandardSeparation (F G1 G2 : Set (Rn 2)) : Prop :=
  diam F ≤ 1/10 ∧ diam G1 ≤ 1/10 ∧ diam G2 ≤ 1/10 ∧
  1/2 ≤ setDist G1 G2 ∧ 1/2 ≤ setDist F {0}
\end{Verbatim}
{\footnotesize\textbf{Mathlib support.} G, M.\par
\textbf{Incoming edges.} \nodecard{Igeom} (api)\par}

\end{tcolorbox}\end{minipage}\par\medskip
\begin{minipage}{\linewidth}
\hypertarget{node-L8_4}{}
\begin{tcolorbox}[colback=white,colframe=radialfg,colbacktitle=radialbg,coltitle=ink,boxrule=0.45pt,arc=1.6mm,left=3mm,right=3mm,top=2mm,bottom=2mm,title={\sffamily\bfseries WZ1 Lemma 8.4 \hfill \texttt{L8\_4}}]
{\sffamily\small\color{muted}Thin tubes force a large dot-product image | WZ1 | WZ1: (8.12), pp. 56-57}\par
\textbf{Mathematical statement.} For $\varepsilon>0$, sufficiently small $\delta$, and $K\ge1$, let $F,G_1,G_2$ be $(\delta,1,K)$-Frostman sets with standard separation.
If $(G_1,G_2)$ has discretized $(1,K,1-c)$-thin tubes and $H\subset F\times G_1\times G_2$ is uniformly $2c$-dense, then
$\mathcal E_\delta(\{a\cdot(b_1-b_2):(a,b_1,b_2)\in H\})\gtrsim c^5K^{-2}\delta^{\varepsilon-1}$.
\smallskip\textbf{Lean statement sketch.}
\begin{Verbatim}[fontsize=\small,breaklines=true,breakanywhere=true,tabsize=2]
theorem dot_diff_large_from_thin_tubes (he : 0 < eps) :
  ∃ delta0 > 0, ∀ delta ∈ Set.Ioc 0 delta0,
  ThinExpansionHyp delta K c F G1 G2 H →
    c_thin*c^5*K^(-2:ℤ)*delta^(eps-1) ≤
      coverNumR delta (dotDiffImage H)
\end{Verbatim}
{\footnotesize\textbf{Mathlib support.} F, G, M, Q.\par
\textbf{Incoming edges.} \nodecard{T8_2} (core); \nodecard{DthinD} (core); \nodecard{Dhyp} (core); \nodecard{DFrost} (core); \nodecard{Dsep} (core)\par}
{\footnotesize\textbf{Interface note.} Uniform 2c-density is stronger than the discarded-mass c and makes intersection with the good thin-tube pairs possible.}
\end{tcolorbox}\end{minipage}\par\medskip
\begin{minipage}{\linewidth}
\hypertarget{node-P8_6}{}
\begin{tcolorbox}[colback=white,colframe=radialfg,colbacktitle=radialbg,coltitle=ink,boxrule=0.45pt,arc=1.6mm,left=3mm,right=3mm,top=2mm,bottom=2mm,title={\sffamily\bfseries WZ1 Proposition 8.6 \hfill \texttt{P8\_6}}]
{\sffamily\small\color{muted}Radial bootstrap for measures | WZ1 | WZ1: pp. 57-58; [35, Lemma 2.9]}\par
\textbf{Mathematical statement.} Given $\beta,\varepsilon>0$, there are $\tau=\tau(\beta,\varepsilon)>0$ and $M=M(\beta,\varepsilon)\ge1$.
Let $\beta\le s\le1-\varepsilon$, $0<c<1/10$, $K,C\ge1$, and let separated probability measures satisfy $\nu_i(B(x,r))\le Cr$.
If both ordered pairs have $(s,K,1-c)$-thin tubes, both have $(s+\tau,K',1-3c)$-thin tubes, where
$K'=\max\{K,C^2M/c\}^{M}$.
\smallskip\textbf{Lean statement sketch.}
\begin{Verbatim}[fontsize=\small,breaklines=true,breakanywhere=true,tabsize=2]
theorem measure_radial_bootstrap (hb : 0 < beta) (he : 0 < eps) :
  ∃ tau > 0, ∃ M ≥ 1, ∀ s ∈ Set.Icc beta (1-eps),
  MeasureBootstrapHyp beta eps s c K C nu1 nu2 →
    ThinTubes (s+tau) (max K (C^2*M/c)^M) (3*c) nu1 nu2 ∧
    ThinTubes (s+tau) (max K (C^2*M/c)^M) (3*c) nu2 nu1
\end{Verbatim}
{\footnotesize\textbf{Mathlib support.} G, M, Q.\par
\textbf{Incoming edges.} \nodecard{XOSW} (core); \nodecard{Dthin} (core)\par}
{\footnotesize\textbf{Interface note.} s denotes the thin-tube exponent here, avoiding confusion with sigma\_3. The paper explains its compact-range uniformity and the freedom in beta in Remark 8.4; it ultimately fixes beta=1/4.}
\end{tcolorbox}\end{minipage}\par\medskip
\begin{minipage}{\linewidth}
\hypertarget{node-P8_6p}{}
\begin{tcolorbox}[colback=white,colframe=radialfg,colbacktitle=radialbg,coltitle=ink,boxrule=0.45pt,arc=1.6mm,left=3mm,right=3mm,top=2mm,bottom=2mm,title={\sffamily\bfseries WZ1 Proposition 8.6' \hfill \texttt{P8\_6p}}]
{\sffamily\small\color{muted}Radial bootstrap for grid sets | WZ1 | WZ1: p. 58}\par
\textbf{Mathematical statement.} Fix $\beta=1/4$. The same conclusion holds for separated $(\delta,1,C)$-Frostman grid sets and both ordered discretized thin-tube pairs with exponent $s\in[1/4,1-\varepsilon]$.
The resulting exponent is $s+\tau$, discarded mass is $3c$, and
$K'=\max\{4K/(3c),\,4C^2M/(3c^2)\}^{M}$.
\smallskip\textbf{Lean statement sketch.}
\begin{Verbatim}[fontsize=\small,breaklines=true,breakanywhere=true,tabsize=2]
theorem discrete_radial_bootstrap (he : 0 < eps) :
  ∃ tau > 0, ∃ M ≥ 1, ∀ s ∈ Set.Icc (1/4) (1-eps),
  DiscreteBootstrapHyp delta s c K C G1 G2 →
    BothDiscreteThin delta (s+tau)
      (max (4*K/(3*c)) (4*C^2*M/(3*c^2))^M) (3*c) G1 G2
\end{Verbatim}
{\footnotesize\textbf{Mathlib support.} F, G, M, Q.\par
\textbf{Incoming edges.} \nodecard{P8_6} (core); \nodecard{L8_3} (core)\par}

\end{tcolorbox}\end{minipage}\par\medskip
\begin{minipage}{\linewidth}
\hypertarget{node-C8_7}{}
\begin{tcolorbox}[colback=white,colframe=radialfg,colbacktitle=radialbg,coltitle=ink,boxrule=0.45pt,arc=1.6mm,left=3mm,right=3mm,top=2mm,bottom=2mm,title={\sffamily\bfseries WZ1 Corollary 8.7 \hfill \texttt{C8\_7}}]
{\sffamily\small\color{muted}Iterate the bootstrap up to exponent 1 | WZ1 | WZ1: p. 58}\par
\textbf{Mathematical statement.} For every $\varepsilon>0$ there is $N=N(\varepsilon)\ge1$.
If both ordered pairs of separated $(\delta,1,C)$-Frostman sets have discretized $(1/4,K,1-c)$-thin tubes, both have exponent $1-\varepsilon$, constant
$K'=\max\{4K/(3c),4C^2N/(3c^2)\}^{N}$, and discarded mass $3^Nc$.
At scales $r\ge\delta$, these are also exponent-1 thin tubes with constant $\delta^{-\varepsilon}K'$.
\smallskip\textbf{Lean statement sketch.}
\begin{Verbatim}[fontsize=\small,breaklines=true,breakanywhere=true,tabsize=2]
theorem thin_tubes_near_one (he : 0 < eps) :
  ∃ N : ℕ, 1 ≤ N ∧
  ∀ G1 G2, BootstrapSeedHyp delta c K C G1 G2 →
    BothDiscreteThin delta 1
      (delta^(-eps)*(max (4*K/(3*c)) (4*C^2*N/(3*c^2))^N))
      ((3:ℝ)^N*c) G1 G2
\end{Verbatim}
{\footnotesize\textbf{Mathlib support.} F, G, M, Q.\par
\textbf{Incoming edges.} \nodecard{P8_6p} (core); \nodecard{Iparams} (api)\par}
{\footnotesize\textbf{Interface note.} The term 3\textasciicircum{}N c is a power, not 3Nc. If it reaches 1, the probability-mass conclusion must be treated as vacuous or the initial c must be chosen smaller; the application does choose c sufficiently small.}
\end{tcolorbox}\end{minipage}\par\medskip
\begin{minipage}{\linewidth}
\hypertarget{node-L8_8}{}
\begin{tcolorbox}[colback=white,colframe=radialfg,colbacktitle=radialbg,coltitle=ink,boxrule=0.45pt,arc=1.6mm,left=3mm,right=3mm,top=2mm,bottom=2mm,title={\sffamily\bfseries WZ1 Lemma 8.8 \hfill \texttt{L8\_8}}]
{\sffamily\small\color{muted}Obtain a fixed 1/4-thin-tube seed | WZ1 | WZ1: (8.16), pp. 59-60}\par
\textbf{Mathematical statement.} For $\lambda,\zeta,\alpha>0$ and sufficiently small $\delta$, let separated $G_1,G_2$ be $(\delta,1,\delta^{-\lambda})$-Frostman sets.
If $|G_2\cap N_r(\ell)|\le(\delta^{-\lambda}r)^\zeta |G_2|$ for all lines and $r\ge\delta$, then
$(G_1,G_2)$ has discretized $(1/4,\delta^{-3\alpha/\zeta-\lambda},1-\delta^\alpha)$-thin tubes.
Swapping the roles gives the reverse ordered pair when both sets satisfy the strip bound.
\smallskip\textbf{Lean statement sketch.}
\begin{Verbatim}[fontsize=\small,breaklines=true,breakanywhere=true,tabsize=2]
theorem quarter_thin_seed
  (hl : 0 < lambda) (hz : 0 < zeta) (ha : 0 < alpha) :
  ∃ delta0 > 0, ∀ delta ∈ Set.Ioc 0 delta0,
  QuarterSeedHyp delta lambda zeta G1 G2 →
    DiscreteThinTubes delta (1/4)
      (delta^(-3*alpha/zeta-lambda)) (delta^alpha) G1 G2
\end{Verbatim}
{\footnotesize\textbf{Mathlib support.} F, G, M, Q.\par
\textbf{Incoming edges.} \nodecard{DFrost} (core)\par}
{\footnotesize\textbf{Interface note.} The fixed seed exponent is important: starting the iteration directly at zeta would make lambda depend on zeta, contrary to the quantifier contract of Proposition 8.5 (Remark 8.5).}
\end{tcolorbox}\end{minipage}\par\medskip
\begin{minipage}{\linewidth}
\hypertarget{node-P8_5}{}
\begin{tcolorbox}[colback=white,colframe=radialfg,colbacktitle=radialbg,coltitle=ink,boxrule=0.45pt,arc=1.6mm,left=3mm,right=3mm,top=2mm,bottom=2mm,title={\sffamily\bfseries WZ1 Proposition 8.5 \hfill \texttt{P8\_5}}]
{\sffamily\small\color{muted}Strip nonconcentration forces expansion | WZ1 | WZ1: (8.14)-(8.15), p. 57; proof p. 61}\par
\textbf{Mathematical statement.} For $0<\varepsilon,\zeta<1$, choose $\lambda=\lambda(\varepsilon)>0$, $\alpha=\alpha(\varepsilon,\zeta)>0$ and $\delta_0$.
Let $F,G_1,G_2$ be $(\delta,1,\delta^{-\lambda})$-Frostman sets with standard separation. Require
$|G_i\cap N_r(\ell)|\le(\delta^{-\lambda}r)^\zeta|G_i|$ for all lines and $r\ge\delta$.
For uniformly $\delta^\alpha$-dense $H$,
$\mathcal E_\delta(\{a\cdot(b_1-b_2):(a,b_1,b_2)\in H\})\ge\delta^{\varepsilon-1}$.
\smallskip\textbf{Lean statement sketch.}
\begin{Verbatim}[fontsize=\small,breaklines=true,breakanywhere=true,tabsize=2]
theorem strip_nonconcentration_expands (he : 0 < eps ∧ eps < 1) :
  ∃ lambda > 0, ∀ zeta ∈ Set.Ioo 0 1,
  ∃ alpha > 0, ∃ delta0 > 0, ∀ delta ∈ Set.Ioc 0 delta0,
  StripExpansionHyp delta lambda zeta alpha F G1 G2 H →
    delta^(eps-1) ≤ coverNumR delta (dotDiffImage H)
\end{Verbatim}
{\footnotesize\textbf{Mathlib support.} F, G, M, Q.\par
\textbf{Incoming edges.} \nodecard{L8_8} (core); \nodecard{C8_7} (core); \nodecard{L8_4} (core); \nodecard{Iparams} (api)\par}
{\footnotesize\textbf{Interface note.} Lambda is chosen before zeta. This dependency matters in the anisotropic reduction and must not be weakened in a Lean interface.}
\end{tcolorbox}\end{minipage}\par\medskip
\begin{minipage}{\linewidth}
\hypertarget{node-L8_10}{}
\begin{tcolorbox}[colback=white,colframe=radialfg,colbacktitle=radialbg,coltitle=ink,boxrule=0.45pt,arc=1.6mm,left=3mm,right=3mm,top=2mm,bottom=2mm,title={\sffamily\bfseries WZ1 Lemma 8.10 \hfill \texttt{L8\_10}}]
{\sffamily\small\color{muted}Two-ends localization to a strip | WZ1 | WZ1: (8.22), p. 62}\par
\textbf{Mathematical statement.} For a $\delta$-square union $E\subset[0,1]^2$ and $0<\zeta\le1/4$, some width $w>0$ and line $\ell$ give $E'=E\cap N_w(\ell)$ with
$|E'|\ge\tfrac12w^\zeta |E|$ and
$|E'\cap N_r(\ell')|\le(r/w)^\zeta|E'|$
for every line $\ell'$ and every $0<r\le1$.
\smallskip\textbf{Lean statement sketch.}
\begin{Verbatim}[fontsize=\small,breaklines=true,breakanywhere=true,tabsize=2]
theorem two_ends_strip_localization
  (hz : 0 < zeta ∧ zeta ≤ 1/4) (hE : GridUnion delta E) :
  ∃ w > 0, ∃ ell, let E' := E ∩ neighborhood w ell
  (1/2)*w^zeta*vol E ≤ vol E' ∧
  ∀ ell' r, 0 < r → r ≤ 1 →
    vol (E' ∩ neighborhood r ell') ≤ (r/w)^zeta*vol E'
\end{Verbatim}
{\footnotesize\textbf{Mathlib support.} F, G, M, Q.\par
\textbf{Incoming edges.} \nodecard{Ientropy} (api)\par}
{\footnotesize\textbf{Interface note.} The paper calls this standard and omits the proof. It still needs an implementation; it is not an external citation.}
\end{tcolorbox}\end{minipage}\par\medskip
\begin{minipage}{\linewidth}
\hypertarget{node-L8_11}{}
\begin{tcolorbox}[colback=white,colframe=radialfg,colbacktitle=radialbg,coltitle=ink,boxrule=0.45pt,arc=1.6mm,left=3mm,right=3mm,top=2mm,bottom=2mm,title={\sffamily\bfseries WZ1 Lemma 8.11 \hfill \texttt{L8\_11}}]
{\sffamily\small\color{muted}Restore Frostman after anisotropic rescaling | WZ1 | WZ1: (8.23), p. 62}\par
\textbf{Mathematical statement.} For $\varepsilon>0$, sufficiently small $\delta$, $w\in[\delta^{1-\varepsilon},1]$, a $1\times w$ rectangle $R$, a $(\delta,1,C)$-Frostman set $E\subset R$, and the affine map $\phi:R\to[0,1]^2$:
some $E'\subset E$ has $|E'|\ge(\delta/w)^\varepsilon|E|$ such that the union of $\delta/w$-squares meeting $\phi(E')$ is $(\delta/w,1,(w/\delta)^\varepsilon C)$-Frostman.
Each occupied output square has comparable intersection volume with $\phi(E')$.
\smallskip\textbf{Lean statement sketch.}
\begin{Verbatim}[fontsize=\small,breaklines=true,breakanywhere=true,tabsize=2]
theorem anisotropic_Frostman_refinement (he : 0 < eps) :
  ∃ delta0 > 0, ∀ delta ∈ Set.Ioc 0 delta0,
  AffineFrostmanHyp delta eps w C E R phi →
  ∃ E' ⊆ E, (delta/w)^eps*vol E ≤ vol E' ∧
    FrostmanSet (delta/w) 1 ((w/delta)^eps*C)
      (recubulate (delta/w) (phi '' E')) ∧
    BalancedOccupancy (delta/w) (phi '' E')
\end{Verbatim}
{\footnotesize\textbf{Mathlib support.} A, F, G, M, Q.\par
\textbf{Incoming edges.} \nodecard{DFrost} (core); \nodecard{Iaffine} (api); \nodecard{Imass} (api)\par}
{\footnotesize\textbf{Interface note.} Normalized Frostman constants are not inherited unchanged by arbitrary refinement. The occupancy condition is precisely what makes this transport estimate work.}
\end{tcolorbox}\end{minipage}\par\medskip
\begin{minipage}{\linewidth}
\hypertarget{node-L8_12}{}
\begin{tcolorbox}[colback=white,colframe=radialfg,colbacktitle=radialbg,coltitle=ink,boxrule=0.45pt,arc=1.6mm,left=3mm,right=3mm,top=2mm,bottom=2mm,title={\sffamily\bfseries WZ1 Lemma 8.12 \hfill \texttt{L8\_12}}]
{\sffamily\small\color{muted}Frostman coarsening after pigeonholing | WZ1 | WZ1: pp. 62-63}\par
\textbf{Mathematical statement.} For $0<\delta\le\rho\le1$ and a $(\delta,\alpha,C)$-Frostman set $E\subset\mathbb R^2$, there is $E'\subset E$ with $|E'|\ge|\log\delta|^{-1}|E|$ such that the union of $\rho$-squares meeting $E'$ is $(\rho,\alpha,C')$-Frostman, with $C'\asymp|\log\delta|C$.
\smallskip\textbf{Lean statement sketch.}
\begin{Verbatim}[fontsize=\small,breaklines=true,breakanywhere=true,tabsize=2]
theorem coarsen_Frostman
  (hs : 0 < delta ∧ delta < 1 ∧ delta ≤ rho ∧ rho ≤ 1)
  (hE : FrostmanSet delta alpha C E) :
  ∃ E' ⊆ E, |Real.log delta|⁻¹*vol E ≤ vol E' ∧
    FrostmanSet rho alpha (C_coarse*|Real.log delta|*C)
      (recubulate rho E')
\end{Verbatim}
{\footnotesize\textbf{Mathlib support.} F, G, M, Q.\par
\textbf{Incoming edges.} \nodecard{DFrost} (core); \nodecard{Imass} (api)\par}
{\footnotesize\textbf{Interface note.} The printed formula is meaningful only below scale 1. In applications delta is sufficiently small and all logarithmic factors are absorbed into power allowances.}
\end{tcolorbox}\end{minipage}\par\medskip
\begin{minipage}{\linewidth}
\hypertarget{node-L8_13}{}
\begin{tcolorbox}[colback=white,colframe=radialfg,colbacktitle=radialbg,coltitle=ink,boxrule=0.45pt,arc=1.6mm,left=3mm,right=3mm,top=2mm,bottom=2mm,title={\sffamily\bfseries WZ1 Lemma 8.13 \hfill \texttt{L8\_13}}]
{\sffamily\small\color{muted}Comparable orthogonal strips or expansion | WZ1 | WZ1: (8.3), pp. 63-65}\par
\textbf{Mathematical statement.} For sufficiently small $\delta$ and separated $(\delta,1,\delta^{-\eta})$-Frostman $F,G_1,G_2$, let $H$ be uniformly $\delta^\eta$-dense and suppose $G_1\subset N_w(\ell)$.
Either
$\pi_{G_2}H\subset N_{\delta^{-\varepsilon}w}(\ell)$ and
$\pi_FH\subset N_{\delta^{-\varepsilon}w}(\ell^\perp)$,
where $\ell^\perp$ passes through the origin; or the expansion conclusion (8.3) holds at some $\rho\ge\delta$.
\smallskip\textbf{Lean statement sketch.}
\begin{Verbatim}[fontsize=\small,breaklines=true,breakanywhere=true,tabsize=2]
theorem comparable_strip_or_expansion (he : 0 < eps) :
  ∃ eta > 0, ∃ delta0 > 0, ∀ delta ∈ Set.Ioc 0 delta0,
  SeparatedUniformHyp delta eta F G1 G2 H →
  G1 ⊆ neighborhood w ell →
    (projG2 H ⊆ neighborhood (delta^(-eps)*w) ell ∧
     projF H ⊆ neighborhood (delta^(-eps)*w) (originPerp ell)) ∨
    ExpansionConclusion delta eta eps (dotDiffImage H)
\end{Verbatim}
{\footnotesize\textbf{Mathlib support.} A, F, G, M, Q.\par
\textbf{Incoming edges.} \nodecard{L8_1} (core); \nodecard{L8_11} (core); \nodecard{L8_12} (core); \nodecard{T8_2} (core); \nodecard{Dsep} (core); \nodecard{Iaffine} (api)\par}

\end{tcolorbox}\end{minipage}\par\medskip
\begin{minipage}{\linewidth}
\hypertarget{node-P8_9}{}
\begin{tcolorbox}[colback=white,colframe=radialfg,colbacktitle=radialbg,coltitle=ink,boxrule=0.45pt,arc=1.6mm,left=3mm,right=3mm,top=2mm,bottom=2mm,title={\sffamily\bfseries WZ1 Proposition 8.9 \hfill \texttt{P8\_9}}]
{\sffamily\small\color{muted}Separated, uniform version of the dichotomy | WZ1 | WZ1: p. 61; proof pp. 66-67}\par
\textbf{Mathematical statement.} For each $\varepsilon>0$, choose $\eta,\delta_0>0$.
If $F,G_1,G_2$ are $(\delta,1,\delta^{-\eta})$-Frostman sets with standard separation and $H\subset F\times G_1\times G_2$ is uniformly $\delta^\eta$-dense, either the three orthogonal-line concentration bounds (8.2) hold, or the expansion conclusion (8.3) holds for this $H$.
\smallskip\textbf{Lean statement sketch.}
\begin{Verbatim}[fontsize=\small,breaklines=true,breakanywhere=true,tabsize=2]
theorem separated_uniform_dichotomy (he : 0 < eps) :
  ∃ eta > 0, ∃ delta0 > 0, ∀ delta ∈ Set.Ioc 0 delta0,
  SeparatedUniformHyp delta eta F G1 G2 H →
    OrthogonalLineConcentration delta eps F G1 G2 ∨
    ExpansionConclusion delta eta eps (dotDiffImage H)
\end{Verbatim}
{\footnotesize\textbf{Mathlib support.} A, F, G, M, Q.\par
\textbf{Incoming edges.} \nodecard{L8_10} (core); \nodecard{L8_13} (core); \nodecard{L8_11} (core); \nodecard{L8_1} (core); \nodecard{P8_5} (core); \nodecard{Iaffine} (api)\par}
{\footnotesize\textbf{Interface note.} The proof localizes to strips, compares widths, uses dual affine changes for F and the G sets, restores Frostman constants, and applies Proposition 8.5.}
\end{tcolorbox}\end{minipage}\par\medskip
\begin{minipage}{\linewidth}
\hypertarget{node-T5_2p}{}
\begin{tcolorbox}[colback=white,colframe=radialfg,colbacktitle=radialbg,coltitle=ink,boxrule=0.45pt,arc=1.6mm,left=3mm,right=3mm,top=2mm,bottom=2mm,title={\sffamily\bfseries WZ1 Theorem 5.2' \hfill \texttt{T5\_2p}}]
{\sffamily\small\color{muted}Frostman form, without separation or uniformity | WZ1 | WZ1: pp. 51-52; proof pp. 68-69}\par
\textbf{Mathematical statement.} For every $\varepsilon>0$, choose $\eta,\delta_0>0$.
For $(\delta,1,\delta^{-\eta})$-Frostman $F,G_1,G_2\subset[0,1]^2$, either an affine line $\ell$ and its origin-based perpendicular satisfy
$\mathcal E_\delta(F\cap N_\delta\ell),\mathcal E_\delta(G_i\cap N_\delta\ell^\perp)\ge\delta^{\varepsilon-1}$ for $i=1,2$; or
every $H\subset F\times G_1\times G_2$ with $\mathcal E_\delta(H)\ge\delta^{\eta-3}$ satisfies the multiscale expansion conclusion (8.3).
\smallskip\textbf{Lean statement sketch.}
\begin{Verbatim}[fontsize=\small,breaklines=true,breakanywhere=true,tabsize=2]
theorem dot_product_Frostman_dichotomy (he : 0 < eps) :
  ∃ eta > 0, ∃ delta0 > 0, ∀ delta ∈ Set.Ioc 0 delta0,
  FrostmanTripleHyp delta eta F G1 G2 →
    OrthogonalLineConcentration delta eps F G1 G2 ∨
    ∀ H, HyperSubset H F G1 G2 →
      delta^(eta-3) ≤ coverNumR delta H →
      ExpansionConclusion delta eta eps (dotDiffImage H)
\end{Verbatim}
{\footnotesize\textbf{Mathlib support.} A, F, G, M, Q.\par
\textbf{Incoming edges.} \nodecard{P8_9} (core); \nodecard{L8_1} (core); \nodecard{Iaffine} (api)\par}
{\footnotesize\textbf{Interface note.} The edge to Theorem 5.2 includes the small-set/vacuous-expansion case on p. 52. AD-to-Frostman conversion alone is insufficient without a lower size bound.}
\end{tcolorbox}\end{minipage}\par\medskip
\sectitle{Node catalogue: External interfaces}
\begin{minipage}{\linewidth}
\hypertarget{node-XBCT}{}
\begin{tcolorbox}[colback=white,colframe=externalfg,colbacktitle=externalbg,coltitle=ink,boxrule=0.45pt,arc=1.6mm,left=3mm,right=3mm,top=2mm,bottom=2mm,title={\sffamily\bfseries External [1] / [8] \hfill \texttt{XBCT}}]
{\sffamily\small\color{muted}Endpoint multilinear Kakeya | External | WZ1: Direct use: Theorem 4.2, p. 19; bibliography pp. 69-70}\par
\textbf{Mathematical statement.} Imported interface: the determinant-weighted trilinear tube integral is bounded by $(\delta^2\#\mathcal T)^{3/2}$, as recorded in Theorem 4.2.
The paper credits Bennett--Carbery--Tao [1], and points to the endpoint version of Carbery--Valdimarsson [8].
\smallskip\textbf{Lean statement sketch.}
\begin{Verbatim}[fontsize=\small,breaklines=true,breakanywhere=true,tabsize=2]
axiom endpoint_multilinear_kakeya :
  ∀ delta Ts, NormalizedTubes delta Ts →
    MultilinearKakeyaBound delta Ts
-- MultilinearKakeyaBound is the formula in node T4_2.
\end{Verbatim}
{\footnotesize\textbf{Mathlib support.} A, F, G, M.\par
\textbf{Incoming edges.} None; foundational or external interface.\par}
{\footnotesize\textbf{Interface note.} Black-box import. No claim is made about its current formalization status.}
\end{tcolorbox}\end{minipage}\par\medskip
\begin{minipage}{\linewidth}
\hypertarget{node-XExt}{}
\begin{tcolorbox}[colback=white,colframe=externalfg,colbacktitle=externalbg,coltitle=ink,boxrule=0.45pt,arc=1.6mm,left=3mm,right=3mm,top=2mm,bottom=2mm,title={\sffamily\bfseries External [29] / Mathlib \hfill \texttt{XExt}}]
{\sffamily\small\color{muted}Scalar Lipschitz extension | External | WZ1: Direct use: proof of Proposition 4.1, p. 31}\par
\textbf{Mathematical statement.} A real-valued $K$-Lipschitz function defined on a subset of a metric space extends to a $K$-Lipschitz function on the whole space.
Only the scalar extension of the slope $f$ is needed at this citation site. The paper cites Kirszbraun [29]; the scalar interface is already provided by Mathlib.
\smallskip\textbf{Lean statement sketch.}
\begin{Verbatim}[fontsize=\small,breaklines=true,breakanywhere=true,tabsize=2]
-- Existing Mathlib theorem, usable after representing f ambiently:
LipschitzOnWith.extend_real :
  LipschitzOnWith K f S →
  ∃ g, LipschitzWith K g ∧ Set.EqOn f g S
\end{Verbatim}
{\footnotesize\textbf{Mathlib support.} L.\par
\textbf{Incoming edges.} None; foundational or external interface.\par}
{\footnotesize\textbf{Interface note.} This is a support-library fact, not an additional analytic axiom. For a Euclidean-valued extension with the Euclidean norm, use the actual Hilbert-space extension theorem rather than a sup-norm product extension.}
\end{tcolorbox}\end{minipage}\par\medskip
\begin{minipage}{\linewidth}
\hypertarget{node-XOSW}{}
\begin{tcolorbox}[colback=white,colframe=externalfg,colbacktitle=externalbg,coltitle=ink,boxrule=0.45pt,arc=1.6mm,left=3mm,right=3mm,top=2mm,bottom=2mm,title={\sffamily\bfseries External [35, Lemma 2.9] \hfill \texttt{XOSW}}]
{\sffamily\small\color{muted}Radial thin-tube bootstrap | External | WZ1: Direct use: Proposition 8.6, pp. 57-58; bibliography p. 71}\par
\textbf{Mathematical statement.} Imported interface: improve the thin-tube exponent by a fixed positive amount on a compact subrange below 1, for separated 1-Frostman probability measures, with quantitative loss in the thin-tube constant and good-pair mass.
The exact specialized statement is Proposition 8.6. Source: Orponen--Shmerkin--Wang [35, Lemma 2.9].
\smallskip\textbf{Lean statement sketch.}
\begin{Verbatim}[fontsize=\small,breaklines=true,breakanywhere=true,tabsize=2]
axiom OSW_radial_bootstrap :
  ∀ beta eps, 0 < beta → 0 < eps →
    MeasureRadialBootstrapContract beta eps
-- Contract includes tau(beta,eps), M(beta,eps),
-- both ordered pairs, and the constants in node P8_6.
\end{Verbatim}
{\footnotesize\textbf{Mathlib support.} G, M, Q.\par
\textbf{Incoming edges.} None; foundational or external interface.\par}
{\footnotesize\textbf{Interface note.} The paper's Remark 8.4 discusses how its stated interface differs from the cited lemma. That explanation is part of the import boundary.}
\end{tcolorbox}\end{minipage}\par\medskip
\begin{minipage}{\linewidth}
\hypertarget{node-XPYZ}{}
\begin{tcolorbox}[colback=white,colframe=externalfg,colbacktitle=externalbg,coltitle=ink,boxrule=0.45pt,arc=1.6mm,left=3mm,right=3mm,top=2mm,bottom=2mm,title={\sffamily\bfseries External: Pramanik-Yang-Zahl \hfill \texttt{XPYZ}}]
{\sffamily\small\color{muted}Cinematic L3/2 maximal theorem | External | WZ1 [36, Theorem 1.7] = WZ2 [13, Theorem 1.7]; cinematic definition [36, Definition 1.6]}\par
\textbf{Mathematical statement.} Imported interface: a separated one-dimensional Frostman subfamily of a uniformly cinematic, doubling $C^2$ family has the $L^{3/2}$ curve-neighbourhood estimate stated in Theorem 7.2.
Source: Pramanik--Yang--Zahl [36, Theorem 1.7]; the cinematic definition is adapted from [36, Definition 1.6].
\smallskip\textbf{Lean statement sketch.}
\begin{Verbatim}[fontsize=\small,breaklines=true,breakanywhere=true,tabsize=2]
axiom PYZ_cinematic_maximal :
  ∀ I K D Fs, CinematicFamily I K D Fs →
    CinematicMaximalContract I K D Fs
-- Contract is node T7_2, including its small-scale threshold.
\end{Verbatim}
{\footnotesize\textbf{Mathlib support.} C, F, G, M, P.\par
\textbf{Incoming edges.} None; foundational or external interface.\par}

\end{tcolorbox}\end{minipage}\par\medskip
\begin{minipage}{\linewidth}
\hypertarget{node-XFO}{}
\begin{tcolorbox}[colback=white,colframe=externalfg,colbacktitle=externalbg,coltitle=ink,boxrule=0.45pt,arc=1.6mm,left=3mm,right=3mm,top=2mm,bottom=2mm,title={\sffamily\bfseries External: Fässler-Orponen Lem. 3.13 \hfill \texttt{XFO}}]
{\sffamily\small\color{muted}Content-to-Katz-Tao extraction | External | WZ2 Lemma 7.7, PDF p. 45; WZ2 [3, Lemma 3.13]}\par
\textbf{Mathematical statement.} If the $s$-dimensional Hausdorff content of a bounded set is $\kappa>0$, then at mesh $\delta$ it contains a $(\delta,s,O(1))$-Katz-Tao set with covering number $\gtrsim\kappa\delta^{-s}$.
The fixed constant may be taken as $100$ in the WZ2 normalization.
\smallskip\textbf{Lean statement sketch.}
\begin{Verbatim}[fontsize=\small,breaklines=true,breakanywhere=true,tabsize=2]
axiom FO_content_extraction : ContentKatzTaoContract
-- Used after mesh thickening, as in WZ2 Lemmas 7.6-7.7.
\end{Verbatim}
{\footnotesize\textbf{Mathlib support.} F, G, M, Q.\par
\textbf{Incoming edges.} None; foundational or external interface.\par}

\end{tcolorbox}\end{minipage}\par\medskip
\begin{minipage}{\linewidth}
\hypertarget{node-XOSuniform}{}
\begin{tcolorbox}[colback=white,colframe=externalfg,colbacktitle=externalbg,coltitle=ink,boxrule=0.45pt,arc=1.6mm,left=3mm,right=3mm,top=2mm,bottom=2mm,title={\sffamily\bfseries External: Orponen-Shmerkin Lem. 2.15 \hfill \texttt{XOSuniform}}]
{\sffamily\small\color{muted}Uniform refinement black box | External | WZ2 [11, Lemma 2.15], cited as WZ2 Lemma 7.9}\par
\textbf{Mathematical statement.} A set at a prescribed sparse dyadic mesh admits a uniform occupied-cell refinement with the polylogarithmic loss stated in WZ2 Lemma 7.9.
No further proof from the cited paper is opened here.
\smallskip\textbf{Lean statement sketch.}
\begin{Verbatim}[fontsize=\small,breaklines=true,breakanywhere=true,tabsize=2]
axiom OS_uniform_refinement : SparseDyadicUniformRefinementContract
\end{Verbatim}
{\footnotesize\textbf{Mathlib support.} F, G, M, Q.\par
\textbf{Incoming edges.} None; foundational or external interface.\par}

\end{tcolorbox}\end{minipage}\par\medskip
\begin{minipage}{\linewidth}
\hypertarget{node-XSWdecomp}{}
\begin{tcolorbox}[colback=white,colframe=externalfg,colbacktitle=externalbg,coltitle=ink,boxrule=0.45pt,arc=1.6mm,left=3mm,right=3mm,top=2mm,bottom=2mm,title={\sffamily\bfseries External: Shmerkin-Wang Lem. 5.21 \hfill \texttt{XSWdecomp}}]
{\sffamily\small\color{muted}Decompose the scale-entropy profile | External | WZ2 [14, Lemma 5.21]; WZ2 proof of Lemma 7.10, PDF p. 46}\par
\textbf{Mathematical statement.} For a uniform set, decompose the scale range into finitely many intervals $[a_j,a_{j+1}]$ of a fixed minimum length, with per-interval Frostman exponents $t_j$ and weighted entropy
$\sum_jt_j(a_{j+1}-a_j)\ge s-o(1)$.
The quoted quantitative decomposition is used as a black box.
\smallskip\textbf{Lean statement sketch.}
\begin{Verbatim}[fontsize=\small,breaklines=true,breakanywhere=true,tabsize=2]
axiom SW_scale_decomposition : ScaleEntropyDecompositionContract
\end{Verbatim}
{\footnotesize\textbf{Mathlib support.} F, G, M, Q.\par
\textbf{Incoming edges.} None; foundational or external interface.\par}

\end{tcolorbox}\end{minipage}\par\medskip
\begin{minipage}{\linewidth}
\hypertarget{node-XSWmerge}{}
\begin{tcolorbox}[colback=white,colframe=externalfg,colbacktitle=externalbg,coltitle=ink,boxrule=0.45pt,arc=1.6mm,left=3mm,right=3mm,top=2mm,bottom=2mm,title={\sffamily\bfseries External: Shmerkin-Wang Lem. 5.20 \hfill \texttt{XSWmerge}}]
{\sffamily\small\color{muted}Merge adjacent entropy intervals | External | WZ2 [14, Lemma 5.20]; WZ2 proof of Lemma 7.10, PDF p. 46}\par
\textbf{Mathematical statement.} Adjacent intervals with $t_j\ge t_{j+1}$ can be merged; their new exponent is the length-weighted average and their Frostman loss remains controlled.
The weighted entropy and the total scale length are preserved.
\smallskip\textbf{Lean statement sketch.}
\begin{Verbatim}[fontsize=\small,breaklines=true,breakanywhere=true,tabsize=2]
axiom SW_merging_lemma : AdjacentScaleMergingContract
\end{Verbatim}
{\footnotesize\textbf{Mathlib support.} F, G, M, Q.\par
\textbf{Incoming edges.} None; foundational or external interface.\par}

\end{tcolorbox}\end{minipage}\par\medskip
\sectitle{Node catalogue: API}
\begin{minipage}{\linewidth}
\hypertarget{node-Igeom}{}
\begin{tcolorbox}[colback=white,colframe=apifg,colbacktitle=apibg,coltitle=ink,boxrule=0.45pt,arc=1.6mm,left=3mm,right=3mm,top=2mm,bottom=2mm,title={\sffamily\bfseries Geometry API \hfill \texttt{Igeom}}]
{\sffamily\small\color{muted}Lines, tubes, cubes, shadings | WZ2 | WZ2 Section 6.1, PDF pp. 37-38; WZ1 Section 2.1}\par
\textbf{Mathematical statement.} Work in a bounded affine-line chart in $\mathbb R^3$. Throughout the structural proof, a tube is
$T_\delta(\ell)=N_{6\delta}(\ell)\cap[-1,1]^3$ and $Y(T)\subset T$ is a union of axis-aligned $\delta$-cubes.
Put $E_C=\bigcup_TY(T)$ and $\mathcal T_C(p)=\{T:p\in Y(T)\}$.
The original unit-segment tubes and measurable shadings are converted by the bounded-loss cubulation in WZ2 Section 6.1.
\smallskip\textbf{Lean statement sketch.}
\begin{Verbatim}[fontsize=\small,breaklines=true,breakanywhere=true,tabsize=2]
abbrev Rn (n : ℕ) := EuclideanSpace ℝ (Fin n)
structure LineChart (n : ℕ) where
  base : Rn n
  direction : Rn n
  -- unit direction, base in the chart, last coordinate >= 1/2
structure Config (n : ℕ) (delta : ℝ) where
  tubes : Finset (LineChart n)
  shade : LineChart n → Set (Rn n)
  -- grid-union and tube-containment certificates
\end{Verbatim}
{\footnotesize\textbf{Mathlib support.} F, G, M.\par
\textbf{Incoming edges.} None; foundational or external interface.\par}
{\footnotesize\textbf{Interface note.} WZ2 Section 6.1 expressly adopts the WZ1 tube and cube conventions. A finite chart split and bounded dilations are infrastructure, not a direction-separation assumption. The threshold condition for cubulation is interpreted as a cell-density condition; harmless constants are absorbed before the final scale choice.}
\end{tcolorbox}\end{minipage}\par\medskip
\begin{minipage}{\linewidth}
\hypertarget{node-Ientropy}{}
\begin{tcolorbox}[colback=white,colframe=apifg,colbacktitle=apibg,coltitle=ink,boxrule=0.45pt,arc=1.6mm,left=3mm,right=3mm,top=2mm,bottom=2mm,title={\sffamily\bfseries Entropy API \hfill \texttt{Ientropy}}]
{\sffamily\small\color{muted}Finite covering numbers | API | WZ2 Section 6.1 and Definitions 7.2-7.3; WZ1 Section 8}\par
\textbf{Mathematical statement.} $\mathcal E_r(S)$ denotes the covering number of $S$ at radius/mesh $r$.
For bounded grid sets, covering numbers, separated-point counts and occupied-cube counts agree up to dimensional constants. Fix one convention and prove comparison lemmas. Neighbourhoods are $N_r(S)=\{x:\operatorname{dist}(x,S)<r\}$.
\smallskip\textbf{Lean statement sketch.}
\begin{Verbatim}[fontsize=\small,breaklines=true,breakanywhere=true,tabsize=2]
noncomputable def coverNum (r : ℝ) (S : Set X) : ℕ := ...
def neighborhood (r : ℝ) (S : Set X) : Set X := ...
theorem grid_entropy_comparable
  (h : GridUnion delta S) (hr : delta ≤ r) :
  Comparable (coverNum r S) (occupiedCubeCount r S)
\end{Verbatim}
{\footnotesize\textbf{Mathlib support.} F, G.\par
\textbf{Incoming edges.} \nodecard{Igeom} (api)\par}
{\footnotesize\textbf{Interface note.} CoverNum is a proposed project definition, not an assertion about an existing Mathlib declaration. Use Nat counts with explicit real coercions in inequalities.}
\end{tcolorbox}\end{minipage}\par\medskip
\begin{minipage}{\linewidth}
\hypertarget{node-Imass}{}
\begin{tcolorbox}[colback=white,colframe=apifg,colbacktitle=apibg,coltitle=ink,boxrule=0.45pt,arc=1.6mm,left=3mm,right=3mm,top=2mm,bottom=2mm,title={\sffamily\bfseries Refinement API \hfill \texttt{Imass}}]
{\sffamily\small\color{muted}Mass and constant multiplicity | WZ2 | WZ2 Definition 2.1 and Section 6.1, PDF pp. 6, 37-38}\par
\textbf{Mathematical statement.} Write $m(C)=\sum_T|Y(T)|$ and $v(C)=\sum_T|T|\asymp\delta^2\#\mathcal T$.
A $\lambda$-dense shading means $m(C)\ge\lambda v(C)$; uniformly dense means $|Y(T)|\ge\lambda|T|$ for every tube.
A refinement retains $m(C')\gtrsim |\log\delta|^{-A}m(C)$. Constant-multiplicity pigeonholing gives $\sum_T1_{Y'(T)}=\mu 1_{E_{C'}}$.
Power-loss selections are separately recorded as subcollections with explicit mass retention.
\smallskip\textbf{Lean statement sketch.}
\begin{Verbatim}[fontsize=\small,breaklines=true,breakanywhere=true,tabsize=2]
def Dense (lambda : ℝ) (C : Config 3 delta) : Prop :=
  lambda * tubeVolumeMass C ≤ shadeMass C
def UniformlyDense (lambda : ℝ) (C : Config 3 delta) : Prop :=
  ∀ T ∈ C.tubes, lambda * vol (tube T delta) ≤ vol (C.shade T)
def Refinement (C' C : Config 3 delta) : Prop :=
  Subcollection C' C ∧ PolylogMassRetention C' C
-- Constant-multiplicity refinement is the inherited finite-grid operation.
\end{Verbatim}
{\footnotesize\textbf{Mathlib support.} F, M, Q.\par
\textbf{Incoming edges.} \nodecard{Igeom} (api)\par}
{\footnotesize\textbf{Interface note.} Average density is the hypothesis of Theorem 5.2. Uniform density is obtained by selecting popular tubes and then rebuilding the required covers; it is not silently assumed at the root.}
\end{tcolorbox}\end{minipage}\par\medskip
\begin{minipage}{\linewidth}
\hypertarget{node-Icover}{}
\begin{tcolorbox}[colback=white,colframe=apifg,colbacktitle=apibg,coltitle=ink,boxrule=0.45pt,arc=1.6mm,left=3mm,right=3mm,top=2mm,bottom=2mm,title={\sffamily\bfseries Cover API \hfill \texttt{Icover}}]
{\sffamily\small\color{muted}Balanced coarse/fine configurations | WZ2 | WZ2 Definitions 2.10-2.12 and Section 6.1, PDF pp. 8-9, 37-38}\par
\textbf{Mathematical statement.} For geometric covers, $\mathcal B[A]=\{B\in\mathcal B:B\subset A\}$.
A $K$-uniform cover has child counts comparable within $K$. A partitioning cover has disjoint families $\mathcal B[2A]$ for distinct parents.
For the shaded structural argument, use the compatible line-metric cover and require child shadings to lie in the parent shading.
A balanced cover has $|E_C\cap Q|$ independent of the occupied coarse cube $Q$; constant-multiplicity data transfers mass through such a cover.
\smallskip\textbf{Lean statement sketch.}
\begin{Verbatim}[fontsize=\small,breaklines=true,breakanywhere=true,tabsize=2]
def Covers (D : Config n rho) (C : Config n delta) : Prop := ...
def Balanced (D : Config n rho) (C : Config n delta) : Prop :=
  Covers D C ∧ ∃ a, ∀ Q, CoarseOccupied D Q →
    vol (union C ∩ Q) = a
theorem balanced_pullback : Balanced D C →
  ConstantMultiplicity C →
  CoarseSelection D D' → MassTransferBound C D'
\end{Verbatim}
{\footnotesize\textbf{Mathlib support.} F, G, M.\par
\textbf{Incoming edges.} \nodecard{Igeom} (api); \nodecard{Imass} (api)\par}
{\footnotesize\textbf{Interface note.} Geometric partitioning, uniform child counts and balanced shaded volumes are three different requirements. WZ2 Section 6.1 reconciles the stronger WZ1 line-metric cover with geometric containment; do not identify them by definition.}
\end{tcolorbox}\end{minipage}\par\medskip
\begin{minipage}{\linewidth}
\hypertarget{node-Iaffine}{}
\begin{tcolorbox}[colback=white,colframe=apifg,colbacktitle=apibg,coltitle=ink,boxrule=0.45pt,arc=1.6mm,left=3mm,right=3mm,top=2mm,bottom=2mm,title={\sffamily\bfseries Affine transport API \hfill \texttt{Iaffine}}]
{\sffamily\small\color{muted}Jacobian, direction and grid transport | API | WZ2 Section 6.1 and the transferred WZ1 Lemma 3.5; WZ1 pp. 12-14, 18}\par
\textbf{Mathematical statement.} For an invertible affine map $\phi(x)=Ax+b$, volume transforms by $|\phi(S)|=|\det A|\,|S|$.
Lines and their directions are transported by $\phi$ and $v\mapsto Av/|Av|$. Projection normals are transported using the inverse transpose.
After transport, re-cubulate at the chosen scale; prove bounded overlap and containment before replacing geometric sets by grid unions.
\smallskip\textbf{Lean statement sketch.}
\begin{Verbatim}[fontsize=\small,breaklines=true,breakanywhere=true,tabsize=2]
def transportedDirection (A : Rn n ≃L[ℝ] Rn n) (v : Rn n) :=
  (‖A v‖)⁻¹ • A v
theorem affine_volume (hS : MeasurableSet S) :
  vol (phi '' S) = |detLinearPart phi| * vol S
theorem recubulate_transport : ControlledAffine phi →
  GridUnion delta S → ControlledGridCover phi S
\end{Verbatim}
{\footnotesize\textbf{Mathlib support.} A, G, M.\par
\textbf{Incoming edges.} \nodecard{Igeom} (api)\par}
{\footnotesize\textbf{Interface note.} Convex-containment counts transport by the determinant and parent-volume normalization. Re-cubulation and balancing still require witnesses; no arbitrary-subcollection preservation of the multi-scale axiom is asserted.}
\end{tcolorbox}\end{minipage}\par\medskip
\begin{minipage}{\linewidth}
\hypertarget{node-Icalc}{}
\begin{tcolorbox}[colback=white,colframe=apifg,colbacktitle=apibg,coltitle=ink,boxrule=0.45pt,arc=1.6mm,left=3mm,right=3mm,top=2mm,bottom=2mm,title={\sffamily\bfseries Calculus API \hfill \texttt{Icalc}}]
{\sffamily\small\color{muted}C2 bounds and compactly supported patches | API | WZ1: Sections 5.2 and 7, pp. 41-50}\par
\textbf{Mathematical statement.} Represent a function on a compact interval by an ambient function with a $C^2$ extension to a neighbourhood. Use the jet norm $\|f\|_{C^2(I)}=\max_{0\le j\le2}\sup_{t\in I}|f^{(j)}(t)|$.
Finite sums of patches with disjoint support at each scale preserve derivative bounds. Endpoint conventions must be fixed for derivatives and gluing.
\smallskip\textbf{Lean statement sketch.}
\begin{Verbatim}[fontsize=\small,breaklines=true,breakanywhere=true,tabsize=2]
def C2Near (f : ℝ → ℝ) (I : Set ℝ) : Prop :=
  ∃ U, IsOpen U ∧ I ⊆ U ∧ ContDiffOn ℝ 2 f U
noncomputable def c2Norm (f : ℝ → ℝ) (I : Set ℝ) : ℝ := ...
def C2Bounded f I B := C2Near f I ∧ c2Norm f I ≤ B
\end{Verbatim}
{\footnotesize\textbf{Mathlib support.} C, F, M.\par
\textbf{Incoming edges.} None; foundational or external interface.\par}
{\footnotesize\textbf{Interface note.} A custom normed C2 function space or a jet-based metric is needed for Definition 7.1. The signatures below use this proposed API.}
\end{tcolorbox}\end{minipage}\par\medskip
\begin{minipage}{\linewidth}
\hypertarget{node-Iparams}{}
\begin{tcolorbox}[colback=white,colframe=apifg,colbacktitle=apibg,coltitle=ink,boxrule=0.45pt,arc=1.6mm,left=3mm,right=3mm,top=2mm,bottom=2mm,title={\sffamily\bfseries Quantifier API \hfill \texttt{Iparams}}]
{\sffamily\small\color{muted}Small-parameter and finite-iteration contracts | API | WZ2 Proposition 6.2, Lemma 7.12 and proof of Proposition 7.1; transferred WZ1 iterations}\par
\textbf{Mathematical statement.} Choose $\varepsilon$ first, then all auxiliary loss exponents, and finally $\delta_0$.
At a middle scale $s$ with $\delta^a\le s\le\delta^b$, convert an error $\delta^{-\eta}$ to an $s$-error using the appropriate endpoint exponent.
For a fixed finite iteration, reserve enough loss at every call, including density, cover rebuilding and logarithmic refinements.
In Section 7, write $q$ for the parameter cube size, $h=\delta/q$ for its rescaled mesh, and $r=q$ for the final coarse spatial scale.
\smallskip\textbf{Lean statement sketch.}
\begin{Verbatim}[fontsize=\small,breaklines=true,breakanywhere=true,tabsize=2]
def SmallScale (P : ℝ → ℝ → Prop) : Prop :=
  ∀ eps > 0, ∃ eta > 0, ∃ delta0 > 0,
    ∀ delta, 0 < delta → delta ≤ delta0 → P eta delta
theorem absorb_log_loss (a eps : ℝ) (he : 0 < eps) :
  ∃ delta0 > 0, ∀ delta ∈ Set.Ioc 0 delta0,
    Real.log (1/delta) ^ a ≤ delta ^ (-eps)
\end{Verbatim}
{\footnotesize\textbf{Mathlib support.} F, Q.\par
\textbf{Incoming edges.} None; foundational or external interface.\par}
{\footnotesize\textbf{Interface note.} The source repeatedly reuses rho for a cube size and a rescaled mesh. This atlas keeps them separate and uses independent cinematic and geometric loss exponents.}
\end{tcolorbox}\end{minipage}\par\medskip
\begin{minipage}{\linewidth}
\hypertarget{node-Iwolffbalance}{}
\begin{tcolorbox}[colback=white,colframe=apifg,colbacktitle=apibg,coltitle=ink,boxrule=0.45pt,arc=1.6mm,left=3mm,right=3mm,top=2mm,bottom=2mm,title={\sffamily\bfseries WZ2 proof operation \hfill \texttt{Iwolffbalance}}]
{\sffamily\small\color{muted}Restriction and rebuilding of cover witnesses | WZ2 | WZ2 proof of Proposition 6.2, PDF pp. 39-40}\par
\textbf{Mathematical statement.} A $q$-fraction subfamily inherits a single-scale convex count with error at most $A/q$.
For the multi-scale condition, separately rebalance the partitioning tree and each surviving child fibre.
Record uniform parent counts, surviving child counts and the new errors as witnesses.
This proof operation is used for the specific refinements constructed in the argument; it is not an arbitrary-subfamily inheritance theorem.
\smallskip\textbf{Lean statement sketch.}
\begin{Verbatim}[fontsize=\small,breaklines=true,breakanywhere=true,tabsize=2]
structure MWRefinementWitness (C' C : Config 3 delta) (A' : ℝ) where
  subfamily : Subcollection C' C
  mass_retention : MassRetentionWitness C' C
  rebuilt_covers : MultiscaleConvexWolff A' C'
-- Every use supplies its constructed witnesses and its loss budget.
\end{Verbatim}
{\footnotesize\textbf{Mathlib support.} F, G, M, Q.\par
\textbf{Incoming edges.} \nodecard{W2L2_13} (api); \nodecard{W2CWA} (api)\par}
{\footnotesize\textbf{Interface note.} This card makes the inheritance obligation visible. The WZ2 sketch describes repeated pigeonholing to restore the axiom; a full formalization must implement that construction, not assume closure under all refinements.}
\end{tcolorbox}\end{minipage}\par\medskip
\begin{minipage}{\linewidth}
\hypertarget{node-Ienvelope}{}
\begin{tcolorbox}[colback=white,colframe=apifg,colbacktitle=apibg,coltitle=ink,boxrule=0.45pt,arc=1.6mm,left=3mm,right=3mm,top=2mm,bottom=2mm,title={\sffamily\bfseries WZ2 Section 6.2: consequence \hfill \texttt{Ienvelope}}]
{\sffamily\small\color{muted}Uniform lower envelope for admissible unions | WZ2 | WZ2 deficit definition and proof of Proposition 6.2}\par
\textbf{Mathematical statement.} For every $\gamma>0$ there are $\eta,\delta_0>0$ such that every admissible pair at $0<\delta\le\delta_0$ satisfies $|E_C|\ge\delta^{\sigma+\gamma}$.
Otherwise $\sigma+\gamma$ would be a bad exponent, contradicting the supremum.
At constant multiplicity this gives $\mu\lesssim\delta^{2-\sigma-\gamma}\#\mathcal T$.
This normalized estimate replaces the old absolute multiplicity bounds.
\smallskip\textbf{Lean statement sketch.}
\begin{Verbatim}[fontsize=\small,breaklines=true,breakanywhere=true,tabsize=2]
theorem CW_union_lower_envelope (hg : 0 < gamma) :
  ∃ eta > 0, ∃ delta0 > 0, ∀ delta ∈ Set.Ioc 0 delta0,
  ∀ C : Config 3 delta, MWAdmissible eta C →
    delta^(sigmaCW+gamma) ≤ vol (union C)
\end{Verbatim}
{\footnotesize\textbf{Mathlib support.} F, G, M, Q.\par
\textbf{Incoming edges.} \nodecard{W2Sigma} (core)\par}

\end{tcolorbox}\end{minipage}\par\medskip
\sectitle{Dependency edge ledger}
Every row is a displayed or catalogue dependency. Assembly edges transfer an output to a later hypothesis; they need not prove the generic lemma itself. Overview-only compressed paths are excluded from this ledger.
\begin{longtable}{@{}p{.20\linewidth}p{.20\linewidth}p{.09\linewidth}p{.43\linewidth}@{}}\toprule Prerequisite & Consumer & Kind & Proof / usage site\\\midrule\endhead
\nodecard{Igeom} & \nodecard{Ientropy} & api & WZ1: Bounded Euclidean grid sets.\\
\nodecard{Igeom} & \nodecard{Imass} & api & WZ1: Section 2.1, pp. 7-8.\\
\nodecard{Igeom} & \nodecard{Icover} & api & WZ1: Section 2.1, p. 8.\\
\nodecard{Igeom} & \nodecard{Iaffine} & api & WZ1: (3.3)-(3.4), pp. 13-14.\\
\nodecard{Imass} & \nodecard{Icover} & api & WZ1: Balanced mass pullback needs multiplicity control.\\
\nodecard{Ientropy} & \nodecard{W2Sigma} & api & WZ1: Finite discretization and multiscale cover. CW interface transferred through WZ2 Sections 6.2-6.5.\\
\nodecard{Imass} & \nodecard{W2Sigma} & api & WZ1: Definition 2.1(c), p. 8. CW interface transferred through WZ2 Sections 6.2-6.5.\\
\nodecard{Icover} & \nodecard{W2Sigma} & api & WZ1: Definition 2.1(b), p. 8. CW interface transferred through WZ2 Sections 6.2-6.5.\\
\nodecard{W2Sigma} & \nodecard{W2D6_1} & api & WZ1: Definition 3.1, p. 11. CW interface transferred through WZ2 Sections 6.2-6.5.\\
\nodecard{Iaffine} & \nodecard{W2Unit} & api & WZ1: Definition 3.2 and (3.3)-(3.4). CW interface transferred through WZ2 Sections 6.2-6.5.\\
\nodecard{Icover} & \nodecard{W2Unit} & api & WZ1: Children must be covered by the chosen parent. CW interface transferred through WZ2 Sections 6.2-6.5.\\
\nodecard{Igeom} & \nodecard{Dplane} & api & WZ1: Definition 4.1, p. 18.\\
\nodecard{Ientropy} & \nodecard{DAD} & api & WZ1: Definition 4.2, p. 18.\\
\nodecard{Igeom} & \nodecard{DFrost} & api & WZ1: Definition 8.1, p. 51.\\
\nodecard{Imass} & \nodecard{Dhyp} & api & WZ1: Section 8.0.1 identifies points with occupied cells.\\
\nodecard{Icalc} & \nodecard{Dcinematic} & api & WZ1: Definition 7.1, p. 49.\\
\nodecard{Igeom} & \nodecard{Dtwist} & api & WZ1: Definition 6.1, p. 48.\\
\nodecard{Icalc} & \nodecard{Dtwist} & api & WZ1: Twisted lines are C2 curves.\\
\nodecard{Ientropy} & \nodecard{Dtrap} & api & WZ1: Trapezoid lengths, heights and separation.\\
\nodecard{Igeom} & \nodecard{Dthin} & api & WZ1: Definition 8.5, p. 55.\\
\nodecard{Dthin} & \nodecard{DthinD} & api & WZ1: Definition 8.5' is the grid counterpart.\\
\nodecard{Igeom} & \nodecard{Dsep} & api & WZ1: Definition 8.6, p. 56.\\
\nodecard{DAD} & \nodecard{Iadgap} & api & WZ1: Direct use of the AD upper covering inequality.\\
\nodecard{W2Sigma} & \nodecard{W2Exists} & core & Near-supremal bad exponent yields existential scales.\\
\nodecard{W2D6_1} & \nodecard{W2Exists} & core & Target extremality contract.\\
\nodecard{W2Sigma} & \nodecard{L3_3} & core & WZ1: (3.6)-(3.8), pp. 13-14. CW interface transferred through WZ2 Sections 6.2-6.5.\\
\nodecard{W2Unit} & \nodecard{L3_3} & core & WZ1: Unit-rescaling geometry and admissibility, pp. 13-14. CW interface transferred through WZ2 Sections 6.2-6.5.\\
\nodecard{L3_3} & \nodecard{W2P6_2} & core & WZ1: Step 2 invokes Lemma 3.3, p. 15. CW interface transferred through WZ2 Sections 6.2-6.5.\\
\nodecard{W2Sigma} & \nodecard{W2P6_2} & core & WZ2 Section 6.2: normalized near-extremal union comparison.\\
\nodecard{W2D6_1} & \nodecard{W2P6_2} & core & WZ1: Extremality premise and targets, pp. 12-17. CW interface transferred through WZ2 Sections 6.2-6.5.\\
\nodecard{W2P6_2} & \nodecard{L3_4} & core & WZ1: Proof on p. 18, rho=delta\textasciicircum{}epsilon. CW interface transferred through WZ2 Sections 6.2-6.5.\\
\nodecard{W2D6_1} & \nodecard{L3_5} & core & WZ1: Extremality under restriction and transport, p. 18. CW interface transferred through WZ2 Sections 6.2-6.5.\\
\nodecard{T4_2} & \nodecard{L4_3} & core & WZ1: Multilinear bound (4.5), p. 19.\\
\nodecard{L3_4} & \nodecard{L4_3} & core & WZ1: Initial robust-transversality refinement, p. 19.\\
\nodecard{W2P6_2} & \nodecard{L4_4} & core & WZ1: Applied to obtain coarse/fine multiplicity, pp. 20-21. CW interface transferred through WZ2 Sections 6.2-6.5.\\
\nodecard{L3_4} & \nodecard{L4_5} & core & WZ1: Applied at the start of the proof, p. 21.\\
\nodecard{W2P6_2} & \nodecard{L4_5} & core & WZ1: Balanced cover and multiplicity bounds, pp. 21-22. CW interface transferred through WZ2 Sections 6.2-6.5.\\
\nodecard{L4_5} & \nodecard{C4_6} & core & WZ1: Cube separation after Lemma 4.5, p. 22.\\
\nodecard{C4_6} & \nodecard{L4_7} & core & WZ1: Iterated over rho\_j=delta\textasciicircum{}(j/N), p. 23.\\
\nodecard{W2Exists} & \nodecard{L4_8} & core & WZ1: Existence assembly explicitly cited, p. 23. CW interface transferred through WZ2 Sections 6.2-6.5.\\
\nodecard{L4_3} & \nodecard{L4_8} & core & WZ1: Sections 4.1-4.2 combined, p. 23.\\
\nodecard{L4_4} & \nodecard{L4_8} & core & WZ1: Coarsening weak planiness in that assembly.\\
\nodecard{L4_7} & \nodecard{L4_8} & core & WZ1: Lipschitz refinement in that assembly.\\
\nodecard{W2P6_2} & \nodecard{L4_9} & core & WZ1: Two cover scales and multiplicities, p. 24. CW interface transferred through WZ2 Sections 6.2-6.5.\\
\nodecard{L3_4} & \nodecard{L4_9} & core & WZ1: Robust transversality at coarse scale, p. 24.\\
\nodecard{DAD} & \nodecard{L4_9} & core & WZ1: Target grain entropy bound (4.15).\\
\nodecard{L4_9} & \nodecard{L4_10} & core & WZ1: Applied twice, p. 26.\\
\nodecard{W2P6_2} & \nodecard{L4_10} & core & WZ1: Balanced cover and additional coarsening, pp. 26-27. CW interface transferred through WZ2 Sections 6.2-6.5.\\
\nodecard{L4_10} & \nodecard{L4_11} & core & WZ1: Finite iteration in tau, pp. 28-29.\\
\nodecard{L4_11} & \nodecard{L4_12} & core & WZ1: Finite iteration in rho, pp. 29-30.\\
\nodecard{L4_8} & \nodecard{W2P6_3} & core & WZ1: Step 1, p. 30. CW interface transferred through WZ2 Sections 6.2-6.5.\\
\nodecard{L4_12} & \nodecard{W2P6_3} & core & WZ1: Used in Step 1 and at completion, pp. 30,32. CW interface transferred through WZ2 Sections 6.2-6.5.\\
\nodecard{W2P6_2} & \nodecard{W2P6_3} & core & WZ1: Coarse cover and extremal parent rescaling, p. 30. CW interface transferred through WZ2 Sections 6.2-6.5.\\
\nodecard{W2Unit} & \nodecard{W2P6_3} & core & WZ1: Pass from local to global grains, pp. 30-31. CW interface transferred through WZ2 Sections 6.2-6.5.\\
\nodecard{L4_3} & \nodecard{W2P6_3} & core & WZ1: Rebuild a weak local plane map, p. 31. CW interface transferred through WZ2 Sections 6.2-6.5.\\
\nodecard{L4_4} & \nodecard{W2P6_3} & core & WZ1: Convert it to a coarse plane map, p. 31. CW interface transferred through WZ2 Sections 6.2-6.5.\\
\nodecard{L3_5} & \nodecard{W2P6_3} & core & WZ1: Final normalization of Lipschitz constants, p. 32. CW interface transferred through WZ2 Sections 6.2-6.5.\\
\nodecard{XExt} & \nodecard{W2P6_3} & core & WZ1: Scalar slope extension, p. 31. CW interface transferred through WZ2 Sections 6.2-6.5.\\
\nodecard{T5_2} & \nodecard{L5_3} & core & WZ1: Apply the dot-product dichotomy, p. 37.\\
\nodecard{L5_3} & \nodecard{L5_4} & core & WZ1: Apply locally in each popular slab, p. 39.\\
\nodecard{W2P6_2} & \nodecard{L5_4} & core & WZ1: Balanced covers at rho and sqrt(rho), p. 38. CW interface transferred through WZ2 Sections 6.2-6.5.\\
\nodecard{L5_4} & \nodecard{L5_5} & core & WZ1: Iterate over j=1,...,N, pp. 39-40.\\
\nodecard{L5_5} & \nodecard{C5_6} & core & WZ1: Restate the refined conclusion via trapezoids, pp. 40-41.\\
\nodecard{Dtrap} & \nodecard{C5_6} & core & WZ1: Trapezoid representation, p. 40.\\
\nodecard{L5_8} & \nodecard{P5_7} & core & WZ1: Patch construction and finite sum, pp. 42-43.\\
\nodecard{Dtrap} & \nodecard{P5_7} & core & WZ1: Nested input trapezoids, p. 41.\\
\nodecard{W2P6_3} & \nodecard{W2P6_4} & core & WZ1: Start with Proposition 4.1 data, p. 43. CW interface transferred through WZ2 Sections 6.2-6.5.\\
\nodecard{C5_6} & \nodecard{W2P6_4} & core & WZ1: Apply the trapezoid-tree corollary, p. 43. CW interface transferred through WZ2 Sections 6.2-6.5.\\
\nodecard{P5_7} & \nodecard{W2P6_4} & core & WZ1: C2 approximant, p. 43. CW interface transferred through WZ2 Sections 6.2-6.5.\\
\nodecard{L3_5} & \nodecard{W2P6_4} & core & WZ1: Normalize the C2 bound by vertical rescaling, p. 43. CW interface transferred through WZ2 Sections 6.2-6.5.\\
\nodecard{DAD} & \nodecard{L6_2} & core & WZ1: Compare projection count with the AD upper bound, p. 44.\\
\nodecard{L6_2} & \nodecard{W2L6_6} & core & WZ1: Final direction-counting contradiction, p. 47. CW interface transferred through WZ2 Sections 6.2-6.5.\\
\nodecard{W2P6_4} & \nodecard{L6_4} & core & WZ1: Initial C2 grain example, p. 47. CW interface transferred through WZ2 Sections 6.2-6.5.\\
\nodecard{W2P6_2} & \nodecard{L6_4} & core & WZ1: Cover at scale delta\textasciicircum{}epsilon\_2, p. 47. CW interface transferred through WZ2 Sections 6.2-6.5.\\
\nodecard{W2L6_6} & \nodecard{L6_4} & core & Replace WZ1 Lemma 6.3 and omit the direction-selection step.\\
\nodecard{L6_4} & \nodecard{W2P6_5} & core & WZ2 proof of Proposition 6.5: inherited comparable-slope slab normalization.\\
\nodecard{L3_5} & \nodecard{W2P6_5} & core & WZ1: Preserve extremality after anisotropic change, p. 48. CW interface transferred through WZ2 Sections 6.2-6.5.\\
\nodecard{W2P6_5} & \nodecard{W2L6_7} & core & WZ1: Small projected union plus Holder, pp. 48-49. CW interface transferred through WZ2 Sections 6.2-6.5.\\
\nodecard{Dtwist} & \nodecard{W2L6_7} & core & WZ1: Projection definition and projected shading mass. CW interface transferred through WZ2 Sections 6.2-6.5.\\
\nodecard{Dcinematic} & \nodecard{T7_2} & core & WZ1: Hypotheses of the imported theorem, p. 49.\\
\nodecard{XPYZ} & \nodecard{T7_2} & core & WZ1: Special case of [36, Theorem 1.7], p. 49.\\
\nodecard{Dcinematic} & \nodecard{L7_3} & core & WZ1: Check the cinematic-family definition, pp. 49-50.\\
\nodecard{W2L6_7} & \nodecard{W2Zero} & core & WZ2 Section 6.5: small projected union, in the audited loss-budget form.\\
\nodecard{Dhyp} & \nodecard{L8_1} & core & WZ1: Prune all coordinate fibres, p. 53.\\
\nodecard{L8_1} & \nodecard{T8_2} & core & WZ1: Uniform dense refinement at start of proof, p. 54.\\
\nodecard{DFrost} & \nodecard{T8_2} & core & WZ1: Energy estimates for normalized measures, pp. 54-55.\\
\nodecard{Dthin} & \nodecard{L8_3} & core & WZ1: Measure good-pair set, p. 55.\\
\nodecard{DthinD} & \nodecard{L8_3} & core & WZ1: Target grid good-pair set, pp. 55-56.\\
\nodecard{T8_2} & \nodecard{L8_4} & core & WZ1: Kaufman projection estimate applied, p. 57.\\
\nodecard{DthinD} & \nodecard{L8_4} & core & WZ1: Construct many radial directions, p. 56.\\
\nodecard{Dhyp} & \nodecard{L8_4} & core & WZ1: 2c-uniform density of sections, pp. 56-57.\\
\nodecard{DFrost} & \nodecard{L8_4} & core & WZ1: Normalized nonconcentration premise, p. 56.\\
\nodecard{Dsep} & \nodecard{L8_4} & core & WZ1: Standard separation gives |b1-b2|>=1/2.\\
\nodecard{XOSW} & \nodecard{P8_6} & core & WZ1: Imported specialized bootstrap, pp. 57-58.\\
\nodecard{Dthin} & \nodecard{P8_6} & core & WZ1: Both ordered measure pairs must be thin.\\
\nodecard{P8_6} & \nodecard{P8_6p} & core & WZ1: Normalize grid volumes to probability measures, p. 58.\\
\nodecard{L8_3} & \nodecard{P8_6p} & core & WZ1: Re-discretize after the measure bootstrap, p. 58.\\
\nodecard{P8_6p} & \nodecard{C8_7} & core & WZ1: Repeated finite bootstrap, p. 58.\\
\nodecard{DFrost} & \nodecard{L8_8} & core & WZ1: Initial 1/4-thin-tube seed proof, pp. 59-60.\\
\nodecard{L8_8} & \nodecard{P8_5} & core & WZ1: Produce the fixed-exponent seeds, p. 61.\\
\nodecard{C8_7} & \nodecard{P8_5} & core & WZ1: Improve them to exponent 1, p. 61.\\
\nodecard{L8_4} & \nodecard{P8_5} & core & WZ1: Conclude large dot-product image, p. 61.\\
\nodecard{DFrost} & \nodecard{L8_11} & core & WZ1: Input normalized ball bound, p. 62.\\
\nodecard{DFrost} & \nodecard{L8_12} & core & WZ1: Input normalized ball bound, pp. 62-63.\\
\nodecard{L8_1} & \nodecard{L8_13} & core & WZ1: Refine the bipartite graph X, p. 65.\\
\nodecard{L8_11} & \nodecard{L8_13} & core & WZ1: Restore rescaled F, p. 65.\\
\nodecard{L8_12} & \nodecard{L8_13} & core & WZ1: Coarsen the rescaled F, p. 65.\\
\nodecard{T8_2} & \nodecard{L8_13} & core & WZ1: Apply Kaufman to F and radial directions, p. 65.\\
\nodecard{Dsep} & \nodecard{L8_13} & core & WZ1: Separated uniform setup, p. 63.\\
\nodecard{L8_10} & \nodecard{P8_9} & core & WZ1: Two-ends localization for each G set, p. 66.\\
\nodecard{L8_13} & \nodecard{P8_9} & core & WZ1: Compare widths or obtain expansion, p. 66.\\
\nodecard{L8_11} & \nodecard{P8_9} & core & WZ1: Restore all three rescaled Frostman sets, p. 67.\\
\nodecard{L8_1} & \nodecard{P8_9} & core & WZ1: Repeatedly restore uniform hypergraph density, pp. 66-67.\\
\nodecard{P8_5} & \nodecard{P8_9} & core & WZ1: Expand after the strip reduction, p. 67.\\
\nodecard{P8_9} & \nodecard{T5_2p} & core & WZ1: Separated version after localization, pp. 68-69.\\
\nodecard{L8_1} & \nodecard{T5_2p} & core & WZ1: Uniformize the localized H, p. 69.\\
\nodecard{T5_2p} & \nodecard{T5_2} & core & WZ1: Frostman generalization implies the AD version, p. 52.\\
\nodecard{DFrost} & \nodecard{T5_2} & core & WZ1: AD-to-Frostman conversion, with the size dichotomy, p. 52.\\
\nodecard{DAD} & \nodecard{T5_2} & core & WZ1: AD hypotheses and upper counts in the size dichotomy.\\
\nodecard{XBCT} & \nodecard{T4_2} & core & WZ1: The external version identified on p. 19.\\
\nodecard{Imass} & \nodecard{L3_3} & api & WZ1: Constant multiplicity, p. 13.\\
\nodecard{Icover} & \nodecard{W2P6_2} & api & WZ1: Balanced covers and fine/coarse pullback, pp. 15-17. CW interface transferred through WZ2 Sections 6.2-6.5.\\
\nodecard{Imass} & \nodecard{W2P6_2} & api & WZ1: Pigeonholing, fine/coarse mass and multiplicity. CW interface transferred through WZ2 Sections 6.2-6.5.\\
\nodecard{Iaffine} & \nodecard{L3_5} & api & WZ1: Explicit transport obligations for the omitted proof.\\
\nodecard{Imass} & \nodecard{L4_3} & api & WZ1: Constant multiplicity refinement, p. 19.\\
\nodecard{Dplane} & \nodecard{L4_4} & api & WZ1: Exact plane-map conclusion.\\
\nodecard{Dplane} & \nodecard{L4_5} & api & WZ1: Plane-map premise.\\
\nodecard{Iparams} & \nodecard{L4_7} & api & WZ1: Finite scale iteration, p. 23.\\
\nodecard{Imass} & \nodecard{L4_9} & api & WZ1: Multiplicity and popular grain refinements, p. 24.\\
\nodecard{Iparams} & \nodecard{L4_11} & api & WZ1: Finite hierarchy in tau, pp. 28-29.\\
\nodecard{Iparams} & \nodecard{L4_12} & api & WZ1: Finite hierarchy in rho, pp. 29-30.\\
\nodecard{Iadgap} & \nodecard{L5_3} & api & WZ1: Expansion contradicts sub-one-dimensional AD spacing, p. 37.\\
\nodecard{Icover} & \nodecard{L5_4} & api & WZ1: Balanced pullback from coarse selected slabs, pp. 38-39.\\
\nodecard{Imass} & \nodecard{L5_5} & api & WZ1: Popularity and separation refinement, p. 40.\\
\nodecard{Iparams} & \nodecard{L5_5} & api & WZ1: Finite iteration with changing exponent budget.\\
\nodecard{Icalc} & \nodecard{L5_8} & api & WZ1: Controlled smooth cutoff and C2 gluing, pp. 41-42.\\
\nodecard{Icalc} & \nodecard{P5_7} & api & WZ1: Finite patch sum with scale-wise disjoint supports.\\
\nodecard{Icalc} & \nodecard{W2L6_6} & api & WZ1: C2 control spreads a small derivative over the slab. CW interface transferred through WZ2 Sections 6.2-6.5.\\
\nodecard{Iaffine} & \nodecard{W2P6_5} & api & WZ1: Rescale both slope and geometry, p. 48. CW interface transferred through WZ2 Sections 6.2-6.5.\\
\nodecard{Icalc} & \nodecard{L7_3} & api & WZ1: Jet estimates for a+b f+d t f.\\
\nodecard{W2Sigma} & \nodecard{W2Zero} & api & WZ1: Nonnegativity and meaning of sigma\_3. CW interface transferred through WZ2 Sections 6.2-6.5.\\
\nodecard{Iparams} & \nodecard{C8_7} & api & WZ1: Finite iteration; control 3\textasciicircum{}N c and K'.\\
\nodecard{Iparams} & \nodecard{P8_5} & api & WZ1: Lambda(epsilon) independent of zeta, p. 61.\\
\nodecard{Ientropy} & \nodecard{L8_10} & api & WZ1: Strip covering and scale-localized pruning.\\
\nodecard{Iaffine} & \nodecard{L8_11} & api & WZ1: Invert affine maps and re-cubulate, p. 62.\\
\nodecard{Imass} & \nodecard{L8_11} & api & WZ1: Comparable rectangle occupancy by dyadic pigeonholing.\\
\nodecard{Imass} & \nodecard{L8_12} & api & WZ1: Comparable coarse-cell occupancy.\\
\nodecard{Iaffine} & \nodecard{L8_13} & api & WZ1: Dual affine transport preserving dot products, p. 65.\\
\nodecard{Iaffine} & \nodecard{P8_9} & api & WZ1: Anisotropic changes and entropy scales, pp. 66-67.\\
\nodecard{Iaffine} & \nodecard{T5_2p} & api & WZ1: Common translation and separate scalar dilations, pp. 68-69.\\
\nodecard{L4_3} & \nodecard{L4_4} & assembly & WZ1: Weak planiness supplies the hypothesis in the existence assembly.\\
\nodecard{L4_4} & \nodecard{L4_7} & assembly & WZ1: The coarsened plane map is cube-constant.\\
\nodecard{L4_8} & \nodecard{L4_9} & assembly & WZ1: Existence node supplies the generic grain lemma's input.\\
\nodecard{L4_7} & \nodecard{W2P6_3} & assembly & WZ1: Implementation must explicitly supply a Lipschitz W before reapplying Lemma 4.12 in Step 2, pp. 31-32. CW interface transferred through WZ2 Sections 6.2-6.5.\\
\nodecard{W2P6_3} & \nodecard{L5_3} & assembly & WZ1: Grain hypotheses are instantiated after localization. CW interface transferred through WZ2 Sections 6.2-6.5.\\
\nodecard{W2P6_4} & \nodecard{W2L6_6} & assembly & WZ1: The generic slab lemma consumes the C2 and two-grain conclusions. CW interface transferred through WZ2 Sections 6.2-6.5.\\
\nodecard{Igeom} & \nodecard{W2CWA} & api & Geometry and finite tube/multiset containment.\\
\nodecard{Igeom} & \nodecard{W2MW} & api & Essential distinction, coarse tube scales and line charts.\\
\nodecard{Icover} & \nodecard{W2MW} & core & WZ2 Definition 2.12: uniform partitioning covers.\\
\nodecard{W2Unit} & \nodecard{W2MW} & core & WZ2 Definition 2.12(ii): rescaled children.\\
\nodecard{W2CWA} & \nodecard{W2MW} & core & WZ2 Definition 2.12(ii): single-scale convex count.\\
\nodecard{W2MW} & \nodecard{W2L2_13} & core & WZ2 Lemma 2.13: nested partitioning tree.\\
\nodecard{Icover} & \nodecard{W2L2_13} & api & Partitioning forces compatible cluster nesting.\\
\nodecard{W2MW} & \nodecard{W2Sigma} & core & WZ2 Section 6.2 admissible class.\\
\nodecard{Imass} & \nodecard{W2Sigma} & core & Average shading density, not absolute total mass.\\
\nodecard{W2Sigma} & \nodecard{Ienvelope} & core & A stronger counterexample exponent contradicts the supremum.\\
\nodecard{W2MW} & \nodecard{W2D6_1} & core & WZ2 Definition 6.1(1).\\
\nodecard{Imass} & \nodecard{W2D6_1} & core & WZ2 Definition 6.1(2).\\
\nodecard{W2Sigma} & \nodecard{W2D6_1} & core & WZ2 Definition 6.1(3).\\
\nodecard{W2L2_13} & \nodecard{W2P6_2} & core & WZ2 proof p. 39: coarse/fine MW witnesses.\\
\nodecard{Ienvelope} & \nodecard{L3_3} & core & Normalized child multiplicity from the admissible union lower envelope.\\
\nodecard{Ienvelope} & \nodecard{W2P6_2} & core & Lower union bound bounds fine and coarse multiplicities.\\
\nodecard{Iwolffbalance} & \nodecard{W2P6_2} & api & Density selection and cover rebuilding in the WZ2 sketch.\\
\nodecard{W2L2_13} & \nodecard{Iwolffbalance} & api & Tree representation of the retained cluster hierarchy.\\
\nodecard{W2CWA} & \nodecard{Iwolffbalance} & api & Single-scale restriction has an explicit relative-cardinality loss.\\
\nodecard{Iwolffbalance} & \nodecard{L3_5} & api & Transported cover and fibre witnesses require rebuilding.\\
\nodecard{Ienvelope} & \nodecard{L4_3} & core & Near-extremal lower multiplicity in normalized count form.\\
\nodecard{W2P6_4} & \nodecard{W1SlabPrep} & assembly & Input global/local grains to WZ1 Lemma 6.3 Steps 1-3.\\
\nodecard{DAD} & \nodecard{W1SlabPrep} & core & Global/local upper AD bounds.\\
\nodecard{Imass} & \nodecard{W1SlabPrep} & api & Popular fibres, tubes and grains are selected with a loss budget.\\
\nodecard{W1SlabPrep} & \nodecard{W2SlabCount} & core & The popular-square count and nondegenerate projected normal.\\
\nodecard{L6_2} & \nodecard{W2SlabCount} & core & WZ2 equation (6.10): shaded segments force near orthogonality.\\
\nodecard{W2CWA} & \nodecard{W2SlabCount} & core & Count full tubes inside the thin convex slab.\\
\nodecard{W2SlabCount} & \nodecard{W2L6_6} & core & Contradiction when the slope is flat at one point.\\
\nodecard{W1SlabPrep} & \nodecard{W2L6_6} & core & Steps 1-3 of the inherited proof.\\
\nodecard{W2P6_4} & \nodecard{W2P6_5} & assembly & C2/global/local input retained by the normalized slope construction.\\
\nodecard{W2CWA} & \nodecard{W2L6_7} & api & Single-scale axiom inherited from the multi-scale configuration.\\
\nodecard{W2MW} & \nodecard{W2L6_7} & api & A top-scale cover and its normalized child convex counts give the needed global single-scale axiom; constants absorbed.\\
\nodecard{DAD} & \nodecard{W2L6_7} & core & One-dimensional AD slices bound projected length before Fubini.\\
\nodecard{W2CWA} & \nodecard{W2TW} & core & A radius-r convex tube has volume O(r\textasciicircum{}2).\\
\nodecard{Ientropy} & \nodecard{W2DFrost} & api & Normalized mesh covering count.\\
\nodecard{Ientropy} & \nodecard{W2DKT} & api & Absolute mesh covering count.\\
\nodecard{Ientropy} & \nodecard{W2DUnif} & api & Occupied cells in a sparse dyadic hierarchy.\\
\nodecard{W2TW} & \nodecard{W2DFrost} & core & WZ2 Remark 7.4: R4 line-parameter Frostman bound.\\
\nodecard{W2DFrost} & \nodecard{W2L7_6} & core & WZ2 Lemma 7.6 hypothesis.\\
\nodecard{W2DKT} & \nodecard{W2L7_6} & api & Katz-Tao output contract.\\
\nodecard{XFO} & \nodecard{W2L7_6} & core & WZ2 Lemma 7.7: content-to-net black box.\\
\nodecard{W2DUnif} & \nodecard{W2L7_9} & api & Uniform output convention.\\
\nodecard{XOSuniform} & \nodecard{W2L7_9} & core & WZ2 [11, Lemma 2.15].\\
\nodecard{W2DUnif} & \nodecard{W2L7_10} & core & WZ2 Lemma 7.10 requires a uniform set.\\
\nodecard{XSWdecomp} & \nodecard{W2L7_10} & core & Quoted scale-entropy decomposition.\\
\nodecard{XSWmerge} & \nodecard{W2L7_10} & core & Merge inversions before selecting the final Frostman scale.\\
\nodecard{W2DFrost} & \nodecard{W2L7_10} & api & Normalized Frostman conclusion at mesh delta/q.\\
\nodecard{W2DKT} & \nodecard{W2L7_11} & core & The absolute local ball count survives separated replication.\\
\nodecard{W2L7_9} & \nodecard{W2ImageUniform} & core & Uniformize the target set after incidence-level pigeonholing.\\
\nodecard{Dtwist} & \nodecard{W2ImageUniform} & api & Fibre-incidence map and target inverse images.\\
\nodecard{Imass} & \nodecard{W2ImageUniform} & api & Mass and target area match on a selected incidence level.\\
\nodecard{W2DFrost} & \nodecard{W2Slice} & core & At least delta\textasciicircum{}(-2+loss) line-parameter cells.\\
\nodecard{W2L7_9} & \nodecard{W2Slice} & core & Uniform refinement of the selected three-parameter slice.\\
\nodecard{W2ImageUniform} & \nodecard{W2Slice} & assembly & Selected projected-incidence shadings are the slice input.\\
\nodecard{W2Slice} & \nodecard{W2JetCount} & assembly & Uniform rich triple set supplied to the zoom/extraction construction.\\
\nodecard{W2L7_10} & \nodecard{W2JetCount} & core & A parameter cube with Frostman rescaling.\\
\nodecard{W2L7_6} & \nodecard{W2JetCount} & core & Extract near-saturated Katz-Tao triples.\\
\nodecard{W2L7_11} & \nodecard{W2JetCount} & core & Replicate while retaining the saturation deficit.\\
\nodecard{L7_3} & \nodecard{W2JetCount} & core & Parameter/jet bi-Lipschitz comparison and cinematic geometry.\\
\nodecard{W2DKT} & \nodecard{W2JetCount} & api & Absolute one-dimensional ball count for C2 curves.\\
\nodecard{W2JetCount} & \nodecard{W2AreaCompare} & assembly & The separated cinematic set and retained projected mass.\\
\nodecard{T7_2} & \nodecard{W2AreaCompare} & core & Cinematic L3/2 maximal estimate with a separately reserved exponent.\\
\nodecard{Imass} & \nodecard{W2AreaCompare} & api & Uniformly dense tubes give projected area mass.\\
\nodecard{Iparams} & \nodecard{W2AreaCompare} & api & Convert the delta-loss to a mesh-ratio loss.\\
\nodecard{W2AreaCompare} & \nodecard{W2L7_12} & core & Local area ratio plus uniform target-cell comparison.\\
\nodecard{W2ImageUniform} & \nodecard{W2L7_12} & core & Transfer the local comparison through the uniform projected set.\\
\nodecard{W2TW} & \nodecard{W2L7_12} & api & Tube-Wolff hypothesis of the one-step lemma.\\
\nodecard{W2L7_9} & \nodecard{IcoarseTW} & core & Uniformity in line-parameter space makes coarse child counts comparable.\\
\nodecard{W2TW} & \nodecard{IcoarseTW} & core & Transfer the normalized radius-squared count to uniform coarse parents.\\
\nodecard{Imass} & \nodecard{IcoarseTW} & api & Coarse shading density and neighbourhood transport.\\
\nodecard{W2L7_12} & \nodecard{W2Iterate} & core & One-step area ratios telescope.\\
\nodecard{IcoarseTW} & \nodecard{W2Iterate} & core & Admissible coarse family and target-neighbourhood compatibility at every step.\\
\nodecard{Iparams} & \nodecard{W2Iterate} & api & Finite iteration loss hierarchy and terminal-scale choice.\\
\nodecard{W2Iterate} & \nodecard{W2P7_1} & core & The terminal coarse curve has positive neighbourhood area.\\
\nodecard{W2P7_1} & \nodecard{W2P6_8} & core & The Tube-Wolff statement is stronger than the required convex-Wolff statement.\\
\nodecard{W2TW} & \nodecard{W2P6_8} & core & Convert convex Wolff to the weaker Tube-Wolff input.\\
\nodecard{W2P6_8} & \nodecard{W2Zero} & core & Projected-area lower bound at exponent sigma/4.\\
\nodecard{Iparams} & \nodecard{W2Zero} & api & Match density, convex errors and scale thresholds before comparison.\\
\nodecard{W2Zero} & \nodecard{W2T5_2} & core & A failed root estimate would produce a positive bad exponent.\\
\nodecard{W2Sigma} & \nodecard{W2T5_2} & core & Counterexample supremum connects deficit zero to the requested estimate.\\
\nodecard{Igeom} & \nodecard{W2T5_2} & api & Bounded-loss conversion returns from chart/cube conventions to original shadings.\\
\bottomrule\end{longtable}
\sectitle{External interfaces and references}
WZ1 and WZ2 bibliography numbers are kept in separate namespaces. In particular, WZ2 [15] is WZ1 itself; its invoked proof is expanded here and is not treated as an unopened external leaf.
\begin{itemize}
\item WZ1 [1] / [8]: endpoint multilinear Kakeya, in the form of WZ1 Theorem 4.2.
\item WZ1 [29]: scalar Lipschitz extension (also a Mathlib interface).
\item WZ1 [35, Lemma 2.9]: Orponen--Shmerkin--Wang radial thin-tube bootstrap; used through WZ1 Proposition 8.6.
\item WZ1 [36, Theorem 1.7] = WZ2 [13]: Pramanik--Yang--Zahl cinematic maximal theorem; used through WZ1 Theorem 7.2 = WZ2 Theorem 7.13.
\item WZ2 [3, Lemma 3.13]: F\"assler--Orponen content extraction; WZ2 Lemma 7.7.
\item WZ2 [11, Lemma 2.15]: Orponen--Shmerkin sparse-scale uniform refinement; WZ2 Lemma 7.9.
\item WZ2 [14, Lemmas 5.20 and 5.21]: Shmerkin--Wang merging and scale-profile decomposition.
\end{itemize}
The reference atlas's historical Kaufman, Dvir, Bourgain ancestry and introductory AD-gap provenance leaves are omitted: they add no independent mathematical input to the requested proof. The in-paper Kaufman-style theorem WZ1 Theorem 8.2 is retained and expanded through its own finite-grid proof.

\textbf{Visual reference.} The supplied \emph{Sticky Kakeya dependency atlas} and Prompt.txt. The external Carleson blueprint graph at \url{https://florisvandoorn.com/carleson/blueprint/dep_graph_document.html} motivates clickable cards and shape conventions only; its mathematics and formalization-status colours are not imported.

\end{document}